\documentclass[11pt,reqno]{amsart}

\usepackage[T1]{fontenc}
\usepackage{lmodern}
\usepackage{microtype}
\usepackage[a4paper,margin=32mm]{geometry}
\usepackage{mathtools}
\usepackage{amssymb}
\usepackage{mathrsfs}
\usepackage{booktabs}
\usepackage{array}
\usepackage{float}
\usepackage{enumitem}
\usepackage{xcolor}
\usepackage{tikz}
\usetikzlibrary{arrows.meta,positioning,shapes.geometric}
% Author: Dr. Denys Dutykh (Khalifa University of Science and Technology, Abu Dhabi, UAE)
% House palette shared by the figures of the article and by their standalone builds.
\definecolor{navy}{HTML}{17365D}
\definecolor{wine}{HTML}{8E2A3F}
\definecolor{teal}{HTML}{007F86}
\definecolor{ochre}{HTML}{B66A19}
\definecolor{plum}{HTML}{6E4B9E}
\definecolor{lilac}{HTML}{9C86C8}
\definecolor{ink}{HTML}{1C1C1C}
\definecolor{slate}{HTML}{5C6670}
\definecolor{mist}{HTML}{C7CDD4}
\definecolor{sand}{HTML}{F1EBDD}
\definecolor{paper}{HTML}{FBF9F4}

\usepackage[colorlinks=true,linkcolor=blue!55!black,citecolor=green!40!black,urlcolor=blue!65!black]{hyperref}
\usepackage{aliascnt}

\newtheorem{theorem}{Theorem}[section]
\newaliascnt{proposition}{theorem}
\newtheorem{proposition}[proposition]{Proposition}
\aliascntresetthe{proposition}
\newaliascnt{lemma}{theorem}
\newtheorem{lemma}[lemma]{Lemma}
\aliascntresetthe{lemma}
\newaliascnt{corollary}{theorem}
\newtheorem{corollary}[corollary]{Corollary}
\aliascntresetthe{corollary}
\newaliascnt{conjecture}{theorem}
\newtheorem{conjecture}[conjecture]{Conjecture}
\aliascntresetthe{conjecture}
\theoremstyle{definition}
\newaliascnt{definition}{theorem}
\newtheorem{definition}[definition]{Definition}
\aliascntresetthe{definition}
\newaliascnt{remark}{theorem}
\newtheorem{remark}[remark]{Remark}
\aliascntresetthe{remark}

\usepackage[nameinlink,noabbrev]{cleveref}
\crefname{theorem}{theorem}{theorems}
\Crefname{theorem}{Theorem}{Theorems}
\crefname{proposition}{proposition}{propositions}
\Crefname{proposition}{Proposition}{Propositions}
\crefname{lemma}{lemma}{lemmas}
\Crefname{lemma}{Lemma}{Lemmas}
\crefname{corollary}{corollary}{corollaries}
\Crefname{corollary}{Corollary}{Corollaries}
\crefname{conjecture}{conjecture}{conjectures}
\Crefname{conjecture}{Conjecture}{Conjectures}
\crefname{definition}{definition}{definitions}
\Crefname{definition}{Definition}{Definitions}
\crefname{remark}{remark}{remarks}
\Crefname{remark}{Remark}{Remarks}

\newcommand{\Z}{\mathbb{Z}}
\newcommand{\Q}{\mathbb{Q}}
\newcommand{\R}{\mathbb{R}}
\newcommand{\abs}[1]{\lvert#1\rvert}
\newcommand{\norm}[1]{\left\lVert#1\right\rVert}
\newcommand{\cont}[1]{\operatorname{cont}\!\left(#1\right)}
\newcommand{\height}[1]{h_{\mathrm{abs}}\!\left(#1\right)}
\newcommand{\dist}[1]{\left\lVert#1\right\rVert}


\newaliascnt{example}{theorem}
\theoremstyle{definition}
\newtheorem{example}[example]{Example}
\aliascntresetthe{example}
\crefname{example}{example}{examples}
\Crefname{example}{Example}{Examples}
\DeclareMathOperator{\tr}{tr}
\DeclareMathOperator{\SL}{SL}
\DeclareMathOperator{\GL}{GL}
\DeclareMathOperator{\arcosh}{arcosh}
\newcommand{\eps}{\varepsilon}
\newcommand{\mc}[1]{\mathcal{#1}}
\frenchspacing

\title[Rational reflections and palindromic words]{Rational reflections and palindromic words\\in the Cohn monoid}

\author{Denys Dutykh}
\address{Mathematics Department, Khalifa University of Science and Technology, PO Box 127788, Abu Dhabi, United Arab Emirates}
\email{denys.dutykh@ku.ac.ae}

\author{Laurent Vuillon}
\address{Univ. Savoie Mont Blanc, CNRS, LAMA, 73000 Chamb\'ery, France}
\email{laurent.vuillon@univ-smb.fr}


\hypersetup{
 pdftitle={Rational reflections and palindromic words in the Cohn monoid},
 pdfauthor={Denys Dutykh and Laurent Vuillon},
 pdfsubject={Standard morphisms, exact word decoding and palindromic Gram reflections},
 pdfkeywords={palindromic closure, standard morphisms, Thue-Morse morphism, Cohn matrices, Markov equation, rational reflections, Frobenius root}
}
\subjclass[2020]{Primary 68R15; Secondary 11D25, 11A55}
\keywords{palindromic closure, standard morphisms, Thue--Morse morphism, Cohn matrices, Markov equation, rational reflections, Frobenius root}
\begin{document}
\begin{abstract}
We study the rational reflection that a palindrome $p$ defines for the two Cohn matrices. For
the reflection of every palindrome we classify all finite words whose terminal columns are taken
to other word columns: they are exactly the concatenations of the blocks $ab\,p$ and $ba\,p$, and
the reflection exchanges these blocks; the projective relation lifts to exact equality.
Consequently two palindromes have equal Gram labels and a root midpoint at the ratio of the centre
exactly when they are exchanged around it, with no hypothesis on a common prefix. Every nontrivial
such pair has nonprimitive completed letter counts, so that no two distinct central occurrences
collide at such a midpoint, and an explicit quadruple of palindromes shows that the midpoint cannot
be dropped. At a standard centre a second proof, by projective cylinders, gives the quantitative
information used away from the midpoint. The exchange conjugates one Gaussian factor of the common
label, and exchanges around nested centres compose into cubes of palindromes with a common label;
the quadruple is the first square, and the midpoints of its two diagonals are shifted from the
centre by half the width of its short rungs. In the Thue--Morse coding of Reutenauer and Vuillon, the Gram label of a
palindrome is the length of an explicit Christoffel word, and the Gram root is its number of
letters $b$, which also fixes where its standard factorization cuts. An exact enumeration of all
palindromes with label below $10^{26}$ finds no coincidence of labels that these exchanges do not
produce, and every class it finds is a single exchange or a square, in support of a conjecture
whose proof would imply Frobenius--Markov uniqueness. In the same coding the length of the
directing word is half the sum of the partial quotients of the label divided by the root; twins on
a reflection axis have directing words of the same length, and for Markov words the equality of
these lengths under a coincidence of numbers is equivalent to uniqueness. Away from the midpoint we
derive the exact two-sided reflection equation, an integral midpoint defect valid for both label
parities, and a fourth-power bound on the common-ancestor label when the parent imbalance is
bounded. These are proved structural restrictions; they do not resolve the uniqueness conjecture.
\end{abstract}
\maketitle
\section{Introduction}
\label{sec:introduction}

The Markov equation
$$
 x^2+y^2+z^2=3xyz
$$
has a tree of positive integer solutions, and the Frobenius--Markov
uniqueness conjecture asks whether a solution is determined, up to permutation, by
its largest entry. One classical way to encode a solution is by a word:
Christoffel words in two letters parametrize the tree, and under the Cohn
representation the largest entry of a triple becomes a matrix invariant of
the corresponding central word. A hypothetical pair of distinct triples
with the same largest entry would then be a pair of distinct central words
with the same invariant, and one can look for structures in the word
setting that forbid such a coincidence. This article isolates one such
structure and determines it completely.

The structure is a reflection. For the two Cohn matrices
$$
 A=\begin{pmatrix}2&1\\1&1\end{pmatrix},\qquad
 B=\begin{pmatrix}5&2\\2&1\end{pmatrix},
$$
a positive matrix word carries three kinds of information at once: a
literal binary word, an integral Gram matrix, and the arithmetic of a root
of $-1$ modulo its first diagonal entry. Each description makes some
operations transparent and hides others. A reflection acts simply on Gram
roots, but its compatibility with positive words turns out to be far more
restrictive, and the question we answer is exactly how restrictive.

The distinction is already visible in a small example. The palindromes
$baab$ and $abba$ have Gram matrices
$$
 \begin{pmatrix}1130&467\\467&193\end{pmatrix},\qquad
 \begin{pmatrix}1130&437\\437&169\end{pmatrix}.
$$
They have the same label and their normalized roots have midpoint $2/5$.
There is no numerical approximation here. The coincidence itself has been noticed
before, as an equality of matrix entries~\cite[p.~4]{LapointeReutenauer2021} and as
a counterexample to a strengthening of the uniqueness conjecture for the values of
lattice paths~\cite[p.~28]{LagisquetEtAl2021}. Yet their completed letter
counts are $(3,3)$, which rules out central words and hence the usual
Christoffel interpretation. In the Thue--Morse coding of Reutenauer and
Vuillon, recalled in \cref{subsec:thue-morse}, the two palindromes become
$abaaaba$ and $aababaa$, and the equality of their labels becomes the equality
of the lengths of two central words; it is the first such coincidence. The issue is to
determine whether this
failure of primitivity is incidental to the example or forced by the
reflection.

We prove that it is forced, and classify every pair responsible for it.
A palindrome $p$ determines a rational reflection $R_p$, which fixes the
terminal column $e=(2,1)^{\mathsf T}$ and is orthogonal for the Gram form
of $p$. The principal word theorem states that, for all finite words $x$
and $y$,
$$
 [R_p\mu(x)e]=[\mu(y)e]
 \quad\Longleftrightarrow\quad
 x\in\{ab\,p,\ ba\,p\}^*,\quad y=\widehat x,
$$
where brackets denote a projective line and the hat exchanges the
blocks $ab\,p$ and $ba\,p$ (\cref{thm:centre-decoder}). This is a
classification for arbitrary finite words, uniformly in the centre, and
the projective equality always lifts to exact equality of columns. When
$p$ is the central word $g_d$ of a standard directive $d$, the blocks are
the images of $ab$ and $ba$ under the standard morphism $\phi_d$, and the
theorem reads $[R_d\nu_d(x)e]=[\nu_d(y)e]$ if and only if
$x\in\{ab,ba\}^*$ and $y=\widehat x$, with $\nu_d=\mu\circ\phi_d$
(\cref{thm:word-decoder}). For this case we give a second proof by
projective cylinders, which also measures how far a reflected column can
move.

For palindromes, the word theorem has a stronger arithmetic consequence.
Equal labels $M$ and the midpoint condition
$$
 \frac{\rho_X+\rho_Y}{2M}=\frac{\rho_p}{M_p}
$$
are equivalent to a full matrix conjugacy by $R_p$, and the decoder then
forces the two palindromes to be $p\,X$ and $p\,\widehat X$, with $X$ a
concatenation of the blocks (\cref{cor:centre-classification}): no common
prefix needs to be assumed, and the prefix $p$ is forced. If $X$ has $2k$
blocks, the completed counts of both palindromes are $2k+1$ times those of
$p$, so they are not primitive when $k>0$. For a standard ancestor this
gives the preimages $\widetilde w w$ and $\widetilde{\widehat w}\,\widehat w$,
with $w$ a concatenation of $ab$ and $ba$, and completed counts
$(2k+1)L_d(1,1)^{\mathsf T}$ (\cref{thm:palindromic-classification}).
Thus the reflection creates an infinite family of exact palindromic
label equalities while excluding every nontrivial central pair at the
midpoint of any palindrome.

We also retain the information that changes when the midpoint moves. The
conjugacy becomes a two-sided matrix equation with a rank-one term, and
its action on a centered reciprocal cusp coordinate is an affine
reflection. For a hypothetical pair of distinct central words with equal
labels, reduced integral coordinates for the midpoint supply a nonzero
integer defect, and the first-letter cylinders of the standard generators
convert that defect into a quantitative restriction on the common
ancestor. The precise statement is \cref{thm:displaced-ancestor-bound}:
when the two parent traces are comparable, the ancestor label $M_0$ is
bounded by a constant multiple of $M^{1/4}$, the constant being explicit
in the data of the actual pair, with no factorization of $M$ required and
no restriction on the parity of $M$. Along a one-sided directive the two
parents grow apart, the constant degrades, and the bound weakens to a
cubic scale. Neither regime excludes all collisions.

The block exchange around a palindrome $p$ (\cref{prop:any-centre}) has a simple
arithmetic: the label of $p$ divides the common label, and the exchange conjugates
the complementary Gaussian factor (\cref{prop:exchange-arithmetic}). Exchanges
around nested centres compose. When the centre of an exchange is itself one of two
exchanged palindromes, the four words built around them form a square of palindromes
with one label, and $r$ nested levels form a cube with $2^r$ vertices
(\cref{thm:nested-cubes}). The pairs of such a square that are not exchanges have their
root midpoints shifted off every centre, by exactly half the width of a short rung
(\cref{cor:nested-square}). In the
Thue--Morse coding the label is the length of an explicit Christoffel word and the
root its number of letters $b$ (\cref{prop:thue-morse}). An exact enumeration of
all palindromes with label below $10^{26}$ finds no coincidence of labels that the
exchanges do not produce, and every class of equal labels it finds is a single exchange
or a square. This is the content of \cref{conj:structure}, whose proof would imply the
uniqueness conjecture; \cref{sec:outlook} explains how far the reflection goes
towards it. It also measures one consequence of the conjecture: two palindromes with the
same label would have directing words of the same length in the Thue--Morse coding. This
length is half the sum of the partial quotients of the label divided by the root, and for
Markov words the fixed-sum theorem of Lagisquet, Pelantov\'a, Tavenas and
Vuillon~\cite{LagisquetEtAl2021} turns the equality of lengths into an equivalent form of the
uniqueness conjecture (\cref{prop:lengths-uniqueness}). On every reflection axis the
equality holds, since twins there are exchanges.

The prescribed midpoint is essential: a general equal-label pair need not
have it (\cref{rem:quadruple} exhibits one), and away from it the inequalities
above remain necessary conditions rather than exclusions. What the article contributes is the
complete word classification at a specified reflection, together with a
quantitative account of the information that survives off its axis. The
Frobenius--Markov uniqueness conjecture itself is not resolved here.

\subsection{Relation to earlier work}
\label{sec:literature}

The standard morphisms and palindromic extensions used here belong to
classical Sturmian combinatorics. Justin's functional identity relates
iterated palindromic closure to composition of standard
morphisms~\cite[Lemma~2.1]{Justin2005}. Its left and right forms, with the
same conventions used below, appear in Perrin and
Reutenauer~\cite[Section~2]{PerrinReutenauer2023}. The incidence formulas
of Berth\'e, de Luca and Reutenauer describe the completed letter counts
of central words and their standard factorizations~\cite[Section~3]{BertheDeLucaReutenauer2008}.
These results supply the word transport and the primitive-count test in
our application. The involution obtained by reversing a directive in
\cite{BertheDeLucaReutenauer2008} changes the central word by Christoffel
duality; our reflection instead acts on a positive matrix orbit while the
directive remains fixed.

The language that appears in the answer is also familiar. If
$\theta(a)=ab$ and $\theta(b)=ba$, then
$(ab\mid ba)^*=\theta(\{a,b\}^*)$ is the image of the Thue--Morse
substitution. Reutenauer and Vuillon use this substitution together with
palindromic and antipalindromic closure to construct Markov
numbers~\cite{ReutenauerVuillon2017}: in their coding a Markov number is two more than the
length of $\operatorname{Pal}(\theta(\operatorname{Pal}(av)))$, and the injectivity of this
length is equivalent to the uniqueness conjecture~\cite[Theorem~1]{ReutenauerVuillon2017}.
\Cref{prop:thue-morse} identifies that length plus two, the length of the Christoffel word
$a\operatorname{Pal}(\theta(F_a(z)))b$, with our Gram label, for every palindrome $z$.
Abram, Lapointe and Reutenauer
extend the construction to whole Markov triples~\cite{AbramLapointeReutenauer2020}.
Our conclusion concerns the role of this language in a different
problem: it exhausts the finite words whose positive terminal vectors
are carried to another such vector by the prescribed rational reflection.
The proof establishes this exhaustion uniformly for the reflection of every
palindrome, before imposing palindromicity on the reflected words.

Palindromic symmetry is well established in the geometry of two-generator
groups. Kassel and Reutenauer prove that a cyclically reduced basis with
both words of odd length has a unique palindromic conjugate
basis~\cite{KasselReutenauer2007}. Gilman and Keen relate palindromic
words to axes perpendicular to the common perpendicular of the generator
axes~\cite{GilmanKeen2009}. These results concern primitive bases or
geometric axes. The decoder here retains the positive word, the finite
terminal vector and the fixed rational reflection, and applies to every
finite word, including nonprimitive ones.

The distinction between a full matrix and one of its entries matters
when connecting the result to Markov arithmetic. Reutenauer's
Christoffel correspondence is a bijection with whole proper Markov
triples~\cite[Theorem~1]{Reutenauer2009}. More recently, Veselov identifies
Markov fractions with Cohn matrix indices and describes continued
fractions through a mirrored Conway topograph~\cite[Sections~3--4]{Veselov2026}.
Jouteur, Paris-Romaskevich and Thomas characterize modular involutions
exchanging rational pairs and study Springborn operations, including a
$q$-deformed midpoint formula for Farey neighbors~\cite[Sections~5--7]{JouteurParisRomaskevichThomas2026}.
Those constructions retain a tree position, an arithmetic regularity
condition or a whole rational fraction. They do not identify which
arbitrary positive words survive reflection in the orbit considered here.

There are also injectivity theorems for restricted word languages.
Lapointe and Reutenauer prove that a Cohn matrix entry is strictly
increasing in military order on the language of one fixed Sturmian
sequence~\cite[Corollary~1]{LapointeReutenauer2021}. Our comparison allows
arbitrary palindromes, whose completions need not lie in one such
language. Their Gram labels are positive integers, and are Markov
numbers under the additional Christoffel hypothesis. Equal labels and a
specified root midpoint recover a full Gram congruence, which brings the
arbitrary-word decoder into play. This yields an exact classification of
the reflected pairs and shows why every nonempty pair fails the primitive
completed-count condition. No claim is made that an arbitrary equal-label
pair has the prescribed midpoint.


\subsection{Organization}

\Cref{sec:standard-grams} fixes the literal standard morphisms,
completed counts and Gram normalization, and reads the label and the root in the
Thue--Morse coding. \Cref{sec:decoder} proves the finite-word decoder by separating
projective cylinders and reading two letters at a time. \Cref{sec:midpoint} derives
the palindromic classification and the central-word obstruction, extends the
construction, the decoder and the classification to every palindromic centre, shows
that the midpoint hypothesis cannot be removed, and proves that nested exchanges
produce cubes of twins. \Cref{sec:displaced} treats moving midpoints, explains the
weighted scale that must be retained, and proves the common-ancestor bound.
\Cref{sec:outlook} states the conjecture on twins, explains the ladders in which
they appear, reads the lengths of the directing words, and isolates the remaining
mathematical question.

\section{Standard words and their Gram matrices}
\label{sec:standard-grams}

We first fix the word conventions. This is essential because two
standard morphisms can have the same incidence matrix and different
literal images. Throughout, a word is a finite word over $\{a,b\}$,
$\widetilde z$ denotes reversal, and $\eps$ denotes the empty word.
A palindrome satisfies $z=\widetilde z$. Products of matrices follow
the order of letters from left to right.

\subsection{Palindromic closure and standard morphisms}

For a word $z$, let $z^+$ be the shortest palindrome having $z$ as a
prefix. The iterated right palindromization map is defined by
$$
 \operatorname{Pal}(\eps)=\eps,\qquad
 \operatorname{Pal}(dc)=(\operatorname{Pal}(d)c)^+,
 \quad c\in\{a,b\}.
$$
For example, $\operatorname{Pal}(ab)=aba$ and
$\operatorname{Pal}(aba)=abaaba$. We call the values of this map
\emph{central words}. This is the usual binary central-word class;
its relation with Christoffel words is recalled below~\cite[Section~3]{BertheDeLucaReutenauer2008}.
The word $d$ is called a \emph{directive}.

Use the right standard morphisms
\begin{equation}\label{eq:right-standard}
 \phi_a:\ a\mapsto a,\ b\mapsto ba,
 \qquad
 \phi_b:\ a\mapsto ab,\ b\mapsto b,
 \qquad \phi_{dc}=\phi_d\circ\phi_c.
\end{equation}
The corresponding left morphisms are
$\psi_a(a)=a$, $\psi_a(b)=ab$, $\psi_b(a)=ba$, and
$\psi_b(b)=b$. Define, for every word $z$,
$$
 F_c(z)=c\phi_c(z)=\psi_c(z)c,\qquad
 F_{dc}=F_d\circ F_c,\qquad F_\eps(z)=z.
$$
The equality in this definition follows by moving the final $c$
across the image of each letter. Reversing it shows that $F_c$
preserves palindromes. Justin's functional identity~\cite[Lemma~2.1]{Justin2005},
in its right-morphism form~\cite[Section~2]{PerrinReutenauer2023}, identifies
$F_d(\eps)$ with $\operatorname{Pal}(d)$; in the present convention
it gives
\begin{equation}\label{eq:standard-palindrome-image}
 F_d(z)=g_d\phi_d(z),\qquad g_d=\operatorname{Pal}(d),
 \qquad F_d(\operatorname{Pal}(t))=\operatorname{Pal}(dt).
\end{equation}
The first equality also follows directly by induction on $|d|$
once $F_d(\eps)=g_d$ is known. Thus $F_d$ extends the action on
central words to every binary palindrome.

The completed-count formulas below are the incidence form of the
standard central-word construction~\cite[Corollaries~3.1--3.2]{BertheDeLucaReutenauer2008}.
We include the short proof to fix the composition convention.

\begin{lemma}[Literal identities and completed counts]
\label{lem:standard-identities}
For every directive $d$,
\begin{equation}\label{eq:literal-standard-identities}
 \phi_d(ab)=abg_d,\qquad \phi_d(ba)=bag_d.
\end{equation}
Let $L_d$ be the incidence matrix of $\phi_d$, with columns recording
the numbers of $a$'s and $b$'s in the images of $a$ and $b$.
For the \emph{completed count column}
$$
 c(z)=\binom{|z|_a+1}{|z|_b+1},
$$
one has
\begin{equation}\label{eq:completed-count-transport}
 c(F_d(z))=L_dc(z),\qquad
 L_a=\begin{pmatrix}1&1\\0&1\end{pmatrix},\quad
 L_b=\begin{pmatrix}1&0\\1&1\end{pmatrix}.
\end{equation}
In particular, the greatest common divisor of the two completed
counts is unchanged by $F_d$, and every central word has primitive
completed counts.
\end{lemma}

\begin{proof}
The identities hold for the empty directive. Write $u=\phi_d(a)$
and $v=\phi_d(b)$. Appending $a$ changes $(u,v)$ to $(u,vu)$ and
$g_d$ to $g_du$; appending $b$ changes them to $(uv,v)$ and $g_dv$.
If $uv=abg_d$ and $vu=bag_d$, these updates immediately preserve
both identities in \eqref{eq:literal-standard-identities}.
Counting their letters gives
$c(g_d)=L_d(1,1)^{\mathsf T}$. Counting
$F_d(z)=g_d\phi_d(z)$ then proves
\eqref{eq:completed-count-transport}. Each $L_d$ is an integral
unimodular matrix. Such a matrix preserves the ideal generated by
the entries of a column, and therefore their greatest common divisor.
Apply this to $c(\eps)=(1,1)^{\mathsf T}$ for the last assertion.
\end{proof}

Primitivity is only a necessary condition for centrality among
palindromes. We shall use precisely that implication, never its
converse. The maps $F_d$ are injective: each of the two morphisms in
\eqref{eq:right-standard} has a prefix code as its set of letter
images, so is injective; a fixed prefix $g_d$ can then be cancelled.

\subsection{A positive integral Gram model}

The Cohn representation used in this paper is
\begin{equation}\label{eq:cohn-generators}
 \mu(a)=A=\begin{pmatrix}2&1\\1&1\end{pmatrix},\qquad
 \mu(b)=B=\begin{pmatrix}5&2\\2&1\end{pmatrix},\qquad
 P=\begin{pmatrix}2&1\\1&0\end{pmatrix},\quad e=\binom21.
\end{equation}
The matrices $A,B$ are symmetric, positive definite and unimodular.
Consequently $\mu(\widetilde z)=\mu(z)^{\mathsf T}$, and for each
palindrome $z$ the matrix $\mu(z)$ is positive definite.
Indeed, if $z=\widetilde w w$, it is $\mu(w)^{\mathsf T}\mu(w)$;
if $z=\widetilde w c w$, it is
$\mu(w)^{\mathsf T}\mu(c)\mu(w)$.

Define its \emph{Gram label} $M(z)$ and \emph{Gram root} $\rho(z)$ by
\begin{equation}\label{eq:palindromic-gram}
 \mc G(z)=P\mu(z)P
 =\begin{pmatrix}M(z)&\rho(z)\\\rho(z)&H(z)\end{pmatrix}.
\end{equation}
These are positive integers, and
\begin{equation}\label{eq:gram-determinant}
 M(z)H(z)-\rho(z)^2=1.
\end{equation}
In particular, $\rho(z)^2\equiv-1\pmod{M(z)}$.
The label is also a primitive sum of two squares. To see this directly,
put $P_1=\left(\begin{smallmatrix}1&1\\1&0\end{smallmatrix}\right)$
and $P_2=P$. Then $A=P_1^2$, $B=P_2^2$, and the factorizations above
give $\mc G(z)=T^{\mathsf T}T$ with an integral unimodular $T$:
use $T=\mu(w)P$ in the even case and $T=P_c\mu(w)P$ in the odd
case, where $P_a=P_1$ and $P_b=P_2$. The first column of $T$ is
positive and primitive, so its squared Euclidean norm is $M(z)$.

There is also a direct connection with the more familiar trace label.
For $J_*=AP^{-1}=\left(\begin{smallmatrix}1&0\\1&-1\end{smallmatrix}\right)$,
one has
\begin{equation}\label{eq:completed-word-trace}
 \mu(azb)=J_*\mc G(z)P,
 \qquad \tr\mu(azb)=3M(z).
\end{equation}
Both formulas follow by multiplication, using $B=P^2$.
If $z$ is central, $azb$ is a proper lower Christoffel word, and
the classical Cohn correspondence identifies $M(z)$ as a Markov
number~\cite[Theorem~1]{Reutenauer2009}. Here a Markov number is an entry of a positive integral
solution of $x^2+y^2+z^2=3xyz$. For a general palindrome,
\eqref{eq:completed-word-trace} remains an exact identity, but it
does not assert that the label is a Markov number.

Among palindromes with prescribed completed counts, the central word is
distinguished by the label itself. Lagisquet, Pelantov\'a, Tavenas and
Vuillon prove a minimum for the upper-right path value at fixed counts,
with its exact equality cases~\cite[Proposition~21]{LagisquetEtAl2021};
transported to the present Gram normalization, it says that for coprime
completed counts every palindrome satisfies $M(z)\geqslant M(z_0)$,
where $z_0$ is the central word with those counts, and that equality
holds only for $z_0$. We record this because it fixes the direction of
the comparison used below: a competing palindrome with the same counts
can only have a larger label, never a smaller one. The statement is
theirs; nothing in this paper strengthens it, and the noncoprime case is
not used.

\subsection{What the Gram root counts}
\label{subsec:gram-root-counts}

The identity \eqref{eq:gram-determinant} makes $\rho(z)$ a square root
of $-1$ modulo the label. That root also has a word-theoretic meaning,
and the incidence transport of
\eqref{eq:completed-count-transport} is enough to find it. Write
$\widehat d$ for the \emph{reversal complement} of a word $d$, obtained
by reading $d$ backwards and exchanging the two letters.

\begin{lemma}[The reversal complement transposes the incidence matrix]
\label{lem:incidence-transpose}
For every directive $d$, $L_{\widehat d}=L_d^{\mathsf T}$.
\end{lemma}

\begin{proof}
The two matrices of \eqref{eq:completed-count-transport} satisfy
$L_a^{\mathsf T}=L_b$ and $L_b^{\mathsf T}=L_a$, so exchanging the
letters of a directive transposes each factor of its incidence product.
Reading the directive backwards reverses the order of those factors.
Performing both operations therefore transposes the product.
\end{proof}

\begin{proposition}[The doubled directive and its letter counts]
\label{prop:doubled-directive-counts}
Let $d$ be a directive and let
$$
 \binom qs=c(\operatorname{Pal}(d))=L_d\binom11
$$
be its completed counts, so that $\gcd(q,s)=1$. Put
$m=q^2+s^2$. Then the completed counts of the doubled central word
$\operatorname{Pal}(\widehat d\,d)$ are
\begin{equation}\label{eq:doubled-counts}
 c\bigl(\operatorname{Pal}(\widehat d\,d)\bigr)
 =L_d^{\mathsf T}L_d\binom11 ,
\end{equation}
so the completed word $a\operatorname{Pal}(\widehat d\,d)b$ has length
exactly $m$.\footnote{That a Markov number is a sum of two squares was
obtained earlier, by a substitution argument on words resting on Justin's
formula, by L.~Vuillon and A.~De~Luca (personal communication, 2026). The
proof below is independent of theirs and yields the letter counts as well as
the length.} Writing $r_d$ for its number of letters $b$, one has
$0<r_d<m$ and
\begin{equation}\label{eq:doubled-root}
 r_d\,q\equiv s \pmod m,
 \qquad\hbox{hence}\qquad r_d^{\,2}\equiv-1\pmod m .
\end{equation}
The number of letters $a$ is $m-r_d$.
\end{proposition}

\begin{proof}
By \eqref{eq:standard-palindrome-image},
$\operatorname{Pal}(\widehat d\,d)=F_{\widehat d}(\operatorname{Pal}(d))$,
so \eqref{eq:completed-count-transport} and
\cref{lem:incidence-transpose} give \eqref{eq:doubled-counts}. Write
$L_d=\left(\begin{smallmatrix}\alpha&\beta\\\gamma&\delta\end{smallmatrix}\right)$,
so that $q=\alpha+\beta$ and $s=\gamma+\delta$, and
$\alpha\delta-\beta\gamma=1$. Expanding \eqref{eq:doubled-counts},
$$
 m-r_d=\alpha q+\gamma s,\qquad r_d=\beta q+\delta s,
$$
whose sum is $(\alpha+\beta)^2+(\gamma+\delta)^2=q^2+s^2$, which is the
length assertion.

Both $\alpha$ and $\delta$ are at least one, since $L_d$ has nonnegative
entries and determinant one, and $\alpha=0$ would force
$-\beta\gamma=1$. Hence $r_d\geqslant\delta s\geqslant1$ and
$m-r_d\geqslant\alpha q\geqslant1$, which is the strict range.

Finally $\gcd(q,m)=\gcd(q,s^2)=1$, so $q$ is invertible modulo $m$, and
$q^2\equiv-s^2$. Therefore
$$
 r_d\,q=\beta q^2+\delta sq
 \equiv-\beta s^2+\delta sq
 =s\bigl(\delta(\alpha+\beta)-\beta(\gamma+\delta)\bigr)
 =s(\alpha\delta-\beta\gamma)=s \pmod m .
$$
Squaring and cancelling the unit $q^2$ gives $r_d^{\,2}\equiv-1$.
\end{proof}

The two letter counts of that word are not merely arithmetic data: they are
also the two lengths into which it splits. Recall that a proper Christoffel
word has a unique factorization into two Christoffel words, its
\emph{standard factorization}, and that the lengths of the two factors are the
column sums of the incidence matrix of its
directive~\cite[Corollaries~3.1--3.2]{BertheDeLucaReutenauer2008}.

\begin{corollary}[The cut and the counts are the same pair]
\label{cor:cut-at-letter-count}
Let $d$ be a directive, let
$\Omega=a\operatorname{Pal}(\widehat d\,d)b$, and let
$\Omega=\Omega_1\Omega_2$ be its standard factorization. Then
\begin{equation}\label{eq:cut-at-letter-count}
 \abs{\Omega_1}=\abs{\Omega}_a=m-r_d,
 \qquad
 \abs{\Omega_2}=\abs{\Omega}_b=r_d .
\end{equation}
That is, $\Omega$ is cut after exactly as many letters as it contains
letters $a$.
\end{corollary}

\begin{proof}
The directive of $\Omega$ is $\widehat d\,d$, so by
\cref{lem:incidence-transpose} its incidence matrix is
$L_{\widehat d\,d}=L_{\widehat d}L_d=L_d^{\mathsf T}L_d$. This is a Gram
matrix, hence symmetric, and therefore its column sums and its row sums are
the same ordered pair:
$$
 \mathbf 1^{\mathsf T}L_d^{\mathsf T}L_d
 =\bigl(L_d^{\mathsf T}L_d\mathbf 1\bigr)^{\mathsf T},
 \qquad \mathbf 1=(1,1)^{\mathsf T}.
$$
The left side is the pair of factor lengths of the standard factorization,
and the right side is, by \eqref{eq:doubled-counts}, the pair of completed
counts of $\Omega$. The identification of the individual entries is the
order of \eqref{eq:doubled-counts}, and $\abs{\Omega}_a=m-r_d$ was proved in
\cref{prop:doubled-directive-counts}.
\end{proof}

Symmetry of a Gram matrix is the whole content: the cut point and the letter
count are one pair of integers read along the two axes of a single matrix.
Exchanging the two letters of $d$ replaces that matrix by the Gram matrix of
the rows of $L_d$, which has the same total mass and the two entries of
\eqref{eq:cut-at-letter-count} interchanged, so the cut is reflected in the
middle of the word.

The Gram model sees the same root. Recall the factorization
$\mc G(z)=T^{\mathsf T}T$ constructed above, with $T=\mu(w)P$ in the even
case and $T=P_c\mu(w)P$ in the odd case.

\begin{corollary}[The Gram root as a letter count]
\label{cor:gram-root-letter-count}
Let $z$ be a palindrome, let $(q,s)$ be the first column of $T$, and put
$\eps_z=\det T$, which is $+1$ when $\abs z$ is odd and $-1$ when
$\abs z$ is even. Then $M(z)=q^2+s^2$ with $\gcd(q,s)=1$, and
$$
 \rho(z)\,q\equiv\eps_z\,s \pmod{M(z)} .
$$
Let $d$ be the unique directive with completed counts $(q,s)$. Then
$\rho(z)=r_d$ when $\eps_z=+1$, and $\rho(z)=M(z)-r_d$ when
$\eps_z=-1$.
\end{corollary}

\begin{proof}
Let $(\beta,\delta)$ be the second column of $T$. Then $M(z)=q^2+s^2$
and $\rho(z)=\beta q+\delta s$ are the two entries of the first row of
$T^{\mathsf T}T$, and $\det T=q\delta-s\beta=\eps_z$. Since
$\det\mu(w)=1$, $\det P=\det P_1=-1$, the two cases of the construction
give $\eps_z=-1$ and $\eps_z=+1$ respectively. As in the proposition,
$q^2\equiv-s^2$ modulo $M(z)$, so
$$
 \rho(z)\,q=\beta q^2+\delta sq\equiv s(q\delta-s\beta)=\eps_z s .
$$
The first column of $T$ is positive and primitive, so $(q,s)$ is the
completed-count column of exactly one central word, namely
$\operatorname{Pal}(d)$ for a unique directive $d$; this is the
inverse of the unimodular transport
\eqref{eq:completed-count-transport}. Both $\rho(z)$ and $r_d$ lie
strictly between $0$ and $M(z)$, and by
\eqref{eq:doubled-root} they satisfy the same congruence up to the sign
$\eps_z$. A residue class modulo $M(z)$ meets that range once, so they
agree when $\eps_z=+1$ and are complementary when $\eps_z=-1$.
\end{proof}

When $z$ is central, $azb$ is a proper lower Christoffel word and $M(z)$
is a Markov number, so $\rho(z)$ is then the classical root of $-1$
attached to the occurrence, and $q^2+s^2$ is its primitive two-square
representation. \Cref{prop:doubled-directive-counts} says that this root
is not an abstract residue: it is the number of letters $b$ in one
explicit Christoffel word of length $M(z)$, the one directed by the
doubled directive $\widehat d\,d$. The statement itself needs no
centrality, and its proof uses only \eqref{eq:completed-count-transport}
and a determinant. \Cref{cor:cut-at-letter-count} adds that the same integer
also measures where that word splits: the standard factorization cuts after
$M(z)-\rho(z)$ letters when $\eps_z=+1$, and after $\rho(z)$ letters when
$\eps_z=-1$. A root of $-1$ modulo a Markov number is thus visible twice in a
single word, once as a count and once as a position.

\subsection{The Thue--Morse coding}
\label{subsec:thue-morse}

The label and the root have a second reading, in the coding of Markov numbers by Reutenauer
and Vuillon~\cite{ReutenauerVuillon2017}, and a search for equal lengths among all palindromes of
that coding finds the pairs of \cref{sec:midpoint}. Let $\theta$ be
the Thue--Morse morphism, $\theta(a)=ab$ and $\theta(b)=ba$, and recall the matrix
$P_1=\left(\begin{smallmatrix}1&1\\1&0\end{smallmatrix}\right)$ of the Gram factorization, with
$P_1^2=A$.

\begin{proposition}[The Thue--Morse coding]
\label{prop:thue-morse}
For every word $z$,
\begin{equation}\label{eq:thue-morse-incidence}
 L_{\theta(F_a(z))}=P_1\,\mu(z)\,P_1 .
\end{equation}
Consequently, for every palindrome $z$ the Christoffel word
$\Omega_z=a\operatorname{Pal}(\theta(F_a(z)))b$ has length $M(z)$ and contains exactly
$\rho(z)$ letters $b$. Moreover, $F_a$ is a bijection from the set of palindromes onto the
set of palindromes that begin with $a$ and do not contain the factor $bb$.
\end{proposition}

\begin{proof}
The standard morphism of a concatenation of directives is the composition of their
standard morphisms, so $L_{uv}=L_uL_v$; since $\theta$ is a morphism, it suffices to compare
the images of the generators. On the word side, $L_{\theta(a)}=L_aL_b=A$ and
$$
 L_{\theta(\phi_a(b))}=L_{baab}=(L_bL_a)(L_aL_b)
 =\begin{pmatrix}1&1\\1&2\end{pmatrix}\begin{pmatrix}2&1\\1&1\end{pmatrix}
 =\begin{pmatrix}3&2\\4&3\end{pmatrix}.
$$
On the matrix side, $P_1^{-1}AP_1=A$ because $A=P_1^2$, and
$$
 P_1^{-1}BP_1=\begin{pmatrix}0&1\\1&-1\end{pmatrix}\begin{pmatrix}7&5\\3&2\end{pmatrix}
 =\begin{pmatrix}3&2\\4&3\end{pmatrix}.
$$
Since $\theta(\phi_a(a))=\theta(a)$, this proves $L_{\theta(\phi_a(c))}=P_1^{-1}\mu(c)P_1$ for
both letters $c$. The product over the letters of $z$ telescopes to
$L_{\theta(\phi_a(z))}=P_1^{-1}\mu(z)P_1$, and therefore
$L_{\theta(F_a(z))}=L_{\theta(a)}L_{\theta(\phi_a(z))}=P_1^2P_1^{-1}\mu(z)P_1$, which is
\eqref{eq:thue-morse-incidence}.

Every directive $D$ satisfies $c(\operatorname{Pal}(D))=L_D(1,1)^{\mathsf T}$, by
\eqref{eq:standard-palindrome-image} and \eqref{eq:completed-count-transport} applied to the
empty word. Take $D=\theta(F_a(z))$ and write $\mu(z)e=(x,y)^{\mathsf T}$. Since
$P_1(1,1)^{\mathsf T}=e$, the identity \eqref{eq:thue-morse-incidence} gives
$$
 c(\operatorname{Pal}(D))=P_1\mu(z)e=\binom{x+y}{x}.
$$
For a palindrome $z$ one has $M(z)=e^{\mathsf T}\mu(z)e=2x+y$, because the first column of
$P$ is $e$, and $\rho(z)=(1,0)\,\mu(z)e=x$, because the second row of $P$ is $(1,0)$. The
completed counts of $\operatorname{Pal}(D)$ are the letter counts of $\Omega_z$, which are
therefore $M(z)-\rho(z)$ and $\rho(z)$. The word $\Omega_z$ is a Christoffel word because
$\operatorname{Pal}(D)$ is central.

For the last assertion, recall $F_a(z)=a\phi_a(z)=\psi_a(z)a$. The reversal of
$\phi_a(z)$ is $\psi_a(\widetilde z)$, so the reversal of $F_a(z)$ is $F_a(\widetilde z)$. If
$z$ is a palindrome, $F_a(z)$ is thus a palindrome; it begins with $a$, and it has no factor
$bb$ because every $b$ of $\phi_a(z)$ is followed by an $a$. Conversely, let $u=au'$ be a
palindrome without the factor $bb$. The word $u'$ is empty or ends with $a$, so each of its
letters $b$ is followed by the letter $a$, and $u'$ factors uniquely over the code
$\{a,ba\}$: $u'=\phi_a(z)$ and $u=F_a(z)$. Finally $F_a(z)=u=F_a(\widetilde z)$, and the
injectivity of $F_a$ gives $z=\widetilde z$.
\end{proof}

For a central word $z=\operatorname{Pal}(v)$ one has $F_a(z)=\operatorname{Pal}(av)$ by
\eqref{eq:standard-palindrome-image}. Then $\Omega_z$ is the word
$a\operatorname{Pal}(\theta(\operatorname{Pal}(av)))b$ of Reutenauer and Vuillon, and the length
statement is their theorem that $2+\abs{\operatorname{Pal}(\theta(\operatorname{Pal}(av)))}$
is the Markov number of the occurrence~\cite[Theorem~1]{ReutenauerVuillon2017}. The
proposition adds the root. By \eqref{eq:completed-word-trace} the first row of the Cohn
matrix $\mu(azb)=J_*\mc G(z)P$ is $(M,\rho)P=(2M+\rho,\,M)$, so the number of letters $b$ of
their word is $\mu(azb)_{11}-2M$, for every occurrence and with no choice of sign. For an
arbitrary palindrome, $F_a(z)$ is an arbitrary palindrome that begins with $a$ and avoids
$bb$, and a coincidence of the lengths
$\abs{\operatorname{Pal}(\theta(F_a(z)))}$ is exactly an equality of Gram labels. The
pair of the introduction is the first one: with $F_a(baab)=abaaaba$ and $F_a(abba)=aababaa$,
the directives $\theta(abaaaba)$ and $\theta(aababaa)$ give central words whose completed counts
are $(663,467)$ and $(693,437)$, and both completed words have length $1130$.

The Thue--Morse directive is the doubled directive of
\cref{prop:doubled-directive-counts}, which explains the sign in
\cref{cor:gram-root-letter-count}.

\begin{corollary}[The Thue--Morse directive is a doubled directive]
\label{cor:thue-morse-doubled}
Let $z$ be a palindrome and let $d$ be the directive of \cref{cor:gram-root-letter-count},
whose completed counts are the first column of $T$. Then $\theta(F_a(z))=\widehat d\,d$ when
$\abs z$ is odd, and $\theta(F_a(z))$ is obtained from $\widehat d\,d$ by exchanging the two
letters when $\abs z$ is even. In particular the sign $\eps_z$ of
\cref{cor:gram-root-letter-count} is the exchange of letters, and the Gram root is the number
of letters $b$ of $\Omega_z$ in both parities.
\end{corollary}

\begin{proof}
Every matrix of $\SL_2(\Z)$ with nonnegative entries is a product of $L_a$ and $L_b$ in
exactly one way, and the column $L_{d'}(1,1)^{\mathsf T}$ determines the directive $d'$, as in
the proof of \cref{cor:gram-root-letter-count}. In the odd case, $z=\widetilde wcw$, the
matrix $L=P_c\mu(w)P_1$ has determinant $(-1)\cdot(-1)=1$ and nonnegative entries, and
$L(1,1)^{\mathsf T}=P_c\mu(w)e=Te_1$. Hence $L=L_d$, and by \cref{lem:incidence-transpose}
$$
 L_{\widehat d\,d}=L_d^{\mathsf T}L_d=P_1\mu(\widetilde w)P_c^2\mu(w)P_1=P_1\mu(z)P_1,
$$
because $P_c$ is symmetric and $P_c^2=\mu(c)$. In the even case, $z=\widetilde ww$, the
matrix $L=\mu(w)L_a$ has determinant one and nonnegative entries, and
$L(1,1)^{\mathsf T}=\mu(w)e=Te_1$; hence $L=L_d$ and
$L_{\widehat d\,d}=L_b\mu(\widetilde w)\mu(w)L_a=L_b\mu(z)L_a$. With
$S=\left(\begin{smallmatrix}0&1\\1&0\end{smallmatrix}\right)$ one has $P_1=L_aS=SL_b$, so
$P_1\mu(z)P_1=S\,L_{\widehat d\,d}\,S$, and conjugation by $S$ exchanges $L_a$ and $L_b$, that
is, the two letters of a directive. In both cases \eqref{eq:thue-morse-incidence} and the
uniqueness of the factorization prove the claim. Exchanging the letters of a directive
exchanges the two letter counts of its central word, which gives the statement about
$\eps_z$ by comparison with \cref{cor:gram-root-letter-count}.
\end{proof}

The coding turns the question of this article into a question about the lengths of central
words. The pairs classified in \cref{sec:midpoint} are coincidences of these lengths, and
\cref{sec:outlook} asks whether they are the only ones.

\subsection{The reflection attached to an ancestor}

Fix a directive $d$ and abbreviate
\begin{equation}\label{eq:ancestor-data}
 G=\mu(g_d),\quad \nu_d=\mu\circ\phi_d,\quad
 M_0=e^{\mathsf T}Ge,\quad \rho_0=(Ge)_1,\quad r_0=\rho_0/M_0.
\end{equation}
The entries of $Ge$ are positive, so $0<r_0<1/2$.
Equation~\eqref{eq:standard-palindrome-image} becomes
\begin{equation}\label{eq:image-gram}
 \mc G(F_d(z))=PG\nu_d(z)P.
\end{equation}
In particular, $G\nu_d(z)$ is symmetric for a palindrome $z$,
although $\nu_d(z)$ itself need not be symmetric.

The ancestor determines the rational matrix
\begin{equation}\label{eq:ancestor-reflection}
 R_d=\frac{2ee^{\mathsf T}G}{M_0}-I.
\end{equation}
It fixes $e$, squares to $I$, and satisfies
$R_d^{\mathsf T}G=GR_d$ and $R_d^{\mathsf T}GR_d=G$.
Thus it is an orthogonal reflection for the inner product defined by
$G$. For example, if $G=T^{\mathsf T}T$, then $TR_dT^{-1}$ is the
ordinary Euclidean reflection in the line spanned by $Te$.
In the Gram coordinates its form is particularly simple:
\begin{equation}\label{eq:gram-reflection}
 R_d=PJ_0P^{-1},\qquad
 J_0=\begin{pmatrix}1&2r_0\\0&-1\end{pmatrix}.
\end{equation}
The next section determines exactly when this reflection carries a
finite Cohn word column to another such column.

\section{Decoding a rational reflection on finite words}
\label{sec:decoder}

The matrix identities for a standard directive supply many pairs related
by its reflection: one may exchange the blocks $ab$ and $ba$ in any
concatenation. The question is whether there are other pairs. We prove
that there are none, even among arbitrary finite words. The proof reads
two letters at a time. A change of projective coordinate places all
possible first letters in two fixed intervals, and the final letter of
the directive determines which inequalities exclude a wrong continuation.
\Cref{subsec:centre-decoder} proves the same statement for the reflection of
every palindrome by a shorter argument in the slope coordinate. The cylinders
of this section are kept because they also measure how far a reflected column
can move, which is what \cref{sec:displaced} needs.

For a directive $d$, retain the notation
$$
 U=\mu(\phi_d(a)),\qquad V=\mu(\phi_d(b)),\qquad
 G=\mu(\operatorname{Pal}(d)),\qquad
 \nu_d(w)=\mu(\phi_d(w)).
$$
Thus $\nu_d$ evaluates a word by replacing $a,b$ with $U,V$.
Write $M_0=e^{\mathsf T}Ge$, $\rho_0=(Ge)_1$ and
$$
 K=2ee^{\mathsf T}G-M_0I,\qquad R=K/M_0.
$$
Since $(ee^{\mathsf T}G)^2=M_0ee^{\mathsf T}G$, one has
$K^2=M_0^2I$; in particular $R$ is an invertible involution.
For a word $x$ in the two blocks $ab,ba$, let $\widehat{x}$ exchange
these blocks in their original order. The parsing into blocks is
unambiguous because both blocks have length two. For example,
$\widehat{abba}=baab$. For a nonzero column $v$, the notation $[v]$
denotes its line in the real projective line.

\begin{theorem}[Finite-word reflection decoder]\label{thm:word-decoder}
For every finite directive $d$ and every pair of finite words $x,y$,
including the empty word,
\begin{equation}\label{eq:decoder-projective}
 [R\nu_d(x)e]=[\nu_d(y)e]
 \quad\Longleftrightarrow\quad
 x\in\{ab,ba\}^*,\qquad y=\widehat{x}.
\end{equation}
Whenever these conditions hold, the proportionality is exact:
\begin{equation}\label{eq:decoder-exact}
 R\nu_d(x)e=\nu_d(\widehat{x})e.
\end{equation}
\end{theorem}

There is no palindrome or primitive-count hypothesis in this theorem.
Those restrictions will enter only when the decoder is applied to the
midpoint problem.

\subsection{The two block identities}

The literal standard identities proved in the preceding section give
\begin{equation}\label{eq:decoder-standard}
 UV=ABG,\qquad VU=BAG,\qquad G=G^{\mathsf T}.
\end{equation}
Appending $a$ to $d$ replaces $(U,V)$ by $(U,VU)$, and appending
$b$ replaces it by $(UV,V)$. The direction of this update will matter
when we compare the two generators.

Direct multiplication gives $AB+BA=6ee^{\mathsf T}$. For every
symmetric matrix $G$ one also has
$\tr(ABG)=3e^{\mathsf T}Ge$. Consequently
\begin{equation}\label{eq:decoder-intertwiners}
 3K=UV+VU-\tr(UV)I,\qquad
 KUV=VUK,\qquad KVU=UVK,\qquad Ke=M_0e.
\end{equation}
Here the middle identities can be checked without a word argument:
Cayley--Hamilton makes
$$
 (UV+VU-\tr(UV)I)UV
 =VUUV-I
 =VU(UV+VU-\tr(UV)I).
$$
Interchanging $U,V$ proves the other identity. In particular, these
relations prove the forward construction in
\eqref{eq:decoder-exact} by induction on the number of blocks. The
rest of the section proves exhaustion.

\subsection{A common coordinate at the terminal slope}

Matrices act on slopes by
$$
 \begin{pmatrix}a&b\\c&d\end{pmatrix}(t)
 =\frac{at+b}{ct+d},
$$
with the usual interpretation at infinity. The terminal column $e$
has slope two. Set
\begin{equation}\label{eq:decoder-interval}
 \mathcal J=[8/5,29/12],\qquad z=\frac1{t-2}.
\end{equation}
Conjugation by $P=\left(\begin{smallmatrix}2&1\\1&0\end{smallmatrix}\right)$
is exactly this change of coordinate, and sends the terminal slope to
$z=\infty$. The usefulness of the coordinate is that every positive
word already lies in the same interval:
$$
 A(\mathcal J)=[21/13,70/41]\subset\operatorname{int}\mathcal J,\qquad
 B(\mathcal J)=[50/21,169/70]\subset\operatorname{int}\mathcal J.
$$
Both matrices have determinant one, so both maps are increasing on
$\mathcal J$ and these endpoint evaluations prove the inclusions. The
second is the tight one: its upper endpoint clears that of $\mathcal J$ by
$29/12-169/70=1/420$, so the inclusion must be read in exact arithmetic.
Since $2\in\mathcal J$, induction applies to every finite positive word.
The interval $\mathcal J$ is a rational neighbourhood of the interval between the
attracting fixed points of $A$ and $B$ on slopes, the golden ratio $(1+\sqrt5)/2$ and
$1+\sqrt2$; its endpoints $8/5$ and $29/12$ are a Fibonacci and a Pell quotient just
outside them. Moreover $A(t)<2<B(t)$ for every positive $t$, so the first letter of a
nonempty word decides on which side of the terminal slope its column lies. In the
coordinate $z$ the two halves $[8/5,2)$ and $(2,29/12]$ of $\mathcal J$ become the rays
$z\leqslant-5/2$ and $z\geqslant12/5$, and the terminal slope $2$ becomes $z=\infty$.

There are positive integers $p,q$ such that
\begin{equation}\label{eq:decoder-trace-column}
 \tr U=3p,\quad (P^{-1}UP)_{21}=-p,\qquad
 \tr V=3q,\quad (P^{-1}VP)_{21}=q.
\end{equation}
At the empty directive this holds with $(p,q)=(1,2)$.
For the induction step, multiplication with a general symmetric $G$
gives
$$
 \tr(ABG)=\tr(BAG)=3M_0,\qquad
 (P^{-1}ABGP)_{21}=-M_0,\qquad
 (P^{-1}BAGP)_{21}=M_0.
$$
Thus the two directive updates replace $q$ or $p$, respectively,
by the positive integer $M_0$.

Define the positive reciprocal parameters and two coordinate centers by
$$
 \alpha_0=p^{-2},\qquad\beta_0=q^{-2},\qquad
 u=\frac{(P^{-1}UP)_{11}}p,\qquad
 v=\frac{(P^{-1}VP)_{11}}q.
$$
Trace and determinant determine the remaining entries:
\begin{equation}\label{eq:decoder-cusp-matrices}
 \begin{split}
 P^{-1}UP&=
 p\begin{pmatrix}u&u^2-3u+\alpha_0\\-1&3-u\end{pmatrix},\\
 P^{-1}VP&=
 q\begin{pmatrix}v&3v-v^2-\beta_0\\1&3-v\end{pmatrix}.
 \end{split}
\end{equation}
Their projective maps in the new coordinate are therefore
\begin{equation}\label{eq:decoder-fg}
 f(z)=-u-\frac{\alpha_0}{z+u-3},\qquad
 g(z)=v-\frac{\beta_0}{z+3-v}.
\end{equation}

\begin{lemma}[Uniform bounds and the last directive letter]
\label{lem:decoder-dominance}
For every directive,
\begin{equation}\label{eq:decoder-uv}
 3\leqslant u\leqslant41/12,\qquad
 5/2\leqslant v\leqslant21/8.
\end{equation}
If $L=u+v-3$, then
\begin{equation}\label{eq:decoder-curve}
 M_0=pqL,\qquad
 \alpha_0+\beta_0=L(3-L),\qquad 5/2\leqslant L<3.
\end{equation}
In particular $p\geqslant1$ and $q\geqslant2$. For a nonempty
directive the following alternatives hold:
\begin{equation}\label{eq:decoder-dominance}
 \begin{array}{c@{\quad}c@{\quad}l}
 \text{last letter}&\text{trace dominance}&\text{reciprocal bounds}\\[2pt]
 a&q\geqslant5p&\beta_0\leqslant\alpha_0/25,\quad\alpha_0\leqslant1\\
 b&p\geqslant(5/2)q&\alpha_0\leqslant4\beta_0/25,\quad\beta_0\leqslant1/4.
 \end{array}
\end{equation}
\end{lemma}

\begin{proof}
The word $\phi_d(a)$ is either $a$ or begins with $ab$; the word
$\phi_d(b)$ is either $b$ or begins with $ba$. This follows by
induction from the two updates following \eqref{eq:decoder-standard}.
Consequently $U(2)=A(s)$, where $s=2$ if $U=A$ and otherwise $s$ is the slope of a column
of a word beginning with $b$, so that $s\in B(\mathcal J)\subset[2,29/12]$. Likewise
$V(2)=B(s')$ with $s'=2$ or $s'\in A(\mathcal J)\subset[8/5,2]$. Since $A$ and $B$ are
increasing,
$$
 5/3=A(2)\leqslant U(2)\leqslant A(29/12)=70/41,\qquad
 50/21=B(8/5)\leqslant V(2)\leqslant B(2)=12/5.
$$
On the other hand \eqref{eq:decoder-cusp-matrices} gives $U(2)=2-1/u$ and
$V(2)=2+1/v$: the column $e$ has cusp coordinate $\infty$, and $f(\infty)=-u$,
$g(\infty)=v$ by \eqref{eq:decoder-fg}. Solving, $12/41\leqslant1/u\leqslant1/3$ and
$8/21\leqslant1/v\leqslant2/5$, which is \eqref{eq:decoder-uv}.

The lower-left entry of $P^{-1}VUP$ is $pq(u+v-3)$, so it
equals $M_0=pqL$. The matrix
$P^{-1}(UV+VU)P=6P^{-1}ee^{\mathsf T}GP$ has zero second row.
Its lower-right entry, expanded with
\eqref{eq:decoder-cusp-matrices}, gives
$\alpha_0+\beta_0=L(3-L)$. The bounds on $u,v$ give $L\geqslant5/2$;
positivity of $\alpha_0+\beta_0$ then gives $L<3$.

This argument has not required lower bounds on $p,q$ beyond their
positive integrality. It now shows that either update to $pqL$
strictly increases the replaced coordinate. Starting at $(1,2)$,
we obtain $p\geqslant1,q\geqslant2$ at every stage.
For an appended $a$, the new $q$ is
$L_{\rm old}p q_{\rm old}\geqslant5p$.
For an appended $b$, the new $p$ is
$L_{\rm old}p_{\rm old}q\geqslant(5/2)q$.
Taking reciprocal squares gives \eqref{eq:decoder-dominance}.
\end{proof}

\subsection{First-letter intervals and the reflection}

Let $\mathcal C$ be the set of cusp slopes of all $\nu_d(w)e$.
The empty word contributes infinity, and
\begin{equation}\label{eq:decoder-orbit}
 \mathcal C=\{\infty\}\cup f(\mathcal C)\cup g(\mathcal C)
 \subset\mathcal H
 :=\{z\leqslant-5/2\}\cup\{z\geqslant12/5\}\cup\{\infty\}.
\end{equation}
The inclusion follows by transforming $\mathcal J$ under
\eqref{eq:decoder-interval}, as explained above. Each of the maps \eqref{eq:decoder-fg} has
a single pole: $f$ at $z=3-u$ and $g$ at $z=v-3$. By \eqref{eq:decoder-uv} the first lies
in $[-5/12,0]$ and the second in $[-1/2,-3/8]$, so both poles sit in the gap between the
two rays of $\mathcal H$, and their distances to the rays are bounded below as in
\cref{tab:poles}. On the positive ray the denominators $z+u-3$ and $z+3-v$ are these
distances, and on the negative ray they are their negatives; in particular no denominator
vanishes on $\mathcal H$.

\begin{table}[htbp]
\centering
\caption{The poles of the two first-letter maps \eqref{eq:decoder-fg} and lower bounds for
their distances to the two rays of $\mathcal H$, from \eqref{eq:decoder-uv}. Dividing the
numerators $\alpha_0$ and $\beta_0$ by these distances gives the distances from the centres
$-u$ and $v$ to the four endpoints of the cylinders \eqref{eq:decoder-cylinders}.}
\label{tab:poles}
\small
\renewcommand{\arraystretch}{1.15}
\begin{tabular}{@{}lccc@{}}
\toprule
map & pole & distance to $z\geqslant12/5$ & distance to $z\leqslant-5/2$\\
\midrule
$f$ (letter $a$) & $3-u\in[-5/12,\,0]$ & $\geqslant12/5-0=12/5$ & $\geqslant-5/12+5/2=25/12$\\
$g$ (letter $b$) & $v-3\in[-1/2,\,-3/8]$ & $\geqslant12/5+3/8=111/40$ & $\geqslant-1/2+5/2=2$\\
\bottomrule
\end{tabular}
\end{table}

A map $z\mapsto c-\beta/(z-z_0)$ with $\beta>0$ sends every point at distance at least
$\delta$ from its pole $z_0$ to within $\beta/\delta$ of $c$, below $c$ to the right of the
pole and above $c$ to the left of it, and it sends $\infty$ to $c$. For $f$ the centre is
$c=-u$, the numerator $\alpha_0$ and the pole $3-u$; for $g$ they are $v$, $\beta_0$ and
$v-3$. Applied to the four entries of \cref{tab:poles}, this gives
\begin{equation}\label{eq:decoder-cylinders}
 \begin{split}
 f(\mathcal H)&\subset
 \mathcal U:=[-u-5\alpha_0/12,-u+12\alpha_0/25],\\
 g(\mathcal H)&\subset
 \mathcal V:=[v-40\beta_0/111,v+\beta_0/2].
 \end{split}
\end{equation}
For example, the lower end of $\mathcal V$ is $v-\beta_0\cdot40/111$: the point of the
positive ray closest to the pole of $g$ is $12/5$, at distance $12/5-(v-3)\geqslant111/40$.
\Cref{fig:cusp-cylinders} shows the configuration at the empty directive.

\begin{figure}[htbp]
\centering
% Author: Dr. Denys Dutykh (Khalifa University of Science and Technology, Abu Dhabi, UAE)
% Generated by supplement/code/make_figures.py from exact data; do not edit by hand.

\begin{tikzpicture}[font=\footnotesize, >={Stealth[length=4pt]}]
\draw[mist, densely dotted, line width=0.5pt] (2.871,-0.56) -- (2.871,-1.54);
\draw[mist, densely dotted, line width=0.5pt] (12.573,-0.56) -- (12.573,-1.54);
\draw[mist, line width=0.8pt] (2.871,0) -- (12.573,0);
\draw[navy, line width=2.2pt] (2.871,0) -- (0.15,0);
\draw[navy, line width=2.2pt] (12.573,0) -- (13.71,0);
\draw[navy, ->, line width=0.8pt] (0.2,0) -- (-0.05,0);
\draw[navy, ->, line width=0.8pt] (13.66,0) -- (13.91,0);
\node[ink, below] at (0,-0.08) {$\infty$};
\node[ink, below] at (13.86,-0.08) {$\infty$};
\node[ink, below] at (2.871,-0.08) {$-5/2$};
\node[ink, below] at (12.573,-0.08) {$12/5$};
\node[ink, below] at (7.821,-0.08) {$0$};
\draw[ink, line width=0.5pt] (7.821,-0.07) -- (7.821,0.07);
\filldraw[fill=sand, draw=ink, line width=0.4pt] (1.056,0.16) rectangle (2.831,0.58);
\node[ink, above] at (1.944,0.6) {$\mathcal U$};
\filldraw[fill=sand, draw=ink, line width=0.4pt] (12.593,0.16) rectangle (13.018,0.58);
\node[ink, above] at (12.806,0.6) {$\mathcal V$};
\draw[navy, line width=0.25pt] (1.061,0.2) -- (1.061,0.54);
\draw[navy, line width=0.25pt] (1.062,0.2) -- (1.062,0.54);
\draw[navy, line width=0.25pt] (1.063,0.2) -- (1.063,0.54);
\draw[navy, line width=0.25pt] (1.089,0.2) -- (1.089,0.54);
\draw[navy, line width=0.25pt] (1.115,0.2) -- (1.115,0.54);
\draw[navy, line width=0.25pt] (1.116,0.2) -- (1.116,0.54);
\draw[navy, line width=0.25pt] (1.119,0.2) -- (1.119,0.54);
\draw[navy, line width=0.25pt] (1.123,0.2) -- (1.123,0.54);
\draw[navy, line width=0.25pt] (1.124,0.2) -- (1.124,0.54);
\draw[navy, line width=0.25pt] (1.125,0.2) -- (1.125,0.54);
\draw[navy, line width=0.25pt] (1.881,0.2) -- (1.881,0.54);
\draw[navy, line width=0.25pt] (2.461,0.2) -- (2.461,0.54);
\draw[navy, line width=0.25pt] (2.463,0.2) -- (2.463,0.54);
\draw[navy, line width=0.25pt] (2.466,0.2) -- (2.466,0.54);
\draw[navy, line width=0.25pt] (2.541,0.2) -- (2.541,0.54);
\draw[navy, line width=0.25pt] (2.612,0.2) -- (2.612,0.54);
\draw[navy, line width=0.25pt] (2.613,0.2) -- (2.613,0.54);
\draw[navy, line width=0.25pt] (2.624,0.2) -- (2.624,0.54);
\draw[navy, line width=0.25pt] (2.634,0.2) -- (2.634,0.54);
\draw[navy, line width=0.25pt] (2.635,0.2) -- (2.635,0.54);
\draw[navy, line width=0.25pt] (2.637,0.2) -- (2.637,0.54);
\draw[navy, line width=0.25pt] (12.601,0.2) -- (12.601,0.54);
\draw[navy, line width=0.25pt] (12.606,0.2) -- (12.606,0.54);
\draw[navy, line width=0.25pt] (12.611,0.2) -- (12.611,0.54);
\draw[navy, line width=0.25pt] (12.612,0.2) -- (12.612,0.54);
\draw[navy, line width=0.25pt] (12.771,0.2) -- (12.771,0.54);
\draw[navy, line width=0.25pt] (12.941,0.2) -- (12.941,0.54);
\draw[navy, line width=0.25pt] (12.942,0.2) -- (12.942,0.54);
\draw[navy, line width=0.25pt] (12.943,0.2) -- (12.943,0.54);
\draw[navy, line width=0.25pt] (12.969,0.2) -- (12.969,0.54);
\draw[navy, line width=0.25pt] (12.995,0.2) -- (12.995,0.54);
\draw[navy, line width=0.25pt] (12.996,0.2) -- (12.996,0.54);
\draw[navy, line width=0.25pt] (12.999,0.2) -- (12.999,0.54);
\draw[navy, line width=0.25pt] (13.003,0.2) -- (13.003,0.54);
\draw[navy, line width=0.25pt] (13.004,0.2) -- (13.004,0.54);
\draw[navy, line width=0.25pt] (13.005,0.2) -- (13.005,0.54);
\draw[wine, <->, line width=0.6pt] (2.871,-0.78) -- (6.996,-0.78);
\draw[wine, <->, line width=0.6pt] (7.821,-0.78) -- (12.573,-0.78);
\draw[wine, line width=3pt] (6.996,-0.78) -- (7.821,-0.78);
\node[wine, above] at (4.934,-0.78) {$25/12$};
\node[wine, above] at (10.197,-0.78) {$12/5$};
\node[wine, below] at (7.409,-0.8200000000000001) {pole of $f$};
\draw[teal, <->, line width=0.6pt] (2.871,-1.42) -- (6.831,-1.42);
\draw[teal, <->, line width=0.6pt] (7.079,-1.42) -- (12.573,-1.42);
\draw[teal, line width=3pt] (6.831,-1.42) -- (7.079,-1.42);
\node[teal, above] at (4.851,-1.42) {$2$};
\node[teal, above] at (9.826,-1.42) {$111/40$};
\node[teal, below] at (6.955,-1.46) {pole of $g$};
\end{tikzpicture}

\caption{The cusp coordinate $z$ at the empty directive, where $f(z)=-3-1/z$ and
$g(z)=5/2-1/(4z+2)$. Every column of a nonempty word lies on one of the two rays of
$\mathcal H$ (thick), and the terminal column is the point at infinity. The poles of $f$ and
$g$ lie in the gap between the rays; the bars below the axis show the intervals $[-5/12,0]$ and
$[-1/2,-3/8]$ that contain them for every directive, and the arrows show the four distances of
\cref{tab:poles}, measured from the ends of these ranges. The shaded boxes are the cylinders
$\mathcal U$ and $\mathcal V$ of \eqref{eq:decoder-cylinders}, and the ticks above the axis are
the columns of all words of length at most nine.}
\label{fig:cusp-cylinders}
\end{figure}

These two finite intervals lie strictly in the negative and positive components of
$\mathcal H$:
$$
 -u+12\alpha_0/25\leqslant-3+12/25<-5/2,\qquad
 v-40\beta_0/111\geqslant5/2-10/111>12/5.
$$
Consequently a finite slope in $\mathcal C$ determines the first
letter, and inversion of its projective generator determines the
remainder. The terminal value cannot belong to either finite
interval. Induction therefore proves injectivity of finite-word
evaluation at the terminal slope.

Put $\tau=2\rho_0/M_0$. Since both entries of $Ge$ are positive,
$M_0=2\rho_0+(Ge)_2$ gives $0<\tau<1$. Direct conjugation and
\eqref{eq:decoder-intertwiners} give
\begin{equation}\label{eq:decoder-k}
 P^{-1}RP=\begin{pmatrix}1&\tau\\0&-1\end{pmatrix},\qquad
 k(z)=-z-\tau,\qquad
 \tau=u-v+\delta,\quad \delta=\frac{\alpha_0-\beta_0}{L}.
\end{equation}
For the last equality, expand the upper-right entry of
$3P^{-1}KP$ using \eqref{eq:decoder-cusp-matrices}.
The map $k$ fixes infinity and sends the two finite components
of $\mathcal H$ to opposite signs. It need not preserve
$\mathcal H$; whenever both endpoints belong to $\mathcal C$,
their first letters must nevertheless be opposite.

We will also use positivity after an intermediate reflection, even
when that reflected point has not yet been shown to be in the orbit.
In the original coordinate,
\begin{equation}\label{eq:decoder-positive-reflection}
 R(t)=2-\frac{t-2}{1+\tau(t-2)}>0\qquad(t\in\mathcal J).
\end{equation}
Indeed the denominator exceeds $3/5$. If $t<2$, the value exceeds
two; if $t\geqslant2$, it is at least $2-(t-2)\geqslant19/12$.
As $U$ begins with $A$ and $V$ with $B$, they send every positive
input below and above two, respectively. Hence $fk$ has negative
cusp slope on $\mathcal C$, while $gk$ has positive cusp slope.
These statements include the terminal input.

\subsection{Excluding a wrong second letter}

After opposite first letters have been removed, the two residual maps
are
\begin{equation}\label{eq:decoder-residuals}
 \begin{split}
 \mathcal E(z):=g^{-1}kf(z)
 &=\frac{\beta_0}{\delta-\alpha_0/(z+u-3)}+v-3,\\
 \mathcal D(z):=f^{-1}kg(z)
 &=\frac{\alpha_0}{\delta-\beta_0/(z+3-v)}+3-u.
 \end{split}
\end{equation}
These formulas follow by solving \eqref{eq:decoder-fg} for its
input. The identities \eqref{eq:decoder-intertwiners} imply
\begin{equation}\label{eq:decoder-residual-return}
 \mathcal E g=fk,\qquad \mathcal D f=gk.
\end{equation}
It remains to exclude an empty remainder or the wrong next letter.
We prove precisely
\begin{equation}\label{eq:decoder-wrong}
 \mathcal E(\{\infty\}\cup\mathcal U)\cap\mathcal C=\varnothing,
 \qquad
 \mathcal D(\{\infty\}\cup\mathcal V)\cap\mathcal C=\varnothing.
\end{equation}
Only these domains are needed. In particular, no estimate across a
pole of either residual map will be used.

\smallskip\noindent\emph{A directive ending in $a$.}
Here $\beta_0\leqslant\alpha_0/25$ and $\alpha_0\leqslant1$.
Since $L<3$,
$$
 \delta>\frac{8\alpha_0}{25}>0,\qquad
 0<\frac{\beta_0}{\delta}<\frac18.
$$
It follows from $-1/2\leqslant v-3\leqslant-3/8$ that
$\mathcal E(\infty)\in(-1/2,-1/4)$.
For $z\in\mathcal U$, the quantity $z+u-3$ is negative,
so the denominator of $\mathcal E$ is larger than $\delta$.
Therefore
\begin{equation}\label{eq:decoder-a-E}
 \mathcal E(\{\infty\}\cup\mathcal U)\subset(-1/2,-1/4),
\end{equation}
which is disjoint from $\mathcal H$.

For the other residual, direct simplification gives
$$
 \mathcal D(\infty)-v=\frac{L\beta_0}{\alpha_0-\beta_0}
 >\frac52\beta_0>\frac{\beta_0}{2}.
$$
If $z\in\mathcal V$, put $s=z+3-v$. The cylinder bound gives
$s\geqslant3-40\beta_0/111>2$. Its residual denominator $h_0$ satisfies
$$
 h_0:=\delta-\frac{\beta_0}{s}
 >\left(\frac8{25}-\frac1{50}\right)\alpha_0
 =\frac{3\alpha_0}{10}>0,\qquad h_0<\delta<\frac{\alpha_0}{L}.
$$
Using $L\delta=\alpha_0-\beta_0$ now yields
\begin{equation}\label{eq:decoder-a-D}
 \mathcal D(z)-v
 =\frac{\beta_0(1+L/s)}{h_0}
 >\frac{L\beta_0}{\alpha_0}\geqslant\frac52\beta_0>\frac{\beta_0}{2}.
\end{equation}
Both the terminal value and these finite values are positive and
strictly above the entire positive cylinder $\mathcal V$. They
cannot be in $\mathcal C$.

\smallskip\noindent\emph{A directive ending in $b$.}
Now $\alpha_0\leqslant4\beta_0/25$ and $\beta_0\leqslant1/4$. Set
$\eta=-\delta=(\beta_0-\alpha_0)/L$. Then
$$
 \eta>\frac{7\beta_0}{25}>0,\qquad
 0<\frac{\alpha_0}{\eta}<\frac47.
$$
At the terminal input,
$$
 \mathcal E(\infty)+u=-\frac{L\alpha_0}{\beta_0-\alpha_0}
 <-10\alpha_0.
$$
For $z\in\mathcal U$, put $s=3-z-u$.
Since $\alpha_0\leqslant1/25$, the cylinder bound gives
$s\geqslant3-12\alpha_0/25>2$.
The positive denominator
$$
 h_0:=\eta-\frac{\alpha_0}{s}
 >\left(\frac7{25}-\frac2{25}\right)\beta_0
 =\frac{\beta_0}{5},\qquad h_0<\eta<\frac{\beta_0}{L},
$$
allows the exact simplification
\begin{equation}\label{eq:decoder-b-E}
 \mathcal E(z)+u
 =-\frac{\alpha_0(1+L/s)}{h_0}
 <-\frac{L\alpha_0}{\beta_0}\leqslant-10\alpha_0.
\end{equation}
These values and the terminal value are negative and strictly below
$-u-5\alpha_0/12$, the lower endpoint of $\mathcal U$.
Thus they are outside $\mathcal C$.

Finally $\mathcal D(\infty)=-\alpha_0/\eta+3-u$.
For $z\in\mathcal V$, the quantity $s=z+3-v$ is positive and
$$
 \mathcal D(z)=-\frac{\alpha_0}{\eta+\beta_0/s}+3-u.
$$
The bounds $0\leqslant u-3\leqslant5/12$ and
$0<\alpha_0/\eta<4/7$ give
\begin{equation}\label{eq:decoder-b-D}
 \mathcal D(\{\infty\}\cup\mathcal V)
 \subset(-83/84,0),\qquad \frac5{12}+\frac47=\frac{83}{84}.
\end{equation}
This interval is disjoint from $\mathcal H$. Equations
\eqref{eq:decoder-a-E}--\eqref{eq:decoder-b-D} prove all of
\eqref{eq:decoder-wrong} for every nonempty directive.

\subsection{The empty directive and finite termination}

There is no last letter for the empty directive, so we check its
residual exclusions directly in the original coordinate. Here
$U=A$, $V=B$, $M_0=5$, and
$$
 K=\begin{pmatrix}3&4\\4&-3\end{pmatrix},\qquad
 E_0=B^{-1}KA=\begin{pmatrix}0&5\\5&-9\end{pmatrix},\qquad
 D_0=A^{-1}KB=\begin{pmatrix}9&5\\5&0\end{pmatrix}.
$$
The use of $K$ instead of $R$ does not change a projective map.
The terminal input is now two. We have $E_0(2)=5\notin \mathcal J$ and
$$
 E_0A(t)=\frac{5(t+1)}{t-4}<0\qquad(t\in\mathcal J).
$$
Thus an empty source remainder or a second $a$ is excluded.
For the other residual,
$$
 D_0(2)=23/10\in(2,50/21),\qquad
 D_0B(t)=\frac{55t+23}{25t+10}.
$$
The latter map is decreasing on $\mathcal J$ and satisfies
$$
 2<D_0B(t)\leqslant111/50<50/21.
$$
All these values lie between the terminal slope and the least
possible slope of a word starting with $B$. They lie outside the
encoded orbit. This proves exactly the empty-directive version
of \eqref{eq:decoder-wrong}, without a dominance assumption.

\begin{proof}[Proof of \cref{thm:word-decoder}]
Represent \eqref{eq:decoder-projective} in cusp coordinates.
If one word is empty, its slope is infinity; as $k(\infty)=\infty$
and no nonempty word has that slope, both words are empty.
Otherwise their first letters are opposite.

Suppose $x=ax_1$ and $y=by_1$. Removing the target's first
letter gives the residual relation
$y_1(\infty)=\mathcal E(x_1(\infty))$, where words here denote
their cusp maps. The first exclusion in
\eqref{eq:decoder-wrong} says that $x_1$ is neither empty nor
starts with $a$. Hence $x_1=bx_2$, and
\eqref{eq:decoder-residual-return} gives
$$
 y_1(\infty)=fk(x_2(\infty)).
$$
By \eqref{eq:decoder-positive-reflection} this is a finite negative
cusp slope. Therefore $y_1=ay_2$. Inverting $f$ restores
exactly the original relation
$y_2(\infty)=k(x_2(\infty))$.
Thus the first source block is $ab$ and the first target block
is $ba$.

If $x$ begins with $b$, the second exclusion in
\eqref{eq:decoder-wrong} and the identity $\mathcal Df=gk$
give the exchanged conclusion: the source begins with $ba$,
the target with $ab$, and removing those two letters restores
the same reflection relation. The empty-directive calculation
supplies the identical step in its original coordinate.

Each step removes two letters from each word. Induction on their
total length terminates, and the terminal case forces the two
words to end simultaneously. They are therefore precisely
the block-exchanged pair in \eqref{eq:decoder-projective}.
Conversely, repeated use of \eqref{eq:decoder-intertwiners} gives
$K\nu_d(x)=\nu_d(\widehat{x})K$ for every block word.
Applying both sides to $e$ proves \eqref{eq:decoder-exact} and
the converse implication.
\end{proof}

The exhaustion rests on the four strict residual exclusions and on the
termination of the two-letter descent, both established above. The
displayed identities and rational constants are also confirmed by exact
computation in the supplementary material; no enumeration of directives
or of finite words enters the proof.

\section{Complete classification at the ancestor midpoint}
\label{sec:midpoint}

The decoder concerns columns. The original arithmetic question concerns
Gram matrices with the same first diagonal entry. The bridge between
them is a full matrix conjugacy, which we prove before invoking the word
theorem. This avoids any choice of a square root of a Gram matrix or of
a middle letter in an odd palindrome.

\begin{lemma}[The Gram bridge]\label{lem:gram-bridge}
Fix a directive $d$ and two palindromes $z_X,z_Y$. Write
$M_i=M(F_d(z_i))$ and $\rho_i=\rho(F_d(z_i))$.
The following conditions are equivalent:
\begin{align}
 M_X=M_Y=M,\qquad \rho_X+\rho_Y&=2r_0M;
 \label{eq:ancestor-midpoint-condition}\\
 \nu_d(z_Y)&=R_d\nu_d(z_X)R_d.
 \label{eq:full-reflection-conjugacy}
\end{align}
Under these conditions,
$\nu_d(z_Y)e=R_d\nu_d(z_X)e$.
\end{lemma}

\begin{proof}
For a symmetric determinant-one matrix
$\left(\begin{smallmatrix}M&\rho\\\rho&H\end{smallmatrix}\right)$,
conjugation by $J_0$ as a bilinear form gives
$$
 J_0^{\mathsf T}
 \begin{pmatrix}M&\rho\\\rho&H\end{pmatrix}J_0
 =\begin{pmatrix}
 M&2r_0M-\rho\\
 2r_0M-\rho&4r_0^2M-4r_0\rho+H
 \end{pmatrix}.
$$
The lower-right entry is uniquely determined by the other entries and
the determinant, because $M>0$. Hence
\eqref{eq:ancestor-midpoint-condition} is equivalent to
$\mc G(F_d(z_Y))=J_0^{\mathsf T}\mc G(F_d(z_X))J_0$.
Substitute \eqref{eq:image-gram}, multiply by $P^{-1}$ on both sides,
and use $R_d=PJ_0P^{-1}$ to obtain
$$
 G\nu_d(z_Y)=R_d^{\mathsf T}G\nu_d(z_X)R_d
 =GR_d\nu_d(z_X)R_d.
$$
Cancel the invertible matrix $G$. All steps are reversible, proving
the equivalence. Finally $R_de=e$ gives the column identity.
\end{proof}

\begin{theorem}[Palindromic pairs at the ancestor midpoint]
\label{thm:palindromic-classification}
Fix any finite directive $d$. Two binary palindromes $z_X,z_Y$
satisfy \eqref{eq:ancestor-midpoint-condition} if and only if, for
some word $w\in\{ab,ba\}^*$,
\begin{equation}\label{eq:classified-palindromes}
 z_X=\widetilde w w,
 \qquad z_Y=\widetilde{\widehat w}\,\widehat w.
\end{equation}
The word $w$ is uniquely determined by the ordered pair.
If it contains $k$ blocks, then
\begin{equation}\label{eq:classified-counts}
 c(F_d(z_X))=c(F_d(z_Y))
 =(2k+1)L_d\binom11.
\end{equation}
The case $k=0$ consists of two copies of the ancestor $g_d$.
Every case $k>0$ consists of two distinct, noncentral palindromes.
\end{theorem}

\begin{proof}
The Gram bridge and \cref{thm:word-decoder}, applied to the
whole words $z_X,z_Y$, show that
$z_X\in\{ab,ba\}^*$ and $z_Y=\widehat{z_X}$.
Because $z_X$ has even length and is a palindrome, each of its two
letter counts is even. A word of $m$ blocks $ab,ba$ has $m$ copies
of each letter. Thus $m=2k$ is even, and the midpoint of $z_X$
falls between blocks. Its second half is a unique word
$w\in\{ab,ba\}^*$ and $z_X=\widetilde w w$.
Block exchange commutes with reversal: reversing a block interchanges
$ab$ and $ba$, and reversing a sequence reverses the order of its
blocks. Consequently
$z_Y=\widetilde{\widehat w}\,\widehat w$.

Conversely, the block intertwining identities in
\eqref{eq:decoder-intertwiners} imply
$R_d\nu_d(x)R_d=\nu_d(\widehat x)$ for every
$x\in\{ab,ba\}^*$. Apply this with $x=\widetilde w w$ and then
use the Gram bridge. This proves the converse at the matrix level,
including the equality of labels.

Both preimages in \eqref{eq:classified-palindromes} have $2k$ copies
of each letter. Formula~\eqref{eq:completed-count-transport} proves
\eqref{eq:classified-counts}. If $k>0$, these completed counts have
greatest common divisor $2k+1>1$, whereas every central word has
primitive completed counts. The words are distinct because block
exchange changes their first letter, and $F_d$ is injective. The
empty case follows from $F_d(\eps)=g_d$.
\end{proof}

\begin{corollary}[Exclusion for central occurrences]
\label{cor:no-central-midpoint}
Let $d,t_X,t_Y$ be finite directives. If the central words
$\operatorname{Pal}(dt_X)$ and $\operatorname{Pal}(dt_Y)$ have the
same Gram label $M$ and their roots sum to $2M\rho_0/M_0$, then
$t_X=t_Y=\eps$. In particular, a pair of distinct central words
with equal labels cannot have its root midpoint at a common
standard ancestor's root parameter.
\end{corollary}

\begin{proof}
Use $z_i=\operatorname{Pal}(t_i)$ in the theorem. The images are
central, so only $k=0$ is possible. The palindromization recurrence
strictly increases length at every nonempty directive step; hence
$\operatorname{Pal}(t_i)=\eps$ implies $t_i=\eps$.
\end{proof}

\subsection{Exact examples and the role of primitivity}

At the empty directive, $G=I$, $(M_0,\rho_0)=(5,2)$ and
$$
 R_\eps=\frac15\begin{pmatrix}3&4\\4&-3\end{pmatrix}.
$$
Take $w=ab$. Then $z_X=baab$ and $z_Y=abba$, and direct evaluation gives
\begin{equation}\label{eq:small-gram-pair}
 \mc G(baab)=\begin{pmatrix}1130&467\\467&193\end{pmatrix},
 \qquad
 \mc G(abba)=\begin{pmatrix}1130&437\\437&169\end{pmatrix}.
\end{equation}
Both determinants are one and
$467+437=2\cdot1130\cdot(2/5)$.
This is a real equality of integral Gram labels. Its completed
counts are $(3,3)$, so the theorem explains exactly why it is
irrelevant to a collision between distinct central occurrences.

The same preimages give a pair for every directive, and the whole table
is one formula. It holds around every palindrome, standard or not.

\begin{table}[tb]
\centering
\caption{Images of the fixed palindromes $baab$ and $abba$. Here
$(M_0,\rho_0)$ belongs to the standard ancestor $\operatorname{Pal}(d)$,
$M$ is the common image label, and $\rho_-<\rho_+$ are the image
roots. Every row satisfies $(\rho_-+\rho_+)/(2M)=\rho_0/M_0$.
The completed counts have greatest common divisor three in every row.}
\label{tab:pairs}
\small
\setlength{\tabcolsep}{9pt}
\renewcommand{\arraystretch}{1.12}
\begin{tabular}{@{}crrrrr@{}}
\toprule
$d$ & $M_0$ & $\rho_0$ & $M$ & $\rho_-$ & $\rho_+$ \\
\midrule
$\eps$ & 5 & 2 & 1130 & 437 & 467 \\
$a$ & 13 & 5 & $19\,786$ & 7571 & 7649 \\
$b$ & 29 & 12 & $219\,530$ & $90\,753$ & $90\,927$ \\
$ab$ & 194 & 75 & $65\,712\,650$ & $25\,403\,793$ & $25\,404\,957$ \\
$ba$ & 433 & 179 & $730\,645\,066$ & $302\,043\,659$ & $302\,046\,257$ \\
\bottomrule
\end{tabular}
\end{table}

\begin{proposition}[The one-block pair in closed form]
\label{prop:one-block-closed-form}
Let $p$ be a palindrome with Gram label $M_p$ and Gram root $\rho_p$. The words
$p\,ba\,p\,ab\,p$ and $p\,ab\,p\,ba\,p$ have the common Gram label and the two Gram roots
\begin{equation}\label{eq:one-block-closed-form}
 M=M_p\bigl(9M_p^2+1\bigr),
 \qquad
 \rho_\pm=\bigl(9M_p^2+1\bigr)\rho_p\pm3M_p ,
\end{equation}
the sign being $+$ for $p\,ba\,p\,ab\,p$ and $-$ for $p\,ab\,p\,ba\,p$. In particular the two
roots differ by exactly $6M_p$, and their mean is $M\rho_p/M_p$. For $p=\operatorname{Pal}(d)$
the two words are $F_d(baab)$ and $F_d(abba)$, and $(M_p,\rho_p)=(M_0,\rho_0)$.
\end{proposition}

\begin{proof}
Put $W=\mu(p)$ and $\mc G_p=PWP=\left(\begin{smallmatrix}m&r\\r&h\end{smallmatrix}\right)$, a
symmetric matrix with $mh-r^2=1$, where $m=M_p$ and $r=\rho_p$. Then
$\mu(p\,ba\,p\,ab\,p)=W\,BA\,W\,AB\,W$, and inserting $P^{-1}P$ between the factors gives
$$
 \mc G(p\,ba\,p\,ab\,p)=\mc G_p\,K\,\mc G_p\,K^{\mathsf T}\mc G_p,
 \qquad
 K=P^{-1}BAP^{-1}=\begin{pmatrix}3&-1\\1&0\end{pmatrix},
$$
where $K^{\mathsf T}=P^{-1}ABP^{-1}$ because $A$, $B$ and $P$ are symmetric. A direct multiplication gives
$$
 K\mc G_pK^{\mathsf T}=\begin{pmatrix}9m-6r+h&3m-r\\3m-r&m\end{pmatrix},
 \qquad
 \mc G_pK\mc G_pK^{\mathsf T}=\begin{pmatrix}9m^2-3mr+1&3m^2\\9mr-3r^2+3&3mr+1\end{pmatrix},
$$
where $mh=r^2+1$ was used in the second product. Multiplying by the first column
$(m,r)^{\mathsf T}$ of $\mc G_p$, the first column of the Gram matrix is
$$
 \binom{(9m^2-3mr+1)m+3m^2r}{(9mr-3r^2+3)m+(3mr+1)r}
 =\binom{m(9m^2+1)}{(9m^2+1)r+3m}.
$$
The second word is obtained by exchanging the blocks, so its Gram matrix is
$$
 \mc G(p\,ab\,p\,ba\,p)=\mc G_p\,K^{\mathsf T}\mc G_p\,K\,\mc G_p ,
$$
and in the same way
$$
 K^{\mathsf T}\mc G_pK=\begin{pmatrix}9m+6r+h&-3m-r\\-3m-r&m\end{pmatrix},
 \qquad
 \mc G_pK^{\mathsf T}\mc G_pK=\begin{pmatrix}9m^2+3mr+1&-3m^2\\9mr+3r^2-3&1-3mr\end{pmatrix},
$$
and the first column of its Gram matrix is
$$
 \binom{(9m^2+3mr+1)m-3m^2r}{(9mr+3r^2-3)m+(1-3mr)r}
 =\binom{m(9m^2+1)}{(9m^2+1)r-3m}.
$$
The mean of the two roots is $(9m^2+1)r=Mr/m$. For $p=\operatorname{Pal}(d)$, \eqref{eq:literal-standard-identities}
gives $\phi_d(ba)=ba\,p$ and $\phi_d(ab)=ab\,p$, so that
$F_d(baab)=p\,\phi_d(ba)\phi_d(ab)=p\,ba\,p\,ab\,p$, and likewise for $abba$.
\end{proof}

Every row of \cref{tab:pairs} is \eqref{eq:one-block-closed-form}
evaluated at the ancestor of that row. The entries were also obtained
directly from the literal words, and the determinant, midpoint and count
identities are confirmed by exact integer arithmetic in the
supplementary material; the theorem, including its exhaustion statement,
depends on neither.

The theorem uses neither an odd-label hypothesis nor a factorization
of the common label. It also includes all three possible middle types
of a palindrome: empty, $a$, and $b$. Their apparent difference
disappears because the full-word decoder forces every reflected
preimage in the midpoint case to have the even form
\eqref{eq:classified-palindromes}. This is the main structural gain
of passing from square-root factors to whole Gram matrices.

\subsection{Every palindrome is a centre}
\label{subsec:any-centre}

\Cref{prop:one-block-closed-form} uses nothing about the centre beyond the symmetry and the
determinant of its Gram matrix, and the same is true of the construction half of
\cref{thm:palindromic-classification}. Fix a palindrome $p$ and put
$$
 G=\mu(p),\quad M_p=e^{\mathsf T}Ge,\quad \rho_p=(Ge)_1,\quad
 K_p=2ee^{\mathsf T}G-M_pI,\quad R_p=K_p/M_p .
$$
For blocks $W_1,\dots,W_n\in\{ab,ba\}$ write
$$
 x_p=p\,W_1\,p\,W_2\,p\cdots p\,W_n\,p ,
$$
and let $\widehat x_p$ be the word obtained by exchanging every block. The word $x_p$ is a
palindrome exactly when $n=2k$ is even and $W_{2k+1-i}$ is the exchange of $W_i$ for every
$i$, since a middle block would have to equal its own exchange. We then say that $x_p$ is
\emph{structured around} $p$, and $\widehat x_p$ is structured around $p$ as well.

\begin{proposition}[Block exchange around an arbitrary palindrome]
\label{prop:any-centre}
For every palindrome $p$ and all blocks $W_1,\dots,W_n\in\{ab,ba\}$,
$$
 \begin{gathered}
 K_p\,\mu(W_1p\cdots W_np)=\mu(\widehat W_1p\cdots\widehat W_np)\,K_p ,\\
 R_p\,\mu(W_1p\cdots W_np)\,e=\mu(\widehat W_1p\cdots\widehat W_np)\,e .
 \end{gathered}
$$
Moreover $e^{\mathsf T}\mu(x_p)e=e^{\mathsf T}\mu(\widehat x_p)e$; for palindromic block words
these are the Gram labels. If $x_p$ is structured around $p$ with $2k$ blocks and common label $M$, then
$\mu(\widehat x_p)=R_p^{\mathsf T}\mu(x_p)R_p$, the roots satisfy
$\rho(x_p)+\rho(\widehat x_p)=2M\rho_p/M_p$, and both words have completed counts
$(2k+1)\,c(p)$. For $k\geqslant1$ the two words are distinct and neither is central.
\end{proposition}

\begin{proof}
The derivation of \eqref{eq:decoder-intertwiners} uses only that $G$ is symmetric, through
$\tr(ABG)=3e^{\mathsf T}Ge$, and that the products have determinant one, through the
Cayley--Hamilton identity. With $G=\mu(p)$ it gives $K_p\,ABG=BAG\,K_p$ and
$K_p\,BAG=ABG\,K_p$. Hence $K_p\Pi=\widehat\Pi K_p$ for $\Pi=\mu(W_1p)\cdots\mu(W_np)$ and its
exchange $\widehat\Pi$, which is the displayed identity; applying it to $e$ and using
$K_pe=M_pe$ gives the column identity. Moreover $K_pe=2e(e^{\mathsf T}Ge)-M_pe=M_pe$ and
$e^{\mathsf T}GK_p=M_pe^{\mathsf T}G$. Since $\mu(x_p)=G\Pi$,
$$
 M_p\,e^{\mathsf T}\mu(\widehat x_p)e=e^{\mathsf T}G\widehat\Pi K_pe=e^{\mathsf T}GK_p\Pi e
 =M_p\,e^{\mathsf T}\mu(x_p)e .
$$
Next, $(ee^{\mathsf T}G)^2=M_pee^{\mathsf T}G$ gives $K_p^2=M_p^2I$, so
$\widehat\Pi=R_p\Pi R_p$; and $GR_pG^{-1}=R_p^{\mathsf T}$ because $G$ is symmetric. Therefore
$\mu(\widehat x_p)=GR_p\Pi R_p=R_p^{\mathsf T}\mu(x_p)R_p$. Since $P$ is symmetric, this gives
$\mc G(\widehat x_p)=(P^{-1}R_pP)^{\mathsf T}\,\mc G(x_p)\,(P^{-1}R_pP)$. Since $P^{-1}e=e_1$ and
$e^{\mathsf T}GP=e_1^{\mathsf T}PGP=(M_p,\rho_p)$,
$$
 P^{-1}R_pP=\frac{2e_1(M_p,\rho_p)-M_pI}{M_p}=\begin{pmatrix}1&2\rho_p/M_p\\0&-1\end{pmatrix},
$$
and the computation in the proof of \cref{lem:gram-bridge} shows that the congruence by this
matrix takes a Gram matrix with first column $(M,\rho)$ to one with first column
$(M,\,2M\rho_p/M_p-\rho)$. This is the root relation. The word $x_p$ contains $2k+1$ copies of
$p$ and $2k$ blocks, each with one letter of each kind, which gives the completed counts.
For $k\geqslant1$ their greatest common divisor is at least $2k+1$, while central words have
coprime completed counts by \cref{lem:standard-identities}; and the two words differ at the
first letter of the first block.
\end{proof}

The exchange also has a transparent arithmetic. Recall that the roots $r$ of $-1$ modulo a
label $M$ correspond bijectively to the primitive Gaussian integers $x+iy$ of norm $M$ taken
up to units, through $r\equiv xy^{-1}\pmod M$, and that replacing $r$ by $-r$ replaces $x+iy$
by its conjugate.

\begin{proposition}[The arithmetic of an exchange]
\label{prop:exchange-arithmetic}
Let $x_p$ be structured around $p$ with $2k\geqslant2$ blocks and common label $M$. Then
$M_p$ divides $M$, the quotient $N=M/M_p$ is prime to $M_p$, and
$$
 \rho(x_p)\equiv\rho(\widehat x_p)\equiv\rho_p\pmod{M_p},
 \qquad
 \rho(x_p)+\rho(\widehat x_p)\equiv0\pmod N .
$$
The exchange therefore keeps the Gaussian factor of norm $M_p$ and conjugates the factor of
norm $N$.
\end{proposition}

\begin{proof}
As in the proof of \cref{prop:one-block-closed-form}, inserting $P^{-1}P$ between the factors
gives $\mc G(x_p)=\mc G_pY_1\mc G_pY_2\cdots Y_n\mc G_p$, where $\mc G_p=\mc G(p)$ and $Y_i$ is
$K$ when $W_i=ba$ and $K^{\mathsf T}$ when $W_i=ab$. Directly from these two matrices,
$e_2^{\mathsf T}K=e_1^{\mathsf T}$, $e_2^{\mathsf T}K^{\mathsf T}=-e_1^{\mathsf T}$,
$Ke_2=-e_1$ and $K^{\mathsf T}e_2=e_1$; and modulo $M_p$, entrywise,
$e_1^{\mathsf T}\mc G_p\equiv\rho_pe_2^{\mathsf T}$ and $\mc G_pe_1\equiv\rho_pe_2$, because the
first row of $\mc G_p$ is $(M_p,\rho_p)$. Reading $M(x_p)=e_1^{\mathsf T}\mc G(x_p)e_1$ from the
left, each factor $\mc G_pY_i$ turns a multiple of $e_1^{\mathsf T}$ into $\pm\rho_p$ times
$e_1^{\mathsf T}$, so $M(x_p)\equiv\pm\rho_p^{\,n}\,e_1^{\mathsf T}\mc G_pe_1=\pm\rho_p^{\,n}M_p\equiv0$.
Reading $\rho(x_p)=e_2^{\mathsf T}\mc G(x_p)e_1$ from the right, the column starts as
$\rho_pe_2$, and each factor $\mc G_pY_i$ multiplies a multiple of $e_2$ by $\rho_p$ and by
$-1$ if $W_i=ba$, $+1$ if $W_i=ab$. Hence
$\rho(x_p)\equiv\rho_p^{\,2k+1}(-1)^{j}$, where $j$ is the number of blocks $ba$. For a
structured word the blocks pair off as $W_i$ and its exchange, so $j=k$; and
$\rho_p^2\equiv-1$ because $\det\mc G_p=1$. Hence $\rho(x_p)\equiv(-1)^k(-1)^k\rho_p=\rho_p$,
and the same holds for $\widehat x_p$.

For the coprimality write $\rho(x_p)=\rho_p+M_ps$ and $\rho(\widehat x_p)=\rho_p+M_ps'$. By
\cref{prop:any-centre} their sum is $2\rho_pN$, so $M_p(s+s')=2\rho_p(N-1)$. A prime dividing
both $M_p$ and $N$ divides neither $N-1$ nor $\rho_p$, which is prime to $M_p$, so it must
divide $2$; then $4$ divides $M=M_pN$. But $M$ is a sum of two coprime squares, the squared
norm of the first column of $T$, and such a sum is $1$ or $2$ modulo $4$. Hence
$\gcd(M_p,N)=1$, and $\rho(x_p)+\rho(\widehat x_p)=2\rho_pN\equiv0\pmod N$. By the Chinese
remainder theorem the two roots agree modulo $M_p$ and are opposite modulo $N$; a root and
its negative select conjugate Gaussian primes, which gives the last assertion.
\end{proof}

With the centre $p=baab$, whose label is $1130$ and root $467$, the closed form
\eqref{eq:one-block-closed-form} gives the label $1130\cdot11\,492\,101=12\,986\,074\,130$ and
the roots $5\,366\,814\,557$ and $5\,366\,807\,777$. Since $baab$ is not central, this pair lies
outside \cref{tab:pairs}, and the same label shows that the midpoint condition of
\cref{thm:palindromic-classification} cannot be removed.

\begin{remark}[The midpoint hypothesis cannot be dropped]
\label{rem:quadruple}
Exactly four palindromes carry the label of the example above,
$M=12\,986\,074\,130$, namely
$$
 \begin{aligned}
 Z_1&=baab\,ba\,baab\,ab\,baab, &\qquad Z_2&=baab\,ab\,baab\,ba\,baab,\\
 Z_3&=abba\,ba\,abba\,ab\,abba, &\qquad Z_4&=abba\,ab\,abba\,ba\,abba.
 \end{aligned}
$$
The pairs $(Z_1,Z_2)$ and $(Z_3,Z_4)$ are exchanges around $baab$ and $abba$. Read in blocks of
two letters, all four words lie in $\{ab,ba\}^*$, and $(Z_1,Z_4)$ and $(Z_2,Z_3)$ are exchanges
around the empty word, with eight blocks each. The remaining pairs $(Z_1,Z_3)$ and $(Z_2,Z_4)$
have equal labels and lie in the image of the empty directive, but neither word is the block
exchange of the other: their root midpoints are $2/5\pm3390/M$ instead of the ancestor
ratio $2/5$. Without the midpoint hypothesis, the conclusion of
\cref{thm:palindromic-classification} fails. A label can also be carried by more than two
palindromes, and two palindromes with the same label need not be exchanges of each other.
\end{remark}

The four words form a single class under exchanges, each around its own centre, and the two
pairs that are not exchanges have their midpoints shifted off every centre. The next subsection
shows that a midpoint at the ratio of a palindrome forces an exchange around it, at every centre,
and \cref{subsec:nested} explains the shifted midpoints.

\subsection{Decoding at every palindromic centre}
\label{subsec:centre-decoder}

\Cref{prop:any-centre} produces pairs of word columns related by the reflection $R_p$ of any
palindrome $p$. We now prove that there are no others, for every palindrome and all finite words,
and deduce the classification of \cref{thm:palindromic-classification} at every centre, without
the hypotheses of a common standard image and of a common prefix. The argument is shorter than that
of \cref{sec:decoder}. It works with the letters $a$ and $b$ in the slope coordinate, reads the
first two letters of each word, and then reads the whole centre letter by letter through the
transposed reflection. Beyond the symmetry of $\mu(p)$, the only number attached to the centre
that enters the estimates is the window in which its ratio lies; the letters of the centre are read
one at a time, by estimates valid for every word.

We write $t(v)=v_1/v_2$ for the slope of a column $v$, with $t(v)=\infty$ when $v_2=0$, and $t(w)$
for the slope of $\mu(w)e$; a matrix acts on slopes by fractional linear transformation. Two
nonzero columns span the same line exactly when their slopes agree. By \cref{sec:decoder}, every
word slope lies in $\mathcal J=[8/5,29/12]$, with $t(\varepsilon)=2$, $t(aw)\in A(\mathcal J)=[21/13,70/41]$
and $t(bw)\in B(\mathcal J)=[50/21,169/70]$, because $t(cw)=\mu(c)(t(w))$ and the two maps are
increasing on $\mathcal J$. So $t(w)<2$, $t(w)=2$ or $t(w)>2$ according as $w$ begins with $a$, is
empty, or begins with $b$, and no word slope is infinite, nonpositive, smaller than $21/13$ or
larger than $169/70$. We call this the \emph{first-letter rule}.

Fix a palindrome $p$ and write $G=\mu(p)$, $M=M_p$, $\rho=\rho_p$, $K=K_p$ and $\tau=2\rho/M$.
Since $Ge=(\rho,M-2\rho)^{\mathsf T}$ and $e^{\mathsf T}G=(Ge)^{\mathsf T}$,
\begin{equation}\label{eq:centre-K}
 K=\begin{pmatrix}4\rho-M&4M-8\rho\\2\rho&M-4\rho\end{pmatrix},\qquad
 K^{\mathsf T}=\begin{pmatrix}4\rho-M&2\rho\\4M-8\rho&M-4\rho\end{pmatrix}.
\end{equation}
Five facts follow by direct computation. First, $K^2=M^2I$, $Ke=Me$ and $GK=K^{\mathsf T}G$, as in
the proof of \cref{prop:any-centre}; hence $K^{\mathsf T}Ge=GKe=M\,Ge$, and
$K^{\mathsf T}(1,-2)^{\mathsf T}=2Ge\,e^{\mathsf T}(1,-2)^{\mathsf T}-M(1,-2)^{\mathsf T}=-M(1,-2)^{\mathsf T}$.
Second, multiplying \eqref{eq:centre-K} by $AB$, $BA$ and $A$,
\begin{equation}\label{eq:centre-intertwining}
 KAB=BAK^{\mathsf T},\qquad KBA=ABK^{\mathsf T},\qquad
 B^{-1}KA=M\begin{pmatrix}0&1\\1&-1-\tau\end{pmatrix};
\end{equation}
for instance both upper-left entries in the first identity equal $12(4\rho-M)+7(4M-8\rho)=16M-8\rho$.
Third, the slope $t(p)=\rho/(M-2\rho)$ lies in $\mathcal J$ and $\tau=2t(p)/(2t(p)+1)$ increases
with it, so $16/21\leqslant\tau\leqslant29/35$, the values at $8/5$ and at $29/12$. Fourth, for a
real $t$, with $\eta=1+\tau(t-2)$,
\begin{equation}\label{eq:centre-slopes}
 K\binom t1=M\binom{2\eta+2-t}{\eta},\qquad
 K^{\mathsf T}\binom t1=M\binom{(2\tau-1)t+\tau}{(4-4\tau)t+1-2\tau}.
\end{equation}
Fifth, a relation $[T\mu(x)e]=[\mu(y)e]$ with $T^2$ a nonzero multiple of $I$ is symmetric in $x$
and $y$: applying $T$ to both sides gives $[\mu(x)e]=[T\mu(y)e]$.

\begin{lemma}[The first two letters]
\label{lem:centre-first-letters}
Let $p$ be a palindrome and $x,y$ words with $[K_p\mu(x)e]=[\mu(y)e]$. Then $x$ is empty if and
only if $y$ is empty; if they are not empty, their first letters differ; and if $x$ begins with
$a$, then $x=ab\,x_1$ and $y=ba\,y_1$ for words $x_1,y_1$ with $[K_p^{\mathsf T}\mu(x_1)e]=[\mu(y_1)e]$.
\end{lemma}

\begin{proof}
\emph{First letters.} Let $t=t(x)$. Since $t-2\geqslant-2/5$ and $\tau\leqslant29/35$, the number
$\eta=1+\tau(t-2)$ is at least $1-58/175=117/175>0$. By \eqref{eq:centre-slopes} the column
$K\mu(x)e$ is a positive multiple of $(2\eta+2-t,\eta)^{\mathsf T}$, so
$t(y)=2+(2-t(x))/\eta$, and $t(y)-2$ has the sign of $2-t(x)$, zero included. The first-letter
rule gives the first two assertions.

\emph{The second letter of $x$.} Let $x=ax'$, so that $y=by'$. Multiplying
$[KA\mu(x')e]=[B\mu(y')e]$ by $B^{-1}$ and using \eqref{eq:centre-intertwining}, the column
$B^{-1}KA\mu(x')e$ is a positive multiple of $(1,\,t(x')-1-\tau)^{\mathsf T}$, so
$t(y')=1/(t(x')-1-\tau)$, read as $\infty$ when the denominator vanishes. If $x'$ were empty, then
$t(y')=1/(1-\tau)\geqslant21/5>169/70$. If $x'$ began with $a$, then
$t(x')\leqslant70/41<37/21\leqslant1+\tau$, and $t(y')$ would be negative. Neither is a word slope,
so $x'=bx_1$.

\emph{The second letter of $y$.} Now $[KAB\mu(x_1)e]=[B\mu(y')e]$, and
\eqref{eq:centre-intertwining} turns it into $[AK^{\mathsf T}\mu(x_1)e]=[\mu(y')e]$. By
\eqref{eq:centre-slopes}, with $t=t(x_1)\geqslant8/5$, the column $K^{\mathsf T}\mu(x_1)e$ is a
positive multiple of $((2\tau-1)t+\tau,\,(4-4\tau)t+1-2\tau)^{\mathsf T}$. Its first entry is
positive because $\tau>1/2$, and its second entry increases with $t$ and is at least
$(37-42\tau)/5\geqslant11/25$ at $t=8/5$. So $K^{\mathsf T}\mu(x_1)e$ has a positive slope $s$, and
$t(y')=A(s)$ lies strictly between $1$ and $2$. By the first-letter rule $y'=ay_1$, and cancelling
$A$ gives $[K^{\mathsf T}\mu(x_1)e]=[\mu(y_1)e]$.
\end{proof}

The transposed reflection $K^{\mathsf T}$ fixes the column $Ge$ of the centre and reverses
$(1,-2)^{\mathsf T}$. The next lemma shows that such a reflection can relate two word columns only if
both words begin as the centre does.

\begin{lemma}[One letter of a fixed word]
\label{lem:centre-one-letter}
Let $T$ be a real matrix and $\kappa\neq0$ with $Tu=\kappa u$ and $Tv=-\kappa v$, where $u=\mu(q)e$
for a nonempty word $q$ with first letter $c$, and $v$ is a column whose slope $\omega$ lies in
$[-1,-2/5]$. If $x$ and $y$ are words with $[T\mu(x)e]=[\mu(y)e]$, then both begin with $c$.
\end{lemma}

\begin{proof}
Put $\theta=t(q)$, which lies in $A(\mathcal J)$ if $c=a$ and in $B(\mathcal J)$ if $c=b$, and
$\Omega=-\omega\in[2/5,1]$. For a real $t\neq\omega$ let $\sigma(t)=(t-\theta)/(t-\omega)$. Writing
$(t,1)^{\mathsf T}=\alpha(\theta,1)^{\mathsf T}+\beta(\omega,1)^{\mathsf T}$ with
$\alpha=(t-\omega)/(\theta-\omega)$ and $\beta=(\theta-t)/(\theta-\omega)$, one has
$T(t,1)^{\mathsf T}=\kappa(\alpha(\theta,1)^{\mathsf T}-\beta(\omega,1)^{\mathsf T})$, and when
its slope $t'$ is finite, $\sigma(t')=\beta/\alpha=-\sigma(t)$. The relation gives $t(y)=t'$ for
$t=t(x)$, so $\sigma(t(y))=-\sigma(t(x))$. On $(\omega,\infty)$, which contains every word slope,
$\sigma$ is increasing, with derivative $(\theta-\omega)/(t-\omega)^2$, and $\sigma<1$.

\emph{The case $c=a$.} Suppose that $x$ does not begin with $a$, so $t(x)\geqslant2$ and
$\sigma(t(x))\geqslant\sigma_0=(2-\theta)/(2+\Omega)>0$. Solving $\sigma(t(y))=-\sigma(t(x))$ for
$t(y)$ and using that $\sigma\mapsto\sigma/(1+\sigma)$ increases,
$$
 t(y)=\theta-\frac{\sigma(t(x))}{1+\sigma(t(x))}(\theta+\Omega)
 \leqslant\theta-\frac{\sigma_0}{1+\sigma_0}(\theta+\Omega)
 =\theta-(2-\theta)\,\frac{\theta+\Omega}{4-\theta+\Omega}.
$$
The last quotient increases with $\Omega$, because $\theta<2$, and with $\theta$, so it is at least
its value $131/181$ at $\theta=21/13$, $\Omega=2/5$. With $\theta\leqslant70/41$ and
$2-\theta\geqslant12/41$ this gives
$t(y)\leqslant70/41-(12/41)(131/181)=11098/7421<21/13$, which is the slope of no word.

\emph{The case $c=b$.} Suppose that $x$ does not begin with $b$, so $t(x)\leqslant2<\theta$ and
$-\sigma(t(x))\geqslant\sigma_0=(\theta-2)/(2+\Omega)>0$; also $\sigma_0\leqslant29/168<1$. Since
$t(y)$ is a word slope, it lies in $(\omega,\infty)$, so $-\sigma(t(x))=\sigma(t(y))<1$. Solving as
before, and using that $\sigma\mapsto\sigma/(1-\sigma)$ increases on $[0,1)$,
$$
 t(y)\geqslant\theta+\frac{\sigma_0}{1-\sigma_0}(\theta+\Omega)
 =\theta+(\theta-2)\,\frac{\theta+\Omega}{4-\theta+\Omega}.
$$
The last quotient now decreases with $\Omega$, because $\theta>2$, and increases with $\theta$, so it
is at least its value $71/55$ at $\theta=50/21$, $\Omega=1$. With $\theta\geqslant50/21$ and
$\theta-2\geqslant8/21$ this gives $t(y)\geqslant50/21+(8/21)(71/55)=158/55>169/70$, which is the
slope of no word.

In both cases $x$ begins with $c$. The columns $u$ and $v$ have different slopes, so $T^2=\kappa^2I$,
the relation is symmetric, and the same argument shows that $y$ begins with $c$.
\end{proof}

\begin{lemma}[Reading the centre]
\label{lem:centre-reading}
Let $p$ be a palindrome. If two words $x$ and $y$ satisfy
$[K_p^{\mathsf T}\mu(x)e]=[\mu(y)e]$, then $x=p\,x_0$ and $y=p\,y_0$, where the words $x_0$ and
$y_0$ satisfy $[K_p\mu(x_0)e]=[\mu(y_0)e]$.
\end{lemma}

\begin{proof}
Write $p=c_1\cdots c_n$, and for $0\leqslant j\leqslant n$ put $\Pi_j=\mu(c_1\cdots c_j)$ and
$T_j=\Pi_j^{-1}K^{\mathsf T}\Pi_j$. Then $T_0=K^{\mathsf T}$ and $T_n=G^{-1}K^{\mathsf T}G=K$. For
$j<n$ let $q_j=c_{j+1}\cdots c_n$, a nonempty word with first letter $c_{j+1}$ and
$\Pi_j\mu(q_j)=G$. Then $T_j\mu(q_j)e=\Pi_j^{-1}K^{\mathsf T}Ge=M\mu(q_j)e$, and the column
$v_j=\Pi_j^{-1}(1,-2)^{\mathsf T}$ satisfies $T_jv_j=-Mv_j$. Its slope $\omega_j$ lies in
$[-1,-2/5]$. Indeed $\omega_0=-1/2$ and $\omega_j=\mu(c_j)^{-1}(\omega_{j-1})$; for $t<0$,
$$
 A^{-1}(t)=\frac{t-1}{2-t}\in\Bigl(-1,-\frac12\Bigr),\qquad
 B^{-1}(t)=\frac{t-2}{5-2t}\in\Bigl(-\frac12,-\frac25\Bigr),
$$
because $A^{-1}(t)+1=1/(2-t)>0$, $A^{-1}(t)+1/2=t/(2(2-t))<0$, $B^{-1}(t)+1/2=1/(2(5-2t))>0$ and
$B^{-1}(t)+2/5=t/(5(5-2t))<0$; induction on $j$ concludes.

Put $x^{(0)}=x$ and $y^{(0)}=y$, and suppose that $[T_j\mu(x^{(j)})e]=[\mu(y^{(j)})e]$ for some
$j<n$. \Cref{lem:centre-one-letter}, applied to $T_j$ with $q=q_j$ and $\kappa=M$, shows that
$x^{(j)}$ and $y^{(j)}$ both begin with $c_{j+1}$. Writing $x^{(j)}=c_{j+1}x^{(j+1)}$ and
$y^{(j)}=c_{j+1}y^{(j+1)}$ and multiplying by $\mu(c_{j+1})^{-1}$ gives
$[T_{j+1}\mu(x^{(j+1)})e]=[\mu(y^{(j+1)})e]$. After $n$ steps $x=p\,x^{(n)}$ and $y=p\,y^{(n)}$, and
the relation for $T_n=K$ is the claim with $x_0=x^{(n)}$ and $y_0=y^{(n)}$.
\end{proof}

\begin{theorem}[Decoding at every palindromic centre]
\label{thm:centre-decoder}
For every palindrome $p$ and all finite words $x$ and $y$, the empty word included, the lines
$[R_p\mu(x)e]$ and $[\mu(y)e]$ coincide if and only if $x\in\{ab\,p,\ ba\,p\}^*$ and $y$ is
obtained from $x$ by exchanging every block. In that case $R_p\mu(x)e=\mu(y)e$ exactly.
\end{theorem}

\begin{proof}
The implication from right to left and the exact equality are \cref{prop:any-centre}: for a block
word $x$ and its exchange $\widehat x$, $K\mu(x)=\mu(\widehat x)K$, and $Ke=Me$. For the converse we
argue by induction on $\abs x+\abs y$. Suppose $[K\mu(x)e]=[\mu(y)e]$. If one of the words is empty,
both are, by \cref{lem:centre-first-letters}, and the empty word is a block word equal to its
exchange. Otherwise their first letters differ. If $x$ begins with $a$,
\cref{lem:centre-first-letters} gives $x=ab\,x_1$, $y=ba\,y_1$ and
$[K^{\mathsf T}\mu(x_1)e]=[\mu(y_1)e]$, and \cref{lem:centre-reading} gives $x_1=p\,x''$,
$y_1=p\,y''$ and $[K\mu(x'')e]=[\mu(y'')e]$. Since $\abs{x''}+\abs{y''}<\abs x+\abs y$, the
induction hypothesis shows that $x''$ is a block word and $y''$ its exchange, so $x=ab\,p\,x''$ is a
block word and $y=ba\,p\,y''$ is its exchange. If $x$ begins with $b$, the relation is symmetric,
because $K^2=M^2I$; the previous case applied to $(y,x)$ gives $y=ab\,p\,y''$ and $x=ba\,p\,x''$ with
$[K\mu(y'')e]=[\mu(x'')e]$, hence $[K\mu(x'')e]=[\mu(y'')e]$, and the induction hypothesis concludes
in the same way.
\end{proof}

For a standard centre $p=\operatorname{Pal}(d)$ one has $ab\,p=\phi_d(ab)$ and $ba\,p=\phi_d(ba)$ by
\eqref{eq:literal-standard-identities}, so the theorem contains \cref{thm:word-decoder} and shows
moreover that no word outside the image of $\phi_d$ is related to a word column by $R_d$. The proof
uses neither the standard letters $U$, $V$ nor the cylinders of \cref{sec:decoder}; those cylinders
are still needed in \cref{sec:displaced}, where they measure how far a column moves when the
midpoint is displaced. The theorem has the following consequence for palindromes.

\begin{corollary}[Classification at every centre]
\label{cor:centre-classification}
Let $p$, $Z$ and $Z'$ be palindromes. Then $M(Z)=M(Z')$ and $\rho(Z)+\rho(Z')=2M(Z)\rho_p/M_p$ if and
only if $Z=p\,X$ and $Z'=p\,\widehat X$ for a block word $X\in\{ab\,p,\ ba\,p\}^*$. In that case either
$Z=Z'=p$, or $Z$ is structured around $p$ and $Z'$ is its exchange, and the two words have the same
numbers of letters $a$ and $b$. In particular two distinct central words never have equal labels
with root midpoint at the ratio of a palindrome.
\end{corollary}

\begin{proof}
Put $J_\tau=\left(\begin{smallmatrix}1&\tau\\0&-1\end{smallmatrix}\right)$ with $\tau=2\rho_p/M_p$;
by the proof of \cref{prop:any-centre}, $P^{-1}R_pP=J_\tau$. The computation in the proof of
\cref{lem:gram-bridge} shows that $J_\tau^{\mathsf T}\mathcal G(Z)J_\tau$ is symmetric of
determinant one with first row $(M(Z),\,\tau M(Z)-\rho(Z))$, and a symmetric matrix of determinant
one with positive upper-left entry is determined by its first row. So the two conditions on labels
and roots hold if and only if $\mathcal G(Z')=J_\tau^{\mathsf T}\mathcal G(Z)J_\tau$, that is, since $P$
is symmetric, $\mu(Z')=R_p^{\mathsf T}\mu(Z)R_p$. If this holds, then
$\mu(Z')e=R_p^{\mathsf T}\mu(Z)e$ because $R_pe=e$, and \cref{lem:centre-reading} gives $Z=p\,X$ and
$Z'=p\,Y$ with $[K_p\mu(X)e]=[\mu(Y)e]$; by \cref{thm:centre-decoder}, $X$ is a block word and
$Y=\widehat X$. Conversely, if $Z=p\,X$ and $Z'=p\,\widehat X$, then $K_p\mu(X)=\mu(\widehat X)K_p$ and
$K_p^2=M_p^2I$ give $\mu(\widehat X)=R_p\mu(X)R_p$, so
$\mu(Z')=G\,R_p\mu(X)R_p=R_p^{\mathsf T}G\mu(X)R_p=R_p^{\mathsf T}\mu(Z)R_p$. Since $Z$ is a palindrome
and every block has length $\abs p+2$, comparing $Z$ with its reversal shows that $X$ has an even
number $2k$ of blocks with $W_{2k+1-i}=\widehat W_i$; so either $k=0$, or $Z$ is structured around
$p$ and $Z'$ is its exchange. Exchanging blocks does not change the letter counts. Finally a word
structured with $k\geqslant1$ has completed counts $(2k+1)\,c(p)$, which are not coprime, so it is
not central.
\end{proof}

The centre in the corollary is determined by the pair: the fraction $\rho_p/M_p$ is reduced,
because $\rho_p^2\equiv-1\pmod{M_p}$, so it gives $M_p$ and $\rho_p$, hence the Gram matrix of $p$
and, the Cohn coding being injective, the palindrome $p$ itself. The corollary contains
\cref{thm:palindromic-classification}: when $Z=F_d(z_X)$ and $Z'=F_d(z_Y)$ and $p=\operatorname{Pal}(d)$,
the block word $X$ lies in $\phi_d(\{ab,ba\}^*)$, and the injectivity of $F_d$ returns the pair
$(\widetilde ww,\widetilde{\widehat w}\,\widehat w)$. \Cref{rem:quadruple} shows that the hypothesis on
the midpoint cannot be dropped.

\subsection{Nested centres and shifted midpoints}
\label{subsec:nested}

The quadruple of \cref{rem:quadruple} raises two questions: why do four palindromes share a
label, and why are the midpoints of two of its pairs shifted off every centre? The answer is that
exchanges around nested centres compose. When the centre of an exchange is itself one of two
exchanged palindromes, the words built around the two of them form a square of twins, and nesting
$r$ levels gives a cube with $2^r$ vertices. The shift of a midpoint is then the geometry of two
reflections of the ratio axis, which compose to a translation.

It is convenient to name the words of \cref{subsec:any-centre} by their blocks. A \emph{block
sequence} is a finite sequence $B=(B_1,\dots,B_n)$ with $B_i\in\{ab,ba\}$; we write
$\widehat B=(\widehat B_1,\dots,\widehat B_n)$, and we call $B$ \emph{structured} when $n=2k\geqslant2$
and $B_{n+1-i}=\widehat B_i$ for every $i$. For a word $q$ put
$$
 q\langle B\rangle=q\,B_1\,q\,B_2\,q\cdots q\,B_n\,q .
$$
The words structured around a palindrome $p$ are exactly the words $p\langle B\rangle$ with $B$
structured, and the exchange of $p\langle B\rangle$ around $p$ is $p\langle\widehat B\rangle$. For
block sequences $X=(X_1,\dots,X_n)$ and $Y=(Y_1,\dots,Y_m)$ let
$$
 X\ast Y=(X,\,Y_1,\,X,\,Y_2,\,\dots,\,X,\,Y_m,\,X)
$$
be the block sequence in which the blocks of $X$ are written $m+1$ times, separated by
$Y_1,\dots,Y_m$; it has $(m+1)n+m$ blocks.

\begin{lemma}[Reading through an inner centre]
\label{lem:nested-reading}
For every word $q$ and all block sequences $X$ and $Y$,
$$
 \bigl(q\langle X\rangle\bigr)\langle Y\rangle=q\langle X\ast Y\rangle,
 \qquad \widehat{X\ast Y}=\widehat X\ast\widehat Y .
$$
If $X$ and $Y$ are structured, then so is $X\ast Y$.
\end{lemma}

\begin{proof}
The word $Q=q\langle X\rangle$ begins and ends with $q$. In $Q\langle Y\rangle=Q\,Y_1\,Q\cdots Y_m\,Q$
every junction reads $\cdots X_n\,q\,Y_j\,q\,X_1\cdots$, so the whole word is $q$ followed by the
blocks $X_1,\dots,X_n,Y_1,X_1,\dots,X_n,Y_2,\dots,Y_m,X_1,\dots,X_n$, each of them followed by one
copy of $q$. This is $q\langle X\ast Y\rangle$. The second identity holds block by block. For the
last assertion let $X$ have $n=2k$ blocks and $Y$ have $m=2j$ blocks. Reading $X\ast Y$ backwards
gives $(\overleftarrow X,Y_m,\overleftarrow X,\dots,Y_1,\overleftarrow X)$, where
$\overleftarrow X=(X_n,\dots,X_1)$. Structure of $X$ means $\overleftarrow X=\widehat X$, and
structure of $Y$ means $Y_{m+1-i}=\widehat Y_i$, so the reversed sequence is
$(\widehat X,\widehat Y_1,\widehat X,\dots,\widehat Y_m,\widehat X)=\widehat X\ast\widehat Y$, which
is $\widehat{X\ast Y}$ by the second identity. Its number of blocks, $(m+1)n+m$, is even and at
least two.
\end{proof}

\begin{theorem}[Cubes of nested exchanges]
\label{thm:nested-cubes}
Let $p$ be a palindrome, let $r\geqslant1$, and let $B_1,\dots,B_r$ be structured block sequences,
$B_i$ with $2k_i$ blocks. For $s\in\{0,1\}^r$ put $B_i^s=B_i$ if $s_i=0$ and $B_i^s=\widehat B_i$
if $s_i=1$, and define
$$
 Q_0^s=p,\qquad Q_i^s=Q_{i-1}^s\langle B_i^s\rangle\quad(1\leqslant i\leqslant r),\qquad Z^s=Q_r^s .
$$
For $1\leqslant i\leqslant r$ let $f_i\in\{0,1\}^r$ have its entries equal to one exactly in the
positions $i,\dots,r$, and write $s+f_i$ for the sum modulo two. Then:
\begin{enumerate}[label=\textup{(\roman*)}]
\item the $2^r$ words $Z^s$ are pairwise distinct palindromes, all with completed counts
 $(2k_1+1)\cdots(2k_r+1)\,c(p)$;
\item for every $s$ and every $i$, the words $Z^s$ and $Z^{s+f_i}$ are exchanges of each other
 around the palindrome $Q_{i-1}^s$;
\item all $2^r$ words have the same Gram label $M$, any two of them are joined by at most $r$
 exchanges of the kind in \textup{(ii)}, and the roots of the pair in \textup{(ii)} satisfy
 $\rho(Z^s)+\rho(Z^{s+f_i})=2M\rho(Q_{i-1}^s)/M(Q_{i-1}^s)$.
\end{enumerate}
\end{theorem}

\begin{proof}
\emph{Every level is a centre.} Define block sequences $R_r^s=B_r^s$ and
$R_i^s=B_i^s\ast R_{i+1}^s$ for $1\leqslant i<r$. We claim that $Z^s=Q_{i-1}^s\langle R_i^s\rangle$
for every $i$, and prove it by downward induction on $i$. For $i=r$ it is the definition
$Z^s=Q_r^s=Q_{r-1}^s\langle B_r^s\rangle$. If $Z^s=Q_i^s\langle R_{i+1}^s\rangle$, then
$Q_i^s=Q_{i-1}^s\langle B_i^s\rangle$ and \cref{lem:nested-reading} give
$Z^s=Q_{i-1}^s\langle B_i^s\ast R_{i+1}^s\rangle=Q_{i-1}^s\langle R_i^s\rangle$. By the last
assertion of \cref{lem:nested-reading} and induction, every $R_i^s$ is structured. The same
argument applied to the word $Q_{i-1}^s$ in place of $Z^s$ writes it as $p\langle R\rangle$ with $R$
structured, or as $p$ itself when $i=1$; in both cases $Q_{i-1}^s$ is a palindrome. Hence $Z^s$ is
structured around the palindrome $Q_{i-1}^s$.

\emph{The exchange is a flip of the last coordinates.} The exchange of $Z^s$ around $Q_{i-1}^s$ is
$Q_{i-1}^s\langle\widehat{R_i^s}\rangle$. Applying the second identity of \cref{lem:nested-reading}
repeatedly, $\widehat{R_i^s}=\widehat{B_i^s}\ast\widehat{R_{i+1}^s}=\cdots$ is the sequence $R_i^{s'}$
built from $s'=s+f_i$, the vector obtained by changing $s_i,\dots,s_r$. The first $i-1$ coordinates
of $s'$ are those of $s$, so $Q_{i-1}^{s'}=Q_{i-1}^s$, and the exchange is
$Q_{i-1}^{s'}\langle R_i^{s'}\rangle=Z^{s'}$. This proves (ii).

\emph{Palindromes, distinctness and counts.} The case $i=1$ of the first step gives
$Z^s=p\langle R_1^s\rangle$ with $R_1^s$ structured, so $Z^s$ is a palindrome. Let $s\neq s'$, let
$i$ be the first coordinate in which they differ, and put $Q=Q_{i-1}^s=Q_{i-1}^{s'}$. Every
$Q_j^s$ begins with $Q_{j-1}^s$, so $Z^s$ begins with $Q\langle B_i^s\rangle$ and $Z^{s'}$ begins
with $Q\langle\widehat{B_i^s}\rangle$; the two words differ in the letter that follows the first
copy of $Q$, which is the first letter of a block in one and of its exchange in the other. By
\cref{prop:any-centre}, a word structured around a centre with $2k$ blocks has completed counts
$2k+1$ times those of the centre; applied at the levels $1,\dots,r$ this gives the counts in (i).

\emph{Labels and roots.} By \cref{prop:any-centre}, two words exchanged around a palindrome have the
same label, and their roots sum to twice the label times the ratio of the centre. The vectors
$f_1,\dots,f_r$ are linearly independent modulo two, since the matrix with these rows is triangular
with ones on the diagonal. Every vector of $\{0,1\}^r$ is therefore a sum of at most $r$ of them,
and applying the corresponding exchanges one after the other joins $Z^s$ to any other $Z^{s'}$.
The label is therefore common to all $2^r$ words, which proves (iii).
\end{proof}

For $r=2$ the cube is a square, and its two pairs that are not exchanges are the pairs of the
quadruple whose midpoints are shifted.

\begin{corollary}[The square and its shifted midpoints]
\label{cor:nested-square}
Let $p$ be a palindrome and $W$, $V$ structured block sequences; put $P_1=p\langle W\rangle$,
$P_2=p\langle\widehat W\rangle$ and
$$
 Z_1=P_1\langle V\rangle,\qquad Z_2=P_1\langle\widehat V\rangle,\qquad
 Z_3=P_2\langle V\rangle,\qquad Z_4=P_2\langle\widehat V\rangle ,
$$
with common label $M$ and roots $\rho_1,\dots,\rho_4$. The pairs $(Z_1,Z_2)$ and $(Z_3,Z_4)$ are
exchanges around $P_1$ and $P_2$, and the pairs $(Z_1,Z_4)$ and $(Z_2,Z_3)$ are exchanges around
$p$. The two remaining pairs satisfy $\rho_1\neq\rho_2$ and
$$
 \frac{\rho_1+\rho_3}{2M}=\frac{\rho_p}{M_p}+\frac{\rho_1-\rho_2}{2M},\qquad
 \frac{\rho_2+\rho_4}{2M}=\frac{\rho_p}{M_p}-\frac{\rho_1-\rho_2}{2M}.
$$
If $W$ and $V$ have two blocks each, then $M(P_1)=M_p(9M_p^2+1)$, $M=M(P_1)\bigl(9M(P_1)^2+1\bigr)$
and $\abs{\rho_1-\rho_2}=6M(P_1)$.
\end{corollary}

\begin{proof}
Take $r=2$, $B_1=W$ and $B_2=V$ in \cref{thm:nested-cubes}; then $Z_1,Z_2,Z_3,Z_4$ are
$Z^{(0,0)},Z^{(0,1)},Z^{(1,0)},Z^{(1,1)}$. The flip $f_2=(0,1)$ joins $Z_1$ to $Z_2$ around
$Q_1^{(0,0)}=P_1$ and $Z_3$ to $Z_4$ around $Q_1^{(1,0)}=P_2$. The flip $f_1=(1,1)$ joins $Z_1$ to
$Z_4$ and $Z_2$ to $Z_3$, both around $Q_0=p$. By part (iii) of the theorem,
$\rho_1+\rho_4=2M\rho_p/M_p$ and $\rho_2+\rho_3=2M\rho_p/M_p$. Hence
$$
 \rho_1+\rho_3=\rho_1-\rho_2+2M\frac{\rho_p}{M_p},\qquad
 \rho_2+\rho_4=\rho_2-\rho_1+2M\frac{\rho_p}{M_p},
$$
which are the displayed identities after division by $2M$. The words $Z_1$ and $Z_2$ are distinct
palindromes with the same label, and a palindrome is determined by its label and root, because
they determine the Gram matrix through its determinant and the Cohn coding is injective; so
$\rho_1\neq\rho_2$. If $W$ and $V$ have two blocks each, \cref{prop:one-block-closed-form}
applied around $p$ gives the label of $P_1$, and applied around $P_1$ it gives $M$ and
$\rho_1-\rho_2=\pm6M(P_1)$.
\end{proof}

For $p=\varepsilon$ and $W=V=(ba,ab)$ one has $P_1=baab$, $P_2=abba$, $M(P_1)=1130$, and the four
words are those of \cref{rem:quadruple}, with $M=1130\cdot11\,492\,101$. The two pairs that are not
exchanges have their midpoints at $2/5\pm3390/M$: the shift is half the width $6\cdot1130$ of the
rungs around $baab$ and $abba$, divided by the label. \Cref{fig:nested-square} draws this square
above the rung of $abba$ and $baab$.

\begin{figure}[htbp]
\centering
% Author: Dr. Denys Dutykh (Khalifa University of Science and Technology, Abu Dhabi, UAE)
% Generated by supplement/code/make_figures.py from exact data; do not edit by hand.

\begin{tikzpicture}[font=\footnotesize]
\draw[ink, line width=0.5pt] (0,0) -- (10.8,0);
\draw[ink, line width=0.5pt] (1.35,0) -- (1.35,-0.08);
\node[ink, below] at (1.35,-0.08) {$0.385$};
\draw[ink, line width=0.5pt] (2.7,0) -- (2.7,-0.08);
\node[ink, below] at (2.7,-0.08) {$0.390$};
\draw[ink, line width=0.5pt] (4.05,0) -- (4.05,-0.08);
\node[ink, below] at (4.05,-0.08) {$0.395$};
\draw[ink, line width=0.5pt] (5.4,0) -- (5.4,-0.08);
\node[ink, below] at (5.4,-0.08) {$0.400$};
\draw[ink, line width=0.5pt] (6.75,0) -- (6.75,-0.08);
\node[ink, below] at (6.75,-0.08) {$0.405$};
\draw[ink, line width=0.5pt] (8.1,0) -- (8.1,-0.08);
\node[ink, below] at (8.1,-0.08) {$0.410$};
\draw[ink, line width=0.5pt] (9.45,0) -- (9.45,-0.08);
\node[ink, below] at (9.45,-0.08) {$0.415$};
\node[ink, below] at (5.4,-0.48) {$\rho/M$};
\node[ink, left] at (-0.1,0.55) {$M=5$};
\node[ink, left] at (-0.1,2.35) {$M=1130$};
\node[ink, left] at (-0.1,4.75) {$M=12\,986\,074\,130$};
\draw[slate, densely dotted, line width=0.4pt] (5.4,0.55) -- (5.4,4.75);
\draw[slate, densely dotted, line width=0.4pt] (8.984,2.35) -- (8.984,4.75);
\draw[slate, densely dotted, line width=0.4pt] (1.816,2.35) -- (1.816,4.75);
\draw[navy, line width=0.8pt] (1.816,2.35) -- (8.984,2.35);
\draw[navy, line width=0.8pt] (2.366,4.75) .. controls (2.366,5.5) and (8.434,5.5) .. (8.434,4.75);
\draw[navy, line width=0.8pt] (1.266,4.75) .. controls (1.266,5.9) and (9.534,5.9) .. (9.534,4.75);
\draw[wine, line width=0.9pt] (8.434,4.75) -- (9.534,4.75);
\draw[wine, line width=0.9pt] (1.266,4.75) -- (2.366,4.75);
\draw[ochre, fill=white, line width=0.6pt] (4.85,4.43) circle (1.5pt);
\draw[ochre, fill=white, line width=0.6pt] (5.95,4.43) circle (1.5pt);
\draw[slate, line width=0.35pt] (4.8,4.37) -- (4.436,4.13);
\node[ink, right, align=left, inner sep=1pt] at (2.136,3.43) {midpoints $2/5\pm3/N$\\of the diagonals\\$(Z_1,Z_3)$ and $(Z_2,Z_4)$};
\draw[wine, fill=white, line width=0.45pt] (9.534,4.75) circle (1.3pt);
\node[ink, below] at (9.534,4.69) {$Z_1$};
\draw[wine, fill=white, line width=0.45pt] (8.434,4.75) circle (1.3pt);
\node[ink, below] at (8.434,4.69) {$Z_2$};
\draw[wine, fill=white, line width=0.45pt] (2.366,4.75) circle (1.3pt);
\node[ink, below] at (2.366,4.69) {$Z_3$};
\draw[wine, fill=white, line width=0.45pt] (1.266,4.75) circle (1.3pt);
\node[ink, below] at (1.266,4.69) {$Z_4$};
\filldraw[ochre, draw=ink, line width=0.25pt] (5.4,0.55) +(0,1.9pt) -- +(1.9pt,0) -- +(0,-1.9pt) -- +(-1.9pt,0) -- cycle;
\node[ink, right=2.2pt] at (5.4,0.55) {$\varepsilon$};
\filldraw[ochre, draw=ink, line width=0.25pt] (8.984,2.35) +(0,1.9pt) -- +(1.9pt,0) -- +(0,-1.9pt) -- +(-1.9pt,0) -- cycle; \draw[wine, line width=0.45pt] (8.984,2.35) circle (2.7pt);
\node[ink, below right=1pt] at (8.984,2.35) {$baab$};
\filldraw[ochre, draw=ink, line width=0.25pt] (1.816,2.35) +(0,1.9pt) -- +(1.9pt,0) -- +(0,-1.9pt) -- +(-1.9pt,0) -- cycle; \draw[wine, line width=0.45pt] (1.816,2.35) circle (2.7pt);
\node[ink, below left=1pt] at (1.816,2.35) {$abba$};
\draw[navy, line width=0.8pt] (0.2,6.8) -- (0.8,6.8); \node[ink, right] at (0.85,6.8) {exchange around $\varepsilon$};
\draw[wine, line width=0.9pt] (4.6,6.8) -- (5.2,6.8); \node[ink, right] at (5.25,6.8) {exchange around $baab$ or $abba$};
\filldraw[ochre, draw=ink, line width=0.25pt] (0.5,6.35) +(0,1.9pt) -- +(1.9pt,0) -- +(0,-1.9pt) -- +(-1.9pt,0) -- cycle; \node[ink, right] at (0.85,6.35) {centre of an exchange};
\draw[ochre, fill=white, line width=0.6pt] (4.9,6.35) circle (1.5pt); \node[ink, right] at (5.25,6.35) {midpoint of a diagonal};
\end{tikzpicture}

\caption{The square of \cref{cor:nested-square} for $p=\varepsilon$ and $W=V=(ba,ab)$, in the plane
of the ratio $\rho/M$ and the label, at the labels $5$, $1130$ and $12\,986\,074\,130$. The two lower
levels are drawn to scale. At the top level $Z_1$ and $Z_2$ lie at $467/1130\pm3/N$ and $Z_3$ and
$Z_4$ at $437/1130\pm3/N$, where $N=9\cdot1130^2+1=11\,492\,101$; these offsets, about
$2.6\cdot10^{-7}$, are drawn enlarged, all by the same factor. The short rungs are the exchanges
around $baab$ and $abba$, each centred above its centre; the arcs are the exchanges around the empty
word, centred at $2/5$ like the rung below them. The diagonals $(Z_1,Z_3)$ and $(Z_2,Z_4)$ are not
exchanges: their midpoints $2/5\pm3/N$ are the inner centre shifted by half the width of a short
rung.}
\label{fig:nested-square}
\end{figure}

When every block sequence has two blocks, the arithmetic of a cube is completely explicit. Put
$M_i=M(Q_i^s)$, which by \cref{prop:one-block-closed-form} does not depend on $s$, and
$N_i=9M_i^2+1$ for $0\leqslant i<r$, so that $M_0=M_p$ and $M_i=M_{i-1}N_{i-1}$. The same
proposition gives $\rho(Q_i^s)=N_{i-1}\rho(Q_{i-1}^s)\pm3M_{i-1}$, with the sign $+$ exactly when the
first block of $B_i^s$ is $ba$. Since $N_{i-1}\equiv1$ modulo $M_{i-1}=M_pN_0\cdots N_{i-2}$, the
numbers $M_p,N_0,\dots,N_{r-1}$ are pairwise coprime with product $M$. Reducing the recursion,
$\rho(Q_i^s)\equiv\rho(Q_{i-1}^s)$ modulo $M_{i-1}$ and $\rho(Q_i^s)\equiv\pm3M_{i-1}$ modulo
$N_{i-1}$; since every later level keeps the residue modulo the label of the previous one,
$$
 \rho(Z^s)\equiv\rho_p\pmod{M_p},\qquad \rho(Z^s)\equiv\pm3M_{i-1}\pmod{N_{i-1}}\quad(1\leqslant i\leqslant r),
$$
and different $s$ give different patterns of signs. In the correspondence between roots of $-1$
and Gaussian integers recalled before \cref{prop:exchange-arithmetic}, $N_{i-1}=(3M_{i-1})^2+1$ is
the norm of $3M_{i-1}\pm i$, whose roots are $\pm3M_{i-1}$. The $2^r$ twins of a cube therefore carry
the Gaussian integers $\gamma_p\prod_{i=1}^r(3M_{i-1}\pm i)$, with $\gamma_p$ the Gaussian integer
attached to $\rho_p$, for all
choices of signs, and the exchange around $Q_{i-1}^s$ conjugates the factors of the levels
$i,\dots,r$. This is the sign change of \cref{prop:exchange-arithmetic}, made explicit. Whether every
coincidence of labels among palindromes arises from such cubes is the question of
\cref{sec:outlook}.

\section{Displaced reflections and the size of a common ancestor}
\label{sec:displaced}

The ancestor-midpoint hypothesis has an exact algebraic meaning: it is the unique position where the reflection fixes the terminal column on both sides of the word equation. Moving the midpoint introduces a rank-one displacement. We first retain this displacement exactly, then combine the two first-letter cylinders with an integral rational-midpoint coordinate. The resulting bound applies to both parities of the common label. It is a necessary restriction on an equal-label pair, not a classification of all displaced pairs.

\subsection{The full two-sided equation}

Fix a directive $d$, put $w=\operatorname{Pal}(d)$, and retain the right standard map $F_d(z)=w\phi_d(z)$. Write
$$
 G=\mu(w),\quad \nu=\mu\circ\phi_d,\quad
 M_0=e^{\mathsf T}Ge,\quad \rho_0=(Ge)_1,\quad r_0=\rho_0/M_0.
$$
The Gram matrix of the ancestor is $PGP$, with first row $(M_0,\rho_0)$. Define two primitive integral columns
\begin{equation}
 \omega=Ge=\binom{\rho_0}{M_0-2\rho_0},\qquad
 f=\binom1{-2},\qquad a_0=G^{-1}f=\binom{M_0-2\rho_0}{-\rho_0}.
 \label{eq:displaced-columns}
\end{equation}
In particular $\omega^{\mathsf T}a_0=0$ and $f^{\mathsf T}a_0=M_0$.

\begin{proposition}[The moving reflection]
\label{prop:displaced-equation}
Suppose the palindrome images $F_d(z_X),F_d(z_Y)$ have the same Gram label $M$ and roots $\rho_X,\rho_Y$. Put
$$
 \tau=\frac{\rho_X+\rho_Y}{M},\quad \delta=\tau-2r_0,
 \qquad R_0=\frac{2ee^{\mathsf T}G-M_0I}{M_0}.
$$
Then, with $\nu_i=\nu(z_i)$,
\begin{align}
 \nu_Y&=L_\tau\nu_XR_\tau,
 \label{eq:displaced-full-word}\\
 R_\tau&=R_0+\delta ef^{\mathsf T},\qquad
 L_\tau=R_0+\delta a_0\omega^{\mathsf T}.
 \label{eq:displaced-left-right}
\end{align}
The columns $v_i=\nu_i e$ satisfy the exact weighted-scale and displacement identities
\begin{equation}
 \omega^{\mathsf T}v_X=\omega^{\mathsf T}v_Y=M,
 \qquad v_Y=R_0v_X+\delta M a_0.
 \label{eq:displaced-weighted-columns}
\end{equation}
The right factor fixes $e$, whereas $L_\tau e=e+\delta M_0a_0$.
\end{proposition}

\begin{proof}
The image Gram matrices have determinant one and the form
$$
 H_i=PG\nu_iP=
 \begin{pmatrix}M&\rho_i\\\rho_i&(\rho_i^2+1)/M\end{pmatrix}.
$$
For $J_\tau=\left(\begin{smallmatrix}1&\tau\\0&-1\end{smallmatrix}\right)$, direct multiplication gives $H_Y=J_\tau^{\mathsf T}H_XJ_\tau$. Cancelling the invertible factors gives \eqref{eq:displaced-full-word} with
$$
 R_\tau=PJ_\tau P^{-1},\qquad L_\tau=G^{-1}R_\tau^{\mathsf T}G.
$$
At $\tau=2r_0$ both are $R_0$; varying $\tau$ gives the two rank-one terms in \eqref{eq:displaced-left-right}. Since $R_\tau e=e$, the full equation gives $v_Y=L_\tau v_X$. The first Gram entry is $\omega^{\mathsf T}v_X=M$, proving \eqref{eq:displaced-weighted-columns}. The terminal statement follows from $R_0e=e$ and $\omega^{\mathsf T}e=M_0$.
\end{proof}

For a positive column $v$ of slope $t=v_1/v_2$, use the reciprocal centered cusp coordinate
\begin{equation}
 \eta(t)=\frac{t-2}{1+r_0(t-2)}
 =\frac{M_0f^{\mathsf T}v}{\omega^{\mathsf T}v}.
 \label{eq:displaced-eta}
\end{equation}
The denominator is positive because $0<r_0<1/2$ and $t>0$. The homogeneous expression also defines the coordinate whenever $\omega^{\mathsf T}v\ne0$. The identities
$$
 \omega^{\mathsf T}L_\tau=\omega^{\mathsf T},\qquad
 f^{\mathsf T}R_0=-f^{\mathsf T}
$$
give
\begin{equation}
 \eta([L_\tau v])=\xi-\eta([v]),\qquad
 \xi=\delta M_0^2.
 \label{eq:displaced-terminal}
\end{equation}
Thus the terminal $\eta(e)=0$ moves to $\xi$. For the two actual positive image columns,
$$
 \eta(t_X)+\eta(t_Y)=\xi.
$$
This formula alone cannot exclude a pair of positive slopes: any such pair determines a value of $\xi$. The word equation must retain the common weighted scale in \eqref{eq:displaced-weighted-columns}. Replacing $L_\tau$ by $R_\tau$ would erase precisely the information lost at a moving terminal.

\subsection{The actual common prefix and an integral midpoint direction}

Among palindromes in general, equal labels with a displaced midpoint do occur
(\cref{rem:quadruple}), so centrality has to enter any restriction on displaced pairs. We now
suppose $Z_X\ne Z_Y$ are central words with equal Gram label $M$. Let $d$ be the longest common prefix of their unique directives. The following lemma ensures that this directive really supplies a symmetric matrix prefix.

\begin{lemma}[The maximal common prefix]
\label{lem:common-prefix}
Let $Z_X\ne Z_Y$ be central words with equal Gram label $M$, let $w_*$ be
their maximal common symmetric prefix, so that $Z_i=w_*h_i\widetilde w_*$,
and let $d$ be the longest common prefix of their directives. Then both
$h_i$ have length at least two and begin with different letters, and
\begin{equation}
 w_*=\operatorname{Pal}(d)=w,\qquad
 Z_i=w h_iw=F_d(z_i),
 \label{eq:displaced-genuine-fork}
\end{equation}
where $z_X,z_Y$ are nonempty central words beginning with different
letters.
\end{lemma}

\begin{proof}
Every palindrome of length at least two has matrix at least $A^2$
entrywise: remove its two outer letters and use that every letter matrix
is at least $I$. The matrix $A^2$ dominates $I,A,B$, strictly in at least
one entry. Against the strictly positive column $\mu(w_*)^{\mathsf T}e$
its quadratic form therefore exceeds those of all length-zero and
length-one cores, and the three short cores themselves have strictly
ordered forms. Equal labels consequently cannot stop at a short core, and
agreeing next letters would permit another symmetric stripping step;
this proves the first assertion.

It follows that $w_*$ is the full longest common word prefix. A directive
extension $dx\cdots$ produces a central word beginning with
$\operatorname{Pal}(d)x$, and neither directive can be a prefix of the
other, since that would make one whole central word a prefix of the
other, contrary to what has just been proved. This gives
\eqref{eq:displaced-genuine-fork}, in which $z_i$ is the central preimage
of the remaining directive and need not equal the literal core $h_i$.
\end{proof}

Write
\begin{equation}
 PG=\begin{pmatrix}u&v\\s&t\end{pmatrix},\qquad
 ut-vs=-1,\quad u=2s+t,\quad
 M_0=2u+v,\quad\rho_0=u.
 \label{eq:displaced-prefix-frame}
\end{equation}
The four entries are nonnegative, with $u,v,s>0$; the empty directive has $t=0$. The symmetric factorization in \eqref{eq:displaced-genuine-fork} gives
$$
 H_i=(PG)\mu(h_i)(PG)^{\mathsf T}.
$$
The two components of $\mu(h_i)(u,v)^{\mathsf T}$ are positive. Hence $\rho_i/M$ is a strict positive weighted average of $s/u$ and $t/v$, proving
\begin{equation}
 \frac tv<\frac{\rho_i}{M}<\frac su,
 \qquad\frac su-\frac tv=\frac1{uv}.
 \label{eq:displaced-farey-interval}
\end{equation}
The Gram label and root determine the complete Gram matrix, and the Cohn coding is injective. Thus distinct $Z_i$ have distinct roots. Order them as $\rho_X<\rho_Y$, and set $\Sigma=\rho_X+\rho_Y$ and $\Delta\rho=\rho_Y-\rho_X>0$. Both roots lie in $(0,M/2)$, so $0<\Sigma<M$; none of these endpoints occurs.

\begin{lemma}[A reduced midpoint without a parity assumption]
\label{lem:displaced-integer-direction}
Define
\begin{equation}
 c_0=\gcd(\Sigma,2M),\qquad Q=\frac{2M}{c_0},\qquad
 \varepsilon_Q=\gcd(Q,2),\qquad
 \kappa=\frac{\varepsilon_Q\Delta\rho}{Q}.
 \label{eq:displaced-reduced-midpoint}
\end{equation}
Then $\kappa$ is a positive integer and there are unique coprime positive integers $\alpha,\beta$ satisfying
\begin{equation}
 Q=\alpha u+\beta v,\qquad
 \frac\Sigma{c_0}=\alpha s+\beta t.
 \label{eq:displaced-primitive-direction}
\end{equation}
They obey
\begin{align}
 c_0&>\frac{2\kappa}{\varepsilon_Q}uv,
 \label{eq:displaced-prefix-separation}\\
 M_0\Sigma-2\rho_0M&=c_0n,\qquad n=\alpha-2\beta.
 \label{eq:displaced-integer-defect}
\end{align}
For the distinct central words under consideration, $n\ne0$.
\end{lemma}

\begin{proof}
The determinant-one Gram equations imply $\rho_X^2\equiv\rho_Y^2\equiv-1\pmod M$, hence $M\mid\Sigma\Delta\rho$. Since $\gcd(\Sigma/c_0,Q)=1$, this gives $Q\mid2\Delta\rho$. Equivalently $Q/\varepsilon_Q\mid\Delta\rho$, proving positive integrality of $\kappa$ without factoring $M$ or assuming it odd.

The strict midpoint interval in \eqref{eq:displaced-farey-interval} gives
$$
 \alpha=\frac{v\Sigma-2Mt}{c_0}>0,\qquad
 \beta=\frac{2Ms-u\Sigma}{c_0}>0.
$$
They are integers because $c_0$ divides both numerators. Inverting the unimodular frame proves \eqref{eq:displaced-primitive-direction}, including uniqueness and coprimality. Moreover
$$
 \frac{\Delta\rho}{M}=\frac{2\kappa}{\varepsilon_Q c_0}<\frac1{uv},
$$
which is \eqref{eq:displaced-prefix-separation}. Substituting \eqref{eq:displaced-primitive-direction} into the left side of \eqref{eq:displaced-integer-defect}, and using \eqref{eq:displaced-prefix-frame}, gives $c_0(\alpha-2\beta)$.

Finally, $n=0$ is precisely $\Sigma/(2M)=\rho_0/M_0$, the ancestor-midpoint case. \Cref{cor:no-central-midpoint} excludes that case for distinct central words. This last exclusion uses centrality; it is not asserted for arbitrary palindromic pairs.
\end{proof}

Put $r=\alpha/\beta$. Since $Q=\beta M_0\{1+r_0(r-2)\}$, the integer-defect identity gives
\begin{equation}
 \xi=\frac{2(r-2)}{1+r_0(r-2)}=2\eta(r),\qquad
 2\eta(r)=\eta(t_X)+\eta(t_Y).
 \label{eq:displaced-mean}
\end{equation}
Thus the primitive direction is a nonlinear mean of the two actual column slopes, rather than a vector of letter counts. The root order also has a direct column interpretation. From $M=(u,v)v_i$, $\rho_i=(s,t)v_i$ and the frame identities, one obtains
$$
 \eta(t_i)=\frac{M_0^2\rho_i}{M}-M_0\rho_0.
$$
An $a$ preimage has $t_i<2$ and hence $\rho_i/M<r_0$; a $b$ preimage has $t_i>2$ and hence $\rho_i/M>r_0$. Therefore the $a$ branch is precisely the smaller-root branch, consistently with the ordering above.

\subsection{Cancellation of the two parent cylinders}

Let $U=\nu(a)$ and $V=\nu(b)$, with $\tr U=3p$, $\tr V=3q$. We recall the notation of the decoder, but write its cusp centers as $u_*,v_*$ to distinguish them from the prefix entries in \eqref{eq:displaced-prefix-frame}. Put
$$
 a_*=p^{-2},\quad b_*=q^{-2},\quad
 L=u_*+v_*-3=\frac{M_0}{pq},\quad
 \delta_*=(a_*-b_*)/L.
$$
The standard-basis formulas in \cref{sec:decoder} give
\begin{equation}
 p\geqslant1,\quad q\geqslant2,\quad
 \frac52\leqslant L<3,\quad
 a_*+b_*=L(3-L),\quad 2r_0=u_*-v_*+\delta_*.
 \label{eq:displaced-parent-parameters}
\end{equation}
The preimages in \eqref{eq:displaced-genuine-fork} begin in opposite letters. Their cusp slopes $\zeta_i=1/(t_i-2)$ therefore lie in the two cylinders \eqref{eq:decoder-cylinders}. Naming the $a$ branch first, we may write
\begin{align}
 \zeta_X&=-u_*+e_U,&-\frac{5a_*}{12}\leqslant e_U\leqslant\frac{12a_*}{25},\nonumber\\
 \zeta_Y&=v_*+e_V,&-\frac{40b_*}{111}\leqslant e_V\leqslant\frac{b_*}{2}.
 \label{eq:displaced-cylinder-errors}
\end{align}
The centered coordinates $y_i=\zeta_i+r_0$ have the form
$$
 y_X=-\frac{L+3}{2}+\frac{\delta_*}{2}+e_U,
 \qquad y_Y=\frac{L+3}{2}+\frac{\delta_*}{2}+e_V.
$$
The leading centers cancel in their reciprocal sum:
\begin{equation}
 \xi=\frac1{y_X}+\frac1{y_Y}
 =\frac{\delta_*+e_U+e_V}{y_Xy_Y}.
 \label{eq:displaced-parent-cancellation}
\end{equation}
This cancellation explains why both parent labels enter the bound.

The ranges in \eqref{eq:displaced-parent-parameters}--\eqref{eq:displaced-cylinder-errors} give
$$
 -\frac{a_*}{12}-\frac{422b_*}{555}
 \leqslant\delta_*+e_U+e_V
 \leqslant\frac{22a_*}{25}+\frac{b_*}{6}.
$$
Every absolute coefficient here is at most $22/25$. The decoder's cusp orbit has $\zeta_X\leqslant-5/2$ and $\zeta_Y\geqslant12/5$, while $0<r_0<1/2$. Hence $y_X<-2$, $y_Y>12/5$, and
$$
 |\xi|<\frac{11}{60}(a_*+b_*)\leqslant\frac{11}{48}.
$$
Inverting \eqref{eq:displaced-mean} gives $r-2=\xi/(2-r_0\xi)$. Its denominator exceeds $181/96>0$, and therefore
\begin{equation}
 \left|\frac\alpha\beta-2\right|
 <\frac{88}{905}(p^{-2}+q^{-2})
 <\frac{p^{-2}+q^{-2}}{10}\leqslant\frac18.
 \label{eq:displaced-parent-cone}
\end{equation}
No parity condition has entered either the cancellation or the rational direction.

\begin{theorem}[A parent-sensitive bound in both parities]
\label{thm:displaced-ancestor-bound}
For two distinct central words with common Gram label $M$, let $d$ be their maximal common directive and use the data above. Let
$$
 \chi=\frac pq+\frac qp,\qquad n=\alpha-2\beta\ne0.
$$
Then
\begin{equation}
 \beta>\frac{10|n|p^2q^2}{p^2+q^2}
 =\frac{10|n|M_0}{\chi L},
 \qquad
 M_0^4<\frac{8\varepsilon_Q\chi(7+5\sqrt2)}{25\kappa|n|}\,M.
 \label{eq:displaced-quartic-bound}
\end{equation}
In particular, bounded parent imbalance gives a fourth-power restriction on the common-ancestor label.
\end{theorem}

\begin{proof}
Since $|r-2|=|n|/\beta$, the first bound follows from \eqref{eq:displaced-parent-cone}. Set $\gamma=1+r_0(r-2)$. The same estimate gives $\gamma>15/16$, whence
$$
 Q=\beta\gamma M_0>
 \frac{10\gamma|n|}{\chi L}M_0^2
 >\frac{25|n|}{8\chi}M_0^2.
$$
The prefix ratio $u/v$ lies in $(\phi,1+\sqrt2)$, where $\phi=(1+\sqrt5)/2$: start at $2$, and apply the two increasing Cohn maps, which preserve this interval. The function $x/(2x+1)^2$ decreases for $x>1/2$. Consequently
$$
 \frac{uv}{M_0^2}=\frac{u/v}{(2u/v+1)^2}>5\sqrt2-7.
$$
Combining this with \eqref{eq:displaced-prefix-separation} and $M=c_0Q/2$ gives
\begin{equation}
 M_0^4<
 \frac{\varepsilon_Q\chi L(7+5\sqrt2)}{10\kappa\gamma|n|}\,M
 <\frac{8\varepsilon_Q\chi(7+5\sqrt2)}{25\kappa|n|}\,M,
 \label{eq:displaced-quartic-refined}
\end{equation}
using $(5\sqrt2-7)^{-1}=7+5\sqrt2$. This also records the stronger data-dependent constant before its uniform simplification.
\end{proof}

The imbalance cannot be dropped. Along a constant $a$ directive, one parent label stays equal to one, while $M_0=Lp q$ and $5/2\leqslant L<3$. Thus $\chi=q+1/q$ grows proportionally to $M_0$; the fourth-power bound then has the earlier cubic scale. The theorem distinguishes balanced parent scales from long one-sided directives without excluding either.

There is also a direct trace meaning for the integer defect. Put $H_0=(\rho_0^2+1)/M_0$, the lower-right entry of the ancestor Gram matrix. Since $(PGP)^{-1}H_i=P^{-1}\nu(z_i)P$, we have
$$
 \tr\nu(z_i)=\tr\bigl((PGP)^{-1}H_i\bigr)
 =H_0M+M_0\frac{\rho_i^2+1}{M}-2\rho_0\rho_i.
$$
Completing a square gives
$$
 \tr\nu(z_i)=\frac M{M_0}+\frac{M_0}M
 +\frac{M_0}M\left(\rho_i-\frac{M\rho_0}{M_0}\right)^2.
$$
Subtracting yields
$$
 \tr\nu(z_Y)-\tr\nu(z_X)
 =M_0\Delta\rho\,\delta
 =\frac{2\kappa n}{\varepsilon_Q}.
$$
Thus the relative words do not themselves have equal traces when $n\ne0$. Neither cancellation of the ancestor nor a projective moving-center argument supplies a new equal-label descent. The remaining problem is to combine literal centrality with the exact weighted-scale equality; the inequalities proved here are necessary conditions only.

\section{Twins, ladders and what remains beyond the reflection axis}
\label{sec:outlook}

The classification has a simple mechanism. At the ancestor midpoint, the Gram reflection
becomes a conjugacy that fixes the terminal column; the finite word is then recoverable from
a rigid succession of paired blocks, and palindromicity turns the paired blocks into an odd
common divisor of the completed counts. \Cref{prop:any-centre} shows that the constructive
half of this mechanism works around every palindrome. This section asks whether the block
exchange is the only source of equal labels among palindromes, records what an exact
computation says about it, and isolates the step that a proof would need.

\subsection{A conjecture on twins}

Call two distinct palindromes with the same Gram label \emph{twins}. By
\cref{prop:thue-morse}, twins are exactly the pairs of palindromes $av\neq av'$ without the
factor $bb$ for which the lengths $\abs{\operatorname{Pal}(\theta(av))}$ and
$\abs{\operatorname{Pal}(\theta(av'))}$ coincide. Reutenauer and Vuillon proved that the
Frobenius--Markov uniqueness conjecture is equivalent to the injectivity of this length on
the central words $\operatorname{Pal}(av)$~\cite[Theorem~1]{ReutenauerVuillon2017}. A search for
coincidences among all palindromes without $bb$ finds only words built around a palindromic
centre, and it suggests the following conjecture.

\begin{conjecture}[Structure of twins]
\label{conj:structure}
\textup{(i)}~A palindrome has a twin if and only if it is structured around some palindrome.
\textup{(ii)}~Two palindromes are twins if and only if a chain of block exchanges, each around
its own centre, leads from one to the other.
\end{conjecture}

Part (ii) implies part (i), and it is the form that \cref{rem:quadruple} imposes, since the
twins of a palindrome need not be its block exchanges. In both parts the \emph{if} direction is \cref{prop:any-centre}. Either part would
imply the Frobenius--Markov uniqueness conjecture. Indeed, two distinct Markov occurrences with
the same Markov number are twin central words, while a structured palindrome has completed
counts divisible by $2k+1\geqslant3$ and is therefore not central.

An exact enumeration by computer of all $16\,743\,538$ palindromes whose label does not exceed
$10^{26}$ supports the conjecture in its stronger form. Exactly $1\,719$ of these labels are
carried by more than one palindrome, $1\,695$ by two and $24$ by four. Every palindrome
involved is structured, every class of twins is connected by exchanges, and every exchange
pair has its root midpoint at its centre. Every class of four is a square of
\cref{cor:nested-square}, so that every pair of twins whose midpoint is shifted off its centres is
a diagonal of such a square, and the two palindromes of every class have the same numbers of
letters $a$ and $b$. The enumeration is part of the supplementary material.

\begin{figure}[htbp]
\centering
% Author: Dr. Denys Dutykh (Khalifa University of Science and Technology, Abu Dhabi, UAE)
% Generated by supplement/code/make_figures.py from exact data; do not edit by hand.

\begin{tikzpicture}[font=\footnotesize]
\draw[ink, line width=0.5pt] (0,7) -- (0,0) -- (12.4,0);
\draw[ink, line width=0.5pt] (0,1.489) -- (-0.08,1.489);
\node[ink, left] at (-0.08,1.489) {$10^{3}$};
\draw[ink, line width=0.5pt] (0,3.277) -- (-0.08,3.277);
\node[ink, left] at (-0.08,3.277) {$10^{6}$};
\draw[ink, line width=0.5pt] (0,5.064) -- (-0.08,5.064);
\node[ink, left] at (-0.08,5.064) {$10^{9}$};
\draw[ink, line width=0.5pt] (0,6.851) -- (-0.08,6.851);
\node[ink, left] at (-0.08,6.851) {$10^{12}$};
\draw[ink, line width=0.5pt] (1.167,0) -- (1.167,-0.08);
\node[ink, below] at (1.167,-0.08) {$0.385$};
\draw[ink, line width=0.5pt] (3.089,0) -- (3.089,-0.08);
\node[ink, below] at (3.089,-0.08) {$0.390$};
\draw[ink, line width=0.5pt] (5.012,0) -- (5.012,-0.08);
\node[ink, below] at (5.012,-0.08) {$0.395$};
\draw[ink, line width=0.5pt] (6.935,0) -- (6.935,-0.08);
\node[ink, below] at (6.935,-0.08) {$0.400$};
\draw[ink, line width=0.5pt] (8.857,0) -- (8.857,-0.08);
\node[ink, below] at (8.857,-0.08) {$0.405$};
\draw[ink, line width=0.5pt] (10.78,0) -- (10.78,-0.08);
\node[ink, below] at (10.78,-0.08) {$0.410$};
\node[ink, below] at (6.2,-0.48) {$\rho(z)/M(z)$, the frequency of $b$ in $\Omega_z$};
\node[ink, rotate=90, above] at (-0.75,3.5) {Gram label $M(z)$};
\fill[slate] (12.4,6.082) circle (0.28pt);
\fill[slate] (12.4,6.33) circle (0.28pt);
\fill[slate] (12.4,6.533) circle (0.28pt);
\fill[slate] (12.4,6.579) circle (0.28pt);
\fill[slate] (12.4,6.828) circle (0.28pt);
\fill[slate] (12.4,5.169) circle (0.28pt);
\fill[slate] (12.4,5.418) circle (0.28pt);
\fill[slate] (12.4,5.621) circle (0.28pt);
\fill[slate] (12.4,6.077) circle (0.28pt);
\fill[slate] (12.4,6.32) circle (0.28pt);
\fill[slate] (12.4,6.533) circle (0.28pt);
\fill[slate] (12.4,6.568) circle (0.28pt);
\fill[slate] (12.4,6.817) circle (0.28pt);
\fill[slate] (12.4,5.667) circle (0.28pt);
\fill[slate] (12.4,5.916) circle (0.28pt);
\fill[slate] (12.4,6.117) circle (0.28pt);
\fill[slate] (12.4,6.573) circle (0.28pt);
\fill[slate] (12.4,6.817) circle (0.28pt);
\fill[slate] (12.4,6.165) circle (0.28pt);
\fill[slate] (12.4,6.414) circle (0.28pt);
\fill[slate] (12.4,6.615) circle (0.28pt);
\fill[slate] (12.4,6.663) circle (0.28pt);
\fill[slate] (12.4,4.257) circle (0.28pt);
\fill[slate] (12.4,4.506) circle (0.28pt);
\fill[slate] (12.4,4.709) circle (0.28pt);
\fill[slate] (12.4,5.165) circle (0.28pt);
\fill[slate] (12.4,5.408) circle (0.28pt);
\fill[slate] (12.4,6.32) circle (0.28pt);
\fill[slate] (12.4,6.568) circle (0.28pt);
\fill[slate] (12.4,6.817) circle (0.28pt);
\fill[slate] (12.4,5.656) circle (0.28pt);
\fill[slate] (12.4,5.905) circle (0.28pt);
\fill[slate] (12.4,6.108) circle (0.28pt);
\fill[slate] (12.4,6.563) circle (0.28pt);
\fill[slate] (12.4,6.807) circle (0.28pt);
\fill[slate] (12.4,6.154) circle (0.28pt);
\fill[slate] (12.4,6.403) circle (0.28pt);
\fill[slate] (12.4,6.604) circle (0.28pt);
\fill[slate] (12.4,6.652) circle (0.28pt);
\fill[slate] (12.4,4.755) circle (0.28pt);
\fill[slate] (12.4,5.004) circle (0.28pt);
\fill[slate] (12.4,5.205) circle (0.28pt);
\fill[slate] (12.4,5.661) circle (0.28pt);
\fill[slate] (12.4,5.904) circle (0.28pt);
\fill[slate] (12.4,6.117) circle (0.28pt);
\fill[slate] (12.4,6.573) circle (0.28pt);
\fill[slate] (12.4,6.816) circle (0.28pt);
\fill[slate] (12.4,6.153) circle (0.28pt);
\fill[slate] (12.4,6.402) circle (0.28pt);
\fill[slate] (12.4,6.604) circle (0.28pt);
\fill[slate] (12.4,6.651) circle (0.28pt);
\fill[slate] (12.4,5.253) circle (0.28pt);
\fill[slate] (12.4,5.502) circle (0.28pt);
\fill[slate] (12.4,5.703) circle (0.28pt);
\fill[slate] (12.4,6.159) circle (0.28pt);
\fill[slate] (12.4,6.402) circle (0.28pt);
\fill[slate] (12.4,6.615) circle (0.28pt);
\fill[slate] (12.4,6.65) circle (0.28pt);
\fill[slate] (12.4,5.751) circle (0.28pt);
\fill[slate] (12.4,6) circle (0.28pt);
\fill[slate] (12.4,6.201) circle (0.28pt);
\fill[slate] (12.4,6.657) circle (0.28pt);
\fill[slate] (12.4,6.249) circle (0.28pt);
\fill[slate] (12.4,6.498) circle (0.28pt);
\fill[slate] (12.4,6.699) circle (0.28pt);
\fill[slate] (12.4,6.747) circle (0.28pt);
\fill[slate] (12.4,3.345) circle (0.28pt);
\fill[slate] (12.4,3.594) circle (0.28pt);
\fill[slate] (12.4,3.796) circle (0.28pt);
\fill[slate] (12.4,4.496) circle (0.28pt);
\fill[slate] (12.4,5.408) circle (0.28pt);
\fill[slate] (12.4,6.077) circle (0.28pt);
\fill[slate] (12.4,6.533) circle (0.28pt);
\fill[slate] (12.4,6.568) circle (0.28pt);
\fill[slate] (12.4,6.817) circle (0.28pt);
\fill[slate] (12.4,5.656) circle (0.28pt);
\fill[slate] (12.4,5.905) circle (0.28pt);
\fill[slate] (12.4,6.107) circle (0.28pt);
\fill[slate] (12.4,6.563) circle (0.28pt);
\fill[slate] (12.4,6.807) circle (0.28pt);
\fill[slate] (12.4,6.154) circle (0.28pt);
\fill[slate] (12.4,6.403) circle (0.28pt);
\fill[slate] (12.4,6.604) circle (0.28pt);
\fill[slate] (12.4,6.652) circle (0.28pt);
\fill[slate] (12.4,4.744) circle (0.28pt);
\fill[slate] (12.4,4.993) circle (0.28pt);
\fill[slate] (12.4,5.195) circle (0.28pt);
\fill[slate] (12.4,5.651) circle (0.28pt);
\fill[slate] (12.4,5.895) circle (0.28pt);
\fill[slate] (12.4,6.107) circle (0.28pt);
\fill[slate] (12.4,6.563) circle (0.28pt);
\fill[slate] (12.4,6.807) circle (0.28pt);
\fill[slate] (12.4,6.143) circle (0.28pt);
\fill[slate] (12.4,6.392) circle (0.28pt);
\fill[slate] (12.4,6.594) circle (0.28pt);
\fill[slate] (12.4,6.641) circle (0.28pt);
\fill[slate] (12.4,5.242) circle (0.28pt);
\fill[slate] (12.4,5.491) circle (0.28pt);
\fill[slate] (12.4,5.692) circle (0.28pt);
\fill[slate] (12.4,6.148) circle (0.28pt);
\fill[slate] (12.4,6.391) circle (0.28pt);
\fill[slate] (12.4,6.604) circle (0.28pt);
\fill[slate] (12.4,6.64) circle (0.28pt);
\fill[slate] (12.4,5.74) circle (0.28pt);
\fill[slate] (12.4,5.989) circle (0.28pt);
\fill[slate] (12.4,6.19) circle (0.28pt);
\fill[slate] (12.4,6.645) circle (0.28pt);
\fill[slate] (12.4,6.238) circle (0.28pt);
\fill[slate] (12.4,6.487) circle (0.28pt);
\fill[slate] (12.4,6.688) circle (0.28pt);
\fill[slate] (12.4,6.736) circle (0.28pt);
\fill[slate] (12.4,3.843) circle (0.28pt);
\fill[slate] (12.4,4.092) circle (0.28pt);
\fill[slate] (12.4,4.293) circle (0.28pt);
\fill[slate] (12.4,4.749) circle (0.28pt);
\fill[slate] (12.4,4.992) circle (0.28pt);
\fill[slate] (12.4,5.205) circle (0.28pt);
\fill[slate] (12.4,5.661) circle (0.28pt);
\fill[slate] (12.4,5.904) circle (0.28pt);
\fill[slate] (12.4,6.117) circle (0.28pt);
\fill[slate] (12.4,6.573) circle (0.28pt);
\fill[slate] (12.4,6.816) circle (0.28pt);
\fill[slate] (12.4,6.152) circle (0.28pt);
\fill[slate] (12.4,6.401) circle (0.28pt);
\fill[slate] (12.4,6.604) circle (0.28pt);
\fill[slate] (12.4,6.65) circle (0.28pt);
\fill[slate] (12.4,5.24) circle (0.28pt);
\fill[slate] (12.4,5.489) circle (0.28pt);
\fill[slate] (12.4,5.692) circle (0.28pt);
\fill[slate] (12.4,6.148) circle (0.28pt);
\fill[slate] (12.4,6.391) circle (0.28pt);
\fill[slate] (12.4,6.604) circle (0.28pt);
\fill[slate] (12.4,6.639) circle (0.28pt);
\fill[slate] (12.4,5.738) circle (0.28pt);
\fill[slate] (12.4,5.987) circle (0.28pt);
\fill[slate] (12.4,6.188) circle (0.28pt);
\fill[slate] (12.4,6.644) circle (0.28pt);
\fill[slate] (12.4,6.236) circle (0.28pt);
\fill[slate] (12.4,6.485) circle (0.28pt);
\fill[slate] (12.4,6.686) circle (0.28pt);
\fill[slate] (12.4,6.734) circle (0.28pt);
\fill[slate] (12.4,4.341) circle (0.28pt);
\fill[slate] (12.4,4.59) circle (0.28pt);
\fill[slate] (12.4,4.791) circle (0.28pt);
\fill[slate] (12.4,5.246) circle (0.28pt);
\fill[slate] (12.4,5.49) circle (0.28pt);
\fill[slate] (12.4,5.703) circle (0.28pt);
\fill[slate] (12.4,6.159) circle (0.28pt);
\fill[slate] (12.4,6.402) circle (0.28pt);
\fill[slate] (12.4,6.615) circle (0.28pt);
\fill[slate] (12.4,6.65) circle (0.28pt);
\fill[slate] (12.4,5.738) circle (0.28pt);
\fill[slate] (12.4,5.987) circle (0.28pt);
\fill[slate] (12.4,6.19) circle (0.28pt);
\fill[slate] (12.4,6.645) circle (0.28pt);
\fill[slate] (12.4,6.236) circle (0.28pt);
\fill[slate] (12.4,6.485) circle (0.28pt);
\fill[slate] (12.4,6.686) circle (0.28pt);
\fill[slate] (12.4,6.734) circle (0.28pt);
\fill[slate] (12.4,4.839) circle (0.28pt);
\fill[slate] (12.4,5.088) circle (0.28pt);
\fill[slate] (12.4,5.289) circle (0.28pt);
\fill[slate] (12.4,5.744) circle (0.28pt);
\fill[slate] (12.4,5.988) circle (0.28pt);
\fill[slate] (12.4,6.2) circle (0.28pt);
\fill[slate] (12.4,6.657) circle (0.28pt);
\fill[slate] (12.4,6.236) circle (0.28pt);
\fill[slate] (12.4,6.485) circle (0.28pt);
\fill[slate] (12.4,6.688) circle (0.28pt);
\fill[slate] (12.4,6.734) circle (0.28pt);
\fill[slate] (12.4,5.337) circle (0.28pt);
\fill[slate] (12.4,5.586) circle (0.28pt);
\fill[slate] (12.4,5.787) circle (0.28pt);
\fill[slate] (12.4,6.242) circle (0.28pt);
\fill[slate] (12.4,6.486) circle (0.28pt);
\fill[slate] (12.4,6.699) circle (0.28pt);
\fill[slate] (12.4,6.734) circle (0.28pt);
\fill[slate] (12.4,5.835) circle (0.28pt);
\fill[slate] (12.4,6.084) circle (0.28pt);
\fill[slate] (12.4,6.285) circle (0.28pt);
\fill[slate] (12.4,6.74) circle (0.28pt);
\fill[slate] (12.4,6.333) circle (0.28pt);
\fill[slate] (12.4,6.582) circle (0.28pt);
\fill[slate] (12.4,6.783) circle (0.28pt);
\fill[slate] (12.4,6.831) circle (0.28pt);
\fill[slate] (12.389,2.433) circle (0.28pt);
\fill[slate] (12.389,2.682) circle (0.28pt);
\fill[slate] (12.391,3.584) circle (0.28pt);
\fill[slate] (12.391,4.708) circle (0.28pt);
\fill[slate] (12.391,5.165) circle (0.28pt);
\fill[slate] (12.391,5.621) circle (0.28pt);
\fill[slate] (12.391,6.077) circle (0.28pt);
\fill[slate] (12.391,6.32) circle (0.28pt);
\fill[slate] (12.391,6.533) circle (0.28pt);
\fill[slate] (12.391,6.568) circle (0.28pt);
\fill[slate] (12.391,6.817) circle (0.28pt);
\fill[slate] (12.391,5.656) circle (0.28pt);
\fill[slate] (12.391,5.905) circle (0.28pt);
\fill[slate] (12.391,6.107) circle (0.28pt);
\fill[slate] (12.391,6.807) circle (0.28pt);
\fill[slate] (12.391,6.154) circle (0.28pt);
\fill[slate] (12.391,6.403) circle (0.28pt);
\fill[slate] (12.391,6.604) circle (0.28pt);
\fill[slate] (12.391,6.652) circle (0.28pt);
\fill[slate] (12.391,4.744) circle (0.28pt);
\fill[slate] (12.391,4.993) circle (0.28pt);
\fill[slate] (12.391,5.895) circle (0.28pt);
\fill[slate] (12.391,6.563) circle (0.28pt);
\fill[slate] (12.391,6.143) circle (0.28pt);
\fill[slate] (12.391,6.392) circle (0.28pt);
\fill[slate] (12.391,6.594) circle (0.28pt);
\fill[slate] (12.391,6.641) circle (0.28pt);
\fill[slate] (12.391,5.242) circle (0.28pt);
\fill[slate] (12.391,5.491) circle (0.28pt);
\fill[slate] (12.391,5.692) circle (0.28pt);
\fill[slate] (12.391,6.147) circle (0.28pt);
\fill[slate] (12.391,6.391) circle (0.28pt);
\fill[slate] (12.391,6.603) circle (0.28pt);
\fill[slate] (12.391,6.639) circle (0.28pt);
\fill[slate] (12.391,5.74) circle (0.28pt);
\fill[slate] (12.391,5.989) circle (0.28pt);
\fill[slate] (12.391,6.189) circle (0.28pt);
\fill[slate] (12.391,6.645) circle (0.28pt);
\fill[slate] (12.391,6.238) circle (0.28pt);
\fill[slate] (12.391,6.487) circle (0.28pt);
\fill[slate] (12.391,6.687) circle (0.28pt);
\fill[slate] (12.391,6.736) circle (0.28pt);
\fill[slate] (12.391,3.832) circle (0.28pt);
\fill[slate] (12.391,4.081) circle (0.28pt);
\fill[slate] (12.391,4.283) circle (0.28pt);
\fill[slate] (12.391,4.739) circle (0.28pt);
\fill[slate] (12.391,4.983) circle (0.28pt);
\fill[slate] (12.391,5.651) circle (0.28pt);
\fill[slate] (12.391,5.895) circle (0.28pt);
\fill[slate] (12.391,6.107) circle (0.28pt);
\fill[slate] (12.391,6.563) circle (0.28pt);
\fill[slate] (12.391,6.807) circle (0.28pt);
\fill[slate] (12.391,6.143) circle (0.28pt);
\fill[slate] (12.391,6.392) circle (0.28pt);
\fill[slate] (12.391,6.594) circle (0.28pt);
\fill[slate] (12.391,6.641) circle (0.28pt);
\fill[slate] (12.391,5.231) circle (0.28pt);
\fill[slate] (12.391,5.48) circle (0.28pt);
\fill[slate] (12.391,5.682) circle (0.28pt);
\fill[slate] (12.391,6.138) circle (0.28pt);
\fill[slate] (12.391,6.382) circle (0.28pt);
\fill[slate] (12.391,6.594) circle (0.28pt);
\fill[slate] (12.391,6.63) circle (0.28pt);
\fill[slate] (12.391,5.729) circle (0.28pt);
\fill[slate] (12.391,5.978) circle (0.28pt);
\fill[slate] (12.391,6.179) circle (0.28pt);
\fill[slate] (12.391,6.635) circle (0.28pt);
\fill[slate] (12.391,6.227) circle (0.28pt);
\fill[slate] (12.391,6.476) circle (0.28pt);
\fill[slate] (12.391,6.676) circle (0.28pt);
\fill[slate] (12.391,6.725) circle (0.28pt);
\fill[slate] (12.391,4.33) circle (0.28pt);
\fill[slate] (12.391,4.579) circle (0.28pt);
\fill[slate] (12.391,4.78) circle (0.28pt);
\fill[slate] (12.391,5.236) circle (0.28pt);
\fill[slate] (12.391,5.479) circle (0.28pt);
\fill[slate] (12.391,5.692) circle (0.28pt);
\fill[slate] (12.391,6.148) circle (0.28pt);
\fill[slate] (12.391,6.391) circle (0.28pt);
\fill[slate] (12.391,6.604) circle (0.28pt);
\fill[slate] (12.391,6.639) circle (0.28pt);
\fill[slate] (12.391,5.727) circle (0.28pt);
\fill[slate] (12.391,5.976) circle (0.28pt);
\fill[slate] (12.391,6.179) circle (0.28pt);
\fill[slate] (12.391,6.635) circle (0.28pt);
\fill[slate] (12.391,6.225) circle (0.28pt);
\fill[slate] (12.391,6.474) circle (0.28pt);
\fill[slate] (12.391,6.675) circle (0.28pt);
\fill[slate] (12.391,6.723) circle (0.28pt);
\fill[slate] (12.391,4.828) circle (0.28pt);
\fill[slate] (12.391,5.077) circle (0.28pt);
\fill[slate] (12.391,5.277) circle (0.28pt);
\fill[slate] (12.391,5.733) circle (0.28pt);
\fill[slate] (12.391,5.977) circle (0.28pt);
\fill[slate] (12.391,6.189) circle (0.28pt);
\fill[slate] (12.391,6.645) circle (0.28pt);
\fill[slate] (12.391,6.225) circle (0.28pt);
\fill[slate] (12.391,6.474) circle (0.28pt);
\fill[slate] (12.391,6.676) circle (0.28pt);
\fill[slate] (12.391,6.723) circle (0.28pt);
\fill[slate] (12.391,5.326) circle (0.28pt);
\fill[slate] (12.391,5.575) circle (0.28pt);
\fill[slate] (12.391,5.775) circle (0.28pt);
\fill[slate] (12.391,6.231) circle (0.28pt);
\fill[slate] (12.391,6.475) circle (0.28pt);
\fill[slate] (12.391,6.687) circle (0.28pt);
\fill[slate] (12.391,6.723) circle (0.28pt);
\fill[slate] (12.391,5.824) circle (0.28pt);
\fill[slate] (12.391,6.073) circle (0.28pt);
\fill[slate] (12.391,6.273) circle (0.28pt);
\fill[slate] (12.391,6.729) circle (0.28pt);
\fill[slate] (12.391,6.322) circle (0.28pt);
\fill[slate] (12.391,6.571) circle (0.28pt);
\fill[slate] (12.391,6.771) circle (0.28pt);
\fill[slate] (12.391,6.82) circle (0.28pt);
\fill[slate] (12.389,2.931) circle (0.28pt);
\fill[slate] (12.389,3.18) circle (0.28pt);
\fill[slate] (12.389,3.381) circle (0.28pt);
\fill[slate] (12.389,3.837) circle (0.28pt);
\fill[slate] (12.389,4.08) circle (0.28pt);
\fill[slate] (12.389,4.293) circle (0.28pt);
\fill[slate] (12.389,4.749) circle (0.28pt);
\fill[slate] (12.389,4.992) circle (0.28pt);
\fill[slate] (12.389,5.205) circle (0.28pt);
\fill[slate] (12.389,5.661) circle (0.28pt);
\fill[slate] (12.389,5.904) circle (0.28pt);
\fill[slate] (12.389,6.117) circle (0.28pt);
\fill[slate] (12.389,6.573) circle (0.28pt);
\fill[slate] (12.389,6.816) circle (0.28pt);
\fill[slate] (12.389,6.152) circle (0.28pt);
\fill[slate] (12.389,6.401) circle (0.28pt);
\fill[slate] (12.389,6.604) circle (0.28pt);
\fill[slate] (12.389,6.65) circle (0.28pt);
\fill[slate] (12.389,5.24) circle (0.28pt);
\fill[slate] (12.389,5.489) circle (0.28pt);
\fill[slate] (12.389,5.691) circle (0.28pt);
\fill[slate] (12.389,6.147) circle (0.28pt);
\fill[slate] (12.389,6.391) circle (0.28pt);
\fill[slate] (12.389,6.603) circle (0.28pt);
\fill[slate] (12.389,6.639) circle (0.28pt);
\fill[slate] (12.389,5.738) circle (0.28pt);
\fill[slate] (12.389,5.987) circle (0.28pt);
\fill[slate] (12.389,6.188) circle (0.28pt);
\fill[slate] (12.389,6.644) circle (0.28pt);
\fill[slate] (12.389,6.236) circle (0.28pt);
\fill[slate] (12.389,6.485) circle (0.28pt);
\fill[slate] (12.389,6.686) circle (0.28pt);
\fill[slate] (12.389,6.734) circle (0.28pt);
\fill[slate] (12.389,4.328) circle (0.28pt);
\fill[slate] (12.389,4.577) circle (0.28pt);
\fill[slate] (12.389,4.78) circle (0.28pt);
\fill[slate] (12.389,5.236) circle (0.28pt);
\fill[slate] (12.389,5.479) circle (0.28pt);
\fill[slate] (12.389,5.692) circle (0.28pt);
\fill[slate] (12.389,6.148) circle (0.28pt);
\fill[slate] (12.389,6.391) circle (0.28pt);
\fill[slate] (12.389,6.604) circle (0.28pt);
\fill[slate] (12.389,6.639) circle (0.28pt);
\fill[slate] (12.389,5.727) circle (0.28pt);
\fill[slate] (12.389,5.976) circle (0.28pt);
\fill[slate] (12.389,6.179) circle (0.28pt);
\fill[slate] (12.389,6.635) circle (0.28pt);
\fill[slate] (12.389,6.225) circle (0.28pt);
\fill[slate] (12.389,6.474) circle (0.28pt);
\fill[slate] (12.389,6.675) circle (0.28pt);
\fill[slate] (12.389,6.723) circle (0.28pt);
\fill[slate] (12.389,4.826) circle (0.28pt);
\fill[slate] (12.389,5.075) circle (0.28pt);
\fill[slate] (12.389,5.276) circle (0.28pt);
\fill[slate] (12.389,5.732) circle (0.28pt);
\fill[slate] (12.389,5.976) circle (0.28pt);
\fill[slate] (12.389,6.188) circle (0.28pt);
\fill[slate] (12.389,6.644) circle (0.28pt);
\fill[slate] (12.389,6.224) circle (0.28pt);
\fill[slate] (12.389,6.473) circle (0.28pt);
\fill[slate] (12.389,6.675) circle (0.28pt);
\fill[slate] (12.389,6.722) circle (0.28pt);
\fill[slate] (12.389,5.324) circle (0.28pt);
\fill[slate] (12.389,5.573) circle (0.28pt);
\fill[slate] (12.389,5.774) circle (0.28pt);
\fill[slate] (12.389,6.23) circle (0.28pt);
\fill[slate] (12.389,6.473) circle (0.28pt);
\fill[slate] (12.389,6.686) circle (0.28pt);
\fill[slate] (12.389,6.722) circle (0.28pt);
\fill[slate] (12.389,5.822) circle (0.28pt);
\fill[slate] (12.389,6.071) circle (0.28pt);
\fill[slate] (12.389,6.272) circle (0.28pt);
\fill[slate] (12.389,6.728) circle (0.28pt);
\fill[slate] (12.389,6.32) circle (0.28pt);
\fill[slate] (12.389,6.569) circle (0.28pt);
\fill[slate] (12.389,6.77) circle (0.28pt);
\fill[slate] (12.389,6.818) circle (0.28pt);
\fill[slate] (12.389,3.429) circle (0.28pt);
\fill[slate] (12.389,3.678) circle (0.28pt);
\fill[slate] (12.389,3.878) circle (0.28pt);
\fill[slate] (12.389,4.334) circle (0.28pt);
\fill[slate] (12.389,4.578) circle (0.28pt);
\fill[slate] (12.389,4.79) circle (0.28pt);
\fill[slate] (12.389,5.246) circle (0.28pt);
\fill[slate] (12.389,5.49) circle (0.28pt);
\fill[slate] (12.389,5.703) circle (0.28pt);
\fill[slate] (12.389,6.159) circle (0.28pt);
\fill[slate] (12.389,6.402) circle (0.28pt);
\fill[slate] (12.389,6.615) circle (0.28pt);
\fill[slate] (12.389,6.65) circle (0.28pt);
\fill[slate] (12.389,5.738) circle (0.28pt);
\fill[slate] (12.389,5.987) circle (0.28pt);
\fill[slate] (12.389,6.189) circle (0.28pt);
\fill[slate] (12.389,6.645) circle (0.28pt);
\fill[slate] (12.389,6.236) circle (0.28pt);
\fill[slate] (12.389,6.485) circle (0.28pt);
\fill[slate] (12.389,6.686) circle (0.28pt);
\fill[slate] (12.389,6.734) circle (0.28pt);
\fill[slate] (12.389,4.826) circle (0.28pt);
\fill[slate] (12.389,5.075) circle (0.28pt);
\fill[slate] (12.389,5.277) circle (0.28pt);
\fill[slate] (12.389,5.733) circle (0.28pt);
\fill[slate] (12.389,5.977) circle (0.28pt);
\fill[slate] (12.389,6.189) circle (0.28pt);
\fill[slate] (12.389,6.645) circle (0.28pt);
\fill[slate] (12.389,6.225) circle (0.28pt);
\fill[slate] (12.389,6.474) circle (0.28pt);
\fill[slate] (12.389,6.676) circle (0.28pt);
\fill[slate] (12.389,6.723) circle (0.28pt);
\fill[slate] (12.389,5.324) circle (0.28pt);
\fill[slate] (12.389,5.573) circle (0.28pt);
\fill[slate] (12.389,5.774) circle (0.28pt);
\fill[slate] (12.389,6.23) circle (0.28pt);
\fill[slate] (12.389,6.473) circle (0.28pt);
\fill[slate] (12.389,6.686) circle (0.28pt);
\fill[slate] (12.389,6.722) circle (0.28pt);
\fill[slate] (12.389,5.822) circle (0.28pt);
\fill[slate] (12.389,6.071) circle (0.28pt);
\fill[slate] (12.389,6.272) circle (0.28pt);
\fill[slate] (12.389,6.727) circle (0.28pt);
\fill[slate] (12.389,6.32) circle (0.28pt);
\fill[slate] (12.389,6.569) circle (0.28pt);
\fill[slate] (12.389,6.77) circle (0.28pt);
\fill[slate] (12.389,6.818) circle (0.28pt);
\fill[slate] (12.389,3.927) circle (0.28pt);
\fill[slate] (12.389,4.176) circle (0.28pt);
\fill[slate] (12.389,4.376) circle (0.28pt);
\fill[slate] (12.389,4.832) circle (0.28pt);
\fill[slate] (12.389,5.076) circle (0.28pt);
\fill[slate] (12.389,5.288) circle (0.28pt);
\fill[slate] (12.389,5.744) circle (0.28pt);
\fill[slate] (12.389,5.988) circle (0.28pt);
\fill[slate] (12.389,6.201) circle (0.28pt);
\fill[slate] (12.389,6.657) circle (0.28pt);
\fill[slate] (12.389,6.236) circle (0.28pt);
\fill[slate] (12.389,6.485) circle (0.28pt);
\fill[slate] (12.389,6.687) circle (0.28pt);
\fill[slate] (12.389,6.734) circle (0.28pt);
\fill[slate] (12.389,5.324) circle (0.28pt);
\fill[slate] (12.389,5.573) circle (0.28pt);
\fill[slate] (12.389,5.775) circle (0.28pt);
\fill[slate] (12.389,6.231) circle (0.28pt);
\fill[slate] (12.389,6.475) circle (0.28pt);
\fill[slate] (12.389,6.687) circle (0.28pt);
\fill[slate] (12.389,6.723) circle (0.28pt);
\fill[slate] (12.389,5.822) circle (0.28pt);
\fill[slate] (12.389,6.071) circle (0.28pt);
\fill[slate] (12.389,6.272) circle (0.28pt);
\fill[slate] (12.389,6.728) circle (0.28pt);
\fill[slate] (12.389,6.32) circle (0.28pt);
\fill[slate] (12.389,6.569) circle (0.28pt);
\fill[slate] (12.389,6.77) circle (0.28pt);
\fill[slate] (12.389,6.818) circle (0.28pt);
\fill[slate] (12.389,4.425) circle (0.28pt);
\fill[slate] (12.389,4.674) circle (0.28pt);
\fill[slate] (12.389,4.874) circle (0.28pt);
\fill[slate] (12.389,5.33) circle (0.28pt);
\fill[slate] (12.389,5.574) circle (0.28pt);
\fill[slate] (12.389,5.786) circle (0.28pt);
\fill[slate] (12.389,6.242) circle (0.28pt);
\fill[slate] (12.389,6.486) circle (0.28pt);
\fill[slate] (12.389,6.699) circle (0.28pt);
\fill[slate] (12.389,6.734) circle (0.28pt);
\fill[slate] (12.389,5.822) circle (0.28pt);
\fill[slate] (12.389,6.071) circle (0.28pt);
\fill[slate] (12.389,6.273) circle (0.28pt);
\fill[slate] (12.389,6.729) circle (0.28pt);
\fill[slate] (12.389,6.32) circle (0.28pt);
\fill[slate] (12.389,6.569) circle (0.28pt);
\fill[slate] (12.389,6.77) circle (0.28pt);
\fill[slate] (12.389,6.818) circle (0.28pt);
\fill[slate] (12.389,4.923) circle (0.28pt);
\fill[slate] (12.389,5.172) circle (0.28pt);
\fill[slate] (12.389,5.372) circle (0.28pt);
\fill[slate] (12.389,5.828) circle (0.28pt);
\fill[slate] (12.389,6.072) circle (0.28pt);
\fill[slate] (12.389,6.284) circle (0.28pt);
\fill[slate] (12.389,6.74) circle (0.28pt);
\fill[slate] (12.389,6.32) circle (0.28pt);
\fill[slate] (12.389,6.569) circle (0.28pt);
\fill[slate] (12.389,6.771) circle (0.28pt);
\fill[slate] (12.389,6.818) circle (0.28pt);
\fill[slate] (12.389,5.421) circle (0.28pt);
\fill[slate] (12.389,5.67) circle (0.28pt);
\fill[slate] (12.389,5.87) circle (0.28pt);
\fill[slate] (12.389,6.326) circle (0.28pt);
\fill[slate] (12.389,6.57) circle (0.28pt);
\fill[slate] (12.389,6.782) circle (0.28pt);
\fill[slate] (12.389,6.818) circle (0.28pt);
\fill[slate] (12.389,5.919) circle (0.28pt);
\fill[slate] (12.389,6.168) circle (0.28pt);
\fill[slate] (12.389,6.368) circle (0.28pt);
\fill[slate] (12.389,6.824) circle (0.28pt);
\fill[slate] (12.389,6.417) circle (0.28pt);
\fill[slate] (12.389,6.666) circle (0.28pt);
\fill[slate] (12.032,1.77) circle (0.28pt);
\fill[slate] (12.086,3.34) circle (0.28pt);
\fill[slate] (12.086,3.796) circle (0.28pt);
\fill[slate] (12.086,4.253) circle (0.28pt);
\fill[slate] (12.086,4.496) circle (0.28pt);
\fill[slate] (12.086,4.709) circle (0.28pt);
\fill[slate] (12.086,5.165) circle (0.28pt);
\fill[slate] (12.086,5.408) circle (0.28pt);
\fill[slate] (12.086,5.621) circle (0.28pt);
\fill[slate] (12.086,6.077) circle (0.28pt);
\fill[slate] (12.086,6.32) circle (0.28pt);
\fill[slate] (12.086,6.533) circle (0.28pt);
\fill[slate] (12.086,6.568) circle (0.28pt);
\fill[slate] (12.086,6.817) circle (0.28pt);
\fill[slate] (12.086,5.656) circle (0.28pt);
\fill[slate] (12.086,5.905) circle (0.28pt);
\fill[slate] (12.086,6.107) circle (0.28pt);
\fill[slate] (12.086,6.563) circle (0.28pt);
\fill[slate] (12.086,6.807) circle (0.28pt);
\fill[slate] (12.086,6.154) circle (0.28pt);
\fill[slate] (12.086,6.403) circle (0.28pt);
\fill[slate] (12.086,6.604) circle (0.28pt);
\fill[slate] (12.086,6.652) circle (0.28pt);
\fill[slate] (12.086,4.744) circle (0.28pt);
\fill[slate] (12.086,4.993) circle (0.28pt);
\fill[slate] (12.086,5.651) circle (0.28pt);
\fill[slate] (12.086,5.895) circle (0.28pt);
\fill[slate] (12.086,6.107) circle (0.28pt);
\fill[slate] (12.086,6.563) circle (0.28pt);
\fill[slate] (12.086,6.807) circle (0.28pt);
\fill[slate] (12.086,6.143) circle (0.28pt);
\fill[slate] (12.086,6.392) circle (0.28pt);
\fill[slate] (12.086,6.594) circle (0.28pt);
\fill[slate] (12.086,6.641) circle (0.28pt);
\fill[slate] (12.086,5.242) circle (0.28pt);
\fill[slate] (12.086,5.491) circle (0.28pt);
\fill[slate] (12.086,5.692) circle (0.28pt);
\fill[slate] (12.086,6.148) circle (0.28pt);
\fill[slate] (12.086,6.391) circle (0.28pt);
\fill[slate] (12.086,6.604) circle (0.28pt);
\fill[slate] (12.086,6.639) circle (0.28pt);
\fill[slate] (12.086,5.74) circle (0.28pt);
\fill[slate] (12.086,5.989) circle (0.28pt);
\fill[slate] (12.086,6.189) circle (0.28pt);
\fill[slate] (12.086,6.645) circle (0.28pt);
\fill[slate] (12.086,6.238) circle (0.28pt);
\fill[slate] (12.086,6.487) circle (0.28pt);
\fill[slate] (12.086,6.687) circle (0.28pt);
\fill[slate] (12.086,6.736) circle (0.28pt);
\fill[slate] (12.086,3.832) circle (0.28pt);
\fill[slate] (12.086,4.081) circle (0.28pt);
\fill[slate] (12.086,5.651) circle (0.28pt);
\fill[slate] (12.086,6.107) circle (0.28pt);
\fill[slate] (12.086,6.563) circle (0.28pt);
\fill[slate] (12.086,6.807) circle (0.28pt);
\fill[slate] (12.086,6.143) circle (0.28pt);
\fill[slate] (12.086,6.392) circle (0.28pt);
\fill[slate] (12.086,6.641) circle (0.28pt);
\fill[slate] (12.086,5.231) circle (0.28pt);
\fill[slate] (12.086,5.48) circle (0.28pt);
\fill[slate] (12.086,5.682) circle (0.28pt);
\fill[slate] (12.086,6.382) circle (0.28pt);
\fill[slate] (12.086,6.594) circle (0.28pt);
\fill[slate] (12.086,6.63) circle (0.28pt);
\fill[slate] (12.086,5.729) circle (0.28pt);
\fill[slate] (12.086,5.978) circle (0.28pt);
\fill[slate] (12.086,6.179) circle (0.28pt);
\fill[slate] (12.086,6.634) circle (0.28pt);
\fill[slate] (12.086,6.227) circle (0.28pt);
\fill[slate] (12.086,6.476) circle (0.28pt);
\fill[slate] (12.086,6.676) circle (0.28pt);
\fill[slate] (12.086,6.725) circle (0.28pt);
\fill[slate] (12.086,4.33) circle (0.28pt);
\fill[slate] (12.086,4.579) circle (0.28pt);
\fill[slate] (12.086,4.78) circle (0.28pt);
\fill[slate] (12.086,5.235) circle (0.28pt);
\fill[slate] (12.086,5.479) circle (0.28pt);
\fill[slate] (12.086,5.691) circle (0.28pt);
\fill[slate] (12.086,6.148) circle (0.28pt);
\fill[slate] (12.086,6.391) circle (0.28pt);
\fill[slate] (12.086,6.604) circle (0.28pt);
\fill[slate] (12.086,6.639) circle (0.28pt);
\fill[slate] (12.086,5.727) circle (0.28pt);
\fill[slate] (12.086,5.976) circle (0.28pt);
\fill[slate] (12.086,6.178) circle (0.28pt);
\fill[slate] (12.086,6.634) circle (0.28pt);
\fill[slate] (12.086,6.225) circle (0.28pt);
\fill[slate] (12.086,6.474) circle (0.28pt);
\fill[slate] (12.086,6.675) circle (0.28pt);
\fill[slate] (12.086,6.723) circle (0.28pt);
\fill[slate] (12.086,4.828) circle (0.28pt);
\fill[slate] (12.086,5.077) circle (0.28pt);
\fill[slate] (12.086,5.277) circle (0.28pt);
\fill[slate] (12.086,5.733) circle (0.28pt);
\fill[slate] (12.086,5.977) circle (0.28pt);
\fill[slate] (12.086,6.189) circle (0.28pt);
\fill[slate] (12.086,6.645) circle (0.28pt);
\fill[slate] (12.086,6.225) circle (0.28pt);
\fill[slate] (12.086,6.474) circle (0.28pt);
\fill[slate] (12.086,6.676) circle (0.28pt);
\fill[slate] (12.086,6.723) circle (0.28pt);
\fill[slate] (12.086,5.326) circle (0.28pt);
\fill[slate] (12.086,5.575) circle (0.28pt);
\fill[slate] (12.086,5.775) circle (0.28pt);
\fill[slate] (12.086,6.231) circle (0.28pt);
\fill[slate] (12.086,6.475) circle (0.28pt);
\fill[slate] (12.086,6.687) circle (0.28pt);
\fill[slate] (12.086,6.723) circle (0.28pt);
\fill[slate] (12.086,5.824) circle (0.28pt);
\fill[slate] (12.086,6.073) circle (0.28pt);
\fill[slate] (12.086,6.273) circle (0.28pt);
\fill[slate] (12.086,6.729) circle (0.28pt);
\fill[slate] (12.086,6.322) circle (0.28pt);
\fill[slate] (12.086,6.571) circle (0.28pt);
\fill[slate] (12.086,6.771) circle (0.28pt);
\fill[slate] (12.086,6.82) circle (0.28pt);
\fill[slate] (12.084,3.169) circle (0.28pt);
\fill[slate] (12.085,4.283) circle (0.28pt);
\fill[slate] (12.085,4.739) circle (0.28pt);
\fill[slate] (12.085,5.195) circle (0.28pt);
\fill[slate] (12.085,5.651) circle (0.28pt);
\fill[slate] (12.085,5.895) circle (0.28pt);
\fill[slate] (12.085,6.108) circle (0.28pt);
\fill[slate] (12.085,6.564) circle (0.28pt);
\fill[slate] (12.085,6.807) circle (0.28pt);
\fill[slate] (12.085,6.143) circle (0.28pt);
\fill[slate] (12.085,6.392) circle (0.28pt);
\fill[slate] (12.085,6.594) circle (0.28pt);
\fill[slate] (12.085,6.641) circle (0.28pt);
\fill[slate] (12.085,5.231) circle (0.28pt);
\fill[slate] (12.085,5.48) circle (0.28pt);
\fill[slate] (12.085,6.138) circle (0.28pt);
\fill[slate] (12.085,6.594) circle (0.28pt);
\fill[slate] (12.085,6.63) circle (0.28pt);
\fill[slate] (12.085,5.729) circle (0.28pt);
\fill[slate] (12.085,5.978) circle (0.28pt);
\fill[slate] (12.085,6.178) circle (0.28pt);
\fill[slate] (12.085,6.634) circle (0.28pt);
\fill[slate] (12.085,6.227) circle (0.28pt);
\fill[slate] (12.085,6.476) circle (0.28pt);
\fill[slate] (12.085,6.676) circle (0.28pt);
\fill[slate] (12.085,6.725) circle (0.28pt);
\fill[slate] (12.085,4.568) circle (0.28pt);
\fill[slate] (12.085,5.682) circle (0.28pt);
\fill[slate] (12.085,6.138) circle (0.28pt);
\fill[slate] (12.085,6.382) circle (0.28pt);
\fill[slate] (12.085,6.594) circle (0.28pt);
\fill[slate] (12.085,6.63) circle (0.28pt);
\fill[slate] (12.085,5.967) circle (0.28pt);
\fill[slate] (12.085,6.216) circle (0.28pt);
\fill[slate] (12.085,6.465) circle (0.28pt);
\fill[slate] (12.085,6.666) circle (0.28pt);
\fill[slate] (12.085,6.714) circle (0.28pt);
\fill[slate] (12.085,4.817) circle (0.28pt);
\fill[slate] (12.085,5.066) circle (0.28pt);
\fill[slate] (12.085,5.267) circle (0.28pt);
\fill[slate] (12.085,5.966) circle (0.28pt);
\fill[slate] (12.085,6.179) circle (0.28pt);
\fill[slate] (12.085,6.635) circle (0.28pt);
\fill[slate] (12.085,6.214) circle (0.28pt);
\fill[slate] (12.085,6.463) circle (0.28pt);
\fill[slate] (12.085,6.666) circle (0.28pt);
\fill[slate] (12.085,6.712) circle (0.28pt);
\fill[slate] (12.085,5.315) circle (0.28pt);
\fill[slate] (12.085,5.564) circle (0.28pt);
\fill[slate] (12.085,5.764) circle (0.28pt);
\fill[slate] (12.085,6.22) circle (0.28pt);
\fill[slate] (12.085,6.464) circle (0.28pt);
\fill[slate] (12.085,6.676) circle (0.28pt);
\fill[slate] (12.085,6.712) circle (0.28pt);
\fill[slate] (12.085,5.813) circle (0.28pt);
\fill[slate] (12.085,6.062) circle (0.28pt);
\fill[slate] (12.085,6.262) circle (0.28pt);
\fill[slate] (12.085,6.718) circle (0.28pt);
\fill[slate] (12.085,6.311) circle (0.28pt);
\fill[slate] (12.085,6.56) circle (0.28pt);
\fill[slate] (12.085,6.76) circle (0.28pt);
\fill[slate] (12.085,6.809) circle (0.28pt);
\fill[slate] (12.084,3.418) circle (0.28pt);
\fill[slate] (12.084,3.667) circle (0.28pt);
\fill[slate] (12.084,3.868) circle (0.28pt);
\fill[slate] (12.084,4.567) circle (0.28pt);
\fill[slate] (12.084,4.78) circle (0.28pt);
\fill[slate] (12.084,5.236) circle (0.28pt);
\fill[slate] (12.084,5.479) circle (0.28pt);
\fill[slate] (12.084,5.692) circle (0.28pt);
\fill[slate] (12.084,6.148) circle (0.28pt);
\fill[slate] (12.084,6.391) circle (0.28pt);
\fill[slate] (12.084,6.604) circle (0.28pt);
\fill[slate] (12.084,6.639) circle (0.28pt);
\fill[slate] (12.084,5.976) circle (0.28pt);
\fill[slate] (12.084,6.178) circle (0.28pt);
\fill[slate] (12.084,6.634) circle (0.28pt);
\fill[slate] (12.084,6.225) circle (0.28pt);
\fill[slate] (12.084,6.474) circle (0.28pt);
\fill[slate] (12.084,6.675) circle (0.28pt);
\fill[slate] (12.084,6.723) circle (0.28pt);
\fill[slate] (12.084,4.815) circle (0.28pt);
\fill[slate] (12.084,5.064) circle (0.28pt);
\fill[slate] (12.084,5.267) circle (0.28pt);
\fill[slate] (12.084,5.966) circle (0.28pt);
\fill[slate] (12.084,6.179) circle (0.28pt);
\fill[slate] (12.084,6.635) circle (0.28pt);
\fill[slate] (12.084,6.214) circle (0.28pt);
\fill[slate] (12.084,6.463) circle (0.28pt);
\fill[slate] (12.084,6.666) circle (0.28pt);
\fill[slate] (12.084,6.712) circle (0.28pt);
\fill[slate] (12.084,5.313) circle (0.28pt);
\fill[slate] (12.084,5.562) circle (0.28pt);
\fill[slate] (12.084,5.763) circle (0.28pt);
\fill[slate] (12.084,6.219) circle (0.28pt);
\fill[slate] (12.084,6.463) circle (0.28pt);
\fill[slate] (12.084,6.675) circle (0.28pt);
\fill[slate] (12.084,6.711) circle (0.28pt);
\fill[slate] (12.084,5.811) circle (0.28pt);
\fill[slate] (12.084,6.06) circle (0.28pt);
\fill[slate] (12.084,6.261) circle (0.28pt);
\fill[slate] (12.084,6.717) circle (0.28pt);
\fill[slate] (12.084,6.309) circle (0.28pt);
\fill[slate] (12.084,6.558) circle (0.28pt);
\fill[slate] (12.084,6.759) circle (0.28pt);
\fill[slate] (12.084,6.807) circle (0.28pt);
\fill[slate] (12.084,3.916) circle (0.28pt);
\fill[slate] (12.084,4.165) circle (0.28pt);
\fill[slate] (12.084,4.365) circle (0.28pt);
\fill[slate] (12.084,4.821) circle (0.28pt);
\fill[slate] (12.084,5.065) circle (0.28pt);
\fill[slate] (12.084,5.277) circle (0.28pt);
\fill[slate] (12.084,5.733) circle (0.28pt);
\fill[slate] (12.084,5.977) circle (0.28pt);
\fill[slate] (12.084,6.19) circle (0.28pt);
\fill[slate] (12.084,6.646) circle (0.28pt);
\fill[slate] (12.084,6.225) circle (0.28pt);
\fill[slate] (12.084,6.474) circle (0.28pt);
\fill[slate] (12.084,6.676) circle (0.28pt);
\fill[slate] (12.084,6.723) circle (0.28pt);
\fill[slate] (12.084,5.313) circle (0.28pt);
\fill[slate] (12.084,5.562) circle (0.28pt);
\fill[slate] (12.084,5.764) circle (0.28pt);
\fill[slate] (12.084,6.22) circle (0.28pt);
\fill[slate] (12.084,6.464) circle (0.28pt);
\fill[slate] (12.084,6.676) circle (0.28pt);
\fill[slate] (12.084,6.712) circle (0.28pt);
\fill[slate] (12.084,5.811) circle (0.28pt);
\fill[slate] (12.084,6.06) circle (0.28pt);
\fill[slate] (12.084,6.261) circle (0.28pt);
\fill[slate] (12.084,6.717) circle (0.28pt);
\fill[slate] (12.084,6.309) circle (0.28pt);
\fill[slate] (12.084,6.558) circle (0.28pt);
\fill[slate] (12.084,6.759) circle (0.28pt);
\fill[slate] (12.084,6.807) circle (0.28pt);
\fill[slate] (12.084,4.414) circle (0.28pt);
\fill[slate] (12.084,4.663) circle (0.28pt);
\fill[slate] (12.084,4.863) circle (0.28pt);
\fill[slate] (12.084,5.319) circle (0.28pt);
\fill[slate] (12.084,5.563) circle (0.28pt);
\fill[slate] (12.084,5.775) circle (0.28pt);
\fill[slate] (12.084,6.231) circle (0.28pt);
\fill[slate] (12.084,6.475) circle (0.28pt);
\fill[slate] (12.084,6.688) circle (0.28pt);
\fill[slate] (12.084,6.723) circle (0.28pt);
\fill[slate] (12.084,5.811) circle (0.28pt);
\fill[slate] (12.084,6.06) circle (0.28pt);
\fill[slate] (12.084,6.262) circle (0.28pt);
\fill[slate] (12.084,6.718) circle (0.28pt);
\fill[slate] (12.084,6.309) circle (0.28pt);
\fill[slate] (12.084,6.558) circle (0.28pt);
\fill[slate] (12.084,6.759) circle (0.28pt);
\fill[slate] (12.084,6.807) circle (0.28pt);
\fill[slate] (12.084,4.912) circle (0.28pt);
\fill[slate] (12.084,5.161) circle (0.28pt);
\fill[slate] (12.084,5.361) circle (0.28pt);
\fill[slate] (12.084,5.817) circle (0.28pt);
\fill[slate] (12.084,6.061) circle (0.28pt);
\fill[slate] (12.084,6.273) circle (0.28pt);
\fill[slate] (12.084,6.729) circle (0.28pt);
\fill[slate] (12.084,6.309) circle (0.28pt);
\fill[slate] (12.084,6.558) circle (0.28pt);
\fill[slate] (12.084,6.76) circle (0.28pt);
\fill[slate] (12.084,6.807) circle (0.28pt);
\fill[slate] (12.084,5.41) circle (0.28pt);
\fill[slate] (12.084,5.659) circle (0.28pt);
\fill[slate] (12.084,5.859) circle (0.28pt);
\fill[slate] (12.084,6.315) circle (0.28pt);
\fill[slate] (12.084,6.559) circle (0.28pt);
\fill[slate] (12.084,6.771) circle (0.28pt);
\fill[slate] (12.084,6.807) circle (0.28pt);
\fill[slate] (12.084,5.908) circle (0.28pt);
\fill[slate] (12.084,6.157) circle (0.28pt);
\fill[slate] (12.084,6.357) circle (0.28pt);
\fill[slate] (12.084,6.813) circle (0.28pt);
\fill[slate] (12.084,6.406) circle (0.28pt);
\fill[slate] (12.084,6.655) circle (0.28pt);
\fill[slate] (12.031,2.019) circle (0.28pt);
\fill[slate] (12.031,2.268) circle (0.28pt);
\fill[slate] (12.039,2.469) circle (0.28pt);
\fill[slate] (12.039,3.168) circle (0.28pt);
\fill[slate] (12.039,3.381) circle (0.28pt);
\fill[slate] (12.039,3.837) circle (0.28pt);
\fill[slate] (12.039,4.08) circle (0.28pt);
\fill[slate] (12.039,4.293) circle (0.28pt);
\fill[slate] (12.039,4.749) circle (0.28pt);
\fill[slate] (12.039,4.992) circle (0.28pt);
\fill[slate] (12.039,5.205) circle (0.28pt);
\fill[slate] (12.039,5.661) circle (0.28pt);
\fill[slate] (12.039,5.904) circle (0.28pt);
\fill[slate] (12.039,6.117) circle (0.28pt);
\fill[slate] (12.039,6.573) circle (0.28pt);
\fill[slate] (12.039,6.816) circle (0.28pt);
\fill[slate] (12.039,6.152) circle (0.28pt);
\fill[slate] (12.039,6.401) circle (0.28pt);
\fill[slate] (12.039,6.604) circle (0.28pt);
\fill[slate] (12.039,6.65) circle (0.28pt);
\fill[slate] (12.039,5.24) circle (0.28pt);
\fill[slate] (12.039,5.489) circle (0.28pt);
\fill[slate] (12.039,5.692) circle (0.28pt);
\fill[slate] (12.039,6.148) circle (0.28pt);
\fill[slate] (12.039,6.391) circle (0.28pt);
\fill[slate] (12.039,6.604) circle (0.28pt);
\fill[slate] (12.039,6.639) circle (0.28pt);
\fill[slate] (12.039,5.738) circle (0.28pt);
\fill[slate] (12.039,5.987) circle (0.28pt);
\fill[slate] (12.039,6.188) circle (0.28pt);
\fill[slate] (12.039,6.644) circle (0.28pt);
\fill[slate] (12.039,6.236) circle (0.28pt);
\fill[slate] (12.039,6.485) circle (0.28pt);
\fill[slate] (12.039,6.686) circle (0.28pt);
\fill[slate] (12.039,6.734) circle (0.28pt);
\fill[slate] (12.039,4.577) circle (0.28pt);
\fill[slate] (12.039,4.779) circle (0.28pt);
\fill[slate] (12.039,5.235) circle (0.28pt);
\fill[slate] (12.039,5.479) circle (0.28pt);
\fill[slate] (12.039,5.691) circle (0.28pt);
\fill[slate] (12.039,6.148) circle (0.28pt);
\fill[slate] (12.039,6.391) circle (0.28pt);
\fill[slate] (12.039,6.604) circle (0.28pt);
\fill[slate] (12.039,6.639) circle (0.28pt);
\fill[slate] (12.039,5.976) circle (0.28pt);
\fill[slate] (12.039,6.178) circle (0.28pt);
\fill[slate] (12.039,6.634) circle (0.28pt);
\fill[slate] (12.039,6.225) circle (0.28pt);
\fill[slate] (12.039,6.474) circle (0.28pt);
\fill[slate] (12.039,6.675) circle (0.28pt);
\fill[slate] (12.039,6.723) circle (0.28pt);
\fill[slate] (12.039,4.826) circle (0.28pt);
\fill[slate] (12.039,5.075) circle (0.28pt);
\fill[slate] (12.039,5.276) circle (0.28pt);
\fill[slate] (12.039,5.975) circle (0.28pt);
\fill[slate] (12.039,6.188) circle (0.28pt);
\fill[slate] (12.039,6.644) circle (0.28pt);
\fill[slate] (12.039,6.224) circle (0.28pt);
\fill[slate] (12.039,6.472) circle (0.28pt);
\fill[slate] (12.039,6.675) circle (0.28pt);
\fill[slate] (12.039,6.721) circle (0.28pt);
\fill[slate] (12.039,5.324) circle (0.28pt);
\fill[slate] (12.039,5.573) circle (0.28pt);
\fill[slate] (12.039,5.774) circle (0.28pt);
\fill[slate] (12.039,6.23) circle (0.28pt);
\fill[slate] (12.039,6.473) circle (0.28pt);
\fill[slate] (12.039,6.686) circle (0.28pt);
\fill[slate] (12.039,6.721) circle (0.28pt);
\fill[slate] (12.039,5.822) circle (0.28pt);
\fill[slate] (12.039,6.071) circle (0.28pt);
\fill[slate] (12.039,6.272) circle (0.28pt);
\fill[slate] (12.039,6.727) circle (0.28pt);
\fill[slate] (12.039,6.32) circle (0.28pt);
\fill[slate] (12.039,6.569) circle (0.28pt);
\fill[slate] (12.039,6.77) circle (0.28pt);
\fill[slate] (12.039,6.818) circle (0.28pt);
\fill[slate] (12.039,3.416) circle (0.28pt);
\fill[slate] (12.039,3.665) circle (0.28pt);
\fill[slate] (12.039,3.868) circle (0.28pt);
\fill[slate] (12.039,4.567) circle (0.28pt);
\fill[slate] (12.039,4.78) circle (0.28pt);
\fill[slate] (12.039,5.236) circle (0.28pt);
\fill[slate] (12.039,5.479) circle (0.28pt);
\fill[slate] (12.039,5.692) circle (0.28pt);
\fill[slate] (12.039,6.148) circle (0.28pt);
\fill[slate] (12.039,6.391) circle (0.28pt);
\fill[slate] (12.039,6.604) circle (0.28pt);
\fill[slate] (12.039,6.639) circle (0.28pt);
\fill[slate] (12.039,5.976) circle (0.28pt);
\fill[slate] (12.039,6.179) circle (0.28pt);
\fill[slate] (12.039,6.634) circle (0.28pt);
\fill[slate] (12.039,6.225) circle (0.28pt);
\fill[slate] (12.039,6.474) circle (0.28pt);
\fill[slate] (12.039,6.675) circle (0.28pt);
\fill[slate] (12.039,6.723) circle (0.28pt);
\fill[slate] (12.039,4.815) circle (0.28pt);
\fill[slate] (12.039,5.064) circle (0.28pt);
\fill[slate] (12.039,5.267) circle (0.28pt);
\fill[slate] (12.039,5.966) circle (0.28pt);
\fill[slate] (12.039,6.179) circle (0.28pt);
\fill[slate] (12.039,6.635) circle (0.28pt);
\fill[slate] (12.039,6.214) circle (0.28pt);
\fill[slate] (12.039,6.463) circle (0.28pt);
\fill[slate] (12.039,6.666) circle (0.28pt);
\fill[slate] (12.039,6.712) circle (0.28pt);
\fill[slate] (12.039,5.313) circle (0.28pt);
\fill[slate] (12.039,5.562) circle (0.28pt);
\fill[slate] (12.039,5.763) circle (0.28pt);
\fill[slate] (12.039,6.219) circle (0.28pt);
\fill[slate] (12.039,6.463) circle (0.28pt);
\fill[slate] (12.039,6.675) circle (0.28pt);
\fill[slate] (12.039,6.711) circle (0.28pt);
\fill[slate] (12.039,5.811) circle (0.28pt);
\fill[slate] (12.039,6.06) circle (0.28pt);
\fill[slate] (12.039,6.261) circle (0.28pt);
\fill[slate] (12.039,6.717) circle (0.28pt);
\fill[slate] (12.039,6.309) circle (0.28pt);
\fill[slate] (12.039,6.558) circle (0.28pt);
\fill[slate] (12.039,6.759) circle (0.28pt);
\fill[slate] (12.039,6.807) circle (0.28pt);
\fill[slate] (12.039,3.914) circle (0.28pt);
\fill[slate] (12.039,4.163) circle (0.28pt);
\fill[slate] (12.039,4.364) circle (0.28pt);
\fill[slate] (12.039,4.82) circle (0.28pt);
\fill[slate] (12.039,5.064) circle (0.28pt);
\fill[slate] (12.039,5.276) circle (0.28pt);
\fill[slate] (12.039,5.732) circle (0.28pt);
\fill[slate] (12.039,5.975) circle (0.28pt);
\fill[slate] (12.039,6.188) circle (0.28pt);
\fill[slate] (12.039,6.644) circle (0.28pt);
\fill[slate] (12.039,6.224) circle (0.28pt);
\fill[slate] (12.039,6.472) circle (0.28pt);
\fill[slate] (12.039,6.675) circle (0.28pt);
\fill[slate] (12.039,6.721) circle (0.28pt);
\fill[slate] (12.039,5.312) circle (0.28pt);
\fill[slate] (12.039,5.561) circle (0.28pt);
\fill[slate] (12.039,5.763) circle (0.28pt);
\fill[slate] (12.039,6.219) circle (0.28pt);
\fill[slate] (12.039,6.463) circle (0.28pt);
\fill[slate] (12.039,6.675) circle (0.28pt);
\fill[slate] (12.039,6.711) circle (0.28pt);
\fill[slate] (12.039,5.81) circle (0.28pt);
\fill[slate] (12.039,6.059) circle (0.28pt);
\fill[slate] (12.039,6.259) circle (0.28pt);
\fill[slate] (12.039,6.715) circle (0.28pt);
\fill[slate] (12.039,6.308) circle (0.28pt);
\fill[slate] (12.039,6.557) circle (0.28pt);
\fill[slate] (12.039,6.757) circle (0.28pt);
\fill[slate] (12.039,6.806) circle (0.28pt);
\fill[slate] (12.039,4.412) circle (0.28pt);
\fill[slate] (12.039,4.661) circle (0.28pt);
\fill[slate] (12.039,4.862) circle (0.28pt);
\fill[slate] (12.039,5.318) circle (0.28pt);
\fill[slate] (12.039,5.561) circle (0.28pt);
\fill[slate] (12.039,5.774) circle (0.28pt);
\fill[slate] (12.039,6.23) circle (0.28pt);
\fill[slate] (12.039,6.473) circle (0.28pt);
\fill[slate] (12.039,6.686) circle (0.28pt);
\fill[slate] (12.039,6.721) circle (0.28pt);
\fill[slate] (12.039,5.81) circle (0.28pt);
\fill[slate] (12.039,6.058) circle (0.28pt);
\fill[slate] (12.039,6.261) circle (0.28pt);
\fill[slate] (12.039,6.717) circle (0.28pt);
\fill[slate] (12.039,6.307) circle (0.28pt);
\fill[slate] (12.039,6.556) circle (0.28pt);
\fill[slate] (12.039,6.757) circle (0.28pt);
\fill[slate] (12.039,6.805) circle (0.28pt);
\fill[slate] (12.039,4.91) circle (0.28pt);
\fill[slate] (12.039,5.159) circle (0.28pt);
\fill[slate] (12.039,5.36) circle (0.28pt);
\fill[slate] (12.039,5.816) circle (0.28pt);
\fill[slate] (12.039,6.059) circle (0.28pt);
\fill[slate] (12.039,6.272) circle (0.28pt);
\fill[slate] (12.039,6.728) circle (0.28pt);
\fill[slate] (12.039,6.308) circle (0.28pt);
\fill[slate] (12.039,6.556) circle (0.28pt);
\fill[slate] (12.039,6.759) circle (0.28pt);
\fill[slate] (12.039,6.805) circle (0.28pt);
\fill[slate] (12.039,5.408) circle (0.28pt);
\fill[slate] (12.039,5.657) circle (0.28pt);
\fill[slate] (12.039,5.858) circle (0.28pt);
\fill[slate] (12.039,6.314) circle (0.28pt);
\fill[slate] (12.039,6.557) circle (0.28pt);
\fill[slate] (12.039,6.77) circle (0.28pt);
\fill[slate] (12.039,6.806) circle (0.28pt);
\fill[slate] (12.039,5.906) circle (0.28pt);
\fill[slate] (12.039,6.155) circle (0.28pt);
\fill[slate] (12.039,6.356) circle (0.28pt);
\fill[slate] (12.039,6.812) circle (0.28pt);
\fill[slate] (12.039,6.404) circle (0.28pt);
\fill[slate] (12.039,6.653) circle (0.28pt);
\fill[slate] (12.031,2.517) circle (0.28pt);
\fill[slate] (12.031,2.766) circle (0.28pt);
\fill[slate] (12.032,2.966) circle (0.28pt);
\fill[slate] (12.032,3.422) circle (0.28pt);
\fill[slate] (12.032,3.666) circle (0.28pt);
\fill[slate] (12.032,3.878) circle (0.28pt);
\fill[slate] (12.032,4.334) circle (0.28pt);
\fill[slate] (12.032,4.578) circle (0.28pt);
\fill[slate] (12.032,4.791) circle (0.28pt);
\fill[slate] (12.032,5.247) circle (0.28pt);
\fill[slate] (12.032,5.49) circle (0.28pt);
\fill[slate] (12.032,5.703) circle (0.28pt);
\fill[slate] (12.032,6.159) circle (0.28pt);
\fill[slate] (12.032,6.402) circle (0.28pt);
\fill[slate] (12.032,6.615) circle (0.28pt);
\fill[slate] (12.032,6.65) circle (0.28pt);
\fill[slate] (12.032,5.738) circle (0.28pt);
\fill[slate] (12.032,5.987) circle (0.28pt);
\fill[slate] (12.032,6.189) circle (0.28pt);
\fill[slate] (12.032,6.645) circle (0.28pt);
\fill[slate] (12.032,6.236) circle (0.28pt);
\fill[slate] (12.032,6.485) circle (0.28pt);
\fill[slate] (12.032,6.686) circle (0.28pt);
\fill[slate] (12.032,6.734) circle (0.28pt);
\fill[slate] (12.032,4.826) circle (0.28pt);
\fill[slate] (12.032,5.075) circle (0.28pt);
\fill[slate] (12.032,5.277) circle (0.28pt);
\fill[slate] (12.032,5.733) circle (0.28pt);
\fill[slate] (12.032,5.977) circle (0.28pt);
\fill[slate] (12.032,6.189) circle (0.28pt);
\fill[slate] (12.032,6.645) circle (0.28pt);
\fill[slate] (12.032,6.225) circle (0.28pt);
\fill[slate] (12.032,6.676) circle (0.28pt);
\fill[slate] (12.032,6.723) circle (0.28pt);
\fill[slate] (12.032,5.324) circle (0.28pt);
\fill[slate] (12.032,5.573) circle (0.28pt);
\fill[slate] (12.032,5.774) circle (0.28pt);
\fill[slate] (12.032,6.23) circle (0.28pt);
\fill[slate] (12.032,6.473) circle (0.28pt);
\fill[slate] (12.032,6.686) circle (0.28pt);
\fill[slate] (12.032,6.721) circle (0.28pt);
\fill[slate] (12.032,5.822) circle (0.28pt);
\fill[slate] (12.032,6.071) circle (0.28pt);
\fill[slate] (12.032,6.271) circle (0.28pt);
\fill[slate] (12.032,6.727) circle (0.28pt);
\fill[slate] (12.032,6.32) circle (0.28pt);
\fill[slate] (12.032,6.569) circle (0.28pt);
\fill[slate] (12.032,6.769) circle (0.28pt);
\fill[slate] (12.032,6.818) circle (0.28pt);
\fill[slate] (12.032,3.914) circle (0.28pt);
\fill[slate] (12.032,4.163) circle (0.28pt);
\fill[slate] (12.032,4.365) circle (0.28pt);
\fill[slate] (12.032,4.821) circle (0.28pt);
\fill[slate] (12.032,5.065) circle (0.28pt);
\fill[slate] (12.032,5.277) circle (0.28pt);
\fill[slate] (12.032,5.734) circle (0.28pt);
\fill[slate] (12.032,5.977) circle (0.28pt);
\fill[slate] (12.032,6.19) circle (0.28pt);
\fill[slate] (12.032,6.646) circle (0.28pt);
\fill[slate] (12.032,6.225) circle (0.28pt);
\fill[slate] (12.032,6.676) circle (0.28pt);
\fill[slate] (12.032,6.723) circle (0.28pt);
\fill[slate] (12.032,5.313) circle (0.28pt);
\fill[slate] (12.032,5.562) circle (0.28pt);
\fill[slate] (12.032,5.764) circle (0.28pt);
\fill[slate] (12.032,6.22) circle (0.28pt);
\fill[slate] (12.032,6.464) circle (0.28pt);
\fill[slate] (12.032,6.676) circle (0.28pt);
\fill[slate] (12.032,6.712) circle (0.28pt);
\fill[slate] (12.032,5.811) circle (0.28pt);
\fill[slate] (12.032,6.06) circle (0.28pt);
\fill[slate] (12.032,6.261) circle (0.28pt);
\fill[slate] (12.032,6.717) circle (0.28pt);
\fill[slate] (12.032,6.309) circle (0.28pt);
\fill[slate] (12.032,6.558) circle (0.28pt);
\fill[slate] (12.032,6.759) circle (0.28pt);
\fill[slate] (12.032,6.807) circle (0.28pt);
\fill[slate] (12.032,4.412) circle (0.28pt);
\fill[slate] (12.032,4.661) circle (0.28pt);
\fill[slate] (12.032,4.862) circle (0.28pt);
\fill[slate] (12.032,5.318) circle (0.28pt);
\fill[slate] (12.032,5.561) circle (0.28pt);
\fill[slate] (12.032,5.774) circle (0.28pt);
\fill[slate] (12.032,6.23) circle (0.28pt);
\fill[slate] (12.032,6.473) circle (0.28pt);
\fill[slate] (12.032,6.686) circle (0.28pt);
\fill[slate] (12.032,6.721) circle (0.28pt);
\fill[slate] (12.032,5.81) circle (0.28pt);
\fill[slate] (12.032,6.058) circle (0.28pt);
\fill[slate] (12.032,6.261) circle (0.28pt);
\fill[slate] (12.032,6.717) circle (0.28pt);
\fill[slate] (12.032,6.307) circle (0.28pt);
\fill[slate] (12.032,6.556) circle (0.28pt);
\fill[slate] (12.032,6.757) circle (0.28pt);
\fill[slate] (12.032,6.805) circle (0.28pt);
\fill[slate] (12.032,4.91) circle (0.28pt);
\fill[slate] (12.032,5.159) circle (0.28pt);
\fill[slate] (12.032,5.36) circle (0.28pt);
\fill[slate] (12.032,5.816) circle (0.28pt);
\fill[slate] (12.032,6.059) circle (0.28pt);
\fill[slate] (12.032,6.272) circle (0.28pt);
\fill[slate] (12.032,6.728) circle (0.28pt);
\fill[slate] (12.032,6.307) circle (0.28pt);
\fill[slate] (12.032,6.556) circle (0.28pt);
\fill[slate] (12.032,6.759) circle (0.28pt);
\fill[slate] (12.032,6.805) circle (0.28pt);
\fill[slate] (12.032,5.408) circle (0.28pt);
\fill[slate] (12.032,5.657) circle (0.28pt);
\fill[slate] (12.032,5.858) circle (0.28pt);
\fill[slate] (12.032,6.313) circle (0.28pt);
\fill[slate] (12.032,6.557) circle (0.28pt);
\fill[slate] (12.032,6.77) circle (0.28pt);
\fill[slate] (12.032,6.805) circle (0.28pt);
\fill[slate] (12.032,5.906) circle (0.28pt);
\fill[slate] (12.032,6.155) circle (0.28pt);
\fill[slate] (12.032,6.356) circle (0.28pt);
\fill[slate] (12.032,6.811) circle (0.28pt);
\fill[slate] (12.032,6.404) circle (0.28pt);
\fill[slate] (12.032,6.653) circle (0.28pt);
\fill[slate] (12.031,3.015) circle (0.28pt);
\fill[slate] (12.031,3.264) circle (0.28pt);
\fill[slate] (12.031,3.464) circle (0.28pt);
\fill[slate] (12.031,3.92) circle (0.28pt);
\fill[slate] (12.031,4.164) circle (0.28pt);
\fill[slate] (12.031,4.376) circle (0.28pt);
\fill[slate] (12.031,4.832) circle (0.28pt);
\fill[slate] (12.031,5.076) circle (0.28pt);
\fill[slate] (12.031,5.289) circle (0.28pt);
\fill[slate] (12.031,5.745) circle (0.28pt);
\fill[slate] (12.031,5.988) circle (0.28pt);
\fill[slate] (12.031,6.201) circle (0.28pt);
\fill[slate] (12.031,6.657) circle (0.28pt);
\fill[slate] (12.031,6.236) circle (0.28pt);
\fill[slate] (12.031,6.485) circle (0.28pt);
\fill[slate] (12.031,6.687) circle (0.28pt);
\fill[slate] (12.031,6.734) circle (0.28pt);
\fill[slate] (12.031,5.324) circle (0.28pt);
\fill[slate] (12.031,5.573) circle (0.28pt);
\fill[slate] (12.031,5.775) circle (0.28pt);
\fill[slate] (12.031,6.231) circle (0.28pt);
\fill[slate] (12.031,6.475) circle (0.28pt);
\fill[slate] (12.031,6.687) circle (0.28pt);
\fill[slate] (12.031,6.723) circle (0.28pt);
\fill[slate] (12.031,5.822) circle (0.28pt);
\fill[slate] (12.031,6.071) circle (0.28pt);
\fill[slate] (12.031,6.272) circle (0.28pt);
\fill[slate] (12.031,6.727) circle (0.28pt);
\fill[slate] (12.031,6.32) circle (0.28pt);
\fill[slate] (12.031,6.569) circle (0.28pt);
\fill[slate] (12.031,6.769) circle (0.28pt);
\fill[slate] (12.031,6.818) circle (0.28pt);
\fill[slate] (12.031,4.412) circle (0.28pt);
\fill[slate] (12.031,4.661) circle (0.28pt);
\fill[slate] (12.031,4.863) circle (0.28pt);
\fill[slate] (12.031,5.319) circle (0.28pt);
\fill[slate] (12.031,5.563) circle (0.28pt);
\fill[slate] (12.031,5.775) circle (0.28pt);
\fill[slate] (12.031,6.231) circle (0.28pt);
\fill[slate] (12.031,6.475) circle (0.28pt);
\fill[slate] (12.031,6.688) circle (0.28pt);
\fill[slate] (12.031,6.723) circle (0.28pt);
\fill[slate] (12.031,5.811) circle (0.28pt);
\fill[slate] (12.031,6.06) circle (0.28pt);
\fill[slate] (12.031,6.262) circle (0.28pt);
\fill[slate] (12.031,6.718) circle (0.28pt);
\fill[slate] (12.031,6.309) circle (0.28pt);
\fill[slate] (12.031,6.558) circle (0.28pt);
\fill[slate] (12.031,6.759) circle (0.28pt);
\fill[slate] (12.031,6.807) circle (0.28pt);
\fill[slate] (12.031,4.91) circle (0.28pt);
\fill[slate] (12.031,5.159) circle (0.28pt);
\fill[slate] (12.031,5.36) circle (0.28pt);
\fill[slate] (12.031,5.816) circle (0.28pt);
\fill[slate] (12.031,6.059) circle (0.28pt);
\fill[slate] (12.031,6.272) circle (0.28pt);
\fill[slate] (12.031,6.728) circle (0.28pt);
\fill[slate] (12.031,6.308) circle (0.28pt);
\fill[slate] (12.031,6.556) circle (0.28pt);
\fill[slate] (12.031,6.759) circle (0.28pt);
\fill[slate] (12.031,6.805) circle (0.28pt);
\fill[slate] (12.031,5.408) circle (0.28pt);
\fill[slate] (12.031,5.657) circle (0.28pt);
\fill[slate] (12.031,5.858) circle (0.28pt);
\fill[slate] (12.031,6.313) circle (0.28pt);
\fill[slate] (12.031,6.557) circle (0.28pt);
\fill[slate] (12.031,6.77) circle (0.28pt);
\fill[slate] (12.031,6.805) circle (0.28pt);
\fill[slate] (12.031,5.906) circle (0.28pt);
\fill[slate] (12.031,6.155) circle (0.28pt);
\fill[slate] (12.031,6.356) circle (0.28pt);
\fill[slate] (12.031,6.811) circle (0.28pt);
\fill[slate] (12.031,6.404) circle (0.28pt);
\fill[slate] (12.031,6.653) circle (0.28pt);
\fill[slate] (12.031,3.513) circle (0.28pt);
\fill[slate] (12.031,3.762) circle (0.28pt);
\fill[slate] (12.031,3.962) circle (0.28pt);
\fill[slate] (12.031,4.418) circle (0.28pt);
\fill[slate] (12.031,4.662) circle (0.28pt);
\fill[slate] (12.031,4.874) circle (0.28pt);
\fill[slate] (12.031,5.33) circle (0.28pt);
\fill[slate] (12.031,5.574) circle (0.28pt);
\fill[slate] (12.031,5.787) circle (0.28pt);
\fill[slate] (12.031,6.243) circle (0.28pt);
\fill[slate] (12.031,6.486) circle (0.28pt);
\fill[slate] (12.031,6.699) circle (0.28pt);
\fill[slate] (12.031,6.734) circle (0.28pt);
\fill[slate] (12.031,5.822) circle (0.28pt);
\fill[slate] (12.031,6.071) circle (0.28pt);
\fill[slate] (12.031,6.273) circle (0.28pt);
\fill[slate] (12.031,6.729) circle (0.28pt);
\fill[slate] (12.031,6.32) circle (0.28pt);
\fill[slate] (12.031,6.569) circle (0.28pt);
\fill[slate] (12.031,6.77) circle (0.28pt);
\fill[slate] (12.031,6.818) circle (0.28pt);
\fill[slate] (12.031,4.91) circle (0.28pt);
\fill[slate] (12.031,5.159) circle (0.28pt);
\fill[slate] (12.031,5.361) circle (0.28pt);
\fill[slate] (12.031,5.817) circle (0.28pt);
\fill[slate] (12.031,6.061) circle (0.28pt);
\fill[slate] (12.031,6.273) circle (0.28pt);
\fill[slate] (12.031,6.729) circle (0.28pt);
\fill[slate] (12.031,6.309) circle (0.28pt);
\fill[slate] (12.031,6.558) circle (0.28pt);
\fill[slate] (12.031,6.76) circle (0.28pt);
\fill[slate] (12.031,6.807) circle (0.28pt);
\fill[slate] (12.031,5.408) circle (0.28pt);
\fill[slate] (12.031,5.657) circle (0.28pt);
\fill[slate] (12.031,5.858) circle (0.28pt);
\fill[slate] (12.031,6.314) circle (0.28pt);
\fill[slate] (12.031,6.557) circle (0.28pt);
\fill[slate] (12.031,6.77) circle (0.28pt);
\fill[slate] (12.031,6.806) circle (0.28pt);
\fill[slate] (12.031,5.906) circle (0.28pt);
\fill[slate] (12.031,6.155) circle (0.28pt);
\fill[slate] (12.031,6.356) circle (0.28pt);
\fill[slate] (12.031,6.811) circle (0.28pt);
\fill[slate] (12.031,6.404) circle (0.28pt);
\fill[slate] (12.031,6.653) circle (0.28pt);
\fill[slate] (12.031,4.011) circle (0.28pt);
\fill[slate] (12.031,4.26) circle (0.28pt);
\fill[slate] (12.031,4.46) circle (0.28pt);
\fill[slate] (12.031,4.916) circle (0.28pt);
\fill[slate] (12.031,5.16) circle (0.28pt);
\fill[slate] (12.031,5.372) circle (0.28pt);
\fill[slate] (12.031,5.828) circle (0.28pt);
\fill[slate] (12.031,6.072) circle (0.28pt);
\fill[slate] (12.031,6.285) circle (0.28pt);
\fill[slate] (12.031,6.741) circle (0.28pt);
\fill[slate] (12.031,6.32) circle (0.28pt);
\fill[slate] (12.031,6.569) circle (0.28pt);
\fill[slate] (12.031,6.771) circle (0.28pt);
\fill[slate] (12.031,6.818) circle (0.28pt);
\fill[slate] (12.031,5.408) circle (0.28pt);
\fill[slate] (12.031,5.657) circle (0.28pt);
\fill[slate] (12.031,5.859) circle (0.28pt);
\fill[slate] (12.031,6.315) circle (0.28pt);
\fill[slate] (12.031,6.559) circle (0.28pt);
\fill[slate] (12.031,6.771) circle (0.28pt);
\fill[slate] (12.031,6.807) circle (0.28pt);
\fill[slate] (12.031,5.906) circle (0.28pt);
\fill[slate] (12.031,6.155) circle (0.28pt);
\fill[slate] (12.031,6.356) circle (0.28pt);
\fill[slate] (12.031,6.812) circle (0.28pt);
\fill[slate] (12.031,6.404) circle (0.28pt);
\fill[slate] (12.031,6.653) circle (0.28pt);
\fill[slate] (12.031,4.509) circle (0.28pt);
\fill[slate] (12.031,4.758) circle (0.28pt);
\fill[slate] (12.031,4.958) circle (0.28pt);
\fill[slate] (12.031,5.414) circle (0.28pt);
\fill[slate] (12.031,5.658) circle (0.28pt);
\fill[slate] (12.031,5.87) circle (0.28pt);
\fill[slate] (12.031,6.326) circle (0.28pt);
\fill[slate] (12.031,6.57) circle (0.28pt);
\fill[slate] (12.031,6.783) circle (0.28pt);
\fill[slate] (12.031,6.818) circle (0.28pt);
\fill[slate] (12.031,5.906) circle (0.28pt);
\fill[slate] (12.031,6.155) circle (0.28pt);
\fill[slate] (12.031,6.357) circle (0.28pt);
\fill[slate] (12.031,6.813) circle (0.28pt);
\fill[slate] (12.031,6.404) circle (0.28pt);
\fill[slate] (12.031,6.653) circle (0.28pt);
\fill[slate] (12.031,5.007) circle (0.28pt);
\fill[slate] (12.031,5.256) circle (0.28pt);
\fill[slate] (12.031,5.456) circle (0.28pt);
\fill[slate] (12.031,5.912) circle (0.28pt);
\fill[slate] (12.031,6.156) circle (0.28pt);
\fill[slate] (12.031,6.368) circle (0.28pt);
\fill[slate] (12.031,6.825) circle (0.28pt);
\fill[slate] (12.031,6.404) circle (0.28pt);
\fill[slate] (12.031,6.653) circle (0.28pt);
\fill[slate] (12.031,5.505) circle (0.28pt);
\fill[slate] (12.031,5.754) circle (0.28pt);
\fill[slate] (12.031,5.954) circle (0.28pt);
\fill[slate] (12.031,6.41) circle (0.28pt);
\fill[slate] (12.031,6.654) circle (0.28pt);
\fill[slate] (12.031,6.003) circle (0.28pt);
\fill[slate] (12.031,6.252) circle (0.28pt);
\fill[slate] (12.031,6.453) circle (0.28pt);
\fill[slate] (12.031,6.501) circle (0.28pt);
\fill[slate] (12.031,6.75) circle (0.28pt);
\fill[slate] (1.832,1.977) circle (0.28pt);
\fill[slate] (1.832,2.433) circle (0.28pt);
\fill[slate] (1.83,2.676) circle (0.28pt);
\fill[slate] (1.832,2.889) circle (0.28pt);
\fill[slate] (1.832,3.345) circle (0.28pt);
\fill[slate] (1.832,3.589) circle (0.28pt);
\fill[slate] (1.832,3.801) circle (0.28pt);
\fill[slate] (1.832,4.257) circle (0.28pt);
\fill[slate] (1.832,4.501) circle (0.28pt);
\fill[slate] (1.832,4.713) circle (0.28pt);
\fill[slate] (1.832,5.17) circle (0.28pt);
\fill[slate] (1.832,5.413) circle (0.28pt);
\fill[slate] (1.832,5.626) circle (0.28pt);
\fill[slate] (1.832,6.082) circle (0.28pt);
\fill[slate] (1.832,6.325) circle (0.28pt);
\fill[slate] (1.832,6.538) circle (0.28pt);
\fill[slate] (1.832,6.573) circle (0.28pt);
\fill[slate] (1.832,6.822) circle (0.28pt);
\fill[slate] (1.832,5.661) circle (0.28pt);
\fill[slate] (1.832,5.91) circle (0.28pt);
\fill[slate] (1.832,6.112) circle (0.28pt);
\fill[slate] (1.832,6.568) circle (0.28pt);
\fill[slate] (1.832,6.812) circle (0.28pt);
\fill[slate] (1.832,6.159) circle (0.28pt);
\fill[slate] (1.832,6.408) circle (0.28pt);
\fill[slate] (1.832,6.609) circle (0.28pt);
\fill[slate] (1.832,6.657) circle (0.28pt);
\fill[slate] (1.832,4.749) circle (0.28pt);
\fill[slate] (1.832,4.998) circle (0.28pt);
\fill[slate] (1.832,5.2) circle (0.28pt);
\fill[slate] (1.832,5.656) circle (0.28pt);
\fill[slate] (1.832,5.9) circle (0.28pt);
\fill[slate] (1.832,6.112) circle (0.28pt);
\fill[slate] (1.832,6.568) circle (0.28pt);
\fill[slate] (1.832,6.812) circle (0.28pt);
\fill[slate] (1.832,6.148) circle (0.28pt);
\fill[slate] (1.832,6.397) circle (0.28pt);
\fill[slate] (1.832,6.599) circle (0.28pt);
\fill[slate] (1.832,6.646) circle (0.28pt);
\fill[slate] (1.832,5.247) circle (0.28pt);
\fill[slate] (1.832,5.496) circle (0.28pt);
\fill[slate] (1.832,5.697) circle (0.28pt);
\fill[slate] (1.832,6.152) circle (0.28pt);
\fill[slate] (1.832,6.396) circle (0.28pt);
\fill[slate] (1.832,6.609) circle (0.28pt);
\fill[slate] (1.832,6.644) circle (0.28pt);
\fill[slate] (1.832,5.745) circle (0.28pt);
\fill[slate] (1.832,5.994) circle (0.28pt);
\fill[slate] (1.832,6.194) circle (0.28pt);
\fill[slate] (1.832,6.65) circle (0.28pt);
\fill[slate] (1.832,6.243) circle (0.28pt);
\fill[slate] (1.832,6.492) circle (0.28pt);
\fill[slate] (1.832,6.692) circle (0.28pt);
\fill[slate] (1.832,6.741) circle (0.28pt);
\fill[slate] (1.832,3.837) circle (0.28pt);
\fill[slate] (1.832,4.086) circle (0.28pt);
\fill[slate] (1.832,4.288) circle (0.28pt);
\fill[slate] (1.832,4.744) circle (0.28pt);
\fill[slate] (1.832,4.988) circle (0.28pt);
\fill[slate] (1.832,5.2) circle (0.28pt);
\fill[slate] (1.832,5.656) circle (0.28pt);
\fill[slate] (1.832,5.899) circle (0.28pt);
\fill[slate] (1.832,6.112) circle (0.28pt);
\fill[slate] (1.832,6.568) circle (0.28pt);
\fill[slate] (1.832,6.812) circle (0.28pt);
\fill[slate] (1.832,6.148) circle (0.28pt);
\fill[slate] (1.832,6.396) circle (0.28pt);
\fill[slate] (1.832,6.599) circle (0.28pt);
\fill[slate] (1.832,6.645) circle (0.28pt);
\fill[slate] (1.832,5.236) circle (0.28pt);
\fill[slate] (1.832,5.485) circle (0.28pt);
\fill[slate] (1.832,5.687) circle (0.28pt);
\fill[slate] (1.832,6.143) circle (0.28pt);
\fill[slate] (1.832,6.387) circle (0.28pt);
\fill[slate] (1.832,6.599) circle (0.28pt);
\fill[slate] (1.832,6.635) circle (0.28pt);
\fill[slate] (1.832,5.734) circle (0.28pt);
\fill[slate] (1.832,5.983) circle (0.28pt);
\fill[slate] (1.832,6.183) circle (0.28pt);
\fill[slate] (1.832,6.639) circle (0.28pt);
\fill[slate] (1.832,6.232) circle (0.28pt);
\fill[slate] (1.832,6.481) circle (0.28pt);
\fill[slate] (1.832,6.681) circle (0.28pt);
\fill[slate] (1.832,6.73) circle (0.28pt);
\fill[slate] (1.832,4.335) circle (0.28pt);
\fill[slate] (1.832,4.584) circle (0.28pt);
\fill[slate] (1.832,4.784) circle (0.28pt);
\fill[slate] (1.832,5.24) circle (0.28pt);
\fill[slate] (1.832,5.484) circle (0.28pt);
\fill[slate] (1.832,5.696) circle (0.28pt);
\fill[slate] (1.832,6.152) circle (0.28pt);
\fill[slate] (1.832,6.396) circle (0.28pt);
\fill[slate] (1.832,6.609) circle (0.28pt);
\fill[slate] (1.832,6.644) circle (0.28pt);
\fill[slate] (1.832,5.732) circle (0.28pt);
\fill[slate] (1.832,5.981) circle (0.28pt);
\fill[slate] (1.832,6.183) circle (0.28pt);
\fill[slate] (1.832,6.639) circle (0.28pt);
\fill[slate] (1.832,6.23) circle (0.28pt);
\fill[slate] (1.832,6.479) circle (0.28pt);
\fill[slate] (1.832,6.68) circle (0.28pt);
\fill[slate] (1.832,6.728) circle (0.28pt);
\fill[slate] (1.832,4.833) circle (0.28pt);
\fill[slate] (1.832,5.082) circle (0.28pt);
\fill[slate] (1.832,5.282) circle (0.28pt);
\fill[slate] (1.832,5.738) circle (0.28pt);
\fill[slate] (1.832,5.982) circle (0.28pt);
\fill[slate] (1.832,6.194) circle (0.28pt);
\fill[slate] (1.832,6.65) circle (0.28pt);
\fill[slate] (1.832,6.23) circle (0.28pt);
\fill[slate] (1.832,6.479) circle (0.28pt);
\fill[slate] (1.832,6.681) circle (0.28pt);
\fill[slate] (1.832,6.728) circle (0.28pt);
\fill[slate] (1.832,5.331) circle (0.28pt);
\fill[slate] (1.832,5.58) circle (0.28pt);
\fill[slate] (1.832,5.78) circle (0.28pt);
\fill[slate] (1.832,6.236) circle (0.28pt);
\fill[slate] (1.832,6.48) circle (0.28pt);
\fill[slate] (1.832,6.692) circle (0.28pt);
\fill[slate] (1.832,6.728) circle (0.28pt);
\fill[slate] (1.832,5.829) circle (0.28pt);
\fill[slate] (1.832,6.078) circle (0.28pt);
\fill[slate] (1.832,6.278) circle (0.28pt);
\fill[slate] (1.832,6.734) circle (0.28pt);
\fill[slate] (1.832,6.327) circle (0.28pt);
\fill[slate] (1.832,6.576) circle (0.28pt);
\fill[slate] (1.832,6.776) circle (0.28pt);
\fill[slate] (1.832,6.825) circle (0.28pt);
\fill[slate] (1.83,3.173) circle (0.28pt);
\fill[slate] (1.83,3.376) circle (0.28pt);
\fill[slate] (1.83,3.832) circle (0.28pt);
\fill[slate] (1.83,4.075) circle (0.28pt);
\fill[slate] (1.83,4.288) circle (0.28pt);
\fill[slate] (1.83,4.744) circle (0.28pt);
\fill[slate] (1.83,4.987) circle (0.28pt);
\fill[slate] (1.83,5.2) circle (0.28pt);
\fill[slate] (1.83,5.656) circle (0.28pt);
\fill[slate] (1.83,5.899) circle (0.28pt);
\fill[slate] (1.83,6.112) circle (0.28pt);
\fill[slate] (1.83,6.568) circle (0.28pt);
\fill[slate] (1.83,6.812) circle (0.28pt);
\fill[slate] (1.83,6.148) circle (0.28pt);
\fill[slate] (1.83,6.396) circle (0.28pt);
\fill[slate] (1.83,6.599) circle (0.28pt);
\fill[slate] (1.83,6.645) circle (0.28pt);
\fill[slate] (1.83,5.235) circle (0.28pt);
\fill[slate] (1.83,5.484) circle (0.28pt);
\fill[slate] (1.83,5.687) circle (0.28pt);
\fill[slate] (1.83,6.143) circle (0.28pt);
\fill[slate] (1.83,6.386) circle (0.28pt);
\fill[slate] (1.83,6.599) circle (0.28pt);
\fill[slate] (1.83,6.634) circle (0.28pt);
\fill[slate] (1.83,5.733) circle (0.28pt);
\fill[slate] (1.83,5.982) circle (0.28pt);
\fill[slate] (1.83,6.183) circle (0.28pt);
\fill[slate] (1.83,6.639) circle (0.28pt);
\fill[slate] (1.83,6.231) circle (0.28pt);
\fill[slate] (1.83,6.48) circle (0.28pt);
\fill[slate] (1.83,6.681) circle (0.28pt);
\fill[slate] (1.83,6.729) circle (0.28pt);
\fill[slate] (1.83,4.572) circle (0.28pt);
\fill[slate] (1.83,4.775) circle (0.28pt);
\fill[slate] (1.83,5.231) circle (0.28pt);
\fill[slate] (1.83,5.474) circle (0.28pt);
\fill[slate] (1.83,5.687) circle (0.28pt);
\fill[slate] (1.83,6.143) circle (0.28pt);
\fill[slate] (1.83,6.386) circle (0.28pt);
\fill[slate] (1.83,6.599) circle (0.28pt);
\fill[slate] (1.83,6.634) circle (0.28pt);
\fill[slate] (1.83,5.971) circle (0.28pt);
\fill[slate] (1.83,6.174) circle (0.28pt);
\fill[slate] (1.83,6.63) circle (0.28pt);
\fill[slate] (1.83,6.22) circle (0.28pt);
\fill[slate] (1.83,6.469) circle (0.28pt);
\fill[slate] (1.83,6.67) circle (0.28pt);
\fill[slate] (1.83,6.718) circle (0.28pt);
\fill[slate] (1.83,4.821) circle (0.28pt);
\fill[slate] (1.83,5.07) circle (0.28pt);
\fill[slate] (1.83,5.271) circle (0.28pt);
\fill[slate] (1.83,5.971) circle (0.28pt);
\fill[slate] (1.83,6.183) circle (0.28pt);
\fill[slate] (1.83,6.639) circle (0.28pt);
\fill[slate] (1.83,6.219) circle (0.28pt);
\fill[slate] (1.83,6.468) circle (0.28pt);
\fill[slate] (1.83,6.67) circle (0.28pt);
\fill[slate] (1.83,6.717) circle (0.28pt);
\fill[slate] (1.83,5.319) circle (0.28pt);
\fill[slate] (1.83,5.568) circle (0.28pt);
\fill[slate] (1.83,5.769) circle (0.28pt);
\fill[slate] (1.83,6.225) circle (0.28pt);
\fill[slate] (1.83,6.469) circle (0.28pt);
\fill[slate] (1.83,6.681) circle (0.28pt);
\fill[slate] (1.83,6.717) circle (0.28pt);
\fill[slate] (1.83,5.817) circle (0.28pt);
\fill[slate] (1.83,6.067) circle (0.28pt);
\fill[slate] (1.83,6.267) circle (0.28pt);
\fill[slate] (1.83,6.723) circle (0.28pt);
\fill[slate] (1.83,6.316) circle (0.28pt);
\fill[slate] (1.83,6.565) circle (0.28pt);
\fill[slate] (1.83,6.765) circle (0.28pt);
\fill[slate] (1.83,6.814) circle (0.28pt);
\fill[slate] (1.83,3.422) circle (0.28pt);
\fill[slate] (1.83,3.671) circle (0.28pt);
\fill[slate] (1.83,3.872) circle (0.28pt);
\fill[slate] (1.83,4.572) circle (0.28pt);
\fill[slate] (1.83,4.784) circle (0.28pt);
\fill[slate] (1.83,5.24) circle (0.28pt);
\fill[slate] (1.83,5.484) circle (0.28pt);
\fill[slate] (1.83,5.696) circle (0.28pt);
\fill[slate] (1.83,6.152) circle (0.28pt);
\fill[slate] (1.83,6.396) circle (0.28pt);
\fill[slate] (1.83,6.609) circle (0.28pt);
\fill[slate] (1.83,6.644) circle (0.28pt);
\fill[slate] (1.83,5.981) circle (0.28pt);
\fill[slate] (1.83,6.183) circle (0.28pt);
\fill[slate] (1.83,6.639) circle (0.28pt);
\fill[slate] (1.83,6.23) circle (0.28pt);
\fill[slate] (1.83,6.479) circle (0.28pt);
\fill[slate] (1.83,6.679) circle (0.28pt);
\fill[slate] (1.83,6.728) circle (0.28pt);
\fill[slate] (1.83,4.82) circle (0.28pt);
\fill[slate] (1.83,5.069) circle (0.28pt);
\fill[slate] (1.83,5.271) circle (0.28pt);
\fill[slate] (1.83,5.971) circle (0.28pt);
\fill[slate] (1.83,6.183) circle (0.28pt);
\fill[slate] (1.83,6.639) circle (0.28pt);
\fill[slate] (1.83,6.219) circle (0.28pt);
\fill[slate] (1.83,6.468) circle (0.28pt);
\fill[slate] (1.83,6.67) circle (0.28pt);
\fill[slate] (1.83,6.717) circle (0.28pt);
\fill[slate] (1.83,5.318) circle (0.28pt);
\fill[slate] (1.83,5.567) circle (0.28pt);
\fill[slate] (1.83,5.768) circle (0.28pt);
\fill[slate] (1.83,6.224) circle (0.28pt);
\fill[slate] (1.83,6.467) circle (0.28pt);
\fill[slate] (1.83,6.68) circle (0.28pt);
\fill[slate] (1.83,6.715) circle (0.28pt);
\fill[slate] (1.83,5.816) circle (0.28pt);
\fill[slate] (1.83,6.065) circle (0.28pt);
\fill[slate] (1.83,6.265) circle (0.28pt);
\fill[slate] (1.83,6.721) circle (0.28pt);
\fill[slate] (1.83,6.314) circle (0.28pt);
\fill[slate] (1.83,6.563) circle (0.28pt);
\fill[slate] (1.83,6.763) circle (0.28pt);
\fill[slate] (1.83,6.812) circle (0.28pt);
\fill[slate] (1.83,3.92) circle (0.28pt);
\fill[slate] (1.83,4.169) circle (0.28pt);
\fill[slate] (1.83,4.37) circle (0.28pt);
\fill[slate] (1.83,4.826) circle (0.28pt);
\fill[slate] (1.83,5.07) circle (0.28pt);
\fill[slate] (1.83,5.282) circle (0.28pt);
\fill[slate] (1.83,5.738) circle (0.28pt);
\fill[slate] (1.83,5.981) circle (0.28pt);
\fill[slate] (1.83,6.194) circle (0.28pt);
\fill[slate] (1.83,6.65) circle (0.28pt);
\fill[slate] (1.83,6.23) circle (0.28pt);
\fill[slate] (1.83,6.478) circle (0.28pt);
\fill[slate] (1.83,6.681) circle (0.28pt);
\fill[slate] (1.83,6.727) circle (0.28pt);
\fill[slate] (1.83,5.318) circle (0.28pt);
\fill[slate] (1.83,5.567) circle (0.28pt);
\fill[slate] (1.83,5.769) circle (0.28pt);
\fill[slate] (1.83,6.225) circle (0.28pt);
\fill[slate] (1.83,6.469) circle (0.28pt);
\fill[slate] (1.83,6.681) circle (0.28pt);
\fill[slate] (1.83,6.717) circle (0.28pt);
\fill[slate] (1.83,5.816) circle (0.28pt);
\fill[slate] (1.83,6.065) circle (0.28pt);
\fill[slate] (1.83,6.265) circle (0.28pt);
\fill[slate] (1.83,6.721) circle (0.28pt);
\fill[slate] (1.83,6.314) circle (0.28pt);
\fill[slate] (1.83,6.563) circle (0.28pt);
\fill[slate] (1.83,6.763) circle (0.28pt);
\fill[slate] (1.83,6.812) circle (0.28pt);
\fill[slate] (1.83,4.418) circle (0.28pt);
\fill[slate] (1.83,4.667) circle (0.28pt);
\fill[slate] (1.83,4.868) circle (0.28pt);
\fill[slate] (1.83,5.324) circle (0.28pt);
\fill[slate] (1.83,5.568) circle (0.28pt);
\fill[slate] (1.83,5.78) circle (0.28pt);
\fill[slate] (1.83,6.236) circle (0.28pt);
\fill[slate] (1.83,6.479) circle (0.28pt);
\fill[slate] (1.83,6.692) circle (0.28pt);
\fill[slate] (1.83,6.727) circle (0.28pt);
\fill[slate] (1.83,5.816) circle (0.28pt);
\fill[slate] (1.83,6.065) circle (0.28pt);
\fill[slate] (1.83,6.267) circle (0.28pt);
\fill[slate] (1.83,6.723) circle (0.28pt);
\fill[slate] (1.83,6.314) circle (0.28pt);
\fill[slate] (1.83,6.563) circle (0.28pt);
\fill[slate] (1.83,6.763) circle (0.28pt);
\fill[slate] (1.83,6.812) circle (0.28pt);
\fill[slate] (1.83,4.916) circle (0.28pt);
\fill[slate] (1.83,5.165) circle (0.28pt);
\fill[slate] (1.83,5.366) circle (0.28pt);
\fill[slate] (1.83,5.822) circle (0.28pt);
\fill[slate] (1.83,6.066) circle (0.28pt);
\fill[slate] (1.83,6.278) circle (0.28pt);
\fill[slate] (1.83,6.734) circle (0.28pt);
\fill[slate] (1.83,6.314) circle (0.28pt);
\fill[slate] (1.83,6.563) circle (0.28pt);
\fill[slate] (1.83,6.765) circle (0.28pt);
\fill[slate] (1.83,6.812) circle (0.28pt);
\fill[slate] (1.83,5.414) circle (0.28pt);
\fill[slate] (1.83,5.663) circle (0.28pt);
\fill[slate] (1.83,5.864) circle (0.28pt);
\fill[slate] (1.83,6.32) circle (0.28pt);
\fill[slate] (1.83,6.564) circle (0.28pt);
\fill[slate] (1.83,6.776) circle (0.28pt);
\fill[slate] (1.83,6.812) circle (0.28pt);
\fill[slate] (1.83,5.913) circle (0.28pt);
\fill[slate] (1.83,6.162) circle (0.28pt);
\fill[slate] (1.83,6.362) circle (0.28pt);
\fill[slate] (1.83,6.818) circle (0.28pt);
\fill[slate] (1.83,6.411) circle (0.28pt);
\fill[slate] (1.83,6.66) circle (0.28pt);
\fill[slate] (1.785,3.376) circle (0.28pt);
\fill[slate] (1.785,3.832) circle (0.28pt);
\fill[slate] (1.785,4.075) circle (0.28pt);
\fill[slate] (1.785,4.288) circle (0.28pt);
\fill[slate] (1.785,4.744) circle (0.28pt);
\fill[slate] (1.785,4.988) circle (0.28pt);
\fill[slate] (1.785,5.2) circle (0.28pt);
\fill[slate] (1.785,5.656) circle (0.28pt);
\fill[slate] (1.785,5.9) circle (0.28pt);
\fill[slate] (1.785,6.112) circle (0.28pt);
\fill[slate] (1.785,6.569) circle (0.28pt);
\fill[slate] (1.785,6.812) circle (0.28pt);
\fill[slate] (1.785,6.148) circle (0.28pt);
\fill[slate] (1.785,6.397) circle (0.28pt);
\fill[slate] (1.785,6.599) circle (0.28pt);
\fill[slate] (1.785,6.646) circle (0.28pt);
\fill[slate] (1.785,5.236) circle (0.28pt);
\fill[slate] (1.785,5.485) circle (0.28pt);
\fill[slate] (1.785,5.687) circle (0.28pt);
\fill[slate] (1.785,6.143) circle (0.28pt);
\fill[slate] (1.785,6.387) circle (0.28pt);
\fill[slate] (1.785,6.599) circle (0.28pt);
\fill[slate] (1.785,6.635) circle (0.28pt);
\fill[slate] (1.785,5.734) circle (0.28pt);
\fill[slate] (1.785,5.983) circle (0.28pt);
\fill[slate] (1.785,6.183) circle (0.28pt);
\fill[slate] (1.785,6.639) circle (0.28pt);
\fill[slate] (1.785,6.232) circle (0.28pt);
\fill[slate] (1.785,6.481) circle (0.28pt);
\fill[slate] (1.785,6.681) circle (0.28pt);
\fill[slate] (1.785,6.73) circle (0.28pt);
\fill[slate] (1.785,4.572) circle (0.28pt);
\fill[slate] (1.785,4.775) circle (0.28pt);
\fill[slate] (1.785,5.231) circle (0.28pt);
\fill[slate] (1.785,5.474) circle (0.28pt);
\fill[slate] (1.785,5.687) circle (0.28pt);
\fill[slate] (1.785,6.143) circle (0.28pt);
\fill[slate] (1.785,6.386) circle (0.28pt);
\fill[slate] (1.785,6.599) circle (0.28pt);
\fill[slate] (1.785,6.634) circle (0.28pt);
\fill[slate] (1.785,5.971) circle (0.28pt);
\fill[slate] (1.785,6.174) circle (0.28pt);
\fill[slate] (1.785,6.63) circle (0.28pt);
\fill[slate] (1.785,6.22) circle (0.28pt);
\fill[slate] (1.785,6.469) circle (0.28pt);
\fill[slate] (1.785,6.67) circle (0.28pt);
\fill[slate] (1.785,6.718) circle (0.28pt);
\fill[slate] (1.785,4.821) circle (0.28pt);
\fill[slate] (1.785,5.07) circle (0.28pt);
\fill[slate] (1.785,5.271) circle (0.28pt);
\fill[slate] (1.785,5.971) circle (0.28pt);
\fill[slate] (1.785,6.183) circle (0.28pt);
\fill[slate] (1.785,6.639) circle (0.28pt);
\fill[slate] (1.785,6.219) circle (0.28pt);
\fill[slate] (1.785,6.468) circle (0.28pt);
\fill[slate] (1.785,6.67) circle (0.28pt);
\fill[slate] (1.785,6.717) circle (0.28pt);
\fill[slate] (1.785,5.319) circle (0.28pt);
\fill[slate] (1.785,5.568) circle (0.28pt);
\fill[slate] (1.785,5.769) circle (0.28pt);
\fill[slate] (1.785,6.225) circle (0.28pt);
\fill[slate] (1.785,6.469) circle (0.28pt);
\fill[slate] (1.785,6.681) circle (0.28pt);
\fill[slate] (1.785,6.717) circle (0.28pt);
\fill[slate] (1.785,5.817) circle (0.28pt);
\fill[slate] (1.785,6.066) circle (0.28pt);
\fill[slate] (1.785,6.267) circle (0.28pt);
\fill[slate] (1.785,6.723) circle (0.28pt);
\fill[slate] (1.785,6.315) circle (0.28pt);
\fill[slate] (1.785,6.564) circle (0.28pt);
\fill[slate] (1.785,6.765) circle (0.28pt);
\fill[slate] (1.785,6.813) circle (0.28pt);
\fill[slate] (1.784,3.661) circle (0.28pt);
\fill[slate] (1.785,4.775) circle (0.28pt);
\fill[slate] (1.785,5.231) circle (0.28pt);
\fill[slate] (1.785,5.474) circle (0.28pt);
\fill[slate] (1.785,5.687) circle (0.28pt);
\fill[slate] (1.785,6.143) circle (0.28pt);
\fill[slate] (1.785,6.387) circle (0.28pt);
\fill[slate] (1.785,6.599) circle (0.28pt);
\fill[slate] (1.785,6.635) circle (0.28pt);
\fill[slate] (1.785,5.971) circle (0.28pt);
\fill[slate] (1.785,6.174) circle (0.28pt);
\fill[slate] (1.785,6.63) circle (0.28pt);
\fill[slate] (1.785,6.22) circle (0.28pt);
\fill[slate] (1.785,6.469) circle (0.28pt);
\fill[slate] (1.785,6.67) circle (0.28pt);
\fill[slate] (1.785,6.718) circle (0.28pt);
\fill[slate] (1.785,5.06) circle (0.28pt);
\fill[slate] (1.785,6.174) circle (0.28pt);
\fill[slate] (1.785,6.63) circle (0.28pt);
\fill[slate] (1.785,6.459) circle (0.28pt);
\fill[slate] (1.785,6.708) circle (0.28pt);
\fill[slate] (1.785,5.309) circle (0.28pt);
\fill[slate] (1.785,5.558) circle (0.28pt);
\fill[slate] (1.785,5.758) circle (0.28pt);
\fill[slate] (1.785,6.214) circle (0.28pt);
\fill[slate] (1.785,6.67) circle (0.28pt);
\fill[slate] (1.785,6.706) circle (0.28pt);
\fill[slate] (1.785,5.807) circle (0.28pt);
\fill[slate] (1.785,6.056) circle (0.28pt);
\fill[slate] (1.785,6.256) circle (0.28pt);
\fill[slate] (1.785,6.712) circle (0.28pt);
\fill[slate] (1.785,6.305) circle (0.28pt);
\fill[slate] (1.785,6.554) circle (0.28pt);
\fill[slate] (1.785,6.754) circle (0.28pt);
\fill[slate] (1.785,6.803) circle (0.28pt);
\fill[slate] (1.784,3.91) circle (0.28pt);
\fill[slate] (1.784,4.159) circle (0.28pt);
\fill[slate] (1.784,4.815) circle (0.28pt);
\fill[slate] (1.784,5.271) circle (0.28pt);
\fill[slate] (1.784,5.727) circle (0.28pt);
\fill[slate] (1.784,5.971) circle (0.28pt);
\fill[slate] (1.784,6.184) circle (0.28pt);
\fill[slate] (1.784,6.64) circle (0.28pt);
\fill[slate] (1.784,6.219) circle (0.28pt);
\fill[slate] (1.784,6.468) circle (0.28pt);
\fill[slate] (1.784,6.67) circle (0.28pt);
\fill[slate] (1.784,6.717) circle (0.28pt);
\fill[slate] (1.784,5.307) circle (0.28pt);
\fill[slate] (1.784,5.556) circle (0.28pt);
\fill[slate] (1.784,6.214) circle (0.28pt);
\fill[slate] (1.784,6.67) circle (0.28pt);
\fill[slate] (1.784,5.805) circle (0.28pt);
\fill[slate] (1.784,6.054) circle (0.28pt);
\fill[slate] (1.784,6.255) circle (0.28pt);
\fill[slate] (1.784,6.711) circle (0.28pt);
\fill[slate] (1.784,6.303) circle (0.28pt);
\fill[slate] (1.784,6.552) circle (0.28pt);
\fill[slate] (1.784,6.753) circle (0.28pt);
\fill[slate] (1.784,6.801) circle (0.28pt);
\fill[slate] (1.784,4.408) circle (0.28pt);
\fill[slate] (1.784,4.657) circle (0.28pt);
\fill[slate] (1.784,4.857) circle (0.28pt);
\fill[slate] (1.784,5.313) circle (0.28pt);
\fill[slate] (1.784,5.557) circle (0.28pt);
\fill[slate] (1.784,5.769) circle (0.28pt);
\fill[slate] (1.784,6.225) circle (0.28pt);
\fill[slate] (1.784,6.469) circle (0.28pt);
\fill[slate] (1.784,6.681) circle (0.28pt);
\fill[slate] (1.784,6.717) circle (0.28pt);
\fill[slate] (1.784,5.805) circle (0.28pt);
\fill[slate] (1.784,6.054) circle (0.28pt);
\fill[slate] (1.784,6.256) circle (0.28pt);
\fill[slate] (1.784,6.712) circle (0.28pt);
\fill[slate] (1.784,6.303) circle (0.28pt);
\fill[slate] (1.784,6.552) circle (0.28pt);
\fill[slate] (1.784,6.753) circle (0.28pt);
\fill[slate] (1.784,6.801) circle (0.28pt);
\fill[slate] (1.784,4.906) circle (0.28pt);
\fill[slate] (1.784,5.155) circle (0.28pt);
\fill[slate] (1.784,5.355) circle (0.28pt);
\fill[slate] (1.784,5.811) circle (0.28pt);
\fill[slate] (1.784,6.055) circle (0.28pt);
\fill[slate] (1.784,6.267) circle (0.28pt);
\fill[slate] (1.784,6.723) circle (0.28pt);
\fill[slate] (1.784,6.303) circle (0.28pt);
\fill[slate] (1.784,6.552) circle (0.28pt);
\fill[slate] (1.784,6.754) circle (0.28pt);
\fill[slate] (1.784,6.801) circle (0.28pt);
\fill[slate] (1.784,5.404) circle (0.28pt);
\fill[slate] (1.784,5.653) circle (0.28pt);
\fill[slate] (1.784,5.853) circle (0.28pt);
\fill[slate] (1.784,6.309) circle (0.28pt);
\fill[slate] (1.784,6.553) circle (0.28pt);
\fill[slate] (1.784,6.765) circle (0.28pt);
\fill[slate] (1.784,6.801) circle (0.28pt);
\fill[slate] (1.784,5.902) circle (0.28pt);
\fill[slate] (1.784,6.151) circle (0.28pt);
\fill[slate] (1.784,6.351) circle (0.28pt);
\fill[slate] (1.784,6.807) circle (0.28pt);
\fill[slate] (1.784,6.4) circle (0.28pt);
\fill[slate] (1.784,6.649) circle (0.28pt);
\fill[slate] (1.784,6.849) circle (0.28pt);
\fill[slate] (1.777,2.511) circle (0.28pt);
\fill[slate] (1.777,2.76) circle (0.28pt);
\fill[slate] (1.778,3.416) circle (0.28pt);
\fill[slate] (1.778,3.872) circle (0.28pt);
\fill[slate] (1.778,4.328) circle (0.28pt);
\fill[slate] (1.778,4.572) circle (0.28pt);
\fill[slate] (1.778,4.785) circle (0.28pt);
\fill[slate] (1.778,5.241) circle (0.28pt);
\fill[slate] (1.778,5.484) circle (0.28pt);
\fill[slate] (1.778,5.697) circle (0.28pt);
\fill[slate] (1.778,6.153) circle (0.28pt);
\fill[slate] (1.778,6.396) circle (0.28pt);
\fill[slate] (1.778,6.609) circle (0.28pt);
\fill[slate] (1.778,6.644) circle (0.28pt);
\fill[slate] (1.778,5.732) circle (0.28pt);
\fill[slate] (1.778,5.981) circle (0.28pt);
\fill[slate] (1.778,6.183) circle (0.28pt);
\fill[slate] (1.778,6.639) circle (0.28pt);
\fill[slate] (1.778,6.23) circle (0.28pt);
\fill[slate] (1.778,6.479) circle (0.28pt);
\fill[slate] (1.778,6.68) circle (0.28pt);
\fill[slate] (1.778,6.728) circle (0.28pt);
\fill[slate] (1.778,4.82) circle (0.28pt);
\fill[slate] (1.778,5.069) circle (0.28pt);
\fill[slate] (1.778,5.271) circle (0.28pt);
\fill[slate] (1.778,5.727) circle (0.28pt);
\fill[slate] (1.778,5.971) circle (0.28pt);
\fill[slate] (1.778,6.183) circle (0.28pt);
\fill[slate] (1.778,6.639) circle (0.28pt);
\fill[slate] (1.778,6.219) circle (0.28pt);
\fill[slate] (1.778,6.468) circle (0.28pt);
\fill[slate] (1.778,6.67) circle (0.28pt);
\fill[slate] (1.778,6.717) circle (0.28pt);
\fill[slate] (1.778,5.318) circle (0.28pt);
\fill[slate] (1.778,5.567) circle (0.28pt);
\fill[slate] (1.778,5.768) circle (0.28pt);
\fill[slate] (1.778,6.224) circle (0.28pt);
\fill[slate] (1.778,6.467) circle (0.28pt);
\fill[slate] (1.778,6.68) circle (0.28pt);
\fill[slate] (1.778,6.715) circle (0.28pt);
\fill[slate] (1.778,5.816) circle (0.28pt);
\fill[slate] (1.778,6.065) circle (0.28pt);
\fill[slate] (1.778,6.265) circle (0.28pt);
\fill[slate] (1.778,6.721) circle (0.28pt);
\fill[slate] (1.778,6.314) circle (0.28pt);
\fill[slate] (1.778,6.563) circle (0.28pt);
\fill[slate] (1.778,6.763) circle (0.28pt);
\fill[slate] (1.778,6.812) circle (0.28pt);
\fill[slate] (1.778,4.815) circle (0.28pt);
\fill[slate] (1.778,5.059) circle (0.28pt);
\fill[slate] (1.778,5.271) circle (0.28pt);
\fill[slate] (1.778,5.728) circle (0.28pt);
\fill[slate] (1.778,5.971) circle (0.28pt);
\fill[slate] (1.778,6.184) circle (0.28pt);
\fill[slate] (1.778,6.64) circle (0.28pt);
\fill[slate] (1.778,6.219) circle (0.28pt);
\fill[slate] (1.778,6.468) circle (0.28pt);
\fill[slate] (1.778,6.67) circle (0.28pt);
\fill[slate] (1.778,6.717) circle (0.28pt);
\fill[slate] (1.778,5.556) circle (0.28pt);
\fill[slate] (1.778,5.758) circle (0.28pt);
\fill[slate] (1.778,6.214) circle (0.28pt);
\fill[slate] (1.778,6.458) circle (0.28pt);
\fill[slate] (1.778,6.67) circle (0.28pt);
\fill[slate] (1.778,6.706) circle (0.28pt);
\fill[slate] (1.778,5.805) circle (0.28pt);
\fill[slate] (1.778,6.054) circle (0.28pt);
\fill[slate] (1.778,6.711) circle (0.28pt);
\fill[slate] (1.778,6.303) circle (0.28pt);
\fill[slate] (1.778,6.552) circle (0.28pt);
\fill[slate] (1.778,6.753) circle (0.28pt);
\fill[slate] (1.778,6.801) circle (0.28pt);
\fill[slate] (1.778,4.406) circle (0.28pt);
\fill[slate] (1.778,4.655) circle (0.28pt);
\fill[slate] (1.778,5.312) circle (0.28pt);
\fill[slate] (1.778,5.768) circle (0.28pt);
\fill[slate] (1.778,6.224) circle (0.28pt);
\fill[slate] (1.778,6.467) circle (0.28pt);
\fill[slate] (1.778,6.68) circle (0.28pt);
\fill[slate] (1.778,6.715) circle (0.28pt);
\fill[slate] (1.778,6.255) circle (0.28pt);
\fill[slate] (1.778,6.711) circle (0.28pt);
\fill[slate] (1.778,6.301) circle (0.28pt);
\fill[slate] (1.778,6.55) circle (0.28pt);
\fill[slate] (1.778,6.799) circle (0.28pt);
\fill[slate] (1.778,4.904) circle (0.28pt);
\fill[slate] (1.778,5.153) circle (0.28pt);
\fill[slate] (1.778,5.354) circle (0.28pt);
\fill[slate] (1.778,5.81) circle (0.28pt);
\fill[slate] (1.778,6.266) circle (0.28pt);
\fill[slate] (1.778,6.722) circle (0.28pt);
\fill[slate] (1.778,6.301) circle (0.28pt);
\fill[slate] (1.778,6.55) circle (0.28pt);
\fill[slate] (1.778,6.753) circle (0.28pt);
\fill[slate] (1.778,6.799) circle (0.28pt);
\fill[slate] (1.778,5.402) circle (0.28pt);
\fill[slate] (1.778,5.651) circle (0.28pt);
\fill[slate] (1.778,5.852) circle (0.28pt);
\fill[slate] (1.778,6.308) circle (0.28pt);
\fill[slate] (1.778,6.551) circle (0.28pt);
\fill[slate] (1.778,6.764) circle (0.28pt);
\fill[slate] (1.778,6.799) circle (0.28pt);
\fill[slate] (1.778,5.9) circle (0.28pt);
\fill[slate] (1.778,6.149) circle (0.28pt);
\fill[slate] (1.778,6.35) circle (0.28pt);
\fill[slate] (1.778,6.806) circle (0.28pt);
\fill[slate] (1.778,6.398) circle (0.28pt);
\fill[slate] (1.778,6.647) circle (0.28pt);
\fill[slate] (1.778,6.848) circle (0.28pt);
\fill[slate] (1.777,3.009) circle (0.28pt);
\fill[slate] (1.777,3.258) circle (0.28pt);
\fill[slate] (1.777,3.458) circle (0.28pt);
\fill[slate] (1.777,3.914) circle (0.28pt);
\fill[slate] (1.777,4.37) circle (0.28pt);
\fill[slate] (1.777,4.826) circle (0.28pt);
\fill[slate] (1.777,5.07) circle (0.28pt);
\fill[slate] (1.777,5.282) circle (0.28pt);
\fill[slate] (1.777,5.738) circle (0.28pt);
\fill[slate] (1.777,5.982) circle (0.28pt);
\fill[slate] (1.777,6.194) circle (0.28pt);
\fill[slate] (1.777,6.651) circle (0.28pt);
\fill[slate] (1.777,6.23) circle (0.28pt);
\fill[slate] (1.777,6.479) circle (0.28pt);
\fill[slate] (1.777,6.681) circle (0.28pt);
\fill[slate] (1.777,6.728) circle (0.28pt);
\fill[slate] (1.777,5.318) circle (0.28pt);
\fill[slate] (1.777,5.567) circle (0.28pt);
\fill[slate] (1.777,5.769) circle (0.28pt);
\fill[slate] (1.777,6.225) circle (0.28pt);
\fill[slate] (1.777,6.469) circle (0.28pt);
\fill[slate] (1.777,6.681) circle (0.28pt);
\fill[slate] (1.777,6.717) circle (0.28pt);
\fill[slate] (1.777,5.816) circle (0.28pt);
\fill[slate] (1.777,6.065) circle (0.28pt);
\fill[slate] (1.777,6.265) circle (0.28pt);
\fill[slate] (1.777,6.721) circle (0.28pt);
\fill[slate] (1.777,6.314) circle (0.28pt);
\fill[slate] (1.777,6.563) circle (0.28pt);
\fill[slate] (1.777,6.763) circle (0.28pt);
\fill[slate] (1.777,6.812) circle (0.28pt);
\fill[slate] (1.777,4.406) circle (0.28pt);
\fill[slate] (1.777,4.655) circle (0.28pt);
\fill[slate] (1.777,4.857) circle (0.28pt);
\fill[slate] (1.777,5.313) circle (0.28pt);
\fill[slate] (1.777,5.557) circle (0.28pt);
\fill[slate] (1.777,5.769) circle (0.28pt);
\fill[slate] (1.777,6.225) circle (0.28pt);
\fill[slate] (1.777,6.469) circle (0.28pt);
\fill[slate] (1.777,6.681) circle (0.28pt);
\fill[slate] (1.777,6.717) circle (0.28pt);
\fill[slate] (1.777,5.805) circle (0.28pt);
\fill[slate] (1.777,6.256) circle (0.28pt);
\fill[slate] (1.777,6.712) circle (0.28pt);
\fill[slate] (1.777,6.303) circle (0.28pt);
\fill[slate] (1.777,6.552) circle (0.28pt);
\fill[slate] (1.777,6.753) circle (0.28pt);
\fill[slate] (1.777,6.801) circle (0.28pt);
\fill[slate] (1.777,4.904) circle (0.28pt);
\fill[slate] (1.777,5.153) circle (0.28pt);
\fill[slate] (1.777,5.354) circle (0.28pt);
\fill[slate] (1.777,5.81) circle (0.28pt);
\fill[slate] (1.777,6.266) circle (0.28pt);
\fill[slate] (1.777,6.722) circle (0.28pt);
\fill[slate] (1.777,6.301) circle (0.28pt);
\fill[slate] (1.777,6.55) circle (0.28pt);
\fill[slate] (1.777,6.753) circle (0.28pt);
\fill[slate] (1.777,6.799) circle (0.28pt);
\fill[slate] (1.777,5.402) circle (0.28pt);
\fill[slate] (1.777,5.651) circle (0.28pt);
\fill[slate] (1.777,5.851) circle (0.28pt);
\fill[slate] (1.777,6.307) circle (0.28pt);
\fill[slate] (1.777,6.551) circle (0.28pt);
\fill[slate] (1.777,6.763) circle (0.28pt);
\fill[slate] (1.777,6.799) circle (0.28pt);
\fill[slate] (1.777,5.9) circle (0.28pt);
\fill[slate] (1.777,6.149) circle (0.28pt);
\fill[slate] (1.777,6.349) circle (0.28pt);
\fill[slate] (1.777,6.805) circle (0.28pt);
\fill[slate] (1.777,6.398) circle (0.28pt);
\fill[slate] (1.777,6.647) circle (0.28pt);
\fill[slate] (1.777,6.847) circle (0.28pt);
\fill[slate] (1.777,3.507) circle (0.28pt);
\fill[slate] (1.777,3.756) circle (0.28pt);
\fill[slate] (1.777,3.956) circle (0.28pt);
\fill[slate] (1.777,4.412) circle (0.28pt);
\fill[slate] (1.777,4.656) circle (0.28pt);
\fill[slate] (1.777,4.868) circle (0.28pt);
\fill[slate] (1.777,5.324) circle (0.28pt);
\fill[slate] (1.777,5.568) circle (0.28pt);
\fill[slate] (1.777,5.78) circle (0.28pt);
\fill[slate] (1.777,6.236) circle (0.28pt);
\fill[slate] (1.777,6.48) circle (0.28pt);
\fill[slate] (1.777,6.692) circle (0.28pt);
\fill[slate] (1.777,6.728) circle (0.28pt);
\fill[slate] (1.777,5.816) circle (0.28pt);
\fill[slate] (1.777,6.065) circle (0.28pt);
\fill[slate] (1.777,6.267) circle (0.28pt);
\fill[slate] (1.777,6.723) circle (0.28pt);
\fill[slate] (1.777,6.314) circle (0.28pt);
\fill[slate] (1.777,6.563) circle (0.28pt);
\fill[slate] (1.777,6.763) circle (0.28pt);
\fill[slate] (1.777,6.812) circle (0.28pt);
\fill[slate] (1.777,4.904) circle (0.28pt);
\fill[slate] (1.777,5.153) circle (0.28pt);
\fill[slate] (1.777,5.355) circle (0.28pt);
\fill[slate] (1.777,5.811) circle (0.28pt);
\fill[slate] (1.777,6.055) circle (0.28pt);
\fill[slate] (1.777,6.267) circle (0.28pt);
\fill[slate] (1.777,6.723) circle (0.28pt);
\fill[slate] (1.777,6.303) circle (0.28pt);
\fill[slate] (1.777,6.552) circle (0.28pt);
\fill[slate] (1.777,6.754) circle (0.28pt);
\fill[slate] (1.777,6.801) circle (0.28pt);
\fill[slate] (1.777,5.402) circle (0.28pt);
\fill[slate] (1.777,5.651) circle (0.28pt);
\fill[slate] (1.777,5.852) circle (0.28pt);
\fill[slate] (1.777,6.308) circle (0.28pt);
\fill[slate] (1.777,6.551) circle (0.28pt);
\fill[slate] (1.777,6.764) circle (0.28pt);
\fill[slate] (1.777,6.799) circle (0.28pt);
\fill[slate] (1.777,5.9) circle (0.28pt);
\fill[slate] (1.777,6.149) circle (0.28pt);
\fill[slate] (1.777,6.349) circle (0.28pt);
\fill[slate] (1.777,6.805) circle (0.28pt);
\fill[slate] (1.777,6.398) circle (0.28pt);
\fill[slate] (1.777,6.647) circle (0.28pt);
\fill[slate] (1.777,6.847) circle (0.28pt);
\fill[slate] (1.777,4.005) circle (0.28pt);
\fill[slate] (1.777,4.254) circle (0.28pt);
\fill[slate] (1.777,4.454) circle (0.28pt);
\fill[slate] (1.777,4.91) circle (0.28pt);
\fill[slate] (1.777,5.154) circle (0.28pt);
\fill[slate] (1.777,5.366) circle (0.28pt);
\fill[slate] (1.777,5.822) circle (0.28pt);
\fill[slate] (1.777,6.066) circle (0.28pt);
\fill[slate] (1.777,6.278) circle (0.28pt);
\fill[slate] (1.777,6.734) circle (0.28pt);
\fill[slate] (1.777,6.314) circle (0.28pt);
\fill[slate] (1.777,6.563) circle (0.28pt);
\fill[slate] (1.777,6.765) circle (0.28pt);
\fill[slate] (1.777,6.812) circle (0.28pt);
\fill[slate] (1.777,5.402) circle (0.28pt);
\fill[slate] (1.777,5.651) circle (0.28pt);
\fill[slate] (1.777,5.853) circle (0.28pt);
\fill[slate] (1.777,6.309) circle (0.28pt);
\fill[slate] (1.777,6.553) circle (0.28pt);
\fill[slate] (1.777,6.765) circle (0.28pt);
\fill[slate] (1.777,6.801) circle (0.28pt);
\fill[slate] (1.777,5.9) circle (0.28pt);
\fill[slate] (1.777,6.149) circle (0.28pt);
\fill[slate] (1.777,6.35) circle (0.28pt);
\fill[slate] (1.777,6.806) circle (0.28pt);
\fill[slate] (1.777,6.398) circle (0.28pt);
\fill[slate] (1.777,6.647) circle (0.28pt);
\fill[slate] (1.777,6.847) circle (0.28pt);
\fill[slate] (1.777,4.503) circle (0.28pt);
\fill[slate] (1.777,4.752) circle (0.28pt);
\fill[slate] (1.777,4.952) circle (0.28pt);
\fill[slate] (1.777,5.408) circle (0.28pt);
\fill[slate] (1.777,5.652) circle (0.28pt);
\fill[slate] (1.777,5.864) circle (0.28pt);
\fill[slate] (1.777,6.32) circle (0.28pt);
\fill[slate] (1.777,6.564) circle (0.28pt);
\fill[slate] (1.777,6.776) circle (0.28pt);
\fill[slate] (1.777,6.812) circle (0.28pt);
\fill[slate] (1.777,5.9) circle (0.28pt);
\fill[slate] (1.777,6.149) circle (0.28pt);
\fill[slate] (1.777,6.351) circle (0.28pt);
\fill[slate] (1.777,6.807) circle (0.28pt);
\fill[slate] (1.777,6.398) circle (0.28pt);
\fill[slate] (1.777,6.647) circle (0.28pt);
\fill[slate] (1.777,6.848) circle (0.28pt);
\fill[slate] (1.777,5.001) circle (0.28pt);
\fill[slate] (1.777,5.25) circle (0.28pt);
\fill[slate] (1.777,5.45) circle (0.28pt);
\fill[slate] (1.777,5.906) circle (0.28pt);
\fill[slate] (1.777,6.15) circle (0.28pt);
\fill[slate] (1.777,6.362) circle (0.28pt);
\fill[slate] (1.777,6.818) circle (0.28pt);
\fill[slate] (1.777,6.398) circle (0.28pt);
\fill[slate] (1.777,6.647) circle (0.28pt);
\fill[slate] (1.777,6.849) circle (0.28pt);
\fill[slate] (1.777,5.499) circle (0.28pt);
\fill[slate] (1.777,5.748) circle (0.28pt);
\fill[slate] (1.777,5.948) circle (0.28pt);
\fill[slate] (1.777,6.404) circle (0.28pt);
\fill[slate] (1.777,6.648) circle (0.28pt);
\fill[slate] (1.777,5.997) circle (0.28pt);
\fill[slate] (1.777,6.246) circle (0.28pt);
\fill[slate] (1.777,6.446) circle (0.28pt);
\fill[slate] (1.777,6.495) circle (0.28pt);
\fill[slate] (1.777,6.744) circle (0.28pt);
\fill[slate] (0.267,2.018) circle (0.28pt);
\fill[slate] (0.268,2.474) circle (0.28pt);
\fill[slate] (0.268,2.93) circle (0.28pt);
\fill[slate] (0.268,3.173) circle (0.28pt);
\fill[slate] (0.268,3.386) circle (0.28pt);
\fill[slate] (0.268,3.842) circle (0.28pt);
\fill[slate] (0.268,4.086) circle (0.28pt);
\fill[slate] (0.268,4.298) circle (0.28pt);
\fill[slate] (0.268,4.754) circle (0.28pt);
\fill[slate] (0.268,4.998) circle (0.28pt);
\fill[slate] (0.268,5.211) circle (0.28pt);
\fill[slate] (0.268,5.667) circle (0.28pt);
\fill[slate] (0.268,5.91) circle (0.28pt);
\fill[slate] (0.268,6.123) circle (0.28pt);
\fill[slate] (0.268,6.579) circle (0.28pt);
\fill[slate] (0.268,6.822) circle (0.28pt);
\fill[slate] (0.268,6.158) circle (0.28pt);
\fill[slate] (0.268,6.407) circle (0.28pt);
\fill[slate] (0.268,6.609) circle (0.28pt);
\fill[slate] (0.268,6.656) circle (0.28pt);
\fill[slate] (0.268,5.246) circle (0.28pt);
\fill[slate] (0.268,5.495) circle (0.28pt);
\fill[slate] (0.268,5.697) circle (0.28pt);
\fill[slate] (0.268,6.153) circle (0.28pt);
\fill[slate] (0.268,6.397) circle (0.28pt);
\fill[slate] (0.268,6.609) circle (0.28pt);
\fill[slate] (0.268,6.645) circle (0.28pt);
\fill[slate] (0.268,5.744) circle (0.28pt);
\fill[slate] (0.268,5.993) circle (0.28pt);
\fill[slate] (0.268,6.194) circle (0.28pt);
\fill[slate] (0.268,6.65) circle (0.28pt);
\fill[slate] (0.268,6.242) circle (0.28pt);
\fill[slate] (0.268,6.491) circle (0.28pt);
\fill[slate] (0.268,6.691) circle (0.28pt);
\fill[slate] (0.268,6.74) circle (0.28pt);
\fill[slate] (0.268,4.334) circle (0.28pt);
\fill[slate] (0.268,4.583) circle (0.28pt);
\fill[slate] (0.268,4.785) circle (0.28pt);
\fill[slate] (0.268,5.241) circle (0.28pt);
\fill[slate] (0.268,5.485) circle (0.28pt);
\fill[slate] (0.268,5.697) circle (0.28pt);
\fill[slate] (0.268,6.153) circle (0.28pt);
\fill[slate] (0.268,6.396) circle (0.28pt);
\fill[slate] (0.268,6.609) circle (0.28pt);
\fill[slate] (0.268,6.645) circle (0.28pt);
\fill[slate] (0.268,5.733) circle (0.28pt);
\fill[slate] (0.268,5.982) circle (0.28pt);
\fill[slate] (0.268,6.184) circle (0.28pt);
\fill[slate] (0.268,6.64) circle (0.28pt);
\fill[slate] (0.268,6.231) circle (0.28pt);
\fill[slate] (0.268,6.48) circle (0.28pt);
\fill[slate] (0.268,6.68) circle (0.28pt);
\fill[slate] (0.268,6.729) circle (0.28pt);
\fill[slate] (0.268,4.832) circle (0.28pt);
\fill[slate] (0.268,5.081) circle (0.28pt);
\fill[slate] (0.268,5.281) circle (0.28pt);
\fill[slate] (0.268,5.737) circle (0.28pt);
\fill[slate] (0.268,5.981) circle (0.28pt);
\fill[slate] (0.268,6.193) circle (0.28pt);
\fill[slate] (0.268,6.65) circle (0.28pt);
\fill[slate] (0.268,6.229) circle (0.28pt);
\fill[slate] (0.268,6.478) circle (0.28pt);
\fill[slate] (0.268,6.68) circle (0.28pt);
\fill[slate] (0.268,6.727) circle (0.28pt);
\fill[slate] (0.268,5.33) circle (0.28pt);
\fill[slate] (0.268,5.579) circle (0.28pt);
\fill[slate] (0.268,5.779) circle (0.28pt);
\fill[slate] (0.268,6.235) circle (0.28pt);
\fill[slate] (0.268,6.479) circle (0.28pt);
\fill[slate] (0.268,6.691) circle (0.28pt);
\fill[slate] (0.268,6.727) circle (0.28pt);
\fill[slate] (0.268,5.828) circle (0.28pt);
\fill[slate] (0.268,6.077) circle (0.28pt);
\fill[slate] (0.268,6.277) circle (0.28pt);
\fill[slate] (0.268,6.733) circle (0.28pt);
\fill[slate] (0.268,6.326) circle (0.28pt);
\fill[slate] (0.268,6.575) circle (0.28pt);
\fill[slate] (0.268,6.775) circle (0.28pt);
\fill[slate] (0.268,6.824) circle (0.28pt);
\fill[slate] (0.267,3.422) circle (0.28pt);
\fill[slate] (0.267,3.671) circle (0.28pt);
\fill[slate] (0.268,3.873) circle (0.28pt);
\fill[slate] (0.268,4.329) circle (0.28pt);
\fill[slate] (0.268,4.572) circle (0.28pt);
\fill[slate] (0.268,4.785) circle (0.28pt);
\fill[slate] (0.268,5.241) circle (0.28pt);
\fill[slate] (0.268,5.484) circle (0.28pt);
\fill[slate] (0.268,5.697) circle (0.28pt);
\fill[slate] (0.268,6.153) circle (0.28pt);
\fill[slate] (0.268,6.396) circle (0.28pt);
\fill[slate] (0.268,6.609) circle (0.28pt);
\fill[slate] (0.268,6.645) circle (0.28pt);
\fill[slate] (0.268,5.732) circle (0.28pt);
\fill[slate] (0.268,5.981) circle (0.28pt);
\fill[slate] (0.268,6.184) circle (0.28pt);
\fill[slate] (0.268,6.64) circle (0.28pt);
\fill[slate] (0.268,6.23) circle (0.28pt);
\fill[slate] (0.268,6.479) circle (0.28pt);
\fill[slate] (0.268,6.68) circle (0.28pt);
\fill[slate] (0.268,6.728) circle (0.28pt);
\fill[slate] (0.268,4.821) circle (0.28pt);
\fill[slate] (0.268,5.07) circle (0.28pt);
\fill[slate] (0.268,5.272) circle (0.28pt);
\fill[slate] (0.268,5.728) circle (0.28pt);
\fill[slate] (0.268,5.971) circle (0.28pt);
\fill[slate] (0.268,6.184) circle (0.28pt);
\fill[slate] (0.268,6.64) circle (0.28pt);
\fill[slate] (0.268,6.22) circle (0.28pt);
\fill[slate] (0.268,6.469) circle (0.28pt);
\fill[slate] (0.268,6.671) circle (0.28pt);
\fill[slate] (0.268,6.718) circle (0.28pt);
\fill[slate] (0.268,5.319) circle (0.28pt);
\fill[slate] (0.268,5.568) circle (0.28pt);
\fill[slate] (0.268,5.768) circle (0.28pt);
\fill[slate] (0.268,6.224) circle (0.28pt);
\fill[slate] (0.268,6.468) circle (0.28pt);
\fill[slate] (0.268,6.68) circle (0.28pt);
\fill[slate] (0.268,6.716) circle (0.28pt);
\fill[slate] (0.268,5.817) circle (0.28pt);
\fill[slate] (0.268,6.066) circle (0.28pt);
\fill[slate] (0.268,6.266) circle (0.28pt);
\fill[slate] (0.268,6.722) circle (0.28pt);
\fill[slate] (0.268,6.315) circle (0.28pt);
\fill[slate] (0.268,6.564) circle (0.28pt);
\fill[slate] (0.268,6.764) circle (0.28pt);
\fill[slate] (0.268,6.813) circle (0.28pt);
\fill[slate] (0.267,3.92) circle (0.28pt);
\fill[slate] (0.267,4.169) circle (0.28pt);
\fill[slate] (0.267,4.369) circle (0.28pt);
\fill[slate] (0.267,4.825) circle (0.28pt);
\fill[slate] (0.267,5.069) circle (0.28pt);
\fill[slate] (0.267,5.281) circle (0.28pt);
\fill[slate] (0.267,5.737) circle (0.28pt);
\fill[slate] (0.267,5.981) circle (0.28pt);
\fill[slate] (0.267,6.193) circle (0.28pt);
\fill[slate] (0.267,6.65) circle (0.28pt);
\fill[slate] (0.267,6.229) circle (0.28pt);
\fill[slate] (0.267,6.478) circle (0.28pt);
\fill[slate] (0.267,6.68) circle (0.28pt);
\fill[slate] (0.267,6.727) circle (0.28pt);
\fill[slate] (0.267,5.317) circle (0.28pt);
\fill[slate] (0.267,5.566) circle (0.28pt);
\fill[slate] (0.267,5.768) circle (0.28pt);
\fill[slate] (0.267,6.224) circle (0.28pt);
\fill[slate] (0.267,6.468) circle (0.28pt);
\fill[slate] (0.267,6.68) circle (0.28pt);
\fill[slate] (0.267,6.716) circle (0.28pt);
\fill[slate] (0.267,5.815) circle (0.28pt);
\fill[slate] (0.267,6.064) circle (0.28pt);
\fill[slate] (0.267,6.265) circle (0.28pt);
\fill[slate] (0.267,6.721) circle (0.28pt);
\fill[slate] (0.267,6.313) circle (0.28pt);
\fill[slate] (0.267,6.562) circle (0.28pt);
\fill[slate] (0.267,6.763) circle (0.28pt);
\fill[slate] (0.267,6.811) circle (0.28pt);
\fill[slate] (0.267,4.418) circle (0.28pt);
\fill[slate] (0.267,4.667) circle (0.28pt);
\fill[slate] (0.267,4.867) circle (0.28pt);
\fill[slate] (0.267,5.323) circle (0.28pt);
\fill[slate] (0.267,5.567) circle (0.28pt);
\fill[slate] (0.267,5.779) circle (0.28pt);
\fill[slate] (0.267,6.235) circle (0.28pt);
\fill[slate] (0.267,6.478) circle (0.28pt);
\fill[slate] (0.267,6.691) circle (0.28pt);
\fill[slate] (0.267,6.727) circle (0.28pt);
\fill[slate] (0.267,5.815) circle (0.28pt);
\fill[slate] (0.267,6.064) circle (0.28pt);
\fill[slate] (0.267,6.266) circle (0.28pt);
\fill[slate] (0.267,6.722) circle (0.28pt);
\fill[slate] (0.267,6.313) circle (0.28pt);
\fill[slate] (0.267,6.562) circle (0.28pt);
\fill[slate] (0.267,6.763) circle (0.28pt);
\fill[slate] (0.267,6.811) circle (0.28pt);
\fill[slate] (0.267,4.916) circle (0.28pt);
\fill[slate] (0.267,5.165) circle (0.28pt);
\fill[slate] (0.267,5.365) circle (0.28pt);
\fill[slate] (0.267,5.821) circle (0.28pt);
\fill[slate] (0.267,6.065) circle (0.28pt);
\fill[slate] (0.267,6.277) circle (0.28pt);
\fill[slate] (0.267,6.733) circle (0.28pt);
\fill[slate] (0.267,6.313) circle (0.28pt);
\fill[slate] (0.267,6.562) circle (0.28pt);
\fill[slate] (0.267,6.764) circle (0.28pt);
\fill[slate] (0.267,6.811) circle (0.28pt);
\fill[slate] (0.267,5.414) circle (0.28pt);
\fill[slate] (0.267,5.663) circle (0.28pt);
\fill[slate] (0.267,5.863) circle (0.28pt);
\fill[slate] (0.267,6.319) circle (0.28pt);
\fill[slate] (0.267,6.563) circle (0.28pt);
\fill[slate] (0.267,6.775) circle (0.28pt);
\fill[slate] (0.267,6.811) circle (0.28pt);
\fill[slate] (0.267,5.912) circle (0.28pt);
\fill[slate] (0.267,6.161) circle (0.28pt);
\fill[slate] (0.267,6.361) circle (0.28pt);
\fill[slate] (0.267,6.817) circle (0.28pt);
\fill[slate] (0.267,6.41) circle (0.28pt);
\fill[slate] (0.267,6.659) circle (0.28pt);
\fill[slate] (0.261,2.961) circle (0.28pt);
\fill[slate] (0.261,3.417) circle (0.28pt);
\fill[slate] (0.261,3.661) circle (0.28pt);
\fill[slate] (0.261,3.873) circle (0.28pt);
\fill[slate] (0.261,4.329) circle (0.28pt);
\fill[slate] (0.261,4.573) circle (0.28pt);
\fill[slate] (0.261,4.785) circle (0.28pt);
\fill[slate] (0.261,5.241) circle (0.28pt);
\fill[slate] (0.261,5.485) circle (0.28pt);
\fill[slate] (0.261,5.697) circle (0.28pt);
\fill[slate] (0.261,6.154) circle (0.28pt);
\fill[slate] (0.261,6.397) circle (0.28pt);
\fill[slate] (0.261,6.61) circle (0.28pt);
\fill[slate] (0.261,6.645) circle (0.28pt);
\fill[slate] (0.261,5.733) circle (0.28pt);
\fill[slate] (0.261,5.982) circle (0.28pt);
\fill[slate] (0.261,6.184) circle (0.28pt);
\fill[slate] (0.261,6.64) circle (0.28pt);
\fill[slate] (0.261,6.231) circle (0.28pt);
\fill[slate] (0.261,6.48) circle (0.28pt);
\fill[slate] (0.261,6.681) circle (0.28pt);
\fill[slate] (0.261,6.729) circle (0.28pt);
\fill[slate] (0.261,4.821) circle (0.28pt);
\fill[slate] (0.261,5.07) circle (0.28pt);
\fill[slate] (0.261,5.272) circle (0.28pt);
\fill[slate] (0.261,5.728) circle (0.28pt);
\fill[slate] (0.261,5.971) circle (0.28pt);
\fill[slate] (0.261,6.184) circle (0.28pt);
\fill[slate] (0.261,6.64) circle (0.28pt);
\fill[slate] (0.261,6.22) circle (0.28pt);
\fill[slate] (0.261,6.469) circle (0.28pt);
\fill[slate] (0.261,6.671) circle (0.28pt);
\fill[slate] (0.261,6.718) circle (0.28pt);
\fill[slate] (0.261,5.319) circle (0.28pt);
\fill[slate] (0.261,5.568) circle (0.28pt);
\fill[slate] (0.261,5.768) circle (0.28pt);
\fill[slate] (0.261,6.224) circle (0.28pt);
\fill[slate] (0.261,6.468) circle (0.28pt);
\fill[slate] (0.261,6.68) circle (0.28pt);
\fill[slate] (0.261,6.716) circle (0.28pt);
\fill[slate] (0.261,5.817) circle (0.28pt);
\fill[slate] (0.261,6.066) circle (0.28pt);
\fill[slate] (0.261,6.266) circle (0.28pt);
\fill[slate] (0.261,6.722) circle (0.28pt);
\fill[slate] (0.261,6.315) circle (0.28pt);
\fill[slate] (0.261,6.564) circle (0.28pt);
\fill[slate] (0.261,6.764) circle (0.28pt);
\fill[slate] (0.261,6.813) circle (0.28pt);
\fill[slate] (0.261,3.909) circle (0.28pt);
\fill[slate] (0.261,4.36) circle (0.28pt);
\fill[slate] (0.261,4.816) circle (0.28pt);
\fill[slate] (0.261,5.06) circle (0.28pt);
\fill[slate] (0.261,5.272) circle (0.28pt);
\fill[slate] (0.261,5.728) circle (0.28pt);
\fill[slate] (0.261,5.971) circle (0.28pt);
\fill[slate] (0.261,6.184) circle (0.28pt);
\fill[slate] (0.261,6.64) circle (0.28pt);
\fill[slate] (0.261,6.22) circle (0.28pt);
\fill[slate] (0.261,6.469) circle (0.28pt);
\fill[slate] (0.261,6.671) circle (0.28pt);
\fill[slate] (0.261,6.718) circle (0.28pt);
\fill[slate] (0.261,5.308) circle (0.28pt);
\fill[slate] (0.261,5.557) circle (0.28pt);
\fill[slate] (0.261,5.759) circle (0.28pt);
\fill[slate] (0.261,6.215) circle (0.28pt);
\fill[slate] (0.261,6.459) circle (0.28pt);
\fill[slate] (0.261,6.671) circle (0.28pt);
\fill[slate] (0.261,6.707) circle (0.28pt);
\fill[slate] (0.261,5.806) circle (0.28pt);
\fill[slate] (0.261,6.055) circle (0.28pt);
\fill[slate] (0.261,6.256) circle (0.28pt);
\fill[slate] (0.261,6.711) circle (0.28pt);
\fill[slate] (0.261,6.304) circle (0.28pt);
\fill[slate] (0.261,6.553) circle (0.28pt);
\fill[slate] (0.261,6.753) circle (0.28pt);
\fill[slate] (0.261,6.802) circle (0.28pt);
\fill[slate] (0.261,4.407) circle (0.28pt);
\fill[slate] (0.261,4.656) circle (0.28pt);
\fill[slate] (0.261,4.857) circle (0.28pt);
\fill[slate] (0.261,5.312) circle (0.28pt);
\fill[slate] (0.261,5.556) circle (0.28pt);
\fill[slate] (0.261,5.768) circle (0.28pt);
\fill[slate] (0.261,6.225) circle (0.28pt);
\fill[slate] (0.261,6.468) circle (0.28pt);
\fill[slate] (0.261,6.681) circle (0.28pt);
\fill[slate] (0.261,6.716) circle (0.28pt);
\fill[slate] (0.261,5.804) circle (0.28pt);
\fill[slate] (0.261,6.256) circle (0.28pt);
\fill[slate] (0.261,6.711) circle (0.28pt);
\fill[slate] (0.261,6.302) circle (0.28pt);
\fill[slate] (0.261,6.551) circle (0.28pt);
\fill[slate] (0.261,6.752) circle (0.28pt);
\fill[slate] (0.261,6.8) circle (0.28pt);
\fill[slate] (0.261,4.905) circle (0.28pt);
\fill[slate] (0.261,5.154) circle (0.28pt);
\fill[slate] (0.261,5.354) circle (0.28pt);
\fill[slate] (0.261,5.81) circle (0.28pt);
\fill[slate] (0.261,6.266) circle (0.28pt);
\fill[slate] (0.261,6.722) circle (0.28pt);
\fill[slate] (0.261,6.302) circle (0.28pt);
\fill[slate] (0.261,6.551) circle (0.28pt);
\fill[slate] (0.261,6.753) circle (0.28pt);
\fill[slate] (0.261,6.8) circle (0.28pt);
\fill[slate] (0.261,5.403) circle (0.28pt);
\fill[slate] (0.261,5.652) circle (0.28pt);
\fill[slate] (0.261,5.852) circle (0.28pt);
\fill[slate] (0.261,6.308) circle (0.28pt);
\fill[slate] (0.261,6.552) circle (0.28pt);
\fill[slate] (0.261,6.764) circle (0.28pt);
\fill[slate] (0.261,6.8) circle (0.28pt);
\fill[slate] (0.261,5.901) circle (0.28pt);
\fill[slate] (0.261,6.15) circle (0.28pt);
\fill[slate] (0.261,6.35) circle (0.28pt);
\fill[slate] (0.261,6.806) circle (0.28pt);
\fill[slate] (0.261,6.399) circle (0.28pt);
\fill[slate] (0.261,6.648) circle (0.28pt);
\fill[slate] (0.261,6.848) circle (0.28pt);
\fill[slate] (0.26,3.257) circle (0.28pt);
\fill[slate] (0.26,3.913) circle (0.28pt);
\fill[slate] (0.26,4.37) circle (0.28pt);
\fill[slate] (0.26,4.826) circle (0.28pt);
\fill[slate] (0.26,5.069) circle (0.28pt);
\fill[slate] (0.26,5.282) circle (0.28pt);
\fill[slate] (0.26,5.738) circle (0.28pt);
\fill[slate] (0.26,5.981) circle (0.28pt);
\fill[slate] (0.26,6.194) circle (0.28pt);
\fill[slate] (0.26,6.65) circle (0.28pt);
\fill[slate] (0.26,6.229) circle (0.28pt);
\fill[slate] (0.26,6.478) circle (0.28pt);
\fill[slate] (0.26,6.681) circle (0.28pt);
\fill[slate] (0.26,6.727) circle (0.28pt);
\fill[slate] (0.26,5.317) circle (0.28pt);
\fill[slate] (0.26,5.566) circle (0.28pt);
\fill[slate] (0.26,5.768) circle (0.28pt);
\fill[slate] (0.26,6.224) circle (0.28pt);
\fill[slate] (0.26,6.468) circle (0.28pt);
\fill[slate] (0.26,6.68) circle (0.28pt);
\fill[slate] (0.26,6.716) circle (0.28pt);
\fill[slate] (0.26,5.815) circle (0.28pt);
\fill[slate] (0.26,6.064) circle (0.28pt);
\fill[slate] (0.26,6.265) circle (0.28pt);
\fill[slate] (0.26,6.721) circle (0.28pt);
\fill[slate] (0.26,6.313) circle (0.28pt);
\fill[slate] (0.26,6.562) circle (0.28pt);
\fill[slate] (0.26,6.763) circle (0.28pt);
\fill[slate] (0.26,6.811) circle (0.28pt);
\fill[slate] (0.26,4.857) circle (0.28pt);
\fill[slate] (0.26,5.312) circle (0.28pt);
\fill[slate] (0.26,5.556) circle (0.28pt);
\fill[slate] (0.26,5.769) circle (0.28pt);
\fill[slate] (0.26,6.225) circle (0.28pt);
\fill[slate] (0.26,6.468) circle (0.28pt);
\fill[slate] (0.26,6.681) circle (0.28pt);
\fill[slate] (0.26,6.716) circle (0.28pt);
\fill[slate] (0.26,5.804) circle (0.28pt);
\fill[slate] (0.26,6.256) circle (0.28pt);
\fill[slate] (0.26,6.711) circle (0.28pt);
\fill[slate] (0.26,6.302) circle (0.28pt);
\fill[slate] (0.26,6.551) circle (0.28pt);
\fill[slate] (0.26,6.752) circle (0.28pt);
\fill[slate] (0.26,6.8) circle (0.28pt);
\fill[slate] (0.26,4.903) circle (0.28pt);
\fill[slate] (0.26,5.152) circle (0.28pt);
\fill[slate] (0.26,5.809) circle (0.28pt);
\fill[slate] (0.26,6.265) circle (0.28pt);
\fill[slate] (0.26,6.721) circle (0.28pt);
\fill[slate] (0.26,6.752) circle (0.28pt);
\fill[slate] (0.26,6.799) circle (0.28pt);
\fill[slate] (0.26,5.401) circle (0.28pt);
\fill[slate] (0.26,5.65) circle (0.28pt);
\fill[slate] (0.26,6.307) circle (0.28pt);
\fill[slate] (0.26,6.55) circle (0.28pt);
\fill[slate] (0.26,6.763) circle (0.28pt);
\fill[slate] (0.26,5.899) circle (0.28pt);
\fill[slate] (0.26,6.148) circle (0.28pt);
\fill[slate] (0.26,6.349) circle (0.28pt);
\fill[slate] (0.26,6.805) circle (0.28pt);
\fill[slate] (0.26,6.397) circle (0.28pt);
\fill[slate] (0.26,6.646) circle (0.28pt);
\fill[slate] (0.26,6.847) circle (0.28pt);
\fill[slate] (0.26,3.506) circle (0.28pt);
\fill[slate] (0.26,3.755) circle (0.28pt);
\fill[slate] (0.26,4.411) circle (0.28pt);
\fill[slate] (0.26,4.655) circle (0.28pt);
\fill[slate] (0.26,4.867) circle (0.28pt);
\fill[slate] (0.26,5.323) circle (0.28pt);
\fill[slate] (0.26,5.567) circle (0.28pt);
\fill[slate] (0.26,5.779) circle (0.28pt);
\fill[slate] (0.26,6.236) circle (0.28pt);
\fill[slate] (0.26,6.479) circle (0.28pt);
\fill[slate] (0.26,6.692) circle (0.28pt);
\fill[slate] (0.26,6.727) circle (0.28pt);
\fill[slate] (0.26,5.815) circle (0.28pt);
\fill[slate] (0.26,6.064) circle (0.28pt);
\fill[slate] (0.26,6.266) circle (0.28pt);
\fill[slate] (0.26,6.722) circle (0.28pt);
\fill[slate] (0.26,6.313) circle (0.28pt);
\fill[slate] (0.26,6.562) circle (0.28pt);
\fill[slate] (0.26,6.763) circle (0.28pt);
\fill[slate] (0.26,6.811) circle (0.28pt);
\fill[slate] (0.26,5.354) circle (0.28pt);
\fill[slate] (0.26,5.81) circle (0.28pt);
\fill[slate] (0.26,6.054) circle (0.28pt);
\fill[slate] (0.26,6.266) circle (0.28pt);
\fill[slate] (0.26,6.722) circle (0.28pt);
\fill[slate] (0.26,6.302) circle (0.28pt);
\fill[slate] (0.26,6.551) circle (0.28pt);
\fill[slate] (0.26,6.753) circle (0.28pt);
\fill[slate] (0.26,6.8) circle (0.28pt);
\fill[slate] (0.26,5.65) circle (0.28pt);
\fill[slate] (0.26,6.307) circle (0.28pt);
\fill[slate] (0.26,6.55) circle (0.28pt);
\fill[slate] (0.26,6.763) circle (0.28pt);
\fill[slate] (0.26,6.798) circle (0.28pt);
\fill[slate] (0.26,5.899) circle (0.28pt);
\fill[slate] (0.26,6.148) circle (0.28pt);
\fill[slate] (0.26,6.804) circle (0.28pt);
\fill[slate] (0.26,6.397) circle (0.28pt);
\fill[slate] (0.26,6.646) circle (0.28pt);
\fill[slate] (0.26,6.847) circle (0.28pt);
\fill[slate] (0.26,4.004) circle (0.28pt);
\fill[slate] (0.26,4.253) circle (0.28pt);
\fill[slate] (0.26,4.453) circle (0.28pt);
\fill[slate] (0.26,4.909) circle (0.28pt);
\fill[slate] (0.26,5.153) circle (0.28pt);
\fill[slate] (0.26,5.365) circle (0.28pt);
\fill[slate] (0.26,5.821) circle (0.28pt);
\fill[slate] (0.26,6.065) circle (0.28pt);
\fill[slate] (0.26,6.277) circle (0.28pt);
\fill[slate] (0.26,6.733) circle (0.28pt);
\fill[slate] (0.26,6.313) circle (0.28pt);
\fill[slate] (0.26,6.562) circle (0.28pt);
\fill[slate] (0.26,6.764) circle (0.28pt);
\fill[slate] (0.26,6.811) circle (0.28pt);
\fill[slate] (0.26,5.65) circle (0.28pt);
\fill[slate] (0.26,5.852) circle (0.28pt);
\fill[slate] (0.26,6.308) circle (0.28pt);
\fill[slate] (0.26,6.552) circle (0.28pt);
\fill[slate] (0.26,6.764) circle (0.28pt);
\fill[slate] (0.26,6.8) circle (0.28pt);
\fill[slate] (0.26,5.899) circle (0.28pt);
\fill[slate] (0.26,6.148) circle (0.28pt);
\fill[slate] (0.26,6.349) circle (0.28pt);
\fill[slate] (0.26,6.805) circle (0.28pt);
\fill[slate] (0.26,6.397) circle (0.28pt);
\fill[slate] (0.26,6.646) circle (0.28pt);
\fill[slate] (0.26,6.847) circle (0.28pt);
\fill[slate] (0.26,4.502) circle (0.28pt);
\fill[slate] (0.26,4.751) circle (0.28pt);
\fill[slate] (0.26,4.951) circle (0.28pt);
\fill[slate] (0.26,5.407) circle (0.28pt);
\fill[slate] (0.26,5.651) circle (0.28pt);
\fill[slate] (0.26,5.863) circle (0.28pt);
\fill[slate] (0.26,6.319) circle (0.28pt);
\fill[slate] (0.26,6.563) circle (0.28pt);
\fill[slate] (0.26,6.775) circle (0.28pt);
\fill[slate] (0.26,6.811) circle (0.28pt);
\fill[slate] (0.26,5.899) circle (0.28pt);
\fill[slate] (0.26,6.148) circle (0.28pt);
\fill[slate] (0.26,6.35) circle (0.28pt);
\fill[slate] (0.26,6.806) circle (0.28pt);
\fill[slate] (0.26,6.397) circle (0.28pt);
\fill[slate] (0.26,6.646) circle (0.28pt);
\fill[slate] (0.26,6.847) circle (0.28pt);
\fill[slate] (0.26,5) circle (0.28pt);
\fill[slate] (0.26,5.249) circle (0.28pt);
\fill[slate] (0.26,5.449) circle (0.28pt);
\fill[slate] (0.26,5.905) circle (0.28pt);
\fill[slate] (0.26,6.149) circle (0.28pt);
\fill[slate] (0.26,6.361) circle (0.28pt);
\fill[slate] (0.26,6.817) circle (0.28pt);
\fill[slate] (0.26,6.397) circle (0.28pt);
\fill[slate] (0.26,6.646) circle (0.28pt);
\fill[slate] (0.26,6.848) circle (0.28pt);
\fill[slate] (0.26,5.498) circle (0.28pt);
\fill[slate] (0.26,5.747) circle (0.28pt);
\fill[slate] (0.26,5.947) circle (0.28pt);
\fill[slate] (0.26,6.403) circle (0.28pt);
\fill[slate] (0.26,6.647) circle (0.28pt);
\fill[slate] (0.26,5.996) circle (0.28pt);
\fill[slate] (0.26,6.245) circle (0.28pt);
\fill[slate] (0.26,6.445) circle (0.28pt);
\fill[slate] (0.26,6.494) circle (0.28pt);
\fill[slate] (0.26,6.743) circle (0.28pt);
\fill[slate] (0.039,2.516) circle (0.28pt);
\fill[slate] (0.038,2.76) circle (0.28pt);
\fill[slate] (0.039,2.972) circle (0.28pt);
\fill[slate] (0.039,3.428) circle (0.28pt);
\fill[slate] (0.039,3.671) circle (0.28pt);
\fill[slate] (0.039,3.884) circle (0.28pt);
\fill[slate] (0.039,4.34) circle (0.28pt);
\fill[slate] (0.039,4.583) circle (0.28pt);
\fill[slate] (0.039,4.796) circle (0.28pt);
\fill[slate] (0.039,5.252) circle (0.28pt);
\fill[slate] (0.039,5.496) circle (0.28pt);
\fill[slate] (0.039,5.708) circle (0.28pt);
\fill[slate] (0.039,6.165) circle (0.28pt);
\fill[slate] (0.039,6.408) circle (0.28pt);
\fill[slate] (0.039,6.621) circle (0.28pt);
\fill[slate] (0.039,6.656) circle (0.28pt);
\fill[slate] (0.039,5.744) circle (0.28pt);
\fill[slate] (0.039,5.993) circle (0.28pt);
\fill[slate] (0.039,6.195) circle (0.28pt);
\fill[slate] (0.039,6.651) circle (0.28pt);
\fill[slate] (0.039,6.242) circle (0.28pt);
\fill[slate] (0.039,6.491) circle (0.28pt);
\fill[slate] (0.039,6.691) circle (0.28pt);
\fill[slate] (0.039,6.74) circle (0.28pt);
\fill[slate] (0.039,4.832) circle (0.28pt);
\fill[slate] (0.039,5.081) circle (0.28pt);
\fill[slate] (0.039,5.283) circle (0.28pt);
\fill[slate] (0.039,5.739) circle (0.28pt);
\fill[slate] (0.039,5.982) circle (0.28pt);
\fill[slate] (0.039,6.195) circle (0.28pt);
\fill[slate] (0.039,6.651) circle (0.28pt);
\fill[slate] (0.039,6.231) circle (0.28pt);
\fill[slate] (0.039,6.48) circle (0.28pt);
\fill[slate] (0.039,6.682) circle (0.28pt);
\fill[slate] (0.039,6.729) circle (0.28pt);
\fill[slate] (0.039,5.33) circle (0.28pt);
\fill[slate] (0.039,5.579) circle (0.28pt);
\fill[slate] (0.039,5.779) circle (0.28pt);
\fill[slate] (0.039,6.235) circle (0.28pt);
\fill[slate] (0.039,6.479) circle (0.28pt);
\fill[slate] (0.039,6.691) circle (0.28pt);
\fill[slate] (0.039,6.727) circle (0.28pt);
\fill[slate] (0.039,5.828) circle (0.28pt);
\fill[slate] (0.039,6.077) circle (0.28pt);
\fill[slate] (0.039,6.277) circle (0.28pt);
\fill[slate] (0.039,6.733) circle (0.28pt);
\fill[slate] (0.039,6.326) circle (0.28pt);
\fill[slate] (0.039,6.575) circle (0.28pt);
\fill[slate] (0.039,6.775) circle (0.28pt);
\fill[slate] (0.039,6.824) circle (0.28pt);
\fill[slate] (0.039,3.92) circle (0.28pt);
\fill[slate] (0.039,4.168) circle (0.28pt);
\fill[slate] (0.039,4.371) circle (0.28pt);
\fill[slate] (0.039,4.827) circle (0.28pt);
\fill[slate] (0.039,5.07) circle (0.28pt);
\fill[slate] (0.039,5.283) circle (0.28pt);
\fill[slate] (0.039,5.739) circle (0.28pt);
\fill[slate] (0.039,5.982) circle (0.28pt);
\fill[slate] (0.039,6.195) circle (0.28pt);
\fill[slate] (0.039,6.651) circle (0.28pt);
\fill[slate] (0.039,6.23) circle (0.28pt);
\fill[slate] (0.039,6.479) circle (0.28pt);
\fill[slate] (0.039,6.682) circle (0.28pt);
\fill[slate] (0.039,6.728) circle (0.28pt);
\fill[slate] (0.039,5.319) circle (0.28pt);
\fill[slate] (0.039,5.567) circle (0.28pt);
\fill[slate] (0.039,5.77) circle (0.28pt);
\fill[slate] (0.039,6.226) circle (0.28pt);
\fill[slate] (0.039,6.469) circle (0.28pt);
\fill[slate] (0.039,6.682) circle (0.28pt);
\fill[slate] (0.039,6.718) circle (0.28pt);
\fill[slate] (0.039,5.816) circle (0.28pt);
\fill[slate] (0.039,6.065) circle (0.28pt);
\fill[slate] (0.039,6.266) circle (0.28pt);
\fill[slate] (0.039,6.722) circle (0.28pt);
\fill[slate] (0.039,6.314) circle (0.28pt);
\fill[slate] (0.039,6.563) circle (0.28pt);
\fill[slate] (0.039,6.764) circle (0.28pt);
\fill[slate] (0.039,6.812) circle (0.28pt);
\fill[slate] (0.039,4.417) circle (0.28pt);
\fill[slate] (0.039,4.666) circle (0.28pt);
\fill[slate] (0.039,4.867) circle (0.28pt);
\fill[slate] (0.039,5.323) circle (0.28pt);
\fill[slate] (0.039,5.567) circle (0.28pt);
\fill[slate] (0.039,5.779) circle (0.28pt);
\fill[slate] (0.039,6.235) circle (0.28pt);
\fill[slate] (0.039,6.479) circle (0.28pt);
\fill[slate] (0.039,6.691) circle (0.28pt);
\fill[slate] (0.039,6.727) circle (0.28pt);
\fill[slate] (0.039,5.815) circle (0.28pt);
\fill[slate] (0.039,6.064) circle (0.28pt);
\fill[slate] (0.039,6.266) circle (0.28pt);
\fill[slate] (0.039,6.722) circle (0.28pt);
\fill[slate] (0.039,6.313) circle (0.28pt);
\fill[slate] (0.039,6.562) circle (0.28pt);
\fill[slate] (0.039,6.763) circle (0.28pt);
\fill[slate] (0.039,6.811) circle (0.28pt);
\fill[slate] (0.039,4.915) circle (0.28pt);
\fill[slate] (0.039,5.164) circle (0.28pt);
\fill[slate] (0.039,5.365) circle (0.28pt);
\fill[slate] (0.039,5.821) circle (0.28pt);
\fill[slate] (0.039,6.064) circle (0.28pt);
\fill[slate] (0.039,6.277) circle (0.28pt);
\fill[slate] (0.039,6.733) circle (0.28pt);
\fill[slate] (0.039,6.313) circle (0.28pt);
\fill[slate] (0.039,6.562) circle (0.28pt);
\fill[slate] (0.039,6.764) circle (0.28pt);
\fill[slate] (0.039,6.811) circle (0.28pt);
\fill[slate] (0.039,5.413) circle (0.28pt);
\fill[slate] (0.039,5.662) circle (0.28pt);
\fill[slate] (0.039,5.863) circle (0.28pt);
\fill[slate] (0.039,6.319) circle (0.28pt);
\fill[slate] (0.039,6.562) circle (0.28pt);
\fill[slate] (0.039,6.775) circle (0.28pt);
\fill[slate] (0.039,6.811) circle (0.28pt);
\fill[slate] (0.039,5.911) circle (0.28pt);
\fill[slate] (0.039,6.16) circle (0.28pt);
\fill[slate] (0.039,6.361) circle (0.28pt);
\fill[slate] (0.039,6.817) circle (0.28pt);
\fill[slate] (0.039,6.409) circle (0.28pt);
\fill[slate] (0.039,6.658) circle (0.28pt);
\fill[slate] (0.038,3.459) circle (0.28pt);
\fill[slate] (0.038,3.915) circle (0.28pt);
\fill[slate] (0.038,4.159) circle (0.28pt);
\fill[slate] (0.038,4.371) circle (0.28pt);
\fill[slate] (0.038,4.827) circle (0.28pt);
\fill[slate] (0.038,5.07) circle (0.28pt);
\fill[slate] (0.038,5.283) circle (0.28pt);
\fill[slate] (0.038,5.739) circle (0.28pt);
\fill[slate] (0.038,5.983) circle (0.28pt);
\fill[slate] (0.038,6.195) circle (0.28pt);
\fill[slate] (0.038,6.651) circle (0.28pt);
\fill[slate] (0.038,6.231) circle (0.28pt);
\fill[slate] (0.038,6.48) circle (0.28pt);
\fill[slate] (0.038,6.682) circle (0.28pt);
\fill[slate] (0.038,6.729) circle (0.28pt);
\fill[slate] (0.038,5.319) circle (0.28pt);
\fill[slate] (0.038,5.567) circle (0.28pt);
\fill[slate] (0.038,5.77) circle (0.28pt);
\fill[slate] (0.038,6.226) circle (0.28pt);
\fill[slate] (0.038,6.469) circle (0.28pt);
\fill[slate] (0.038,6.682) circle (0.28pt);
\fill[slate] (0.038,6.718) circle (0.28pt);
\fill[slate] (0.038,5.816) circle (0.28pt);
\fill[slate] (0.038,6.065) circle (0.28pt);
\fill[slate] (0.038,6.266) circle (0.28pt);
\fill[slate] (0.038,6.722) circle (0.28pt);
\fill[slate] (0.038,6.314) circle (0.28pt);
\fill[slate] (0.038,6.563) circle (0.28pt);
\fill[slate] (0.038,6.764) circle (0.28pt);
\fill[slate] (0.038,6.812) circle (0.28pt);
\fill[slate] (0.038,4.407) circle (0.28pt);
\fill[slate] (0.038,4.656) circle (0.28pt);
\fill[slate] (0.038,4.858) circle (0.28pt);
\fill[slate] (0.038,5.314) circle (0.28pt);
\fill[slate] (0.038,5.558) circle (0.28pt);
\fill[slate] (0.038,5.77) circle (0.28pt);
\fill[slate] (0.038,6.226) circle (0.28pt);
\fill[slate] (0.038,6.469) circle (0.28pt);
\fill[slate] (0.038,6.682) circle (0.28pt);
\fill[slate] (0.038,6.718) circle (0.28pt);
\fill[slate] (0.038,5.806) circle (0.28pt);
\fill[slate] (0.038,6.055) circle (0.28pt);
\fill[slate] (0.038,6.257) circle (0.28pt);
\fill[slate] (0.038,6.713) circle (0.28pt);
\fill[slate] (0.038,6.304) circle (0.28pt);
\fill[slate] (0.038,6.553) circle (0.28pt);
\fill[slate] (0.038,6.753) circle (0.28pt);
\fill[slate] (0.038,6.802) circle (0.28pt);
\fill[slate] (0.038,4.905) circle (0.28pt);
\fill[slate] (0.038,5.154) circle (0.28pt);
\fill[slate] (0.038,5.354) circle (0.28pt);
\fill[slate] (0.038,5.81) circle (0.28pt);
\fill[slate] (0.038,6.054) circle (0.28pt);
\fill[slate] (0.038,6.266) circle (0.28pt);
\fill[slate] (0.038,6.722) circle (0.28pt);
\fill[slate] (0.038,6.302) circle (0.28pt);
\fill[slate] (0.038,6.551) circle (0.28pt);
\fill[slate] (0.038,6.753) circle (0.28pt);
\fill[slate] (0.038,6.8) circle (0.28pt);
\fill[slate] (0.038,5.403) circle (0.28pt);
\fill[slate] (0.038,5.652) circle (0.28pt);
\fill[slate] (0.038,5.852) circle (0.28pt);
\fill[slate] (0.038,6.308) circle (0.28pt);
\fill[slate] (0.038,6.552) circle (0.28pt);
\fill[slate] (0.038,6.764) circle (0.28pt);
\fill[slate] (0.038,6.8) circle (0.28pt);
\fill[slate] (0.038,5.901) circle (0.28pt);
\fill[slate] (0.038,6.15) circle (0.28pt);
\fill[slate] (0.038,6.35) circle (0.28pt);
\fill[slate] (0.038,6.806) circle (0.28pt);
\fill[slate] (0.038,6.399) circle (0.28pt);
\fill[slate] (0.038,6.648) circle (0.28pt);
\fill[slate] (0.038,6.848) circle (0.28pt);
\fill[slate] (0.038,3.955) circle (0.28pt);
\fill[slate] (0.038,4.411) circle (0.28pt);
\fill[slate] (0.038,4.655) circle (0.28pt);
\fill[slate] (0.038,4.867) circle (0.28pt);
\fill[slate] (0.038,5.323) circle (0.28pt);
\fill[slate] (0.038,5.567) circle (0.28pt);
\fill[slate] (0.038,5.78) circle (0.28pt);
\fill[slate] (0.038,6.236) circle (0.28pt);
\fill[slate] (0.038,6.479) circle (0.28pt);
\fill[slate] (0.038,6.692) circle (0.28pt);
\fill[slate] (0.038,6.727) circle (0.28pt);
\fill[slate] (0.038,5.815) circle (0.28pt);
\fill[slate] (0.038,6.064) circle (0.28pt);
\fill[slate] (0.038,6.266) circle (0.28pt);
\fill[slate] (0.038,6.722) circle (0.28pt);
\fill[slate] (0.038,6.313) circle (0.28pt);
\fill[slate] (0.038,6.562) circle (0.28pt);
\fill[slate] (0.038,6.763) circle (0.28pt);
\fill[slate] (0.038,6.811) circle (0.28pt);
\fill[slate] (0.038,4.903) circle (0.28pt);
\fill[slate] (0.038,5.152) circle (0.28pt);
\fill[slate] (0.038,5.354) circle (0.28pt);
\fill[slate] (0.038,5.81) circle (0.28pt);
\fill[slate] (0.038,6.054) circle (0.28pt);
\fill[slate] (0.038,6.266) circle (0.28pt);
\fill[slate] (0.038,6.723) circle (0.28pt);
\fill[slate] (0.038,6.302) circle (0.28pt);
\fill[slate] (0.038,6.551) circle (0.28pt);
\fill[slate] (0.038,6.753) circle (0.28pt);
\fill[slate] (0.038,6.8) circle (0.28pt);
\fill[slate] (0.038,5.65) circle (0.28pt);
\fill[slate] (0.038,5.851) circle (0.28pt);
\fill[slate] (0.038,6.307) circle (0.28pt);
\fill[slate] (0.038,6.55) circle (0.28pt);
\fill[slate] (0.038,6.763) circle (0.28pt);
\fill[slate] (0.038,6.799) circle (0.28pt);
\fill[slate] (0.038,5.899) circle (0.28pt);
\fill[slate] (0.038,6.148) circle (0.28pt);
\fill[slate] (0.038,6.349) circle (0.28pt);
\fill[slate] (0.038,6.805) circle (0.28pt);
\fill[slate] (0.038,6.397) circle (0.28pt);
\fill[slate] (0.038,6.646) circle (0.28pt);
\fill[slate] (0.038,6.847) circle (0.28pt);
\fill[slate] (0.038,4.004) circle (0.28pt);
\fill[slate] (0.038,4.253) circle (0.28pt);
\fill[slate] (0.038,4.909) circle (0.28pt);
\fill[slate] (0.038,5.153) circle (0.28pt);
\fill[slate] (0.038,5.365) circle (0.28pt);
\fill[slate] (0.038,5.821) circle (0.28pt);
\fill[slate] (0.038,6.065) circle (0.28pt);
\fill[slate] (0.038,6.277) circle (0.28pt);
\fill[slate] (0.038,6.733) circle (0.28pt);
\fill[slate] (0.038,6.313) circle (0.28pt);
\fill[slate] (0.038,6.562) circle (0.28pt);
\fill[slate] (0.038,6.764) circle (0.28pt);
\fill[slate] (0.038,6.811) circle (0.28pt);
\fill[slate] (0.038,5.852) circle (0.28pt);
\fill[slate] (0.038,6.308) circle (0.28pt);
\fill[slate] (0.038,6.552) circle (0.28pt);
\fill[slate] (0.038,6.764) circle (0.28pt);
\fill[slate] (0.038,6.8) circle (0.28pt);
\fill[slate] (0.038,6.148) circle (0.28pt);
\fill[slate] (0.038,6.349) circle (0.28pt);
\fill[slate] (0.038,6.805) circle (0.28pt);
\fill[slate] (0.038,6.397) circle (0.28pt);
\fill[slate] (0.038,6.646) circle (0.28pt);
\fill[slate] (0.038,4.502) circle (0.28pt);
\fill[slate] (0.038,4.751) circle (0.28pt);
\fill[slate] (0.038,5.407) circle (0.28pt);
\fill[slate] (0.038,5.651) circle (0.28pt);
\fill[slate] (0.038,5.863) circle (0.28pt);
\fill[slate] (0.038,6.319) circle (0.28pt);
\fill[slate] (0.038,6.563) circle (0.28pt);
\fill[slate] (0.038,6.775) circle (0.28pt);
\fill[slate] (0.038,6.811) circle (0.28pt);
\fill[slate] (0.038,5.899) circle (0.28pt);
\fill[slate] (0.038,6.35) circle (0.28pt);
\fill[slate] (0.038,6.806) circle (0.28pt);
\fill[slate] (0.038,6.847) circle (0.28pt);
\fill[slate] (0.038,5) circle (0.28pt);
\fill[slate] (0.038,5.249) circle (0.28pt);
\fill[slate] (0.038,5.449) circle (0.28pt);
\fill[slate] (0.038,5.905) circle (0.28pt);
\fill[slate] (0.038,6.149) circle (0.28pt);
\fill[slate] (0.038,6.361) circle (0.28pt);
\fill[slate] (0.038,6.817) circle (0.28pt);
\fill[slate] (0.038,6.397) circle (0.28pt);
\fill[slate] (0.038,6.848) circle (0.28pt);
\fill[slate] (0.038,5.498) circle (0.28pt);
\fill[slate] (0.038,5.747) circle (0.28pt);
\fill[slate] (0.038,5.947) circle (0.28pt);
\fill[slate] (0.038,6.403) circle (0.28pt);
\fill[slate] (0.038,6.647) circle (0.28pt);
\fill[slate] (0.038,5.996) circle (0.28pt);
\fill[slate] (0.038,6.245) circle (0.28pt);
\fill[slate] (0.038,6.445) circle (0.28pt);
\fill[slate] (0.038,6.494) circle (0.28pt);
\fill[slate] (0.038,6.743) circle (0.28pt);
\fill[slate] (0.006,3.014) circle (0.28pt);
\fill[slate] (0.006,3.258) circle (0.28pt);
\fill[slate] (0.006,3.47) circle (0.28pt);
\fill[slate] (0.006,3.926) circle (0.28pt);
\fill[slate] (0.006,4.169) circle (0.28pt);
\fill[slate] (0.006,4.382) circle (0.28pt);
\fill[slate] (0.006,4.838) circle (0.28pt);
\fill[slate] (0.006,5.081) circle (0.28pt);
\fill[slate] (0.006,5.294) circle (0.28pt);
\fill[slate] (0.006,5.75) circle (0.28pt);
\fill[slate] (0.006,5.994) circle (0.28pt);
\fill[slate] (0.006,6.206) circle (0.28pt);
\fill[slate] (0.006,6.663) circle (0.28pt);
\fill[slate] (0.006,6.242) circle (0.28pt);
\fill[slate] (0.006,6.491) circle (0.28pt);
\fill[slate] (0.006,6.693) circle (0.28pt);
\fill[slate] (0.006,6.74) circle (0.28pt);
\fill[slate] (0.006,5.33) circle (0.28pt);
\fill[slate] (0.006,5.579) circle (0.28pt);
\fill[slate] (0.006,5.781) circle (0.28pt);
\fill[slate] (0.006,6.237) circle (0.28pt);
\fill[slate] (0.006,6.48) circle (0.28pt);
\fill[slate] (0.006,6.693) circle (0.28pt);
\fill[slate] (0.006,6.729) circle (0.28pt);
\fill[slate] (0.006,5.828) circle (0.28pt);
\fill[slate] (0.006,6.077) circle (0.28pt);
\fill[slate] (0.006,6.277) circle (0.28pt);
\fill[slate] (0.006,6.733) circle (0.28pt);
\fill[slate] (0.006,6.326) circle (0.28pt);
\fill[slate] (0.006,6.575) circle (0.28pt);
\fill[slate] (0.006,6.775) circle (0.28pt);
\fill[slate] (0.006,6.824) circle (0.28pt);
\fill[slate] (0.006,4.418) circle (0.28pt);
\fill[slate] (0.006,4.666) circle (0.28pt);
\fill[slate] (0.006,4.869) circle (0.28pt);
\fill[slate] (0.006,5.325) circle (0.28pt);
\fill[slate] (0.006,5.568) circle (0.28pt);
\fill[slate] (0.006,5.781) circle (0.28pt);
\fill[slate] (0.006,6.237) circle (0.28pt);
\fill[slate] (0.006,6.48) circle (0.28pt);
\fill[slate] (0.006,6.693) circle (0.28pt);
\fill[slate] (0.006,6.728) circle (0.28pt);
\fill[slate] (0.006,5.817) circle (0.28pt);
\fill[slate] (0.006,6.065) circle (0.28pt);
\fill[slate] (0.006,6.268) circle (0.28pt);
\fill[slate] (0.006,6.724) circle (0.28pt);
\fill[slate] (0.006,6.314) circle (0.28pt);
\fill[slate] (0.006,6.563) circle (0.28pt);
\fill[slate] (0.006,6.764) circle (0.28pt);
\fill[slate] (0.006,6.812) circle (0.28pt);
\fill[slate] (0.006,4.915) circle (0.28pt);
\fill[slate] (0.006,5.164) circle (0.28pt);
\fill[slate] (0.006,5.365) circle (0.28pt);
\fill[slate] (0.006,5.821) circle (0.28pt);
\fill[slate] (0.006,6.065) circle (0.28pt);
\fill[slate] (0.006,6.277) circle (0.28pt);
\fill[slate] (0.006,6.733) circle (0.28pt);
\fill[slate] (0.006,6.313) circle (0.28pt);
\fill[slate] (0.006,6.562) circle (0.28pt);
\fill[slate] (0.006,6.764) circle (0.28pt);
\fill[slate] (0.006,6.811) circle (0.28pt);
\fill[slate] (0.006,5.413) circle (0.28pt);
\fill[slate] (0.006,5.662) circle (0.28pt);
\fill[slate] (0.006,5.863) circle (0.28pt);
\fill[slate] (0.006,6.319) circle (0.28pt);
\fill[slate] (0.006,6.562) circle (0.28pt);
\fill[slate] (0.006,6.775) circle (0.28pt);
\fill[slate] (0.006,6.811) circle (0.28pt);
\fill[slate] (0.006,5.911) circle (0.28pt);
\fill[slate] (0.006,6.16) circle (0.28pt);
\fill[slate] (0.006,6.361) circle (0.28pt);
\fill[slate] (0.006,6.817) circle (0.28pt);
\fill[slate] (0.006,6.409) circle (0.28pt);
\fill[slate] (0.006,6.658) circle (0.28pt);
\fill[slate] (0.006,3.506) circle (0.28pt);
\fill[slate] (0.006,3.957) circle (0.28pt);
\fill[slate] (0.006,4.413) circle (0.28pt);
\fill[slate] (0.006,4.657) circle (0.28pt);
\fill[slate] (0.006,4.869) circle (0.28pt);
\fill[slate] (0.006,5.325) circle (0.28pt);
\fill[slate] (0.006,5.568) circle (0.28pt);
\fill[slate] (0.006,5.781) circle (0.28pt);
\fill[slate] (0.006,6.237) circle (0.28pt);
\fill[slate] (0.006,6.481) circle (0.28pt);
\fill[slate] (0.006,6.693) circle (0.28pt);
\fill[slate] (0.006,6.729) circle (0.28pt);
\fill[slate] (0.006,5.817) circle (0.28pt);
\fill[slate] (0.006,6.065) circle (0.28pt);
\fill[slate] (0.006,6.268) circle (0.28pt);
\fill[slate] (0.006,6.724) circle (0.28pt);
\fill[slate] (0.006,6.314) circle (0.28pt);
\fill[slate] (0.006,6.563) circle (0.28pt);
\fill[slate] (0.006,6.764) circle (0.28pt);
\fill[slate] (0.006,6.812) circle (0.28pt);
\fill[slate] (0.006,4.905) circle (0.28pt);
\fill[slate] (0.006,5.154) circle (0.28pt);
\fill[slate] (0.006,5.356) circle (0.28pt);
\fill[slate] (0.006,5.812) circle (0.28pt);
\fill[slate] (0.006,6.056) circle (0.28pt);
\fill[slate] (0.006,6.268) circle (0.28pt);
\fill[slate] (0.006,6.724) circle (0.28pt);
\fill[slate] (0.006,6.304) circle (0.28pt);
\fill[slate] (0.006,6.553) circle (0.28pt);
\fill[slate] (0.006,6.755) circle (0.28pt);
\fill[slate] (0.006,6.802) circle (0.28pt);
\fill[slate] (0.006,5.403) circle (0.28pt);
\fill[slate] (0.006,5.652) circle (0.28pt);
\fill[slate] (0.006,5.852) circle (0.28pt);
\fill[slate] (0.006,6.308) circle (0.28pt);
\fill[slate] (0.006,6.552) circle (0.28pt);
\fill[slate] (0.006,6.764) circle (0.28pt);
\fill[slate] (0.006,6.8) circle (0.28pt);
\fill[slate] (0.006,5.901) circle (0.28pt);
\fill[slate] (0.006,6.15) circle (0.28pt);
\fill[slate] (0.006,6.35) circle (0.28pt);
\fill[slate] (0.006,6.806) circle (0.28pt);
\fill[slate] (0.006,6.399) circle (0.28pt);
\fill[slate] (0.006,6.648) circle (0.28pt);
\fill[slate] (0.006,6.848) circle (0.28pt);
\fill[slate] (0.006,4.453) circle (0.28pt);
\fill[slate] (0.006,4.909) circle (0.28pt);
\fill[slate] (0.006,5.153) circle (0.28pt);
\fill[slate] (0.006,5.365) circle (0.28pt);
\fill[slate] (0.006,5.821) circle (0.28pt);
\fill[slate] (0.006,6.065) circle (0.28pt);
\fill[slate] (0.006,6.278) circle (0.28pt);
\fill[slate] (0.006,6.734) circle (0.28pt);
\fill[slate] (0.006,6.313) circle (0.28pt);
\fill[slate] (0.006,6.562) circle (0.28pt);
\fill[slate] (0.006,6.764) circle (0.28pt);
\fill[slate] (0.006,6.811) circle (0.28pt);
\fill[slate] (0.006,5.401) circle (0.28pt);
\fill[slate] (0.006,5.65) circle (0.28pt);
\fill[slate] (0.006,5.852) circle (0.28pt);
\fill[slate] (0.006,6.308) circle (0.28pt);
\fill[slate] (0.006,6.552) circle (0.28pt);
\fill[slate] (0.006,6.764) circle (0.28pt);
\fill[slate] (0.006,6.8) circle (0.28pt);
\fill[slate] (0.006,5.899) circle (0.28pt);
\fill[slate] (0.006,6.148) circle (0.28pt);
\fill[slate] (0.006,6.349) circle (0.28pt);
\fill[slate] (0.006,6.805) circle (0.28pt);
\fill[slate] (0.006,6.397) circle (0.28pt);
\fill[slate] (0.006,6.646) circle (0.28pt);
\fill[slate] (0.006,6.847) circle (0.28pt);
\fill[slate] (0.006,4.751) circle (0.28pt);
\fill[slate] (0.006,4.951) circle (0.28pt);
\fill[slate] (0.006,5.407) circle (0.28pt);
\fill[slate] (0.006,5.651) circle (0.28pt);
\fill[slate] (0.006,5.863) circle (0.28pt);
\fill[slate] (0.006,6.319) circle (0.28pt);
\fill[slate] (0.006,6.563) circle (0.28pt);
\fill[slate] (0.006,6.775) circle (0.28pt);
\fill[slate] (0.006,6.811) circle (0.28pt);
\fill[slate] (0.006,5.899) circle (0.28pt);
\fill[slate] (0.006,6.148) circle (0.28pt);
\fill[slate] (0.006,6.35) circle (0.28pt);
\fill[slate] (0.006,6.806) circle (0.28pt);
\fill[slate] (0.006,6.397) circle (0.28pt);
\fill[slate] (0.006,6.847) circle (0.28pt);
\fill[slate] (0.006,5) circle (0.28pt);
\fill[slate] (0.006,5.249) circle (0.28pt);
\fill[slate] (0.006,5.905) circle (0.28pt);
\fill[slate] (0.006,6.149) circle (0.28pt);
\fill[slate] (0.006,6.361) circle (0.28pt);
\fill[slate] (0.006,6.817) circle (0.28pt);
\fill[slate] (0.006,6.397) circle (0.28pt);
\fill[slate] (0.006,6.848) circle (0.28pt);
\fill[slate] (0.006,5.498) circle (0.28pt);
\fill[slate] (0.006,5.747) circle (0.28pt);
\fill[slate] (0.006,6.403) circle (0.28pt);
\fill[slate] (0.006,6.647) circle (0.28pt);
\fill[slate] (0.006,5.996) circle (0.28pt);
\fill[slate] (0.006,6.245) circle (0.28pt);
\fill[slate] (0.006,6.445) circle (0.28pt);
\fill[slate] (0.006,6.494) circle (0.28pt);
\fill[slate] (0.006,6.743) circle (0.28pt);
\fill[slate] (0.001,3.512) circle (0.28pt);
\fill[slate] (0.001,3.756) circle (0.28pt);
\fill[slate] (0.001,3.968) circle (0.28pt);
\fill[slate] (0.001,4.424) circle (0.28pt);
\fill[slate] (0.001,4.667) circle (0.28pt);
\fill[slate] (0.001,4.88) circle (0.28pt);
\fill[slate] (0.001,5.336) circle (0.28pt);
\fill[slate] (0.001,5.579) circle (0.28pt);
\fill[slate] (0.001,5.792) circle (0.28pt);
\fill[slate] (0.001,6.248) circle (0.28pt);
\fill[slate] (0.001,6.492) circle (0.28pt);
\fill[slate] (0.001,6.704) circle (0.28pt);
\fill[slate] (0.001,6.74) circle (0.28pt);
\fill[slate] (0.001,5.828) circle (0.28pt);
\fill[slate] (0.001,6.077) circle (0.28pt);
\fill[slate] (0.001,6.279) circle (0.28pt);
\fill[slate] (0.001,6.735) circle (0.28pt);
\fill[slate] (0.001,6.326) circle (0.28pt);
\fill[slate] (0.001,6.575) circle (0.28pt);
\fill[slate] (0.001,6.775) circle (0.28pt);
\fill[slate] (0.001,6.824) circle (0.28pt);
\fill[slate] (0.001,4.916) circle (0.28pt);
\fill[slate] (0.001,5.164) circle (0.28pt);
\fill[slate] (0.001,5.367) circle (0.28pt);
\fill[slate] (0.001,5.823) circle (0.28pt);
\fill[slate] (0.001,6.066) circle (0.28pt);
\fill[slate] (0.001,6.279) circle (0.28pt);
\fill[slate] (0.001,6.735) circle (0.28pt);
\fill[slate] (0.001,6.315) circle (0.28pt);
\fill[slate] (0.001,6.563) circle (0.28pt);
\fill[slate] (0.001,6.766) circle (0.28pt);
\fill[slate] (0.001,6.812) circle (0.28pt);
\fill[slate] (0.001,5.413) circle (0.28pt);
\fill[slate] (0.001,5.662) circle (0.28pt);
\fill[slate] (0.001,5.863) circle (0.28pt);
\fill[slate] (0.001,6.319) circle (0.28pt);
\fill[slate] (0.001,6.563) circle (0.28pt);
\fill[slate] (0.001,6.775) circle (0.28pt);
\fill[slate] (0.001,6.811) circle (0.28pt);
\fill[slate] (0.001,5.911) circle (0.28pt);
\fill[slate] (0.001,6.16) circle (0.28pt);
\fill[slate] (0.001,6.361) circle (0.28pt);
\fill[slate] (0.001,6.817) circle (0.28pt);
\fill[slate] (0.001,6.409) circle (0.28pt);
\fill[slate] (0.001,6.658) circle (0.28pt);
\fill[slate] (0.001,4.004) circle (0.28pt);
\fill[slate] (0.001,4.253) circle (0.28pt);
\fill[slate] (0.001,4.455) circle (0.28pt);
\fill[slate] (0.001,4.911) circle (0.28pt);
\fill[slate] (0.001,5.155) circle (0.28pt);
\fill[slate] (0.001,5.367) circle (0.28pt);
\fill[slate] (0.001,5.823) circle (0.28pt);
\fill[slate] (0.001,6.066) circle (0.28pt);
\fill[slate] (0.001,6.279) circle (0.28pt);
\fill[slate] (0.001,6.735) circle (0.28pt);
\fill[slate] (0.001,6.315) circle (0.28pt);
\fill[slate] (0.001,6.563) circle (0.28pt);
\fill[slate] (0.001,6.766) circle (0.28pt);
\fill[slate] (0.001,6.812) circle (0.28pt);
\fill[slate] (0.001,5.403) circle (0.28pt);
\fill[slate] (0.001,5.652) circle (0.28pt);
\fill[slate] (0.001,5.854) circle (0.28pt);
\fill[slate] (0.001,6.31) circle (0.28pt);
\fill[slate] (0.001,6.554) circle (0.28pt);
\fill[slate] (0.001,6.766) circle (0.28pt);
\fill[slate] (0.001,6.802) circle (0.28pt);
\fill[slate] (0.001,5.901) circle (0.28pt);
\fill[slate] (0.001,6.15) circle (0.28pt);
\fill[slate] (0.001,6.35) circle (0.28pt);
\fill[slate] (0.001,6.806) circle (0.28pt);
\fill[slate] (0.001,6.399) circle (0.28pt);
\fill[slate] (0.001,6.648) circle (0.28pt);
\fill[slate] (0.001,6.848) circle (0.28pt);
\fill[slate] (0.001,4.951) circle (0.28pt);
\fill[slate] (0.001,5.407) circle (0.28pt);
\fill[slate] (0.001,5.651) circle (0.28pt);
\fill[slate] (0.001,5.863) circle (0.28pt);
\fill[slate] (0.001,6.319) circle (0.28pt);
\fill[slate] (0.001,6.563) circle (0.28pt);
\fill[slate] (0.001,6.776) circle (0.28pt);
\fill[slate] (0.001,6.811) circle (0.28pt);
\fill[slate] (0.001,5.899) circle (0.28pt);
\fill[slate] (0.001,6.148) circle (0.28pt);
\fill[slate] (0.001,6.35) circle (0.28pt);
\fill[slate] (0.001,6.806) circle (0.28pt);
\fill[slate] (0.001,6.397) circle (0.28pt);
\fill[slate] (0.001,6.646) circle (0.28pt);
\fill[slate] (0.001,6.847) circle (0.28pt);
\fill[slate] (0.001,5.449) circle (0.28pt);
\fill[slate] (0.001,5.905) circle (0.28pt);
\fill[slate] (0.001,6.149) circle (0.28pt);
\fill[slate] (0.001,6.361) circle (0.28pt);
\fill[slate] (0.001,6.817) circle (0.28pt);
\fill[slate] (0.001,6.397) circle (0.28pt);
\fill[slate] (0.001,6.646) circle (0.28pt);
\fill[slate] (0.001,6.848) circle (0.28pt);
\fill[slate] (0.001,5.498) circle (0.28pt);
\fill[slate] (0.001,5.747) circle (0.28pt);
\fill[slate] (0.001,5.947) circle (0.28pt);
\fill[slate] (0.001,6.403) circle (0.28pt);
\fill[slate] (0.001,6.647) circle (0.28pt);
\fill[slate] (0.001,5.996) circle (0.28pt);
\fill[slate] (0.001,6.245) circle (0.28pt);
\fill[slate] (0.001,6.494) circle (0.28pt);
\fill[slate] (0.001,6.743) circle (0.28pt);
\fill[slate] (0,4.01) circle (0.28pt);
\fill[slate] (0,4.254) circle (0.28pt);
\fill[slate] (0,4.466) circle (0.28pt);
\fill[slate] (0,4.922) circle (0.28pt);
\fill[slate] (0,5.165) circle (0.28pt);
\fill[slate] (0,5.378) circle (0.28pt);
\fill[slate] (0,5.834) circle (0.28pt);
\fill[slate] (0,6.078) circle (0.28pt);
\fill[slate] (0,6.29) circle (0.28pt);
\fill[slate] (0,6.746) circle (0.28pt);
\fill[slate] (0,6.326) circle (0.28pt);
\fill[slate] (0,6.575) circle (0.28pt);
\fill[slate] (0,6.777) circle (0.28pt);
\fill[slate] (0,6.824) circle (0.28pt);
\fill[slate] (0,5.414) circle (0.28pt);
\fill[slate] (0,5.662) circle (0.28pt);
\fill[slate] (0,5.865) circle (0.28pt);
\fill[slate] (0,6.321) circle (0.28pt);
\fill[slate] (0,6.564) circle (0.28pt);
\fill[slate] (0,6.777) circle (0.28pt);
\fill[slate] (0,6.813) circle (0.28pt);
\fill[slate] (0,5.911) circle (0.28pt);
\fill[slate] (0,6.16) circle (0.28pt);
\fill[slate] (0,6.361) circle (0.28pt);
\fill[slate] (0,6.817) circle (0.28pt);
\fill[slate] (0,6.409) circle (0.28pt);
\fill[slate] (0,6.658) circle (0.28pt);
\fill[slate] (0,4.502) circle (0.28pt);
\fill[slate] (0,4.751) circle (0.28pt);
\fill[slate] (0,4.953) circle (0.28pt);
\fill[slate] (0,5.409) circle (0.28pt);
\fill[slate] (0,5.653) circle (0.28pt);
\fill[slate] (0,5.865) circle (0.28pt);
\fill[slate] (0,6.321) circle (0.28pt);
\fill[slate] (0,6.564) circle (0.28pt);
\fill[slate] (0,6.777) circle (0.28pt);
\fill[slate] (0,6.813) circle (0.28pt);
\fill[slate] (0,5.901) circle (0.28pt);
\fill[slate] (0,6.15) circle (0.28pt);
\fill[slate] (0,6.352) circle (0.28pt);
\fill[slate] (0,6.808) circle (0.28pt);
\fill[slate] (0,6.399) circle (0.28pt);
\fill[slate] (0,6.648) circle (0.28pt);
\fill[slate] (0,6.848) circle (0.28pt);
\fill[slate] (0,5) circle (0.28pt);
\fill[slate] (0,5.449) circle (0.28pt);
\fill[slate] (0,5.905) circle (0.28pt);
\fill[slate] (0,6.149) circle (0.28pt);
\fill[slate] (0,6.361) circle (0.28pt);
\fill[slate] (0,6.817) circle (0.28pt);
\fill[slate] (0,6.397) circle (0.28pt);
\fill[slate] (0,6.646) circle (0.28pt);
\fill[slate] (0,6.848) circle (0.28pt);
\fill[slate] (0,5.947) circle (0.28pt);
\fill[slate] (0,6.403) circle (0.28pt);
\fill[slate] (0,6.647) circle (0.28pt);
\fill[slate] (0,6.245) circle (0.28pt);
\fill[slate] (0,6.445) circle (0.28pt);
\fill[slate] (0,6.494) circle (0.28pt);
\fill[slate] (0,6.743) circle (0.28pt);
\fill[slate] (0,4.508) circle (0.28pt);
\fill[slate] (0,4.752) circle (0.28pt);
\fill[slate] (0,4.964) circle (0.28pt);
\fill[slate] (0,5.42) circle (0.28pt);
\fill[slate] (0,5.663) circle (0.28pt);
\fill[slate] (0,5.876) circle (0.28pt);
\fill[slate] (0,6.332) circle (0.28pt);
\fill[slate] (0,6.576) circle (0.28pt);
\fill[slate] (0,6.788) circle (0.28pt);
\fill[slate] (0,6.824) circle (0.28pt);
\fill[slate] (0,5.912) circle (0.28pt);
\fill[slate] (0,6.16) circle (0.28pt);
\fill[slate] (0,6.363) circle (0.28pt);
\fill[slate] (0,6.819) circle (0.28pt);
\fill[slate] (0,6.409) circle (0.28pt);
\fill[slate] (0,6.658) circle (0.28pt);
\fill[slate] (0,5) circle (0.28pt);
\fill[slate] (0,5.249) circle (0.28pt);
\fill[slate] (0,5.451) circle (0.28pt);
\fill[slate] (0,5.907) circle (0.28pt);
\fill[slate] (0,6.151) circle (0.28pt);
\fill[slate] (0,6.363) circle (0.28pt);
\fill[slate] (0,6.819) circle (0.28pt);
\fill[slate] (0,6.399) circle (0.28pt);
\fill[slate] (0,6.648) circle (0.28pt);
\fill[slate] (0,6.85) circle (0.28pt);
\fill[slate] (0,5.498) circle (0.28pt);
\fill[slate] (0,5.747) circle (0.28pt);
\fill[slate] (0,5.947) circle (0.28pt);
\fill[slate] (0,6.403) circle (0.28pt);
\fill[slate] (0,6.647) circle (0.28pt);
\fill[slate] (0,6.445) circle (0.28pt);
\fill[slate] (0,5.006) circle (0.28pt);
\fill[slate] (0,5.25) circle (0.28pt);
\fill[slate] (0,5.462) circle (0.28pt);
\fill[slate] (0,5.918) circle (0.28pt);
\fill[slate] (0,6.161) circle (0.28pt);
\fill[slate] (0,6.374) circle (0.28pt);
\fill[slate] (0,6.83) circle (0.28pt);
\fill[slate] (0,6.41) circle (0.28pt);
\fill[slate] (0,6.658) circle (0.28pt);
\fill[slate] (0,5.498) circle (0.28pt);
\fill[slate] (0,5.747) circle (0.28pt);
\fill[slate] (0,5.949) circle (0.28pt);
\fill[slate] (0,6.405) circle (0.28pt);
\fill[slate] (0,6.649) circle (0.28pt);
\fill[slate] (0,5.996) circle (0.28pt);
\fill[slate] (0,6.245) circle (0.28pt);
\fill[slate] (0,6.445) circle (0.28pt);
\fill[slate] (0,6.494) circle (0.28pt);
\fill[slate] (0,5.504) circle (0.28pt);
\fill[slate] (0,5.748) circle (0.28pt);
\fill[slate] (0,5.96) circle (0.28pt);
\fill[slate] (0,6.416) circle (0.28pt);
\fill[slate] (0,6.659) circle (0.28pt);
\fill[slate] (0,5.996) circle (0.28pt);
\fill[slate] (0,6.245) circle (0.28pt);
\fill[slate] (0,6.447) circle (0.28pt);
\fill[slate] (0,6.494) circle (0.28pt);
\fill[slate] (0,6.743) circle (0.28pt);
\fill[slate] (0,6.002) circle (0.28pt);
\fill[slate] (0,6.246) circle (0.28pt);
\fill[slate] (0,6.458) circle (0.28pt);
\fill[slate] (0,6.494) circle (0.28pt);
\fill[slate] (0,6.743) circle (0.28pt);
\fill[slate] (0,6.5) circle (0.28pt);
\fill[slate] (0,6.744) circle (0.28pt);
\fill[navy] (6.935,0.119) circle (0.75pt);
\fill[navy] (1.019,0.366) circle (0.75pt);
\fill[navy] (0.149,0.614) circle (0.75pt);
\fill[navy] (0.022,0.863) circle (0.75pt);
\fill[navy] (0.003,1.112) circle (0.75pt);
\fill[navy] (0,1.361) circle (0.75pt);
\fill[navy] (0,1.61) circle (0.75pt);
\fill[navy] (0,1.859) circle (0.75pt);
\fill[navy] (0,2.108) circle (0.75pt);
\fill[navy] (0,2.357) circle (0.75pt);
\fill[navy] (0,2.607) circle (0.75pt);
\fill[navy] (0,2.856) circle (0.75pt);
\fill[navy] (0,3.105) circle (0.75pt);
\fill[navy] (0,3.354) circle (0.75pt);
\fill[navy] (0,3.603) circle (0.75pt);
\fill[navy] (0,3.852) circle (0.75pt);
\fill[navy] (0,4.101) circle (0.75pt);
\fill[navy] (0,4.35) circle (0.75pt);
\fill[navy] (0,4.599) circle (0.75pt);
\fill[navy] (0,4.848) circle (0.75pt);
\fill[navy] (0,5.097) circle (0.75pt);
\fill[navy] (0,5.346) circle (0.75pt);
\fill[navy] (0,5.595) circle (0.75pt);
\fill[navy] (0,5.844) circle (0.75pt);
\fill[navy] (0,6.093) circle (0.75pt);
\fill[navy] (0,6.342) circle (0.75pt);
\fill[navy] (0,6.591) circle (0.75pt);
\fill[navy] (0,6.84) circle (0.75pt);
\fill[navy] (0,6.542) circle (0.75pt);
\fill[navy] (0,6.044) circle (0.75pt);
\fill[navy] (0,5.546) circle (0.75pt);
\fill[navy] (0,5.048) circle (0.75pt);
\fill[navy] (0,4.55) circle (0.75pt);
\fill[navy] (0,4.052) circle (0.75pt);
\fill[navy] (0,6.494) circle (0.75pt);
\fill[navy] (0,6.245) circle (0.75pt);
\fill[navy] (0,3.554) circle (0.75pt);
\fill[navy] (0,5.747) circle (0.75pt);
\fill[navy] (0,5.498) circle (0.75pt);
\fill[navy] (0.001,3.056) circle (0.75pt);
\fill[navy] (0.001,5) circle (0.75pt);
\fill[navy] (0.001,4.751) circle (0.75pt);
\fill[navy] (0.001,6.445) circle (0.75pt);
\fill[navy] (0.006,2.558) circle (0.75pt);
\fill[navy] (0.006,4.253) circle (0.75pt);
\fill[navy] (0.006,5.947) circle (0.75pt);
\fill[navy] (0.006,4.004) circle (0.75pt);
\fill[navy] (0.006,5.449) circle (0.75pt);
\fill[navy] (0.038,2.06) circle (0.75pt);
\fill[navy] (0.038,3.506) circle (0.75pt);
\fill[navy] (0.038,4.951) circle (0.75pt);
\fill[navy] (0.038,6.397) circle (0.75pt);
\fill[navy] (0.038,6.148) circle (0.75pt);
\fill[navy] (0.038,3.257) circle (0.75pt);
\fill[navy] (0.038,5.899) circle (0.75pt);
\fill[navy] (0.038,4.453) circle (0.75pt);
\fill[navy] (0.038,5.65) circle (0.75pt);
\fill[navy] (0.038,6.846) circle (0.75pt);
\fill[navy] (0.26,1.562) circle (0.75pt);
\fill[navy] (0.26,2.759) circle (0.75pt);
\fill[navy] (0.26,3.955) circle (0.75pt);
\fill[navy] (0.26,5.152) circle (0.75pt);
\fill[navy] (0.26,6.349) circle (0.75pt);
\fill[navy] (0.26,4.903) circle (0.75pt);
\fill[navy] (0.26,2.51) circle (0.75pt);
\fill[navy] (0.26,4.654) circle (0.75pt);
\fill[navy] (0.26,6.798) circle (0.75pt);
\fill[navy] (0.26,3.458) circle (0.75pt);
\fill[navy] (0.26,6.55) circle (0.75pt);
\fill[navy] (0.26,4.405) circle (0.75pt);
\fill[navy] (0.26,5.353) circle (0.75pt);
\fill[navy] (0.26,6.301) circle (0.75pt);
\fill[navy] (1.781,1.065) circle (0.75pt);
\fill[navy] (1.778,2.013) circle (0.75pt);
\fill[navy] (1.778,2.96) circle (0.75pt);
\fill[navy] (1.778,3.908) circle (0.75pt);
\fill[navy] (1.778,4.856) circle (0.75pt);
\fill[navy] (1.778,5.804) circle (0.75pt);
\fill[navy] (1.778,6.751) circle (0.75pt);
\fill[navy] (1.778,5.555) circle (0.75pt);
\fill[navy] (1.778,3.66) circle (0.75pt);
\fill[navy] (1.778,6.255) circle (0.75pt);
\fill[navy] (1.778,5.307) circle (0.75pt);
\fill[navy] (1.785,1.765) circle (0.75pt);
\fill[navy] (1.784,3.412) circle (0.75pt);
\fill[navy] (1.784,5.059) circle (0.75pt);
\fill[navy] (1.784,6.706) circle (0.75pt);
\fill[navy] (1.784,5.758) circle (0.75pt);
\fill[navy] (1.785,2.464) circle (0.75pt);
\fill[navy] (1.785,4.811) circle (0.75pt);
\fill[navy] (1.785,3.164) circle (0.75pt);
\fill[navy] (1.785,6.21) circle (0.75pt);
\fill[navy] (1.785,3.863) circle (0.75pt);
\fill[navy] (1.785,4.563) circle (0.75pt);
\fill[navy] (1.785,5.262) circle (0.75pt);
\fill[navy] (1.785,5.962) circle (0.75pt);
\fill[navy] (1.785,6.661) circle (0.75pt);
\fill[navy] (12.238,0.573) circle (0.75pt);
\fill[navy] (12.085,1.273) circle (0.75pt);
\fill[navy] (12.085,1.972) circle (0.75pt);
\fill[navy] (12.085,2.672) circle (0.75pt);
\fill[navy] (12.085,3.371) circle (0.75pt);
\fill[navy] (12.085,4.071) circle (0.75pt);
\fill[navy] (12.085,4.77) circle (0.75pt);
\fill[navy] (12.085,5.47) circle (0.75pt);
\fill[navy] (12.085,6.169) circle (0.75pt);
\fill[navy] (12.085,6.625) circle (0.75pt);
\fill[navy] (12.085,5.226) circle (0.75pt);
\fill[navy] (12.085,3.827) circle (0.75pt);
\fill[navy] (12.085,6.382) circle (0.75pt);
\fill[navy] (12.085,5.682) circle (0.75pt);
\fill[navy] (12.086,2.428) circle (0.75pt);
\fill[navy] (12.086,4.283) circle (0.75pt);
\fill[navy] (12.086,6.138) circle (0.75pt);
\fill[navy] (12.086,3.584) circle (0.75pt);
\fill[navy] (12.086,6.594) circle (0.75pt);
\fill[navy] (12.086,4.739) circle (0.75pt);
\fill[navy] (12.086,5.894) circle (0.75pt);
\fill[navy] (12.395,1.029) circle (0.75pt);
\fill[navy] (12.391,2.185) circle (0.75pt);
\fill[navy] (12.391,3.34) circle (0.75pt);
\fill[navy] (12.391,4.496) circle (0.75pt);
\fill[navy] (12.391,5.651) circle (0.75pt);
\fill[navy] (12.391,6.806) circle (0.75pt);
\fill[navy] (12.391,6.107) circle (0.75pt);
\fill[navy] (12.391,3.796) circle (0.75pt);
\fill[navy] (12.391,6.563) circle (0.75pt);
\fill[navy] (12.391,5.408) circle (0.75pt);
\fill[navy] (12.4,1.485) circle (0.75pt);
\fill[navy] (12.4,3.097) circle (0.75pt);
\fill[navy] (12.4,4.708) circle (0.75pt);
\fill[navy] (12.4,6.32) circle (0.75pt);
\fill[navy] (12.4,5.165) circle (0.75pt);
\fill[navy] (12.4,1.942) circle (0.75pt);
\fill[navy] (12.4,4.009) circle (0.75pt);
\fill[navy] (12.4,6.077) circle (0.75pt);
\fill[navy] (12.4,6.533) circle (0.75pt);
\fill[navy] (12.4,2.398) circle (0.75pt);
\fill[navy] (12.4,4.921) circle (0.75pt);
\fill[navy] (12.4,2.854) circle (0.75pt);
\fill[navy] (12.4,5.833) circle (0.75pt);
\fill[navy] (12.4,3.31) circle (0.75pt);
\fill[navy] (12.4,6.746) circle (0.75pt);
\fill[navy] (12.4,3.766) circle (0.75pt);
\fill[navy] (12.4,4.222) circle (0.75pt);
\fill[navy] (12.4,4.678) circle (0.75pt);
\fill[navy] (12.4,5.134) circle (0.75pt);
\fill[navy] (12.4,5.59) circle (0.75pt);
\fill[navy] (12.4,6.046) circle (0.75pt);
\fill[navy] (12.4,6.502) circle (0.75pt);
\draw[slate, densely dotted, line width=0.35pt] (6.935,0.119) -- (6.935,1.521);
\draw[slate, densely dotted, line width=0.35pt] (1.019,0.366) -- (1.019,2.262);
\draw[slate, densely dotted, line width=0.35pt] (12.238,0.573) -- (12.238,2.884);
\draw[slate, densely dotted, line width=0.35pt] (6.935,0.119) -- (6.935,2.92);
\draw[slate, densely dotted, line width=0.35pt] (6.935,0.119) -- (6.935,2.925);
\draw[slate, densely dotted, line width=0.35pt] (0.149,0.614) -- (0.149,3.008);
\draw[slate, densely dotted, line width=0.35pt] (0.022,0.863) -- (0.022,3.755);
\draw[slate, densely dotted, line width=0.35pt] (1.019,0.366) -- (1.019,4.157);
\draw[slate, densely dotted, line width=0.35pt] (1.019,0.366) -- (1.019,4.158);
\draw[slate, densely dotted, line width=0.35pt] (12.395,1.029) -- (12.395,4.252);
\draw[slate, densely dotted, line width=0.35pt] (6.935,0.119) -- (6.935,4.319);
\draw[slate, densely dotted, line width=0.35pt] (6.935,0.119) -- (6.935,4.324);
\draw[slate, densely dotted, line width=0.35pt] (6.935,0.119) -- (6.935,4.324);
\draw[slate, densely dotted, line width=0.35pt] (6.935,0.119) -- (6.935,4.328);
\draw[slate, densely dotted, line width=0.35pt] (1.781,1.065) -- (1.781,4.359);
\draw[slate, densely dotted, line width=0.35pt] (0.003,1.112) -- (0.003,4.502);
\draw[slate, densely dotted, line width=0.35pt] (12.085,1.273) -- (12.085,4.983);
\draw[slate, densely dotted, line width=0.35pt] (12.238,0.573) -- (12.238,5.195);
\draw[slate, densely dotted, line width=0.35pt] (12.238,0.573) -- (12.238,5.195);
\draw[slate, densely dotted, line width=0.35pt] (0,1.361) -- (0,5.249);
\draw[slate, densely dotted, line width=0.35pt] (0.149,0.614) -- (0.149,5.401);
\draw[slate, densely dotted, line width=0.35pt] (0.149,0.614) -- (0.149,5.401);
\draw[slate, densely dotted, line width=0.35pt] (12.4,1.485) -- (12.4,5.621);
\draw[slate, densely dotted, line width=0.35pt] (6.935,0.119) -- (6.935,5.718);
\draw[slate, densely dotted, line width=0.35pt] (6.935,0.119) -- (6.935,5.723);
\draw[slate, densely dotted, line width=0.35pt] (6.935,0.119) -- (6.935,5.723);
\draw[slate, densely dotted, line width=0.35pt] (6.935,0.119) -- (6.935,5.723);
\draw[slate, densely dotted, line width=0.35pt] (6.935,0.119) -- (6.935,5.727);
\draw[slate, densely dotted, line width=0.35pt] (1.83,1.521) -- (1.83,5.727);
\draw[slate, densely dotted, line width=0.35pt] (6.935,0.119) -- (6.935,5.727);
\draw[slate, densely dotted, line width=0.35pt] (6.935,0.119) -- (6.935,5.727);
\draw[slate, densely dotted, line width=0.35pt] (12.039,1.521) -- (12.039,5.727);
\draw[slate, densely dotted, line width=0.35pt] (6.935,0.119) -- (6.935,5.732);
\draw[slate, densely dotted, line width=0.35pt] (0.26,1.562) -- (0.26,5.851);
\draw[slate, densely dotted, line width=0.35pt] (0,1.61) -- (0,5.996);
\draw[slate, densely dotted, line width=0.35pt] (1.019,0.366) -- (1.019,6.052);
\draw[slate, densely dotted, line width=0.35pt] (1.019,0.366) -- (1.019,6.053);
\draw[slate, densely dotted, line width=0.35pt] (1.019,0.366) -- (1.019,6.053);
\draw[slate, densely dotted, line width=0.35pt] (1.019,0.366) -- (1.019,6.054);
\draw[slate, densely dotted, line width=0.35pt] (1.785,1.765) -- (1.785,6.458);
\draw[slate, densely dotted, line width=0.35pt] (12.032,1.77) -- (12.032,6.474);
\draw[slate, densely dotted, line width=0.35pt] (0.022,0.863) -- (0.022,6.646);
\draw[slate, densely dotted, line width=0.35pt] (0.022,0.863) -- (0.022,6.646);
\draw[slate, densely dotted, line width=0.35pt] (0,1.859) -- (0,6.743);
\draw[wine, line width=0.7pt] (1.83,1.521) -- (12.039,1.521);
\draw[wine, line width=0.7pt] (0.261,2.262) -- (1.777,2.262);
\draw[wine, line width=0.7pt] (12.086,2.884) -- (12.391,2.884);
\draw[wine, line width=0.7pt] (1.785,2.92) -- (12.084,2.92);
\draw[wine, line width=0.7pt] (1.83,2.925) -- (12.039,2.925);
\draw[wine, line width=0.7pt] (0.038,3.008) -- (0.26,3.008);
\draw[wine, line width=0.7pt] (0.006,3.755) -- (0.038,3.755);
\draw[wine, line width=0.7pt] (0.26,4.157) -- (1.778,4.157);
\draw[wine, line width=0.7pt] (0.261,4.158) -- (1.777,4.158);
\draw[wine, line width=0.7pt] (12.391,4.252) -- (12.4,4.252);
\draw[wine, line width=0.7pt] (1.785,4.319) -- (12.085,4.319);
\draw[wine, line width=0.7pt] (1.785,4.324) -- (12.084,4.324);
\draw[wine, line width=0.7pt] (1.83,4.324) -- (12.039,4.324);
\draw[wine, line width=0.7pt] (1.83,4.328) -- (12.039,4.328);
\draw[wine, line width=0.7pt] (1.778,4.359) -- (1.784,4.359);
\draw[wine, line width=0.7pt] (0.001,4.502) -- (0.006,4.502);
\draw[wine, line width=0.7pt] (12.085,4.983) -- (12.086,4.983);
\draw[wine, line width=0.7pt] (12.086,5.195) -- (12.391,5.195);
\draw[wine, line width=0.7pt] (12.086,5.195) -- (12.391,5.195);
\draw[wine, line width=0.7pt] (0,5.249) -- (0.001,5.249);
\draw[wine, line width=0.7pt] (0.038,5.401) -- (0.26,5.401);
\draw[wine, line width=0.7pt] (0.038,5.401) -- (0.26,5.401);
\draw[wine, line width=0.7pt] (12.4,5.621) -- (12.4,5.621);
\draw[wine, line width=0.7pt] (1.785,5.718) -- (12.085,5.718);
\draw[wine, line width=0.7pt] (1.785,5.723) -- (12.085,5.723);
\draw[wine, line width=0.7pt] (1.83,5.723) -- (12.039,5.723);
\draw[wine, line width=0.7pt] (1.785,5.723) -- (12.084,5.723);
\draw[wine, line width=0.7pt] (1.785,5.727) -- (12.084,5.727);
\draw[wine, line width=0.7pt] (1.83,5.727) -- (1.83,5.727);
\draw[wine, line width=0.7pt] (1.83,5.727) -- (12.039,5.727);
\draw[wine, line width=0.7pt] (1.83,5.727) -- (12.039,5.727);
\draw[wine, line width=0.7pt] (12.039,5.727) -- (12.039,5.727);
\draw[wine, line width=0.7pt] (1.83,5.732) -- (12.039,5.732);
\draw[wine, line width=0.7pt] (0.26,5.851) -- (0.26,5.851);
\draw[wine, line width=0.7pt] (0,5.996) -- (0,5.996);
\draw[wine, line width=0.7pt] (0.26,6.052) -- (1.778,6.052);
\draw[wine, line width=0.7pt] (0.26,6.053) -- (1.778,6.053);
\draw[wine, line width=0.7pt] (0.261,6.053) -- (1.777,6.053);
\draw[wine, line width=0.7pt] (0.261,6.054) -- (1.777,6.054);
\draw[wine, line width=0.7pt] (1.784,6.458) -- (1.785,6.458);
\draw[wine, line width=0.7pt] (12.032,6.474) -- (12.032,6.474);
\draw[wine, line width=0.7pt] (0.006,6.646) -- (0.038,6.646);
\draw[wine, line width=0.7pt] (0.006,6.646) -- (0.038,6.646);
\draw[wine, line width=0.7pt] (0,6.743) -- (0,6.743);
\draw[wine, fill=white, line width=0.45pt] (0,6.743) circle (1.1pt);
\draw[wine, fill=white, line width=0.45pt] (0,6.743) circle (1.1pt);
\draw[wine, fill=white, line width=0.45pt] (0,5.996) circle (1.1pt);
\draw[wine, fill=white, line width=0.45pt] (0,5.996) circle (1.1pt);
\draw[wine, fill=white, line width=0.45pt] (0,5.249) circle (1.1pt);
\draw[wine, fill=white, line width=0.45pt] (0.001,5.249) circle (1.1pt);
\draw[wine, fill=white, line width=0.45pt] (0.001,4.502) circle (1.1pt);
\draw[wine, fill=white, line width=0.45pt] (0.006,4.502) circle (1.1pt);
\draw[wine, fill=white, line width=0.45pt] (0.006,6.646) circle (1.1pt);
\draw[wine, fill=white, line width=0.45pt] (0.006,3.755) circle (1.1pt);
\draw[wine, fill=white, line width=0.45pt] (0.006,6.646) circle (1.1pt);
\draw[wine, fill=white, line width=0.45pt] (0.038,3.755) circle (1.1pt);
\draw[wine, fill=white, line width=0.45pt] (0.038,6.646) circle (1.1pt);
\draw[wine, fill=white, line width=0.45pt] (0.038,6.646) circle (1.1pt);
\draw[wine, fill=white, line width=0.45pt] (0.038,5.401) circle (1.1pt);
\draw[wine, fill=white, line width=0.45pt] (0.038,3.008) circle (1.1pt);
\draw[wine, fill=white, line width=0.45pt] (0.038,5.401) circle (1.1pt);
\draw[wine, fill=white, line width=0.45pt] (0.26,3.008) circle (1.1pt);
\draw[wine, fill=white, line width=0.45pt] (0.26,5.401) circle (1.1pt);
\draw[wine, fill=white, line width=0.45pt] (0.26,5.401) circle (1.1pt);
\draw[wine, fill=white, line width=0.45pt] (0.26,5.851) circle (1.1pt);
\draw[wine, fill=white, line width=0.45pt] (0.26,5.851) circle (1.1pt);
\draw[wine, fill=white, line width=0.45pt] (0.26,6.052) circle (1.1pt);
\draw[wine, fill=white, line width=0.45pt] (0.26,6.053) circle (1.1pt);
\draw[wine, fill=white, line width=0.45pt] (0.26,4.157) circle (1.1pt);
\draw[wine, fill=white, line width=0.45pt] (0.261,2.262) circle (1.1pt);
\draw[wine, fill=white, line width=0.45pt] (0.261,4.158) circle (1.1pt);
\draw[wine, fill=white, line width=0.45pt] (0.261,6.054) circle (1.1pt);
\draw[wine, fill=white, line width=0.45pt] (0.261,6.053) circle (1.1pt);
\draw[wine, fill=white, line width=0.45pt] (1.777,2.262) circle (1.1pt);
\draw[wine, fill=white, line width=0.45pt] (1.777,6.053) circle (1.1pt);
\draw[wine, fill=white, line width=0.45pt] (1.777,4.158) circle (1.1pt);
\draw[wine, fill=white, line width=0.45pt] (1.777,6.054) circle (1.1pt);
\draw[wine, fill=white, line width=0.45pt] (1.778,4.157) circle (1.1pt);
\draw[wine, fill=white, line width=0.45pt] (1.778,6.053) circle (1.1pt);
\draw[wine, fill=white, line width=0.45pt] (1.778,6.052) circle (1.1pt);
\draw[wine, fill=white, line width=0.45pt] (1.778,4.359) circle (1.1pt);
\draw[wine, fill=white, line width=0.45pt] (1.784,4.359) circle (1.1pt);
\draw[wine, fill=white, line width=0.45pt] (1.784,6.458) circle (1.1pt);
\draw[wine, fill=white, line width=0.45pt] (1.785,6.458) circle (1.1pt);
\draw[wine, fill=white, line width=0.45pt] (1.785,5.718) circle (1.1pt);
\draw[wine, fill=white, line width=0.45pt] (1.785,5.723) circle (1.1pt);
\draw[wine, fill=white, line width=0.45pt] (1.785,4.319) circle (1.1pt);
\draw[wine, fill=white, line width=0.45pt] (1.785,5.727) circle (1.1pt);
\draw[wine, fill=white, line width=0.45pt] (1.785,4.324) circle (1.1pt);
\draw[wine, fill=white, line width=0.45pt] (1.785,2.92) circle (1.1pt);
\draw[wine, fill=white, line width=0.45pt] (1.785,5.723) circle (1.1pt);
\draw[wine, fill=white, line width=0.45pt] (1.83,1.521) circle (1.1pt);
\draw[wine, fill=white, line width=0.45pt] (1.83,5.727) circle (1.1pt);
\draw[wine, fill=white, line width=0.45pt] (1.83,2.925) circle (1.1pt);
\draw[wine, fill=white, line width=0.45pt] (1.83,4.328) circle (1.1pt);
\draw[wine, fill=white, line width=0.45pt] (1.83,5.732) circle (1.1pt);
\draw[wine, fill=white, line width=0.45pt] (1.83,4.324) circle (1.1pt);
\draw[wine, fill=white, line width=0.45pt] (1.83,5.727) circle (1.1pt);
\draw[wine, fill=white, line width=0.45pt] (1.83,5.723) circle (1.1pt);
\draw[wine, fill=white, line width=0.45pt] (12.032,6.474) circle (1.1pt);
\draw[wine, fill=white, line width=0.45pt] (12.032,6.474) circle (1.1pt);
\draw[wine, fill=white, line width=0.45pt] (12.039,1.521) circle (1.1pt);
\draw[wine, fill=white, line width=0.45pt] (12.039,5.723) circle (1.1pt);
\draw[wine, fill=white, line width=0.45pt] (12.039,5.727) circle (1.1pt);
\draw[wine, fill=white, line width=0.45pt] (12.039,4.324) circle (1.1pt);
\draw[wine, fill=white, line width=0.45pt] (12.039,2.925) circle (1.1pt);
\draw[wine, fill=white, line width=0.45pt] (12.039,4.328) circle (1.1pt);
\draw[wine, fill=white, line width=0.45pt] (12.039,5.732) circle (1.1pt);
\draw[wine, fill=white, line width=0.45pt] (12.039,5.727) circle (1.1pt);
\draw[wine, fill=white, line width=0.45pt] (12.084,2.92) circle (1.1pt);
\draw[wine, fill=white, line width=0.45pt] (12.084,5.723) circle (1.1pt);
\draw[wine, fill=white, line width=0.45pt] (12.084,4.324) circle (1.1pt);
\draw[wine, fill=white, line width=0.45pt] (12.084,5.727) circle (1.1pt);
\draw[wine, fill=white, line width=0.45pt] (12.085,4.319) circle (1.1pt);
\draw[wine, fill=white, line width=0.45pt] (12.085,5.723) circle (1.1pt);
\draw[wine, fill=white, line width=0.45pt] (12.085,5.718) circle (1.1pt);
\draw[wine, fill=white, line width=0.45pt] (12.085,4.983) circle (1.1pt);
\draw[wine, fill=white, line width=0.45pt] (12.086,4.983) circle (1.1pt);
\draw[wine, fill=white, line width=0.45pt] (12.086,5.195) circle (1.1pt);
\draw[wine, fill=white, line width=0.45pt] (12.086,2.884) circle (1.1pt);
\draw[wine, fill=white, line width=0.45pt] (12.086,5.195) circle (1.1pt);
\draw[wine, fill=white, line width=0.45pt] (12.391,2.884) circle (1.1pt);
\draw[wine, fill=white, line width=0.45pt] (12.391,5.195) circle (1.1pt);
\draw[wine, fill=white, line width=0.45pt] (12.391,5.195) circle (1.1pt);
\draw[wine, fill=white, line width=0.45pt] (12.391,4.252) circle (1.1pt);
\draw[wine, fill=white, line width=0.45pt] (12.4,4.252) circle (1.1pt);
\draw[wine, fill=white, line width=0.45pt] (12.4,5.621) circle (1.1pt);
\draw[wine, fill=white, line width=0.45pt] (12.4,5.621) circle (1.1pt);
\filldraw[ochre, draw=ink, line width=0.25pt] (6.935,0.119) +(0,1.9pt) -- +(1.9pt,0) -- +(0,-1.9pt) -- +(-1.9pt,0) -- cycle;
\filldraw[ochre, draw=ink, line width=0.25pt] (1.019,0.366) +(0,1.9pt) -- +(1.9pt,0) -- +(0,-1.9pt) -- +(-1.9pt,0) -- cycle;
\filldraw[ochre, draw=ink, line width=0.25pt] (0.149,0.614) +(0,1.9pt) -- +(1.9pt,0) -- +(0,-1.9pt) -- +(-1.9pt,0) -- cycle;
\filldraw[ochre, draw=ink, line width=0.25pt] (0.022,0.863) +(0,1.9pt) -- +(1.9pt,0) -- +(0,-1.9pt) -- +(-1.9pt,0) -- cycle;
\filldraw[ochre, draw=ink, line width=0.25pt] (0.003,1.112) +(0,1.9pt) -- +(1.9pt,0) -- +(0,-1.9pt) -- +(-1.9pt,0) -- cycle;
\filldraw[ochre, draw=ink, line width=0.25pt] (0,1.361) +(0,1.9pt) -- +(1.9pt,0) -- +(0,-1.9pt) -- +(-1.9pt,0) -- cycle;
\filldraw[ochre, draw=ink, line width=0.25pt] (0,1.61) +(0,1.9pt) -- +(1.9pt,0) -- +(0,-1.9pt) -- +(-1.9pt,0) -- cycle;
\filldraw[ochre, draw=ink, line width=0.25pt] (0,1.859) +(0,1.9pt) -- +(1.9pt,0) -- +(0,-1.9pt) -- +(-1.9pt,0) -- cycle;
\filldraw[ochre, draw=ink, line width=0.25pt] (0.26,1.562) +(0,1.9pt) -- +(1.9pt,0) -- +(0,-1.9pt) -- +(-1.9pt,0) -- cycle;
\filldraw[ochre, draw=ink, line width=0.25pt] (1.781,1.065) +(0,1.9pt) -- +(1.9pt,0) -- +(0,-1.9pt) -- +(-1.9pt,0) -- cycle;
\filldraw[ochre, draw=ink, line width=0.25pt] (1.785,1.765) +(0,1.9pt) -- +(1.9pt,0) -- +(0,-1.9pt) -- +(-1.9pt,0) -- cycle;
\filldraw[ochre, draw=ink, line width=0.25pt] (1.83,1.521) +(0,1.9pt) -- +(1.9pt,0) -- +(0,-1.9pt) -- +(-1.9pt,0) -- cycle;
\draw[wine, line width=0.45pt] (1.83,1.521) circle (2.7pt);
\filldraw[ochre, draw=ink, line width=0.25pt] (12.238,0.573) +(0,1.9pt) -- +(1.9pt,0) -- +(0,-1.9pt) -- +(-1.9pt,0) -- cycle;
\filldraw[ochre, draw=ink, line width=0.25pt] (12.032,1.77) +(0,1.9pt) -- +(1.9pt,0) -- +(0,-1.9pt) -- +(-1.9pt,0) -- cycle;
\filldraw[ochre, draw=ink, line width=0.25pt] (12.039,1.521) +(0,1.9pt) -- +(1.9pt,0) -- +(0,-1.9pt) -- +(-1.9pt,0) -- cycle;
\draw[wine, line width=0.45pt] (12.039,1.521) circle (2.7pt);
\filldraw[ochre, draw=ink, line width=0.25pt] (12.085,1.273) +(0,1.9pt) -- +(1.9pt,0) -- +(0,-1.9pt) -- +(-1.9pt,0) -- cycle;
\filldraw[ochre, draw=ink, line width=0.25pt] (12.395,1.029) +(0,1.9pt) -- +(1.9pt,0) -- +(0,-1.9pt) -- +(-1.9pt,0) -- cycle;
\filldraw[ochre, draw=ink, line width=0.25pt] (12.4,1.485) +(0,1.9pt) -- +(1.9pt,0) -- +(0,-1.9pt) -- +(-1.9pt,0) -- cycle;
\node[ink, right=2.2pt] at (6.935,0.119) {$\varepsilon$};
\node[ink, right=2.2pt] at (1.019,0.366) {$a$};
\node[ink, left=2.2pt] at (12.238,0.573) {$b$};
\fill[navy] (0.2,7.42) circle (0.75pt); \node[ink, right] at (0.28,7.42) {central word};
\filldraw[ochre, draw=ink, line width=0.25pt] (2.6,7.42) +(0,1.9pt) -- +(1.9pt,0) -- +(0,-1.9pt) -- +(-1.9pt,0) -- cycle; \node[ink, right] at (2.7,7.42) {centre $p$};
\draw[wine, line width=0.7pt] (4.5,7.42) -- (5.1,7.42); \draw[wine, fill=white, line width=0.45pt] (4.5,7.42) circle (1.1pt) (5.1,7.42) circle (1.1pt); \node[ink, right] at (5.2,7.42) {twins and their rung};
\fill[slate] (8.85,7.42) circle (0.5pt); \node[ink, right] at (8.93,7.42) {other palindrome};
\end{tikzpicture}

\caption{All palindromes with Gram label at most $10^{12}$, placed by the ratio $\rho/M$ and
their Gram label $M$. The ratio $\rho/M=t/(2t+1)$ increases with the slope $t$ of \cref{lem:ladder}, and it is
the frequency of the letter $b$ in $\Omega_z$. Each horizontal segment joins two twins exchanged
around a palindrome $p$ (diamond), at the height of their common label; the dotted post shows
that the segment is centred at the ratio $\rho_p/M_p$ of $p$, which by \cref{lem:ladder} is the
lowest point between its ends. Ladders rise above the empty word, above $a$ and above $b$. The
two ringed diamonds, $abba$ and $baab$, are twins themselves and centres of the exchanges at
$12\,986\,074\,130$, the square of \cref{fig:nested-square}. No central word lies on a rung.}
\label{fig:ladders}
\end{figure}

\Cref{fig:ladders} shows the twins in the plane of the ratio $\rho/M$ and the label. They appear
as the rungs of ladders, which are already visible in the alphabetical order of the directive
words, and the following lemma explains why each rung stands on a palindrome of smaller
label.

\begin{lemma}[Ladders]
\label{lem:ladder}
For a word $w$ let $t(w)$ be the slope of the column $\mu(w)e$. For every word $q$, the slopes
of the words beginning with $q$ lie in the interval $\mu(q)(\mathcal J)$, which contains the
slope of no other word. Consequently, if $Z$ and $Z'$ are twins and $q$ is their longest
common prefix, then every word whose slope lies strictly between $t(Z)$ and $t(Z')$ begins
with $q$, the slope $t(q)$ is one of them, and $q$ has strictly the smallest label among
these words.
\end{lemma}

\begin{proof}
By \cref{sec:decoder}, $A(\mathcal J)$ and $B(\mathcal J)$ lie in $\mathcal J$, on either side
of $2$, and both maps are increasing. A word $qw'$ has slope
$\mu(q)(t(w'))\in\mu(q)(\mathcal J)$. A word that is not of this form either first differs
from $q$ at some letter, after a common prefix $u$, and then its slope lies in
$\mu(u)\mu(c)(\mathcal J)$ while $\mu(q)(\mathcal J)\subset\mu(u)\mu(c')(\mathcal J)$ with
$c\neq c'$, two disjoint intervals; or it is a proper prefix $u$ of $q=uy$, and then its
slope $\mu(u)(2)$ avoids $\mu(u)\mu(y)(\mathcal J)$ because $2\notin A(\mathcal J)\cup B(\mathcal J)$.
This proves the first assertion. Next, for a nonempty word $y=cy'$ the column $\mu(y)e-e$
has positive entries: by induction $\mu(y')e-e$ has nonnegative entries, and then
$\mu(y)e-e=\mu(c)(\mu(y')e-e)+(\mu(c)e-e)$, where $\mu(c)$ has nonnegative entries and
$\mu(c)e-e$ is $(3,2)^{\mathsf T}$ or $(10,4)^{\mathsf T}$. Since the row $e^{\mathsf T}\mu(q)$ has
positive entries, $e^{\mathsf T}\mu(qy)e-e^{\mathsf T}\mu(q)e=e^{\mathsf T}\mu(q)(\mu(y)e-e)>0$.
In particular a twin cannot be a proper prefix of another twin. Hence $Z=q\,c\,Z_1$ and $Z'=q\,c'\,Z_1'$ with $c\neq c'$, say $c=a$. Both slopes
lie in the interval $\mu(q)(\mathcal J)$, and so does everything between them; and
$t(Z)<\mu(q)(2)=t(q)<t(Z')$ because $t(aZ_1)<2<t(bZ_1')$. The words in between begin with $q$,
and the same inequality of labels shows that $q$ is the lowest among them.
\end{proof}

For an exchange pair around $p$ the foot of the rung is $p$ itself, and by
\cref{prop:any-centre} the rung is centred on it in the ratio $\rho/M$, which is an increasing
function of the slope $t$, so that the lemma applies to it unchanged. Since the lemma holds for
every pair of twins, the ladders alone cannot decide the conjecture. In this picture, the
conjecture implies that every palindrome with a twin has at least one rung centred on a
palindrome, and \cref{subsec:centre-decoding} shows that this consequence is equivalent to
part~(i).

\subsection{The midpoint statement}
\label{subsec:centre-decoding}

\Cref{thm:centre-decoder} proves the decoder at every palindromic centre, and
\cref{cor:centre-classification} turns it into a classification: two palindromes with equal
labels whose root midpoint is the ratio $\rho_p/M_p$ of a palindrome $p$ are exchanged around $p$.
This makes part~(i) of \cref{conj:structure} equivalent to one statement about midpoints:
\emph{every palindrome with a twin has a twin whose root midpoint is the ratio of a palindrome.}
Indeed, if a palindrome $Z$ has a twin $Z'$ whose root midpoint is $\rho_p/M_p$, then $Z$ is
structured around $p$ by \cref{cor:centre-classification}; conversely a structured palindrome
has its exchange as a twin with that midpoint, by \cref{prop:any-centre}. No condition on a
common prefix is needed: the corollary forces both palindromes to begin with $p$. Part~(ii) is
equivalent in the same way to the statement that any two twins are joined by a chain of twins
in which consecutive words have their root midpoint at the ratio of a palindrome. This is the
displaced problem of \cref{sec:displaced}, now posed for all palindromes. \Cref{rem:quadruple}
shows that the twin has to be chosen, since some twins are displaced, and the enumeration finds
that every displaced pair is a diagonal of a square of \cref{cor:nested-square}, that is, a
composition of two exchanges. For each pair the centre, when it exists, is unique, because the
ratio of a palindrome determines it.

\subsection{Lengths of directing words}
\label{subsec:lengths}

In the Thue--Morse coding a coincidence of labels is a coincidence of the lengths
$\abs{\operatorname{Pal}(\theta(av))}$ and $\abs{\operatorname{Pal}(\theta(av'))}$ of two central
words, and it is natural to ask whether it forces the two directing words $av$ and $av'$ to have
the same length. The question matters because the directing words of one length form a finite
family, so that a positive answer would allow the coincidences to be studied one length at a
time. This subsection gives the length its meaning in the Cohn coding, shows that
\cref{conj:structure} implies a positive answer, and shows that for Markov words a positive answer
is equivalent to the uniqueness conjecture.

\begin{lemma}[The length of a directing word]
\label{lem:length-cf}
For a word $z$ put $M=e^{\mathsf T}\mu(z)e$ and $\rho=(\mu(z)e)_1$, which are the Gram label and
the Gram root when $z$ is a palindrome, and let $\sigma(z)$ be the sequence of integers obtained
from $z$ by replacing every letter $a$ by $1,1$ and every letter $b$ by $2,2$. Then
$\abs{F_a(z)}=\abs z_a+2\abs z_b+1$, and the regular continued fraction of $M/\rho$ is
$$
 \frac M\rho=[\,2;\ \sigma(z),\ 2\,].
$$
In particular the sum of the partial quotients of $M/\rho$ is $2\abs{F_a(z)}+2$, the directive
$\theta(F_a(z))$ has $2\abs{F_a(z)}$ letters, and $\abs{F_a(z)}=\abs{azb}_b+\abs{azb}-2$.
\end{lemma}

\begin{proof}
The morphism $\phi_a$ sends the letter $a$ to one letter and the letter $b$ to two, which gives the
first formula, and $\theta$ doubles every length. For a positive integer $d$ put
$D(d)=\left(\begin{smallmatrix}d&1\\1&0\end{smallmatrix}\right)$, so that $P_1=D(1)$ and $P=P_2=D(2)$,
with $A=P_1^2$ and $B=P_2^2$. Hence $P\mu(z)P=D(d_0)D(d_1)\cdots D(d_N)$, where
$(d_0,\dots,d_N)=(2,\sigma(z),2)$. Let $K$ denote the continuants, defined by $K()=1$, $K(d_0)=d_0$
and $K(d_0,\dots,d_N)=d_0K(d_1,\dots,d_N)+K(d_2,\dots,d_N)$. The first column of
$D(d_0)\cdots D(d_N)$ is $(K(d_0,\dots,d_N),K(d_1,\dots,d_N))^{\mathsf T}$: for a single factor it is
$(d_N,1)^{\mathsf T}=(K(d_N),K())^{\mathsf T}$, and multiplying on the left by $D(d_0)$ replaces a first
column $(x,y)^{\mathsf T}$ by $(d_0x+y,x)^{\mathsf T}$, which is the recurrence. On the other hand the
first column of $P\mu(z)P$ is $P\mu(z)e$, because $Pe_1=e$; with $\mu(z)e=(x,y)^{\mathsf T}$ it equals
$(2x+y,x)^{\mathsf T}=(M,\rho)^{\mathsf T}$. Therefore $M=K(2,\sigma(z),2)$ and $\rho=K(\sigma(z),2)$.
The quotient of these continuants is the continued fraction: $K(d_0)/K()=d_0$, and
$[d_0;d_1,\dots,d_N]=d_0+1/[d_1;\dots,d_N]=d_0+K(d_2,\dots,d_N)/K(d_1,\dots,d_N)$ by induction, which
is $K(d_0,\dots,d_N)/K(d_1,\dots,d_N)$ by the recurrence. All partial quotients are positive and the
last one equals two, so this is the regular continued fraction of $M/\rho$, which is unique. The pair
$1,1$ adds $2$ to the sum of the partial quotients and the pair $2,2$ adds $4$, so the sum is
$2+2\abs z_a+4\abs z_b+2=2\abs{F_a(z)}+2$. Finally $\abs{azb}_b=\abs z_b+1$ and $\abs{azb}=\abs z+2$.
\end{proof}

The length of the directing word is thus half the sum of the partial quotients of the label divided
by the root, minus one: it counts, with multiplicity, the steps of the Euclidean algorithm on the
pair $(M,\rho)$. For a central word $z$ the identity $M=K(2,\sigma(z),2)$ is the continued-fraction
description of Markov numbers given by Frobenius~\cite[formula~(6), p.~472]{Frobenius1913}; the
proof above does not use centrality. The value of a word in the matrices of Lagisquet, Pelantov\'a, Tavenas and
Vuillon~\cite{LagisquetEtAl2021} connects this length with their ordering of Markov numbers.

\begin{proposition}[Lengths of directing words of Markov words]
\label{prop:lengths-uniqueness}
For a word $w$ let $m(w)=(1,0)\,M^w\,(0,1)^{\mathsf T}$, where $M^w$ is the product along $w$ of
$A_L=\left(\begin{smallmatrix}1&1\\1&2\end{smallmatrix}\right)$ for the letter $a$ and
$B_L=\left(\begin{smallmatrix}2&1\\1&1\end{smallmatrix}\right)$ for the letter $b$; this is the
$m$-value of~\cite{LagisquetEtAl2021}. Then:
\begin{enumerate}[label=\textup{(\roman*)}]
\item $m(a\,F_a(z)\,b)=e^{\mathsf T}\mu(z)e$ for every word $z$;
\item if $z$ is central, then $aF_a(z)b$ is the lower Christoffel word with $p=\abs z_b+1$ letters $b$
 and $q=\abs z+2$ letters $a$, the label $M(z)$ is the Markov number of the Farey fraction $p/q$,
 and $\abs{F_a(z)}=p+q-2$;
\item the Frobenius--Markov uniqueness conjecture holds if and only if any two central words with
 the same label have directing words of the same length; the same holds when the length of the
 directing word is replaced by the number of letters $b$ of the word, or by its length.
\end{enumerate}
\end{proposition}

\begin{proof}
\emph{Part (i).} Put $X=\left(\begin{smallmatrix}1&-1\\0&1\end{smallmatrix}\right)$. Direct
multiplication gives
$$
 XAX^{-1}=\begin{pmatrix}1&1\\1&2\end{pmatrix}=A_L,\qquad
 XBX^{-1}=\begin{pmatrix}3&4\\2&3\end{pmatrix}=B_LA_L .
$$
Since $\phi_a(a)=a$ and $\phi_a(b)=ba$, this says $M^{\phi_a(c)}=X\mu(c)X^{-1}$ for both letters $c$,
and both sides are multiplicative in the word, so $M^{\phi_a(w)}=X\mu(w)X^{-1}$ for every word $w$.
As $aF_a(z)b=a\,a\,\phi_a(z)\,b$,
$$
 m(aF_a(z)b)=(1,0)\,A_L^2X\,\mu(z)\,X^{-1}B_L\binom01 .
$$
Now $(1,0)A_L^2=(2,3)$ and $(2,3)X=(2,1)=e^{\mathsf T}$, while $B_L(0,1)^{\mathsf T}=(1,1)^{\mathsf T}$
and $X^{-1}(1,1)^{\mathsf T}=(2,1)^{\mathsf T}=e$. This proves (i).

\emph{Part (ii).} If $z=\operatorname{Pal}(v)$, then
\eqref{eq:standard-palindrome-image} gives $F_a(z)=\operatorname{Pal}(av)$, a central word, so
$aF_a(z)b$ is a lower Christoffel word~\cite[Section~3]{BertheDeLucaReutenauer2008}. The morphism
$\phi_a$ keeps the number of letters $b$ and adds one letter $a$ for each letter of $z$; hence the
word $aF_a(z)b$ has $\abs z_b+1=p$ letters $b$ and $2+\abs z_a+\abs z_b=\abs z+2=q$ letters $a$, and
$p<q$. A lower Christoffel word is determined by its letter counts, which are coprime, so
$aF_a(z)b$ is the word written $c(p,q)$ in~\cite{LagisquetEtAl2021}, and by part (i) and
\cite[Proposition~3]{LagisquetEtAl2021} its value $M(z)$ is the Markov number of the Farey fraction
$p/q$; we write $m_{p/q}$ for that Markov number. By \cref{lem:length-cf}, $\abs{F_a(z)}=p+q-2$.
Conversely, every reduced fraction $0<p/q<1$ arises from exactly one central $z$. Write
$c(p,q)=a\,w\,b$ with $w$ central; $w$ is not empty, since it has $q-1\geqslant1$ letters $a$. If the
directive of $w$ began with $b$, then $w=F_b(w')=b\,\phi_b(w')$ with $\phi_b(a)=ab$, $\phi_b(b)=b$:
every letter $a$ of $w$ would be followed by a letter $b$, and $w$ would have more letters $b$ than
letters $a$. Since $w$ has $q-1>p-1$ letters $a$, its directive begins with $a$, so $w=F_a(z)$ for a
central $z$, unique because $F_a$ is injective.

\emph{Part (iii).} The uniqueness conjecture is equivalent to the injectivity of $p/q\mapsto m_{p/q}$
on the reduced fractions of $(0,1)$, since the Markov tree indexes by these fractions the
normalized Markov triples with distinct entries; by part (ii) this is the injectivity of the label
on central words, and it makes the three statements vacuously true. Conversely suppose that two
distinct central words have the same label, and let $p/q\neq p'/q'$ be their fractions. Write
$m_{n,d}$ for the least value of $m$ on the lattice paths from $(0,0)$ to $(d,n)$; by
\cite[Proposition~21]{LagisquetEtAl2021} it is attained at $c(n,d)$, so $m_{p,q}$ and
$m_{p',q'}$ are the two equal labels. By \cite[Theorem~2]{LagisquetEtAl2021}, for integers $j<k$,
$$
 m_{j,k}<m_{j,k+1},\qquad m_{j,k}<m_{j+1,k},\qquad m_{j+1,k}<m_{j,k+1}.
$$
If $p+q=p'+q'=s$, then $p\neq p'$, say $p'<p$; for $p'\leqslant j\leqslant p-1$ one has
$j<q-1<s-j-1$, so the third inequality applied to $j<s-j-1$ gives $m_{j+1,s-j-1}<m_{j,s-j}$, and
chaining from $j=p-1$ down to $j=p'$ gives $m_{p,q}<m_{p',q'}$, a contradiction. If $p=p'$, then
$q\neq q'$, say $q<q'$, and the first inequality applied to $p<k$ for $q\leqslant k<q'$ gives
$m_{p,q}<m_{p,q'}$. If $q=q'$, then $p\neq p'$, say $p<p'$, and the second inequality applied to
$j<q$ for $p\leqslant j<p'$ gives $m_{p,q}<m_{p',q}$. Each case contradicts the equality of the
labels, so a coincidence of labels among distinct central words changes the length of the
directing word, the number of letters $b$, and the length. Each of the three statements therefore
implies the uniqueness conjecture.
\end{proof}

Nothing in the proposition is deep, and most of it is known. Part~(i) is, in our normalization,
the identity $g(w)=m(aG(w)ab)$ of~\cite[Appendix~B]{LagisquetEtAl2021}, where $G(a)=a$ and $G(b)=ab$,
so that $aG(w)ab=aF_a(w)b$, and where $g$ is the function of Aigner's book, which is the label
$e^{\mathsf T}\mu(w)e$ written in another basis. Part~(iii) is an immediate consequence of the
ordering theorems, and Lee, Li, Rabideau and Schiffler prove more: $m_{p/q}=m_{p'/q'}$ with
$(p,q)\neq(p',q')$ forces $-5/4<(p-p')/(q-q')<-8/7$~\cite[Corollary~1.5]{LeeEtAl2023}. Since
$\abs{F_a(z)}=p+q-2$, a coincidence would change the length of the directing word by more than a
seventh and less than a quarter of the change of $q$, and never by zero. The point of the proposition
is the identification it makes: the directing word of the Thue--Morse coding is the central factor of
the Christoffel word of the fixed-sum theorem, so that the length question is an equivalent form of
the uniqueness conjecture.

Three consequences place the proposition. First, part~(ii) of \cref{conj:structure} implies that
twins have the same numbers of letters $a$ and $b$, because an exchange does not change them; so the
conjecture implies that twins have directing words of the same length, and by the proposition this
alone already implies uniqueness. Second, the equality of letter counts is a theorem on every
reflection axis: by \cref{cor:centre-classification}, two twins whose root midpoint is the ratio of
a palindrome are exchanged around it, so they have the same numbers of letters $a$ and $b$ and
directing words of the same length; in the enumeration up to $10^{26}$ every class of twins has a
single pair of letter counts. Third, in the language
of~\cite{LagisquetEtAl2021}, part~(ii) says that their Christoffel word $c(p,q)$ is $aF_a(z)b$: the
directing word of the Thue--Morse coding is its central factor, two letters shorter. Their fixed-sum
theorem says that two distinct lower Christoffel words with more letters $a$ than letters $b$ and
of the same length carry different Markov numbers; the length question asks for
the complementary statement, that a coincidence of Markov numbers keeps the length, and the two
statements together are the uniqueness conjecture.

\subsection{Off the axis}

Moving the midpoint breaks the reflection in an explicit way. The two matrix factors of
\cref{prop:displaced-equation} remain involutions, but they differ, and the left factor moves
the terminal column by a rank-one displacement. The integer defect and the parent-sensitive
bound of \cref{thm:displaced-ancestor-bound} give coordinates for this remaining problem: they
show that a balanced common ancestor of two distinct central twins must be small compared with
their label. They do not bound the two directive tails or prevent early branching. A proof of
uniqueness along these lines would have to show that central words cannot realize a displaced
midpoint, and centrality must enter there, since displaced twins exist among palindromes. The
displaced twins found by the enumeration are all diagonals of the squares of
\cref{cor:nested-square}: each pair is a composition of two exchanges, and its root is translated by
twice the label times the difference of the two centre ratios.

\raggedbottom
\section*{Supplementary material}
The supplementary material accompanying this article contains the exact computations used in the proofs
of \cref{thm:word-decoder,thm:palindromic-classification,thm:displaced-ancestor-bound},
the worked examples of \cref{tab:pairs}, the identities and constants of the decoder at every
palindromic centre in \cref{subsec:centre-decoder}, the nested cubes of \cref{subsec:nested}, the
enumeration of twins of \cref{sec:outlook}, the identities of \cref{subsec:lengths}, the data of
the three figures, and instructions for repeating the verification.
\section*{Code and data availability}
The source code and the data supporting this article, with the instructions
that repeat every calculation, are contained in the supplementary material of
this article, which is available, together with that of the other articles of
this series:

{\centering
\url{https://github.com/dutykh/fibonacci-pell-nearest-gaps/}
\par}

\section*{Conflict of interest}
The authors declare that they have no conflict of interest.

\bibliographystyle{amsplain}
\bibliography{references}
\end{document}
